\documentclass[fs85,seriesWit]{korfarm-base}
\usepackage{korfarm-witt}

% ============================================================
% passage-note.tex (v11)
% 지문 박스 + 첫 페이지 우측 메모란 (동적 높이)
% 정책: PAGE_BREAK_POLICY.md §3, §4 참조
%
% 변경 (v11):
%  · 지문 박스 폭 확대: enlarge right by=-3.7cm  (메모 3.4cm + gap 3mm)
%  · 메모란 동적 높이: frame.north east ~ frame.south east 사이 자동 채움
%  · 메모 줄선: clip 영역으로 박스 내부만 채움 (박스 높이 모르고도 작동)
%  · 문단 들여쓰기 활성화: tcolorbox 내부 \parindent=1em, \noindent 제거
%
% 사용:
%   \kfPassageNote{제목}{지문 본문}      → 메모란 포함
%   \kfPassageNoteNT{지문 본문}          → 메모란 포함, 제목 없음
%   \kfPassageFull{제목}{본문}           → 메모란 없음 (활동 내 보기)
% ============================================================

\makeatletter

% 메모란 규격 (cm 고정)
\newlength{\kf@memoW}
\setlength{\kf@memoW}{3.4cm}
\newlength{\kf@memoGap}
\setlength{\kf@memoGap}{3mm}

% --- 메인: 제목 있는 버전 ---
\newcommand{\kf@passagenote@with}[2]{%
  \par\ifkfPassageNewPage\ifdim\pagetotal>0.4\textheight\clearpage\fi\fi\smallskip\needspace{6\baselineskip}%
  \begin{tcolorbox}[
    enhanced, breakable,
    width=\dimexpr\linewidth-4cm\relax,
    colback=kfPassageBg, colframe=kfPrimary!70,
    boxrule=0.6pt, arc=3pt,
    left=\kfBoxPad, right=\kfBoxPad, top=\dimexpr\kfBoxPad+8pt\relax, bottom=\kfBoxPad,
    title={\small\bfseries\color{kfPrimary} #1},
    coltitle=kfPrimary,
    colbacktitle=kfPrimaryLight!60,
    attach boxed title to top left={yshift=-2.4mm, xshift=8pt},
    boxed title style={colframe=kfPrimary!50, boxrule=0.5pt, arc=2pt},
    parbox=false,
    overlay unbroken and first={%
      \begin{scope}
        \draw[kfRule, line width=0.4pt, fill=kfNoteBg, rounded corners=2pt]
          ([xshift=\kf@memoGap]frame.north east) rectangle
          ([xshift=\kf@memoGap+\kf@memoW]frame.south east);
        \node[anchor=north east, font=\footnotesize\bfseries\color{kfMute}, inner sep=3pt]
          at ([xshift=\kf@memoGap+\kf@memoW, yshift=-2pt]frame.north east) {메모};
        \node[anchor=north east, fill=kfPrimary!15, text=kfPrimary,
              font=\scriptsize\bfseries, rounded corners=2pt,
              inner xsep=4pt, inner ysep=2pt]
          at ([xshift=\kf@memoGap+\kf@memoW-3pt, yshift=-19pt]frame.north east) {\kf@memocat};
        \begin{scope}
          \path[clip]
            ([xshift=\kf@memoGap]frame.north east) rectangle
            ([xshift=\kf@memoGap+\kf@memoW]frame.south east);
          \foreach \i in {1,...,40}{%
            \draw[kfRule, line width=0.25pt]
              ([xshift=\kf@memoGap+2mm, yshift=-7mm - \i*5mm]frame.north east) --
              ([xshift=\kf@memoGap+\kf@memoW-2mm, yshift=-7mm - \i*5mm]frame.north east);%
          }
        \end{scope}
      \end{scope}
    }
  ]
    \setlength{\parindent}{1em}%
    \hspace*{1em}\ignorespaces#2%
  \end{tcolorbox}\par\smallskip
}

% --- 제목 없는 버전 ---
\newcommand{\kf@passagenote@notitle}[1]{%
  \par\ifkfPassageNewPage\ifdim\pagetotal>0.4\textheight\clearpage\fi\fi\smallskip\needspace{6\baselineskip}%
  \begin{tcolorbox}[
    enhanced, breakable,
    width=\dimexpr\linewidth-4cm\relax,
    colback=kfPassageBg, colframe=kfPrimary!70,
    boxrule=0.6pt, arc=3pt,
    left=\kfBoxPad, right=\kfBoxPad, top=\dimexpr\kfBoxPad+8pt\relax, bottom=\kfBoxPad,
    parbox=false,
    overlay unbroken and first={%
      \begin{scope}
        \draw[kfRule, line width=0.4pt, fill=kfNoteBg, rounded corners=2pt]
          ([xshift=\kf@memoGap]frame.north east) rectangle
          ([xshift=\kf@memoGap+\kf@memoW]frame.south east);
        \node[anchor=north east, font=\footnotesize\bfseries\color{kfMute}, inner sep=3pt]
          at ([xshift=\kf@memoGap+\kf@memoW, yshift=-2pt]frame.north east) {메모};
        \node[anchor=north east, fill=kfPrimary!15, text=kfPrimary,
              font=\scriptsize\bfseries, rounded corners=2pt,
              inner xsep=4pt, inner ysep=2pt]
          at ([xshift=\kf@memoGap+\kf@memoW-3pt, yshift=-19pt]frame.north east) {\kf@memocat};
        \begin{scope}
          \path[clip]
            ([xshift=\kf@memoGap]frame.north east) rectangle
            ([xshift=\kf@memoGap+\kf@memoW]frame.south east);
          \foreach \i in {1,...,40}{%
            \draw[kfRule, line width=0.25pt]
              ([xshift=\kf@memoGap+2mm, yshift=-7mm - \i*5mm]frame.north east) --
              ([xshift=\kf@memoGap+\kf@memoW-2mm, yshift=-7mm - \i*5mm]frame.north east);%
          }
        \end{scope}
      \end{scope}
    }
  ]
    \setlength{\parindent}{1em}%
    \hspace*{1em}\ignorespaces#1%
  \end{tcolorbox}\par\smallskip
}

\newcommand{\kfPassageNote}[2]{%
  \def\kf@tmptitle{#1}%
  \ifx\kf@tmptitle\@empty
    \kf@passagenote@notitle{#2}%
  \else
    \kf@passagenote@with{#1}{#2}%
  \fi
}
\newcommand{\kfPassageNoteNT}[1]{\kf@passagenote@notitle{#1}}
\makeatother

% ============================================================
% 풀폭 지문 박스 (메모란 없음) — 활동 내 보기 박스
% PAGE_BREAK_POLICY.md §3
% ============================================================
\makeatletter
\newcommand{\kf@passagefull@with}[2]{%
  \par\ifkfPassageNewPage\ifdim\pagetotal>0.4\textheight\clearpage\fi\fi\smallskip\needspace{6\baselineskip}%
  \begin{tcolorbox}[
    enhanced, breakable,
    colback=kfPassageBg, colframe=kfPrimary!70,
    boxrule=0.6pt, arc=3pt,
    left=\kfBoxPad, right=\kfBoxPad, top=\dimexpr\kfBoxPad+8pt\relax, bottom=\kfBoxPad,
    title={\small\bfseries\color{kfPrimary} #1},
    coltitle=kfPrimary,
    colbacktitle=kfPrimaryLight!60,
    attach boxed title to top left={yshift=-2.4mm, xshift=8pt},
    boxed title style={colframe=kfPrimary!50, boxrule=0.5pt, arc=2pt},
    parbox=false
  ]
    \setlength{\parindent}{1em}%
    \hspace*{1em}\ignorespaces#2%
  \end{tcolorbox}\par\smallskip
}
\newcommand{\kf@passagefull@notitle}[1]{%
  \par\ifkfPassageNewPage\ifdim\pagetotal>0.4\textheight\clearpage\fi\fi\smallskip\needspace{6\baselineskip}%
  \begin{tcolorbox}[
    enhanced, breakable,
    colback=kfPassageBg, colframe=kfPrimary!70,
    boxrule=0.6pt, arc=3pt,
    left=\kfBoxPad, right=\kfBoxPad, top=\dimexpr\kfBoxPad+8pt\relax, bottom=\kfBoxPad,
    parbox=false
  ]
    \setlength{\parindent}{1em}%
    \hspace*{1em}\ignorespaces#1%
  \end{tcolorbox}\par\smallskip
}
\newcommand{\kfPassageFull}[2]{%
  \def\kf@tmptitle{#1}%
  \ifx\kf@tmptitle\@empty
    \kf@passagefull@notitle{#2}%
  \else
    \kf@passagefull@with{#1}{#2}%
  \fi
}
\newcommand{\kfPassageFullNT}[1]{\kf@passagefull@notitle{#1}}
\makeatother

% ============================================================
% 작은 문장 박스 (문장 독해용) — PAGE_BREAK_POLICY.md §2.5
% ============================================================
\newcommand{\kfSentenceBox}[2]{%
  \par\smallskip\needspace{3\baselineskip}%
  \noindent\begin{tcolorbox}[
    enhanced, breakable=false,
    colback=kfPassageBg, colframe=kfMute,
    boxrule=0.4pt, arc=2pt,
    left=8pt, right=6pt, top=4pt, bottom=4pt,
    title={\footnotesize\bfseries\color{kfPrimary} #1},
    coltitle=kfPrimary, colbacktitle=kfPaper,
    attach boxed title to top left={yshift=-1.6mm, xshift=4pt},
    boxed title style={colframe=kfMute, boxrule=0.3pt, arc=1pt}
  ]%
    \setstretch{1.3}#2%
  \end{tcolorbox}\par\smallskip%
}

% ============================================================
% 지문 본문 아래 안내/정리 박스 (v16 신규)
% 사용: \kfPassageHint{한 줄 정리}{본문 안내 텍스트}
% ============================================================
\newcommand{\kfPassageHint}[2]{%
  \par\smallskip\needspace{4\baselineskip}%
  \begin{tcolorbox}[
    enhanced, breakable=false,
    colback=kfPrimary!5, colframe=kfPrimary!40,
    boxrule=0.4pt, arc=2pt,
    left=8pt, right=6pt, top=4pt, bottom=4pt,
    title={\footnotesize\bfseries\color{kfPrimary} #1},
    coltitle=kfPrimary, colbacktitle=kfPaper,
    attach boxed title to top left={yshift=-1.5mm, xshift=4pt},
    boxed title style={colframe=kfPrimary!40, boxrule=0.3pt, arc=1pt}
  ]%
    \small\setstretch{1.3}#2%
  \end{tcolorbox}\par\smallskip%
}

% ============================================================
% 환경 버전 (호환성) — 메모란 포함
% ============================================================
\makeatletter
\newenvironment{kfPassageNoteEnv}[1]%
  {\par\needspace{6\baselineskip}%
   \def\kf@psrc{#1}%
   \begin{tcolorbox}[
     enhanced, breakable,
     width=\dimexpr\linewidth-4cm\relax,
     colback=kfPassageBg, colframe=kfPrimary!70,
     boxrule=0.6pt, arc=3pt,
     left=\kfBoxPad, right=\kfBoxPad, top=\dimexpr\kfBoxPad+8pt\relax, bottom=\kfBoxPad,
     title={\small\bfseries\color{kfPrimary} \kf@psrc},
     coltitle=kfPrimary, colbacktitle=kfPrimaryLight!60,
     attach boxed title to top left={yshift=-2.4mm, xshift=8pt},
     boxed title style={colframe=kfPrimary!50, boxrule=0.5pt, arc=2pt},
     parbox=false,
     overlay unbroken and first={%
       \begin{scope}
         \draw[kfRule, line width=0.4pt, fill=kfNoteBg, rounded corners=2pt]
           ([xshift=\kf@memoGap]frame.north east) rectangle
           ([xshift=\kf@memoGap+\kf@memoW]frame.south east);
         \node[anchor=north east, font=\footnotesize\bfseries\color{kfMute}, inner sep=3pt]
           at ([xshift=\kf@memoGap+\kf@memoW, yshift=-2pt]frame.north east) {메모};
         \begin{scope}
           \path[clip]
             ([xshift=\kf@memoGap]frame.north east) rectangle
             ([xshift=\kf@memoGap+\kf@memoW]frame.south east);
           \foreach \i in {1,...,40}{%
             \draw[kfRule, line width=0.25pt]
               ([xshift=\kf@memoGap+2mm, yshift=-7mm - \i*5mm]frame.north east) --
               ([xshift=\kf@memoGap+\kf@memoW-2mm, yshift=-7mm - \i*5mm]frame.north east);%
           }
         \end{scope}
       \end{scope}
     }
   ]%
   \setlength{\parindent}{1em}%
  }%
  {\end{tcolorbox}\par\smallskip}
\makeatother

% ============================================================
% condition-box.tex
% <조건> 박스 컴포넌트 (tcolorbox)
% 사용:
%   \begin{kfCondition}
%     \item 첫째 조건
%     \item 둘째 조건
%   \end{kfCondition}
% ============================================================

\newenvironment{kfCondition}%
  {\par\smallskip\needspace{4\baselineskip}%
   \begin{tcolorbox}[
     enhanced, breakable=false,
     colback=kfPrimaryLight!40, colframe=kfPrimary,
     boxrule=0.5pt, arc=2pt,
     left=\kfBoxPad, right=\kfBoxPad, top=\kfBoxPad, bottom=\kfBoxPad,
     title={\small\bfseries\color{kfPrimary} \,조건\,},
     coltitle=kfPrimary, colbacktitle=kfPaper,
     attach boxed title to top left={yshift=-2mm, xshift=6pt},
     boxed title style={colframe=kfPrimary, boxrule=0.4pt, arc=1pt}
   ]%
   \begin{enumerate}[leftmargin=16pt, itemsep=1pt, topsep=1pt,
                    label=\textbf{\arabic*.}]}%
  {\end{enumerate}\end{tcolorbox}\par\smallskip}

% ============================================================
% evidence-box.tex
% <보기> 박스 컴포넌트
% 사용:
%   \begin{kfEvidence}
%     보기 본문 ...
%   \end{kfEvidence}
% ============================================================

\newenvironment{kfEvidence}%
  {\par\smallskip\needspace{4\baselineskip}%
   \begin{tcolorbox}[
     enhanced, breakable=false,
     colback=kfPaper, colframe=kfMute,
     boxrule=0.5pt, arc=2pt,
     left=\kfBoxPad, right=\kfBoxPad, top=\kfBoxPad, bottom=\kfBoxPad,
     title={\small\bfseries\color{kfInk} \,보기\,},
     coltitle=kfInk, colbacktitle=kfPrimaryLight!50,
     attach boxed title to top left={yshift=-2mm, xshift=6pt},
     boxed title style={colframe=kfMute, boxrule=0.4pt, arc=1pt}
   ]%
   \leavevmode\mbox{}\ignorespaces%
  }%
  {\end{tcolorbox}\par\smallskip}

% ============================================================
% choice-list.tex
% 5선지 (① ② ③ ④ ⑤) 자동 들여쓰기 컴포넌트
% 사용:
%   \begin{kfChoices}
%     \kfChoice 첫째 선택지
%     \kfChoice 둘째 선택지
%     \kfChoice 셋째 선택지
%     \kfChoice 넷째 선택지
%     \kfChoice 다섯째 선택지
%   \end{kfChoices}
% ============================================================

\newcounter{kfChoiceCnt}
\newcommand{\kfChoiceMark}[1]{%
  \ifcase#1\relax%
  \or ①\or ②\or ③\or ④\or ⑤\or ⑥\or ⑦\or ⑧\or ⑨%
  \fi
}

% 환경: 자동 번호
\newenvironment{kfChoices}%
  {\setcounter{kfChoiceCnt}{0}%
   \par\smallskip%
   \begin{list}{}{%
     \setlength{\leftmargin}{20pt}%
     \setlength{\itemindent}{0pt}%
     \setlength{\labelwidth}{18pt}%
     \setlength{\labelsep}{6pt}%
     \setlength{\itemsep}{2pt}%
     \setlength{\topsep}{1pt}%
     \setlength{\parsep}{0pt}%
     \renewcommand{\makelabel}[1]{##1\hss}%
   }}%
  {\end{list}\par\smallskip}

\newcommand{\kfChoice}{%
  \stepcounter{kfChoiceCnt}%
  \item[\textbf{\kfChoiceMark{\value{kfChoiceCnt}}}]\ignorespaces%
}

% --- 한 줄 5선지 (짧은 선택지용) ---
\newcommand{\kfChoicesInline}[5]{%
  \par\smallskip%
  \noindent\textbf{①}~#1\quad\textbf{②}~#2\quad\textbf{③}~#3\quad\textbf{④}~#4\quad\textbf{⑤}~#5%
  \par\smallskip%
}

% ============================================================
% answer-note.tex
% 답안 작성란 (노트 스타일, 줄친 영역)
% 사용:
%   \kfAnswerLines{5}        % 5줄 답안란
%   \kfAnswerNote{서술형 답안란}{8} % 라벨+8줄
% ============================================================

% 단일 줄 (필요 시 인라인)
\newcommand{\kfAnswerLine}{%
  \par\noindent\textcolor{kfRule}{\rule{\linewidth}{0.4pt}}\par
}

% N줄 답안란 (간단)
\newcommand{\kfAnswerLines}[1]{%
  \par\smallskip%
  \begin{minipage}{\linewidth}%
  \foreach \i in {1,...,#1}{%
    \textcolor{kfRule}{\rule{\linewidth}{0.4pt}}\\[6pt]%
  }%
  \end{minipage}%
  \par\smallskip%
}

% 라벨이 있는 노트형 답안란
\newcommand{\kfAnswerNote}[2]{%
  \par\smallskip\needspace{#2\baselineskip}%
  \begin{tcolorbox}[
    enhanced, breakable=false,
    colback=kfPaper, colframe=kfRule,
    boxrule=0.4pt, arc=2pt,
    left=\kfBoxPad, right=\kfBoxPad, top=\kfBoxPad, bottom=\kfBoxPad,
    title={\footnotesize\bfseries\color{kfMute} #1},
    coltitle=kfMute, colbacktitle=kfPaper,
    attach boxed title to top left={yshift=-1.5mm, xshift=6pt},
    boxed title style={colframe=kfRule, boxrule=0.3pt, arc=1pt}
  ]%
  \foreach \i in {1,...,#2}{%
    \textcolor{kfRule}{\rule{\linewidth}{0.3pt}}\\[6pt]%
  }%
  \end{tcolorbox}\par\smallskip%
}

% 짧은 빈칸 (단답형)
\newcommand{\kfBlank}[1][3cm]{\underline{\hspace{#1}}}
% 넓은 빈칸 (서술 한 줄)
\newcommand{\kfBlankWide}[1][6cm]{\underline{\hspace{#1}}}
% 괄호 빈칸
\newcommand{\kfBlankParen}{(\,\hspace{1cm}\,)}
% O/X 빈칸
\newcommand{\kfBlankOX}{(\,\hspace{1.5cm}\,)}

% ============================================================
% question-block.tex
% 문제 블록 (번호 배지 + 발문 + 보기/조건 + 선택지 + 답안란)
% 모든 구성을 하나의 minipage로 묶어 단/페이지 단절 금지
%
% 사용:
%   \begin{kfQBlock}{1}{객관식}
%     \kfQStem{다음 글에 대한 설명으로 가장 적절한 것은?}
%     \begin{kfEvidence} ... \end{kfEvidence}    % 선택
%     \begin{kfCondition} ... \end{kfCondition}  % 선택
%     \begin{kfChoices}
%       \kfChoice ...
%     \end{kfChoices}
%     \kfAnswerLines{2}                            % 선택
%   \end{kfQBlock}
% ============================================================

% 무단절 minipage로 묶고, 발문/보기/조건/선택지/답안란을 자유롭게 배치
% 사용자 피드백 #4: [객관식]/[서술형] 등 텍스트 라벨 제거 — 번호 배지만 표시
% #2(유형 라벨)는 호환성 인자로만 받고 출력하지 않음.
\newenvironment{kfQBlock}[2]%
  {\par\smallskip\needspace{6\baselineskip}%
   \noindent\begin{minipage}{\linewidth}%
   \samepage%
   \par\noindent
   \kfQNum{#1}\hspace{4pt}%
   \begin{minipage}[t]{\dimexpr\linewidth-30pt\relax}%
  }%
  {\end{minipage}\end{minipage}\par\smallskip}

% 문제 발문
\newcommand{\kfQStem}[1]{%
  \par\noindent#1\par\vspace{2pt}%
}

% 짧은 소문항 라벨 (1), (2), (3)
\newcounter{kfSubQ}
\newcommand{\kfSubQReset}{\setcounter{kfSubQ}{0}}
\newcommand{\kfSubQ}{%
  \stepcounter{kfSubQ}%
  \par\noindent\textbf{(\arabic{kfSubQ})}~\ignorespaces%
}

% ============================================================
% activity-table.tex
% 활동: 정리표 / 내용 확인표 (마크다운 표 → tabularx 변환)
% 사용:
%   % 일반 tabular (직접 컬럼 명시)
%   \begin{kfActTable}{|l|X|X|}
%     \kfTHead{항목} & \kfTHead{설명} & \kfTHead{학생 작성} \\\hline
%     ① & ... & \kfTBlank \\\hline
%   \end{kfActTable}
%
%   % adjustbox 자동 축소 (정해진 폭으로 강제)
%   \kfActTableWrap{
%     ... (\begin{tabular}...\end{tabular} 코드) ...
%   }
% ============================================================

% 사용 시 본문 마지막 행은 \\\hline 으로 끝나야 함.
% v17.1: 흰색 mask로 Canva 배경 본문 침범 차단
\newenvironment{kfActTable}[1]%
  {\par\smallskip\needspace{4\baselineskip}%
   \renewcommand{\arraystretch}{1.3}%
   \arrayrulecolor{kfRule}%
   \begin{tcolorbox}[blanker, colback=white, left=2pt, right=2pt, top=2pt, bottom=2pt]%
   \noindent\begin{adjustbox}{max width=\linewidth}%
   \begin{tabular}{#1}%
   \hline}%
  {\end{tabular}%
   \end{adjustbox}%
   \end{tcolorbox}%
   \arrayrulecolor{black}%
   \par\smallskip}

% 헤더 행 헬퍼
\newcommand{\kfTHead}[1]{\textbf{\color{kfPrimary} #1}}
\newcommand{\kfTHeadRow}[1]{\rowcolor{kfPrimaryLight!50}\textbf{\color{kfPrimary} #1}}

% 빈 셀 (학생 작성용)
\newcommand{\kfTBlank}{\rule[-1.6em]{0pt}{3.4em}}
\newcommand{\kfTBlankSm}{\rule[-1.0em]{0pt}{2.4em}}

% adjustbox 강제 축소 래퍼 — Python에서 변환된 raw tabular 블록을 감쌀 때 사용
\newcommand{\kfActTableWrap}[1]{%
  \par\smallskip\needspace{4\baselineskip}%
  \begin{tcolorbox}[blanker, colback=white, left=2pt, right=2pt, top=2pt, bottom=2pt]%
  \noindent\begin{adjustbox}{max width=\linewidth, center}%
  #1%
  \end{adjustbox}%
  \end{tcolorbox}\par\smallskip%
}

% ============================================================
% activity-vocab.tex
% 어휘 학습 표 (3열: 번호 - 뜻 - 빈칸)
% 사용:
%   \begin{kfVocab}
%     \kfVocabItem{1}{눈으로 보고 마음으로 깨달음}
%     \kfVocabItem{2}{사물의 이치를 따져 헤아림}
%   \end{kfVocab}
% ============================================================

% adjustbox로 폭 강제 + 일반 tabular + p{} 열 사용
\newenvironment{kfVocab}%
  {\par\smallskip\needspace{4\baselineskip}%
   \renewcommand{\arraystretch}{1.35}%
   \arrayrulecolor{kfRule}%
   \noindent\begin{adjustbox}{max width=\linewidth, center}%
   \begin{tabular}{|>{\centering\arraybackslash}m{1.0cm}|m{8.5cm}|>{\centering\arraybackslash}m{4.0cm}|}%
   \hline
   \rowcolor{kfPrimaryLight!50}%
   \multicolumn{1}{|c|}{\textbf{\color{kfPrimary}번호}} &
   \multicolumn{1}{c|}{\textbf{\color{kfPrimary}뜻풀이}} &
   \multicolumn{1}{c|}{\textbf{\color{kfPrimary}어휘 (정답 작성)}} \\\hline
   \ignorespaces}%
  {\end{tabular}%
   \end{adjustbox}%
   \arrayrulecolor{black}%
   \par\smallskip}

\newcommand{\kfVocabItem}[2]{%
  \textbf{#1} & #2 & \rule[-1.0em]{0pt}{2.4em}\\\hline%
}

% ============================================================
% activity-blank.tex
% 빈칸 채우기 활동
% 사용:
%   \begin{kfBlankList}
%     \kfBlankItem 첫 번째 문장에서 (\kfBlank)은(는) ...
%     \kfBlankItem 둘째 문장에서 ...
%   \end{kfBlankList}
% ============================================================

\newenvironment{kfBlankList}%
  {\par\smallskip%
   \begin{enumerate}[leftmargin=20pt, itemsep=4pt, topsep=2pt,
                    label=\textbf{\arabic*.}, labelsep=6pt]}%
  {\end{enumerate}\par\smallskip}

\newcommand{\kfBlankItem}{\item}

% 텍스트 안에 박스형 빈칸 (단어 1개 채우기)
\newcommand{\kfBlankBox}[1][2.4cm]{%
  \tikz[baseline=-0.5ex]{%
    \draw[kfRule, line width=0.4pt] (0,0) -- ++(#1,0);%
  }%
}

% ============================================================
% activity-ox.tex (v15)
% O/X 표기 활동 — 그래픽 디자인
% 사용:
%   \begin{kfOXList}
%     \kfOXItem 글쓴이는 자연을 예찬하고 있다.\kfOX
%     \kfOXItem 1연과 4연은 수미상관 구조이다.\kfOX
%   \end{kfOXList}
%
% \kfOX = 우측에 [O] [X] 두 박스 (학생이 하나 동그라미)
% ============================================================

\newenvironment{kfOXList}%
  {\par\smallskip%
   \begin{enumerate}[leftmargin=20pt, itemsep=5pt, topsep=2pt,
                    label=\textbf{\arabic*.}, labelsep=6pt]}%
  {\end{enumerate}\par\smallskip}

\newcommand{\kfOXItem}{\item}

% v16: O/X 그래픽 박스 키움 (18×14 → 24×20pt) + 영역색 강조
\newcommand{\kfOX}{%
  \hfill\nolinebreak
  \tikz[baseline=-0.9ex]{%
    \node[draw=kfPrimary, fill=kfPrimary!5, line width=0.6pt, rounded corners=3pt,
          minimum width=24pt, minimum height=20pt, inner sep=0pt,
          font=\normalsize\bfseries\color{kfPrimary}] (ob) {O};%
    \node[draw=kfMute, fill=kfMute!5, line width=0.6pt, rounded corners=3pt,
          minimum width=24pt, minimum height=20pt, inner sep=0pt,
          font=\normalsize\bfseries\color{kfMute},
          right=4pt of ob] (xb) {X};%
  }%
}

% ============================================================
% activity-connect.tex
% 좌우 선 긋기 활동 (2단 양식 + 중앙 여백)
% 사용:
%   \begin{kfConnect}
%     \kfPair{사르트르}{실존은 본질에 앞선다}
%     \kfPair{니체}{초인 사상}
%     \kfPair{하이데거}{현존재}
%   \end{kfConnect}
% ============================================================

\newenvironment{kfConnect}%
  {\par\smallskip%
   \renewcommand{\arraystretch}{1.6}%
   \begin{tabular}{@{}>{\raggedleft\arraybackslash}m{4.5cm}
                     @{\hspace{2.5cm}}
                     >{\raggedright\arraybackslash}m{6.5cm}@{}}}%
  {\end{tabular}\par\smallskip}

\newcommand{\kfPair}[2]{%
  \tikz[baseline]{\node[draw=kfRule, fill=kfPrimaryLight!30, rounded corners=2pt,
        inner sep=4pt, font=\small]{#1};} \tikz[baseline]{\node[circle, draw=kfMute, fill=kfPaper, inner sep=1.5pt]{};}
  &
  \tikz[baseline]{\node[circle, draw=kfMute, fill=kfPaper, inner sep=1.5pt]{};} \tikz[baseline]{\node[draw=kfRule, fill=kfNoteBg, rounded corners=2pt,
        inner sep=4pt, font=\small]{#2};} \\
}

% ============================================================
% activity-writing.tex
% 글쓰기 활동 (큰 노트 영역)
% 사용:
%   \kfWriting{독후감}{12}     % 라벨 + 12줄 큰 노트
%   \kfWritingPrompt{프롬프트 텍스트}
% ============================================================

\newcommand{\kfWritingPrompt}[1]{%
  \par\smallskip%
  \begin{tcolorbox}[
    enhanced, breakable=false,
    colback=kfPrimaryLight!30, colframe=kfPrimary,
    boxrule=0.4pt, arc=3pt,
    left=\kfBoxPad, right=\kfBoxPad, top=\kfBoxPad, bottom=\kfBoxPad,
  ]%
  \small\textbf{\color{kfPrimary}글쓰기 안내}\\[2pt]%
  #1
  \end{tcolorbox}\par\smallskip%
}

\newcommand{\kfWriting}[2]{%
  \par\smallskip\needspace{\numexpr#2+2\relax\baselineskip}%
  \begin{tcolorbox}[
    enhanced, breakable=false,
    colback=kfPaper, colframe=kfRule,
    boxrule=0.4pt, arc=2pt,
    left=\kfBoxPad, right=\kfBoxPad, top=\kfBoxPad, bottom=\kfBoxPad,
    title={\footnotesize\bfseries\color{kfMute} #1},
    coltitle=kfMute, colbacktitle=kfPaper,
    attach boxed title to top left={yshift=-1.5mm, xshift=6pt},
    boxed title style={colframe=kfRule, boxrule=0.3pt, arc=1pt}
  ]%
  \foreach \i in {1,...,#2}{%
    \textcolor{kfRule}{\rule{\linewidth}{0.3pt}}\\[7pt]%
  }%
  \end{tcolorbox}\par\smallskip%
}

% ============================================================
% activity-structure.tex
% 구조도 (트리 + 빈칸)
% 사용:
%   \begin{kfStructure}
%     \kfNode{0}{서론}{문제 제기}
%     \kfNode{1}{본론}{(\,\quad\,)}
%     \kfNode{1}{본론}{근거 2}
%     \kfNode{0}{결론}{주장 강화}
%   \end{kfStructure}
% ============================================================

\newenvironment{kfStructure}%
  {\par\smallskip\needspace{6\baselineskip}%
   \begin{tcolorbox}[
     enhanced, breakable=false,
     colback=kfPaper, colframe=kfRule,
     boxrule=0.4pt, arc=2pt,
     left=\kfBoxPad, right=\kfBoxPad, top=\kfBoxPad, bottom=\kfBoxPad,
     title={\footnotesize\bfseries\color{kfMute} 글의 구조},
     coltitle=kfMute, colbacktitle=kfPrimaryLight!40,
     attach boxed title to top left={yshift=-1.5mm, xshift=6pt},
     boxed title style={colframe=kfRule, boxrule=0.3pt, arc=1pt}
   ]%
   \begin{tabbing}
     \hspace{1.4cm}\=\hspace{1.4cm}\=\hspace{8cm}\kill}%
  {\end{tabbing}\end{tcolorbox}\par\smallskip}

% 노드: #1=들여쓰기 단계 (0~3), #2=라벨, #3=내용
\newcommand{\kfNode}[3]{%
  \ifcase#1\relax%
    \>\>\textbf{\color{kfPrimary} ▶ #2}\quad #3\\
  \or
    \>\>\quad\textbf{\color{kfMute} ◦ #2}\quad #3\\
  \or
    \>\>\quad\quad\textbf{\color{kfMute} - #2}\quad #3\\
  \or
    \>\>\quad\quad\quad\textbf{\color{kfMute} \textperiodcentered\ #2}\quad #3\\
  \fi
}

% 빈칸 노드
\newcommand{\kfNodeBlank}[2]{\kfNode{#1}{#2}{(\hspace{2.5cm})}}

% ============================================================
% activity-sentence-reading.tex
% 문장 독해 활동 — 문장 박스 1개 + 그에 묶인 문제 N개를 한 minipage로
% 사용:
%   \begin{kfSentenceUnit}{1}{"바람은 내 귀에 속삭이며..."}
%     \kfSentQ{(1)}{서술어 '속삭이며'의 주어는?}
%     \begin{kfChoices}
%       \kfChoice 바람은
%       ...
%     \end{kfChoices}
%   \end{kfSentenceUnit}
% ============================================================

% kfSentenceBox 매크로는 passage-note.tex에 정의됨

\newenvironment{kfSentenceUnit}[2]%
  {\par\smallskip\needspace{6\baselineskip}%
   \noindent\begin{minipage}{\linewidth}%
   \samepage%
   \kfSentenceBox{문장 #1}{#2}%
   \par\smallskip%
  }%
  {\end{minipage}\par\smallskip}

% 같은 문장에 묶인 소문항 (문장 안내 + 발문)
\newcommand{\kfSentQ}[2]{%
  \par\noindent\textbf{#1}~#2\par\smallskip%
}

% 헤더 없이 문장만 있는 묶음 (sentence_ref가 없는 일반 문항)
\newenvironment{kfSentencelessUnit}%
  {\par\smallskip\noindent\begin{minipage}{\linewidth}\samepage}%
  {\end{minipage}\par\smallskip}

% ============================================================
% chapter-cover.tex (v6)
% 챕터 진입 페이지 (표지) — 하단에 영역 목차 추가
% 사용:
%   \kfChapterCover{1}{근대성과 자아}{Chapter 1}{이번 챕터에서는 ...}
%   \begin{kfChapterAreaList}
%     \kfChapterAreaItem{어휘}{kfAreaVocab}{핵심 어휘 12개}
%     ...
%   \end{kfChapterAreaList}
%
%   영역 목차가 없을 때는 \kfChapterCoverEnd 호출
% ============================================================

\newcommand{\kfChapterCover}[4]{%
  \clearpage%
  \thispagestyle{kfBlank}%
  \kfMarkChapter{Chapter #1. #2}%
  \kfChapterCoverBg%
  % v18.5: 챕터 표지 우상단에 로고
  \begin{tikzpicture}[remember picture, overlay]
    \node[anchor=north east, inner sep=0pt] at ($(current page.north east)+(-18mm,-18mm)$) {\kfLogoMd};
  \end{tikzpicture}%
  \vspace*{20mm}%
  \noindent
  \begin{tikzpicture}[remember picture, overlay]
    % 좌측 액센트 바
    \fill[kfPrimary] (current page.west) ++(12mm,25mm) rectangle ++(5mm, -50mm);
    % 챕터 번호 배지 (이중 원, 약간 작게)
    \node[anchor=north west, inner sep=0pt] at ($(current page.north west)+(24mm,-45mm)$) {%
      \begin{tikzpicture}
        \draw[kfPrimary, line width=1pt] (0,0) circle (10mm);
        \draw[kfPrimary, line width=0.4pt] (0,0) circle (8mm);
        \node[font=\normalsize\bfseries\color{kfPrimary}] at (0,2mm) {CHAPTER};
        \node[font=\LARGE\bfseries\color{kfPrimary}] at (0,-3mm) {#1};
      \end{tikzpicture}%
    };
  \end{tikzpicture}%
  \vspace{42mm}%
  \par\noindent\hspace{32mm}%
  \begin{minipage}{0.65\linewidth}%
    {\footnotesize\color{kfMute} #3}\\[4pt]
    {\LARGE\bfseries\color{kfPrimary} #2}\\[8pt]
    {\color{kfRule}\rule{50mm}{0.5pt}}\\[6pt]
    {\small\color{kfInk} #4}
  \end{minipage}%
  \par\vspace{14mm}%
}

% v6 / v18.5 / v18.6: 챕터 표지의 영역 목차 — 흰색 반투명 박스
% v18.6: 배경 가로선/도형과 안 겹치도록 TikZ overlay로 우하단 절대 위치
\makeatletter
% 영역 항목들을 모아 둘 박스
\newsavebox{\kf@arealistbox}
\newenvironment{kfChapterAreaList}{%
  \begin{lrbox}{\kf@arealistbox}%
  \begin{minipage}{0.62\linewidth}%
    \begin{tcolorbox}[%
      enhanced, colback=white, colframe=white,
      opacityback=0.92, opacityframe=0,
      boxrule=0pt, arc=4pt,
      width=\linewidth,
      top=8pt, bottom=8pt, left=10pt, right=10pt,
    ]%
      {\footnotesize\bfseries\color{kfMute} 이 챕터의 영역}\par\vspace{4pt}%
      {\color{kfRule}\hrule height 0.3pt}\par\vspace{4pt}%
      \begin{itemize}[leftmargin=14pt, itemsep=2pt, topsep=0pt, label={\color{kfPrimary!70}\textbullet}]%
}{%
      \end{itemize}%
    \end{tcolorbox}%
  \end{minipage}%
  \end{lrbox}%
  % v18.7: 배경 상단 가로선 ~ 하단 가로선 사이의 중간 영역에 배치
  % 가로선이 내용을 가리지 않도록 박스 중심이 두 선 사이 중점에 위치
  \begin{tikzpicture}[remember picture, overlay]%
    \node[anchor=center, inner sep=0pt]
      at ($(current page.south)+(0mm,90mm)$)
      {\usebox{\kf@arealistbox}};%
  \end{tikzpicture}%
  \vfill%
  \noindent\hfill\kfSeriesBadge\hspace{12mm}%
  \clearpage%
}
\makeatother

% 영역 항목: \kfChapterAreaItem{영역명}{kfArea색}{부제}
\makeatletter
\newcommand{\kfChapterAreaItem}[3]{%
  \item {\small\bfseries\color{#2} #1}%
  \def\kf@cci@sub{#3}%
  \ifx\kf@cci@sub\@empty\else
    {\small\color{kfMute}\, \textemdash\, #3}%
  \fi
}
\makeatother

% 영역 목록 없이 cover만 (legacy)
\newcommand{\kfChapterCoverEnd}{%
  \vfill%
  \noindent\hfill\kfSeriesBadge\hspace{12mm}%
  \clearpage%
}

% ============================================================
% area-header.tex (v16)
% 영역 시작 표제 — 시리즈별 차별화 라벨 + 큰 영역명 + 부제
% PAGE_BREAK_POLICY.md §1.4
%
% v16 (디자인 고도화 Phase 1):
%   - 시리즈별 라벨 박스 추가:
%     소쉬르 = AREA START (영역색 채움 박스)
%     프레게 = SECTION OPEN (영역색 테두리)
%     러셀   = RUSSELL (영역색 라이트 박스)
%     비트   = WITTGENSTEIN (회색 라이트 박스)
%   - 라벨 위치: 영역명 위 좌측
% ============================================================

\makeatletter

% 시리즈 라벨 텍스트
\newcommand{\kf@serieslabel}{%
  \ifcase\kf@series\relax
    AREA START%
  \or
    AREA START%
  \or
    SECTION OPEN%
  \or
    RUSSELL%
  \or
    WITTGENSTEIN%
  \fi
}

% 시리즈 라벨 박스 스타일 (#1 = 영역색)
\newcommand{\kf@labelbox}[1]{%
  \ifcase\kf@series\relax
    % 0/1 = 소쉬르: 영역색 채움
    \tikz[baseline=-2pt]{\node[fill=#1, text=white,
      rounded corners=3pt, inner xsep=6pt, inner ysep=2pt,
      font=\scriptsize\bfseries]{\kf@serieslabel};}%
  \or
    % 1 = 소쉬르: 영역색 채움
    \tikz[baseline=-2pt]{\node[fill=#1, text=white,
      rounded corners=3pt, inner xsep=6pt, inner ysep=2pt,
      font=\scriptsize\bfseries]{\kf@serieslabel};}%
  \or
    % 2 = 프레게: 영역색 테두리
    \tikz[baseline=-2pt]{\node[draw=#1, line width=0.6pt, text=#1,
      rounded corners=3pt, inner xsep=6pt, inner ysep=2pt,
      font=\scriptsize\bfseries]{\kf@serieslabel};}%
  \or
    % 3 = 러셀: 영역색 라이트 박스
    \tikz[baseline=-2pt]{\node[fill=#1!12, text=#1,
      rounded corners=2pt, inner xsep=5pt, inner ysep=1.5pt,
      font=\scriptsize\bfseries]{\kf@serieslabel};}%
  \or
    % 4 = 비트: 회색 미니 박스
    \tikz[baseline=-2pt]{\node[fill=kfMute!10, text=kfMute,
      rounded corners=2pt, inner xsep=5pt, inner ysep=1.5pt,
      font=\scriptsize\bfseries]{\kf@serieslabel};}%
  \fi
}

% v17: 시리즈+영역별 Canva 배경 자동 매핑
% \kf@hasbg=1로 set되면 \kfAreaStart에서 라벨 박스 출력 생략
\def\kf@areaVocab{어휘}\def\kf@areaGrammar{문법}\def\kf@areaConcept{개념}
\def\kf@areaLit{문학}\def\kf@areaNonlit{비문학}
\newif\ifkf@bgon
% 파일 있으면 배경 적용 + boolean true. 없으면 nothing.
\newcommand{\kfBgSet}[1]{%
  \IfFileExists{#1.pdf}{\kfPageBg{#1}\kf@bgontrue}{}%
}
\newcommand{\kf@trybg}[1]{%
  % v18.4: 영역 배경 PDF 비활성화 — 큰 도형이 본문 영역을 침범해
  % "영역 헤더 후 빈 페이지" 문제를 일으킴. 라벨 박스로만 영역 표시.
  \kf@bgonfalse%
}

% Note: 러셀(rus_*)/비트(wit_*) 배경 PDF 미존재 시 \kfPageBg가 에러 → \IfFileExists로 가드
% (현재는 파일이 모두 존재한다고 가정. 부분 제공 시 cls 수정 필요)

\newcommand{\kfAreaStart}[3]{%
  \clearpage%
  \thispagestyle{fancy}%
  \kfMarkArea{#1}%
  \kf@trybg{#1}%
  \vspace*{1mm}%
  % 배경 없으면 라텍스 라벨 박스 출력 (배경 있으면 일러스트가 라벨 역할)
  \ifkf@bgon
    \vspace*{4mm}%
  \else
    \noindent\kf@labelbox{#2}\par\vspace{4pt}%
  \fi
  % 영역명 + 부제
  \noindent
  \begin{tikzpicture}
    \node[anchor=west, inner sep=0pt] (bar) at (0,0) {\color{#2}\rule{3pt}{22pt}};
    \node[anchor=base west, inner sep=0pt, xshift=8pt, yshift=2pt] (lbl) at (bar.east |- 0,0) {%
      {\LARGE\bfseries\color{#2} #1}%
    };
    \def\kf@areasub{#3}%
    \ifx\kf@areasub\@empty\else
      \node[anchor=base west, inner sep=0pt, xshift=8pt, font=\normalsize\color{kfMute}]
        at (lbl.base east) {\,\textperiodcentered\, #3};%
    \fi
  \end{tikzpicture}%
  \par\vspace{2pt}
  {\color{#2!70}\hrule height 1pt}%
  \par\vspace{1pt}%
  {\color{kfRule}\hrule height 0.3pt}%
  \par\vspace{4pt}%
  \needspace{6\baselineskip}\nopagebreak[4]%
  \global\kf@areaJustOpenedtrue% v18.5: 첫 컨텐츠 needspace 무시
}
\makeatother

% 시리즈/영역 단축 매크로
\newcommand{\kfStartVocab}[1]{\kfAreaStart{어휘}{kfAreaVocab}{#1}}
\newcommand{\kfStartGrammar}[1]{\kfAreaStart{문법}{kfAreaGrammar}{#1}}
\newcommand{\kfStartConcept}[1]{\kfAreaStart{개념}{kfAreaConcept}{#1}}
\newcommand{\kfStartLit}[1]{\kfAreaStart{문학}{kfAreaLiterature}{#1}}
\newcommand{\kfStartNonlit}[1]{\kfAreaStart{비문학}{kfAreaNonlit}{#1}}
\newcommand{\kfStartWeekly}[1]{\kfAreaStart{실력 확인}{kfAreaWeekly}{#1}}

% ============================================================
% section-header.tex
% 섹션 시작 라벨 헤더 — [지문] / [활동] / [문제] 라벨
% 좌측 액센트 바 + 라벨 + 부제 (영역 액센트 색)
%
% 사용:
%   \kfSecHeader{kfAreaLiterature}{지문}{빼앗긴 들에도 봄은 오는가 / 이상화}
%   \kfSecHeader{kfAreaConcept}{활동}{내용 확인}
%   \kfSecHeader{kfAreaNonlit}{문제}{}
%
% 헤더 직후 페이지 분할 금지: \needspace{4\baselineskip} + \nopagebreak
% ============================================================

\providecommand{\kfSecHeader}[3]{%
  \par\needspace{10\baselineskip}\filbreak%
  \vspace{3pt}%
  \noindent
  \begin{tikzpicture}
    \node[anchor=west, inner sep=0pt] (bar) at (0,0) {\color{#1}\rule{3pt}{14pt}};
    \node[anchor=west, inner sep=0pt, xshift=6pt, font=\normalsize\bfseries\color{#1}]
       (lbl) at (bar.east) {[\,#2\,]};
    \node[anchor=west, inner sep=0pt, xshift=4pt, font=\small\color{kfMute}]
       at (lbl.east) {#3};
  \end{tikzpicture}%
  \par\vspace{1pt}%
  {\color{#1!50}\hrule height 0.4pt}%
  \par\vspace{2pt}\nopagebreak%
}

% 영역별 단축 매크로
\newcommand{\kfSecPassageVocab}[1]{\kfSecHeader{kfAreaVocab}{지문}{#1}}
\newcommand{\kfSecPassageGrammar}[1]{\kfSecHeader{kfAreaGrammar}{지문}{#1}}
\newcommand{\kfSecPassageConcept}[1]{\kfSecHeader{kfAreaConcept}{지문}{#1}}
\newcommand{\kfSecPassageLit}[1]{\kfSecHeader{kfAreaLiterature}{지문}{#1}}
\newcommand{\kfSecPassageNonlit}[1]{\kfSecHeader{kfAreaNonlit}{지문}{#1}}
\newcommand{\kfSecPassageWeekly}[1]{\kfSecHeader{kfAreaWeekly}{지문}{#1}}

\newcommand{\kfSecActVocab}[1]{\kfSecHeader{kfAreaVocab}{활동}{#1}}
\newcommand{\kfSecActGrammar}[1]{\kfSecHeader{kfAreaGrammar}{활동}{#1}}
\newcommand{\kfSecActConcept}[1]{\kfSecHeader{kfAreaConcept}{활동}{#1}}
\newcommand{\kfSecActLit}[1]{\kfSecHeader{kfAreaLiterature}{활동}{#1}}
\newcommand{\kfSecActNonlit}[1]{\kfSecHeader{kfAreaNonlit}{활동}{#1}}
\newcommand{\kfSecActWeekly}[1]{\kfSecHeader{kfAreaWeekly}{활동}{#1}}

\newcommand{\kfSecQVocab}[1]{\kfSecHeader{kfAreaVocab}{문제}{#1}}
\newcommand{\kfSecQGrammar}[1]{\kfSecHeader{kfAreaGrammar}{문제}{#1}}
\newcommand{\kfSecQConcept}[1]{\kfSecHeader{kfAreaConcept}{문제}{#1}}
\newcommand{\kfSecQLit}[1]{\kfSecHeader{kfAreaLiterature}{문제}{#1}}
\newcommand{\kfSecQNonlit}[1]{\kfSecHeader{kfAreaNonlit}{문제}{#1}}
\newcommand{\kfSecQWeekly}[1]{\kfSecHeader{kfAreaWeekly}{문제}{#1}}


\begin{document}

\thispagestyle{kfBlank}
\kfBookCoverBg
\vspace*{\kfBookCoverVspace}
\begin{center}
{\kfLogoLg}\\[14pt]
{\fontsize{72}{84}\selectfont\bfseries\kfBookCoverColor 비트겐슈타인}\\[18pt]
{\fontsize{28}{32}\selectfont\kfBookCoverSubColor \textsc{Level} \textbf{1}}\\[22pt]
{\large\kfBookCoverSubColor 제 1 권 \quad\textbar\quad 챕터 1\textendash{} 4}
\end{center}
\vfill
\hfill\kfSeriesBadge\hspace{15mm}
\clearpage
\kfChapterCover{1}{비트겐슈타인1 ch01}{Chapter 1}{이번 챕터: 문법 → 문학 5작품 → 비문학 5지문 → 패턴 워크북.}
\begin{kfChapterAreaList}
\kfChapterAreaItem{문법}{kfAreaGrammar}{음운 / 음운의 체계}
\kfChapterAreaItem{문학}{kfAreaLiterature}{「자화상」 — 윤동주 외 4편}
\kfChapterAreaItem{비문학}{kfAreaNonlit}{「소크라테스와 자기 인식」 외 4편}
\kfChapterAreaItem{실력 확인}{kfAreaWeekly}{패턴 워크북 — 총 18문항}
\end{kfChapterAreaList}

\kfStartGrammar{음운 / 음운의 체계}


\kfSecHeader{kfAreaGrammar}{문제}{}

\kfBeginMC
\par\needspace{25\baselineskip}
\kfPassageFull{단모음 분류}{%
국어의 단모음은 발음할 때 혀나 입술이 고정되어 움직이지 않는 모음이다. 단모음은 혀의 높낮이, 혀의 전후 위치, 입술 모양이라는 세 가지 기준으로 분류할 수 있다. 혀의 높이에 따라 고모음(ㅣ, ㅟ, ㅡ, ㅜ), 중모음(ㅔ, ㅚ, ㅓ, ㅗ), 저모음(ㅐ, ㅏ)으로 나뉘며, 혀의 전후 위치에 따라 전설 모음(ㅣ, ㅔ, ㅐ, ㅟ, ㅚ)과 후설 모음(ㅡ, ㅓ, ㅏ, ㅜ, ㅗ)으로 구분된다. 또한 입술 모양에 따라 원순 모음(ㅟ, ㅚ, ㅜ, ㅗ)과 평순 모음(ㅣ, ㅔ, ㅐ, ㅡ, ㅓ, ㅏ)으로 나뉜다. 표준 발음법에서는 10개의 단모음 체계를 인정하고 있으나, 'ㅟ'와 'ㅚ'는 이중 모음으로 발음하는 것도 허용된다.
}\par\nobreak\nopagebreak[4]\kfResetAreaFlag
\begin{kfQBlock}{01}{객관식}
\kfQStem{윗글을 바탕으로 할 때, <보기>의 ㉠\textasciitilde{}㉤에 대한 설명으로 적절하지 \underline{않은} 것은?}
\begin{kfEvidence}㉠ ㅣ \\\relax
㉡ ㅔ \\\relax
㉢ ㅜ \\\relax
㉣ ㅗ \\\relax
㉤ ㅏ\end{kfEvidence}
\begin{kfChoices}
\kfChoice ㉠은 전설 모음이면서 고모음이다.
\kfChoice ㉡은 전설 모음이면서 중모음이다.
\kfChoice ㉢은 후설 모음이면서 원순 모음이다.
\kfChoice ㉣은 전설 모음이면서 원순 모음이다.
\kfChoice ㉤은 후설 모음이면서 저모음이다.
\end{kfChoices}
\end{kfQBlock}

\par\needspace{25\baselineskip}
\kfPassageFull{자음 분류}{%
국어의 자음은 공기가 목 안이나 입안에서 장애를 받으면서 나는 소리로, 모두 19개이다. 자음은 소리가 만들어지는 방법(조음 방법)에 따라 파열음, 파찰음, 마찰음, 비음, 유음으로 나뉜다. 파열음은 공기를 막았다가 터뜨리며 내는 소리(ㅂ, ㅃ, ㅍ, ㄷ, ㄸ, ㅌ, ㄱ, ㄲ, ㅋ)이고, 파찰음은 막았다가 서서히 마찰을 일으켜 내는 소리(ㅈ, ㅉ, ㅊ)이다. 마찰음은 좁은 틈 사이로 공기를 내보내 마찰시키는 소리(ㅅ, ㅆ, ㅎ)이며, 비음은 코로 공기를 내보내며 내는 소리(ㅁ, ㄴ, ㅇ), 유음은 혀끝을 잇몸에 대었다가 떼며 내는 소리(ㄹ)이다. 또한 소리가 만들어지는 위치(조음 위치)에 따라 입술소리, 잇몸소리, 센입천장소리, 여린입천장소리, 목청소리로 구분하며, 소리의 세기에 따라 예사소리, 된소리, 거센소리로 나뉜다.
}\par\nobreak\nopagebreak[4]\kfResetAreaFlag
\begin{kfQBlock}{02}{객관식}
\kfQStem{윗글을 바탕으로 할 때, <보기>의 자음에 대한 설명으로 적절한 것은?}
\begin{kfEvidence}ㄱ, ㅈ, ㅅ, ㄴ, ㄹ\end{kfEvidence}
\begin{kfChoices}
\kfChoice 'ㄱ'과 'ㅈ'은 조음 방법이 같다.
\kfChoice 'ㅈ'과 'ㅅ'은 조음 위치가 같다.
\kfChoice 'ㅅ'과 'ㄹ'은 조음 방법이 같다.
\kfChoice 'ㄴ'과 'ㄹ'은 조음 위치가 같다.
\kfChoice 'ㄱ'과 'ㄴ'은 조음 방법이 같다.
\end{kfChoices}
\end{kfQBlock}

\par\needspace{25\baselineskip}
\kfPassageFull{이중모음}{%
국어의 모음은 단모음과 이중 모음으로 나뉜다. 단모음은 발음할 때 혀의 위치나 입술 모양이 변하지 않는 모음이고, 이중 모음은 발음하는 동안 혀의 위치나 입술 모양이 변하는 모음이다. 이중 모음은 반모음과 단모음이 결합하여 만들어진다. 반모음 'ㅣ'와 단모음이 결합한 것으로 'ㅑ, ㅕ, ㅛ, ㅠ, ㅒ, ㅖ'가 있고, 반모음 'ㅗ/ㅜ'와 단모음이 결합한 것으로 'ㅘ, ㅝ, ㅙ, ㅞ'가 있다. 'ㅢ'는 반모음 'ㅣ'가 뒤에 오는 특수한 경우이다. 한편 표준 발음법에서 'ㅟ'와 'ㅚ'는 원칙적으로 단모음이지만, 이중 모음으로 발음하는 것도 허용하고 있다.
}\par\nobreak\nopagebreak[4]\kfResetAreaFlag
\begin{kfQBlock}{03}{객관식}
\kfQStem{윗글을 바탕으로 할 때, <보기>에서 이중 모음만을 모두 고른 것은?}
\begin{kfEvidence}ㄱ. ㅑ \\\relax
ㄴ. ㅓ \\\relax
ㄷ. ㅘ \\\relax
ㄹ. ㅡ \\\relax
ㅁ. ㅖ\end{kfEvidence}
\begin{kfChoices}
\kfChoice ㄱ, ㄴ, ㄷ
\kfChoice ㄱ, ㄷ, ㅁ
\kfChoice ㄴ, ㄷ, ㄹ
\kfChoice ㄱ, ㄹ, ㅁ
\kfChoice ㄴ, ㄷ, ㅁ
\end{kfChoices}
\end{kfQBlock}

\par\needspace{25\baselineskip}
\kfPassageFull{모음 조화}{%
국어에서 모음은 양성 모음과 음성 모음으로 나뉜다. 양성 모음은 'ㅏ, ㅗ'처럼 밝고 가벼운 느낌을 주는 모음이고, 음성 모음은 'ㅓ, ㅜ'처럼 어둡고 무거운 느낌을 주는 모음이다. 모음 조화란 양성 모음은 양성 모음끼리, 음성 모음은 음성 모음끼리 어울리는 현상을 말한다. 이 현상은 특히 의성어·의태어와 용언의 어간에 결합하는 어미에서 잘 나타난다. 예를 들어 '살살'(양성)에 대응하는 '슬슬'(음성), '찰랑찰랑'(양성)에 대응하는 '출렁출렁'(음성)이 있다. 또한 용언 어간에 결합하는 어미에서도 모음 조화가 적용되어, '막다'의 어간 '막-'에는 '막-아'가, '먹다'의 어간 '먹-'에는 '먹-어'가 결합한다. 다만 현대 국어에서는 모음 조화가 엄격하게 지켜지지 않는 경우도 많다.
}\par\nobreak\nopagebreak[4]\kfResetAreaFlag
\begin{kfQBlock}{04}{객관식}
\kfQStem{윗글을 바탕으로 할 때, 모음 조화가 적용된 예로 적절하지 \underline{않은} 것은?}
\begin{kfChoices}
\kfChoice '잡다'의 어간에 어미 '-아'가 결합하여 '잡아'가 된다.
\kfChoice '넣다'의 어간에 어미 '-어'가 결합하여 '넣어'가 된다.
\kfChoice '깔깔'은 양성 모음끼리 어울린 의태어이다.
\kfChoice '담다'의 어간에 어미 '-어'가 결합하여 '담어'가 된다.
\kfChoice '출렁출렁'은 음성 모음끼리 어울린 의태어이다.
\end{kfChoices}
\end{kfQBlock}

\par\needspace{25\baselineskip}
\kfPassageFull{음운과 음성 구분}{%
언어학에서 음성과 음운은 구분되는 개념이다. 음성은 사람의 발음 기관을 통해 나오는 물리적인 소리 그 자체를 의미한다. 반면 음운은 말의 뜻을 구별해 주는 소리의 가장 작은 단위이다. 예를 들어, '불'과 '물'에서 'ㅂ'과 'ㅁ'은 각각 다른 의미를 만들어 내므로 서로 다른 음운이다. 국어의 음운에는 자음 19개, 모음 21개와 같은 분절 음운과, 소리의 길이나 높낮이와 같은 비분절 음운이 있다. 분절 음운은 다른 음운과 분리하여 하나씩 떼어 낼 수 있지만, 비분절 음운은 홀로 존재할 수 없고 분절 음운에 얹혀서만 실현된다. 예를 들어 '눈'(雪)과 '눈:'(眼), '밤'(夜)과 '밤:'(栗)은 소리의 길이로 뜻이 구별되는데, 이때의 길이가 비분절 음운에 해당한다.
}\par\nobreak\nopagebreak[4]\kfResetAreaFlag
\begin{kfQBlock}{05}{객관식}
\kfQStem{윗글을 바탕으로 할 때, <보기>에 대한 이해로 적절하지 \underline{않은} 것은?}
\begin{kfEvidence}'말'이라는 단어는 문맥에 따라 다양한 의미를 가진다. \\\relax
(가) 말(言) — 사람의 생각이나 느낌을 표현하는 음성 기호 \\\relax
(나) 말:(馬) — 동물의 한 종류 \\\relax
(다) 말(斗) — 곡식 등을 되는 단위\end{kfEvidence}
\begin{kfChoices}
\kfChoice (가)와 (나)는 분절 음운이 동일하다.
\kfChoice (가)와 (나)의 의미 차이는 비분절 음운에 의해 발생한다.
\kfChoice (가)와 (다)는 비분절 음운에 의해 의미가 구별된다.
\kfChoice (나)에서 모음의 길이가 길게 발음되는 것은 비분절 음운에 해당한다.
\kfChoice (가)의 'ㅁ', 'ㅏ', 'ㄹ'은 각각 하나의 분절 음운이다.
\end{kfChoices}
\end{kfQBlock}

\kfEndMC


\kfStartLit{「자화상」 — 윤동주 외 4편}


\kfPassageNote{「자화상」 — 윤동주}{%
산모퉁이를 돌아 논가 외딴 우물을 홀로 찾아가선 \\[4pt]
가만히 들여다봅니다. \\[14pt]
우물 속에는 달이 밝고 구름이 흐르고 하늘이 펼치고 \\[4pt]
파아란 바람이 불고 가을이 있습니다. \\[14pt]
그리고 한 사나이가 있습니다. \\[4pt]
어쩐지 그 사나이가 미워져 돌아갑니다. \\[14pt]
돌아가다 생각하니 그 사나이가 가엾어집니다. \\[4pt]
도로 가 들여다보니 사나이는 그대로 있습니다. \\[14pt]
다시 그 사나이가 미워져 돌아갑니다. \\[4pt]
돌아가다 생각하니 그 사나이가 그리워집니다. \\[14pt]
우물 속에는 달이 밝고 구름이 흐르고 하늘이 펼치고 \\[4pt]
파아란 바람이 불고 가을이 있고 추억처럼 사나이가 있습니다.
}\kfResetAreaFlag
\par\kfMaybeNeedspace{50\baselineskip}\noindent{\kfTipMark\,\,}{\small\color{kfInk} 빈칸에 알맞은 말을 채워 넣으시오. (초성 힌트 참고)}\par\nopagebreak[4]\smallskip\nopagebreak[4]
\begin{kfActTable}{|>{\centering\arraybackslash}m{2.6cm}|m{6cm}|m{4cm}|}
\rowcolor{kfPrimaryLight!50}\textbf{\color{kfPrimary}항목} & \textbf{\color{kfPrimary}내용 (빈칸 채우기)} & \textbf{\color{kfPrimary}답안 작성}\\\hline
갈래 & ㅈㅇㅅ, ㅅㅈㅅ ① & \kfTBlank \\\hline
성격 & ㅅㅊㅈ, ㄱㅈㅈ ② & \kfTBlank \\\hline
주요 소재 & ㅇㅁ, ㅅㄴㅇ ③ & \kfTBlank \\\hline
우물의 상징 & 자기 내면을 비추는 ㄱㅇ ④ & \kfTBlank \\\hline
사나이의 정체 & 화자의 또 다른 ㅈㅇ / ㅂㅅ ⑤ & \kfTBlank \\\hline
감정 변화 & ㅁㅇ → ㄱㅇ → ㄱㄹㅇ ⑥ & \kfTBlank \\\hline
표현 기법 1 & 우물을 거울로 활용한 ㅇㅇㅂ ⑦ & \kfTBlank \\\hline
표현 기법 2 & 첫 연과 마지막 연이 유사한 ㅅㅁㅅㄱ ⑧ & \kfTBlank \\\hline
시대 배경 & ㅇㅈ ㄱㅈㄱ (1939년 작) ⑨ & \kfTBlank \\\hline
주제 & ㅈㅇ ㅅㅊ과 ㅈㄱ ㅌㄱ ⑩ & \kfTBlank \\\hline
마지막 연 변주 & 'ㅊㅇ처럼 사나이가 있습니다' ⑪ & \kfTBlank \\\hline
출전 & 시집 『ㅎㄴ과 ㅂㄹ과 ㅂ과 ㅅ』 ⑫ & \kfTBlank \\\hline
\end{kfActTable}
\kfSecHeader{kfAreaLiterature}{문제}{}

\kfBeginMC
\begin{kfQBlock}{01}{객관식}
\kfQStem{윗시에 대한 설명으로 적절하지 \underline{않은} 것은?}
\begin{kfChoices}
\kfChoice 자연물을 통해 자아의 내면을 들여다보는 구조를 취하고 있다.
\kfChoice 화자의 감정이 '미움→가엾음→그리움'으로 변화하고 있다.
\kfChoice 수미상관의 구조를 활용하여 시상을 마무리하고 있다.
\kfChoice 우물 속 '사나이'는 화자와 대립하는 타인을 가리킨다.
\kfChoice 반복과 변주를 통해 화자의 내면 갈등을 드러내고 있다.
\end{kfChoices}
\end{kfQBlock}

\begin{kfQBlock}{02}{객관식}
\kfQStem{윗시의 '우물'이 지닌 상징적 의미로 가장 적절한 것은?}
\begin{kfChoices}
\kfChoice 화자가 현실에서 도피하기 위해 찾아가는 은신처
\kfChoice 화자의 과거 기억이 저장된 시간의 보관소
\kfChoice 화자가 자기 내면을 비추어 보는 성찰의 매개체
\kfChoice 화자가 자연과 합일을 이루기 위한 신성한 공간
\kfChoice 화자의 죄의식을 씻어 내는 정화(淨化)의 장소
\end{kfChoices}
\end{kfQBlock}

\begin{kfQBlock}{03}{객관식}
\kfQStem{화자가 우물 속 사나이를 '미워'하다가 '가엾어'하고 다시 '그리워'하는 까닭을 가장 적절하게 설명한 것은?}
\begin{kfChoices}
\kfChoice 사나이가 시간이 지남에 따라 외모가 변했기 때문이다.
\kfChoice 화자가 자기 자신에 대한 감정을 복합적으로 인식해 가기 때문이다.
\kfChoice 우물 속에 비친 사나이가 여러 명으로 바뀌기 때문이다.
\kfChoice 화자가 사나이와 과거에 겪었던 갈등을 회상하기 때문이다.
\kfChoice 계절의 변화에 따라 우물에 비친 풍경이 점점 달라지기 때문이다.
\end{kfChoices}
\end{kfQBlock}

\begin{kfQBlock}{04}{객관식}
\kfQStem{윗시의 마지막 연에서 '추억처럼 사나이가 있습니다'라는 구절의 효과로 가장 적절한 것은?}
\begin{kfChoices}
\kfChoice 사나이가 이미 죽은 존재임을 암시하여 비극적 분위기를 형성한다.
\kfChoice 화자가 사나이의 존재를 잊고 자연 풍경에 몰입했음을 보여 준다.
\kfChoice 화자가 사나이와의 갈등을 완전히 해소하고 화해했음을 보여 준다.
\kfChoice 사나이를 추억의 정서로 감싸며 수용의 자세를 드러낸다.
\kfChoice 과거에 만났던 실제 인물에 대한 향수를 환기한다.
\end{kfChoices}
\end{kfQBlock}

\begin{kfQBlock}{05}{객관식}
\kfQStem{윗시를 바탕으로 <보기>를 감상한 내용으로 적절하지 \underline{않은} 것은?}
\begin{kfEvidence}윤동주의 시에서 '우물'은 자기 성찰의 공간으로 자주 등장한다. 일제 강점기라는 시대적 상황에서 윤동주는 적극적인 저항 대신 내면적 성찰을 통해 부끄러움 없는 삶을 추구했다. 이러한 태도는 그의 여러 작품에서 '부끄러움'과 '별'과 같은 시어를 통해 반복적으로 나타난다.\end{kfEvidence}
\begin{kfChoices}
\kfChoice '외딴 우물을 홀로 찾아가'는 것은 자기 성찰을 위해 혼자만의 공간을 찾는 행위로 볼 수 있겠군.
\kfChoice 화자가 사나이를 '미워'하는 것은 시대에 적극적으로 저항하지 못하는 자신에 대한 부끄러움과 관련이 있겠군.
\kfChoice '파아란 바람'이나 '가을' 같은 자연물은 화자가 직접적인 항일 의지를 표출하는 수단이겠군.
\kfChoice 마지막 연에서 사나이를 '추억처럼' 받아들이는 것은 자기 자신과의 내면적 화해 시도로 이해할 수 있겠군.
\kfChoice 이 시의 성찰적 태도는 <보기>에서 말한 '부끄러움 없는 삶의 추구'와 맥이 닿아 있겠군.
\end{kfChoices}
\end{kfQBlock}

\kfEndMC

\kfPassageNote{「유리창」 — 정지용}{%
유리에 차고 슬픈 것이 어린거린다. \\[4pt]
열없이 붙어서서 입김을 흐리우니 \\[4pt]
길들은 양 언 날개를 파닥거린다. \\[4pt]
지우고 보고 지우고 보아도 \\[4pt]
새까만 밤이 밀려나가고 밀려와 부딪키고, \\[4pt]
물 먹은 별이, 반짝, 보석 처럼 백인다. \\[4pt]
밤에 홀로 유리를 닦는 것은 \\[4pt]
외로운 황홀한 심사이어니, \\[4pt]
고운 폐혈관이 찢어진 채로 \\[4pt]
아아, 너는 산(山)새처럼 날아갔구나!
}\kfResetAreaFlag
\par\kfMaybeNeedspace{50\baselineskip}\noindent{\kfTipMark\,\,}{\small\color{kfInk} 빈칸에 알맞은 말을 채워 넣으시오. (초성 힌트 참고)}\par\nopagebreak[4]\smallskip\nopagebreak[4]
\begin{kfActTable}{|>{\centering\arraybackslash}m{2.6cm}|m{6cm}|m{4cm}|}
\rowcolor{kfPrimaryLight!50}\textbf{\color{kfPrimary}항목} & \textbf{\color{kfPrimary}내용 (빈칸 채우기)} & \textbf{\color{kfPrimary}답안 작성}\\\hline
갈래 & ㅈㅇㅅ, ㅅㅈㅅ ① & \kfTBlank \\\hline
성격 & ㄱㄱㅈ, ㅇㅅㅈ ② & \kfTBlank \\\hline
주요 소재 & ㅇㄹ(창), ㅇㄱ, ㅂ ③ & \kfTBlank \\\hline
유리창의 상징 & 화자와 대상 사이의 넘을 수 없는 ㄱㄱ ④ & \kfTBlank \\\hline
'너'의 정체 & 화자가 그리워하는 ㅈㅇ ㅇㅇ (ㅈㅎㅇ) ⑤ & \kfTBlank \\\hline
핵심 기법 1 & 시각·촉각 등 ㄱㄱㅈ ㅅㅅ 활용 ⑥ & \kfTBlank \\\hline
핵심 기법 2 & '외로운 황홀한'처럼 모순된 말의 결합 = ㅇㅅㅂ ⑦ & \kfTBlank \\\hline
감각적 표현 & '물 먹은 별이, 반짝, 보석 처럼 ㅂㅇㄷ' ⑧ & \kfTBlank \\\hline
화자의 행위 & 'ㅈㅇ고 보고 ㅈㅇ고 보아도' — 반복 행위 ⑨ & \kfTBlank \\\hline
마지막 행 & '너는 ㅅㅅ처럼 날아갔구나' ⑩ & \kfTBlank \\\hline
주제 & 사랑하는 존재의 ㅅㅅ과 ㄱㄹㅇ, ㅈㅈ ㅇㅅ ⑪ & \kfTBlank \\\hline
출전 & 『ㅈㅈㅇ ㅅㅈ』(1935) ⑫ & \kfTBlank \\\hline
\end{kfActTable}
\kfSecHeader{kfAreaLiterature}{문제}{}

\kfBeginMC
\begin{kfQBlock}{01}{객관식}
\kfQStem{윗시에 대한 설명으로 적절한 것은?}
\begin{kfChoices}
\kfChoice 명령형 어조를 활용하여 독자에게 행동을 촉구하고 있다.
\kfChoice 반어적 표현을 통해 사회 현실을 비판하고 있다.
\kfChoice 대화체를 사용하여 대상과의 친밀감을 드러내고 있다.
\kfChoice 감각적 이미지를 활용하여 화자의 정서를 형상화하고 있다.
\kfChoice 시간의 흐름에 따라 계절의 변화를 순차적으로 묘사하고 있다.
\end{kfChoices}
\end{kfQBlock}

\begin{kfQBlock}{02}{객관식}
\kfQStem{윗시에서 '유리(창)'가 지닌 상징적 의미로 가장 적절한 것은?}
\begin{kfChoices}
\kfChoice 화자를 외부의 위험에서 보호하는 안전장치
\kfChoice 화자와 그리운 대상을 가르는 넘을 수 없는 경계
\kfChoice 화자의 과거와 현재를 연결해 주는 시간의 통로
\kfChoice 화자가 자신의 모습을 비추어 보는 거울
\kfChoice 세상의 아름다움을 선별하여 보여 주는 액자
\end{kfChoices}
\end{kfQBlock}

\begin{kfQBlock}{03}{객관식}
\kfQStem{'외로운 황홀한 심사'에서 '외로운'과 '황홀한'이 함께 쓰인 효과로 가장 적절한 것은?}
\begin{kfChoices}
\kfChoice 서로 모순되는 감정을 결합하여 화자의 복합적 심리를 드러낸다.
\kfChoice 점차 슬픔에서 기쁨으로 감정이 전환됨을 보여 준다.
\kfChoice 화자가 고독을 즐기는 낙관적 세계관을 보여 준다.
\kfChoice 외로움이 극에 달하면 황홀해진다는 과학적 사실을 반영한다.
\kfChoice 앞뒤 구절의 의미를 강조하기 위한 단순한 수식이다.
\end{kfChoices}
\end{kfQBlock}

\begin{kfQBlock}{04}{객관식}
\kfQStem{윗시의 마지막 행 '아아, 너는 산(山)새처럼 날아갔구나!'에 대한 해석으로 가장 적절한 것은?}
\begin{kfChoices}
\kfChoice 그리운 대상이 새처럼 자유롭게 세상을 여행하고 있음을 축복하고 있다.
\kfChoice 화자 자신이 산새가 되어 대상을 찾아 날아가겠다는 의지를 표출하고 있다.
\kfChoice 그리운 대상이 화자의 곁을 떠나 돌이킬 수 없이 사라졌음을 탄식하고 있다.
\kfChoice 대상이 화자에게 돌아올 것이라는 확신을 감탄의 형식으로 표현하고 있다.
\kfChoice 자연 속 새의 움직임을 객관적으로 묘사하여 시상을 마무리하고 있다.
\end{kfChoices}
\end{kfQBlock}

\begin{kfQBlock}{05}{객관식}
\kfQStem{윗시의 '지우고 보고 지우고 보아도'에 나타난 화자의 행위가 의미하는 바로 가장 적절한 것은?}
\begin{kfChoices}
\kfChoice 유리창의 얼룩을 깨끗이 닦아 시야를 확보하려는 실용적 행위
\kfChoice 유리 위에 그림을 그렸다 지우는 예술적 행위
\kfChoice 과거의 잘못을 반복적으로 반성하는 참회의 과정
\kfChoice 유리에 비친 자신의 모습을 거부하려는 자기 부정의 행위
\kfChoice 잊으려 해도 잊히지 않는 그리움이 반복되는 심리적 상태
\end{kfChoices}
\end{kfQBlock}

\kfEndMC

\kfPassageNote{「별 헤는 밤」 — 윤동주}{%
계절이 지나가는 하늘에는 \\[4pt]
가을로 가득 차 있습니다. \\[14pt]
나는 아무 걱정도 없이 \\[4pt]
가을 속의 별들을 다 헤일 듯합니다. \\[14pt]
가슴 속에 하나 둘 새겨지는 별을 \\[4pt]
이제 다 못 헤는 것은 \\[4pt]
쉬이 아침이 오는 까닭이요, \\[4pt]
내일 밤이 남은 까닭이요, \\[4pt]
아직 나의 청춘이 다하지 않은 까닭입니다. \\[14pt]
별 하나에 추억과 \\[4pt]
별 하나에 사랑과 \\[4pt]
별 하나에 쓸쓸함과 \\[4pt]
별 하나에 동경과 \\[4pt]
별 하나에 시와 \\[4pt]
별 하나에 어머니, 어머니, \\[14pt]
어머님, 나는 별 하나에 아름다운 말 한마디씩 불러 봅니다. 소학교 때 책상을 같이했던 아이들의 이름과, 패, 경, 옥, 이런 이국 소녀의 이름과, 벌써 아기 어머니 된 계집애들의 이름과, 가난한 이웃 사람들의 이름과, 비둘기, 강아지, 토끼, 노새, 노루, '프란시스 잠', '라이너 마리아 릴케', 이런 시인의 이름을 불러 봅니다. \\[14pt]
이네들은 너무나 멀리 있습니다. \\[4pt]
별이 아스라이 멀듯이. \\[14pt]
어머님, \\[4pt]
그리고 당신은 멀리 북간도에 계십니다. \\[14pt]
나는 무엇인지 그리워 \\[4pt]
이 많은 별빛이 내린 언덕 위에 \\[4pt]
내 이름자를 써 보고, \\[4pt]
흙으로 덮어 버리었습니다. \\[14pt]
딴은 밤을 새워 우는 벌레는 \\[4pt]
부끄러운 이름을 슬퍼하는 까닭입니다. \\[14pt]
그러나 겨울이 지나고 나의 별에도 봄이 오면 \\[4pt]
무덤 위에 파란 잔디가 피어나듯이 \\[4pt]
내 이름자 묻힌 언덕 위에도 \\[4pt]
자랑처럼 풀이 무성할 거외다.
}\kfResetAreaFlag
\par\kfMaybeNeedspace{50\baselineskip}\noindent{\kfTipMark\,\,}{\small\color{kfInk} 빈칸에 알맞은 말을 채워 넣으시오. (초성 힌트 참고)}\par\nopagebreak[4]\smallskip\nopagebreak[4]
\begin{kfActTable}{|>{\centering\arraybackslash}m{2.6cm}|m{6cm}|m{4cm}|}
\rowcolor{kfPrimaryLight!50}\textbf{\color{kfPrimary}항목} & \textbf{\color{kfPrimary}내용 (빈칸 채우기)} & \textbf{\color{kfPrimary}답안 작성}\\\hline
갈래 & ㅈㅇㅅ, ㅅㅈㅅ ① & \kfTBlank \\\hline
성격 & ㅅㅊㅈ, ㄱㅂㅈ ② & \kfTBlank \\\hline
주요 소재 & ㅂ, ㅂ, ㅇㄷ ③ & \kfTBlank \\\hline
별의 상징 & 화자에게 소중한 ㅊㅇ, ㅅㄹ, ㅇㅅ 등의 ㄱㅊ ④ & \kfTBlank \\\hline
구조적 특징 & '별 하나에 ○○과'의 ㅂㅂ과 ㅇㄱ ⑤ & \kfTBlank \\\hline
그리운 대상 & ㅇㅁㄴ, 소학교 때 아이들, 시인 등 ⑥ & \kfTBlank \\\hline
이름 행위 & 이름자를 쓰고 ㅎ으로 덮어 버림 ⑦ & \kfTBlank \\\hline
이름에 대한 감정 & 'ㅂㄲㄹㅇ 이름을 슬퍼하는 까닭' ⑧ & \kfTBlank \\\hline
마지막 연 이미지 & 'ㅈㄹ처럼 풀이 ㅁㅅ할 거외다' ⑨ & \kfTBlank \\\hline
마지막 연 의미 & ㅂㄲㄹㅇ을 넘어 떳떳한 존재가 되리라는 ㅎㅁ ⑩ & \kfTBlank \\\hline
시대 배경 & ㅇㅈ ㄱㅈㄱ (1941년 작) ⑪ & \kfTBlank \\\hline
주제 & ㅈㅇ ㅌㅅ과 ㅅㅅ ㅈㅊ, 이상적 자아에 대한 ㄱㅁ ⑫ & \kfTBlank \\\hline
\end{kfActTable}
\kfSecHeader{kfAreaLiterature}{문제}{}

\kfBeginMC
\begin{kfQBlock}{01}{객관식}
\kfQStem{윗시에 대한 설명으로 적절하지 \underline{않은} 것은?}
\begin{kfChoices}
\kfChoice 열거법을 활용하여 화자에게 소중한 것들을 나열하고 있다.
\kfChoice '\textasciitilde{}까닭입니다'의 반복을 통해 시적 리듬을 형성하고 있다.
\kfChoice 과거의 인물들과 현재의 화자를 대비하여 긴장감을 조성하고 있다.
\kfChoice 계절의 순환 이미지를 통해 미래에 대한 희망을 암시하고 있다.
\kfChoice 고백적 어조로 화자의 내면 심리를 직접적으로 드러내고 있다.
\end{kfChoices}
\end{kfQBlock}

\begin{kfQBlock}{02}{객관식}
\kfQStem{윗시에서 '별'이 지닌 상징적 의미로 가장 적절한 것은?}
\begin{kfChoices}
\kfChoice 화자가 도달할 수 없는 현실적 부(富)와 권력
\kfChoice 화자에게 소중한 추억, 사랑, 이상 등의 가치
\kfChoice 화자를 감시하는 권력의 시선
\kfChoice 시간의 흐름 속에서 변하지 않는 자연의 법칙
\kfChoice 화자가 극복해야 할 시련과 고난
\end{kfChoices}
\end{kfQBlock}

\begin{kfQBlock}{03}{객관식}
\kfQStem{화자가 '내 이름자를 써 보고, / 흙으로 덮어 버리'는 행위의 의미로 가장 적절한 것은?}
\begin{kfChoices}
\kfChoice 자신의 이름이 널리 알려지기를 바라는 명예욕의 표현
\kfChoice 현재의 부끄러운 자아를 묻고 새로운 자아로 거듭나고 싶은 소망
\kfChoice 장래에 이름을 남길 문학 작품을 쓰겠다는 결심
\kfChoice 일제에 의해 빼앗긴 이름을 되찾으려는 적극적 독립 운동의 의지
\kfChoice 죽음을 예감하고 자신의 무덤을 미리 만들어 보는 행위
\end{kfChoices}
\end{kfQBlock}

\begin{kfQBlock}{04}{객관식}
\kfQStem{윗시의 마지막 연에서 '무덤 위에 파란 잔디가 피어나듯이 / 내 이름자 묻힌 언덕 위에도 / 자랑처럼 풀이 무성할 거외다'가 지닌 의미로 가장 적절한 것은?}
\begin{kfChoices}
\kfChoice 자연이 인간의 죽음을 매우 무관심하게 덮어 버린다는 허무감
\kfChoice 봄이 되면 무덤을 찾아가겠다는 화자의 구체적 계획
\kfChoice 자신이 죽은 뒤에야 작품이 인정받을 것이라는 예언
\kfChoice 식물의 성장 과정을 관찰하려는 과학적 탐구심
\kfChoice 현재의 부끄러움을 넘어 미래에는 떳떳한 존재가 되리라는 희망
\end{kfChoices}
\end{kfQBlock}

\begin{kfQBlock}{05}{객관식}
\kfQStem{윗시를 바탕으로 <보기>를 감상한 내용으로 적절하지 \underline{않은} 것은?}
\begin{kfEvidence}윤동주의 시 세계는 '부끄러움의 윤리학'으로 요약되기도 한다. 그는 일제 강점기라는 암울한 시대에 적극적 항쟁 대신 자기 내면의 성찰을 통해 양심을 지키고자 했다. 별, 하늘, 바람 등의 자연 이미지는 순수하고 이상적인 세계를 표상하며, 화자는 그 이상에 비추어 자신의 현실적 삶을 반성한다.\end{kfEvidence}
\begin{kfChoices}
\kfChoice 화자가 별에 추억, 사랑 등을 대응시키는 것은 이상적 가치에 대한 동경으로 볼 수 있겠군.
\kfChoice '부끄러운 이름'은 시대에 적극적으로 저항하지 못하는 자신에 대한 반성과 관련이 있겠군.
\kfChoice 화자가 이름을 흙으로 덮는 것은 현실의 부끄러운 자아를 직시하는 성찰의 행위이겠군.
\kfChoice 마지막 연의 '자랑처럼 풀이 무성할 거외다'는 적극적인 무장 투쟁의 의지를 표명한 것이겠군.
\kfChoice '이네들은 너무나 멀리 있습니다'는 화자가 소중한 가치들과 단절된 현실을 인식하는 것이겠군.
\end{kfChoices}
\end{kfQBlock}

\kfEndMC

\kfPassageNote{「수오재기」 — 정약용}{%
나의 서재(書齋)를 '수오재(守吾齋)'라 이름 짓고 그 뜻을 기록한다.

무릇 지킨다[守]고 하는 것은 그 잃어버릴 염려가 있기 때문이다. 집을 지킨다고 하는 것은 도둑이 있기 때문이요, 몸을 지킨다고 하는 것은 질병이 있기 때문이요, 나를 지킨다고 하는 것은 나를 잃어버릴 염려가 있기 때문이다.

집이란 것은 담장과 문으로 둘러싸인 것이니, 하인을 시켜 지키게 할 수 있다. 몸이란 것은 혈기(血氣)와 골육(骨肉)으로 이루어진 것이니, 의원을 불러 지키게 할 수 있다. 그러나 나란 것은 이 몸에 붙어 있으면서도 이 몸과 같지 않은 것이니, 다른 사람이 대신 지켜 줄 수 없고 오직 내가 스스로 지킬 수밖에 없다.

그러므로 집을 지키는 것은 쉽고, 몸을 지키는 것은 그보다 어렵고, 나를 지키는 것은 더욱 어렵다. 집을 잃으면 다시 살 곳을 구할 수 있고, 몸이 상하면 약을 써서 고칠 수 있다. 그러나 나를 잃으면 다시 구할 길이 없으니, 어찌 두렵지 않겠는가?

대저 나란 것은 무엇인가? 인의예지(仁義禮智)가 이것이요, 효제충신(孝悌忠信)이 이것이다. 이것을 굳게 지키면 비록 곤궁한 처지에 빠지더라도 떳떳한 사람이 될 수 있거니와, 이것을 잃으면 비록 부귀영화를 누리더라도 사람 노릇을 한다고 할 수 없다.

내가 지금 바닷가 궁벽한 곳에 귀양살이를 하고 있어 집도 잃고 벼슬도 잃었으나, 나를 아직 잃지 않았으니 진실로 다행한 일이다. 이에 서재 이름을 '수오재'라 하여 스스로를 경계하노라.
}\kfResetAreaFlag
\par\kfMaybeNeedspace{50\baselineskip}\noindent{\kfTipMark\,\,}{\small\color{kfInk} 빈칸에 알맞은 말을 채워 넣으시오. (초성 힌트 참고)}\par\nopagebreak[4]\smallskip\nopagebreak[4]
\begin{kfActTable}{|>{\centering\arraybackslash}m{2.6cm}|m{6cm}|m{4cm}|}
\rowcolor{kfPrimaryLight!50}\textbf{\color{kfPrimary}항목} & \textbf{\color{kfPrimary}내용 (빈칸 채우기)} & \textbf{\color{kfPrimary}답안 작성}\\\hline
갈래 & ㄱㅈ ㅅㅍ — ㄱ(記)ㅁ ① & \kfTBlank \\\hline
서재 이름 & ㅅㅇㅈ(守吾齋) = 'ㄴ를 ㅈㅋ는 집' ② & \kfTBlank \\\hline
지킴 1단계 & ㅈ 지키기(守家) — ㅎㅇ이 대신 가능 ③ & \kfTBlank \\\hline
지킴 2단계 & ㅁ 지키기(守身) — ㅇㅇ이 도울 수 있음 ④ & \kfTBlank \\\hline
지킴 3단계 & ㄴ 지키기(守吾) — 오직 ㅈㄱ ㅈㅅ만 가능 ⑤ & \kfTBlank \\\hline
논리 전개 & 외면→내면으로 심화하는 ㅈㅊㅂ ⑥ & \kfTBlank \\\hline
'나'의 정의 & ㅇㅇㅇㅈ(仁義禮智), ㅎㅈㅊㅅ(孝悌忠信) ⑦ & \kfTBlank \\\hline
글쓴이 상황 & 바닷가 궁벽한 곳에서 ㅇㅂ 생활 ⑧ & \kfTBlank \\\hline
잃은 것 & ㅈ과 ㅂㅅ ⑨ & \kfTBlank \\\hline
지킨 것 & ㄴ (도덕적 ㅈㅇ) ⑩ & \kfTBlank \\\hline
주제 & ㅈㄱ ㅅㅇ과 ㅈㅈ ㅇㅁ에 대한 유배지 ㅅㅊ ⑪ & \kfTBlank \\\hline
출전 & 『ㅇㅇㄷ ㅈㅅ(與猶堂全書)』 ⑫ & \kfTBlank \\\hline
\end{kfActTable}
\kfSecHeader{kfAreaLiterature}{문제}{}

\kfBeginMC
\begin{kfQBlock}{01}{객관식}
\kfQStem{윗글에 대한 설명으로 적절한 것은?}
\begin{kfChoices}
\kfChoice 서재의 이름에 담긴 뜻을 밝히는 기(記)문의 형식을 취하고 있다.
\kfChoice 다른 인물과의 대화를 통해 논지를 전개하고 있다.
\kfChoice 역사적 사건을 시간 순서대로 나열하여 교훈을 이끌어 내고 있다.
\kfChoice 비유와 상징을 활용한 우화(寓話)의 형식으로 쓰여 있다.
\kfChoice 자연 경치를 묘사하며 풍류를 즐기는 태도를 드러내고 있다.
\end{kfChoices}
\end{kfQBlock}

\begin{kfQBlock}{02}{객관식}
\kfQStem{윗글에서 글쓴이가 '지킴'의 대상을 세 단계로 나눈 순서와 그 의도로 가장 적절한 것은?}
\begin{kfChoices}
\kfChoice 몸 → 집 → 나: 개인에서 사회로 논의를 확장하기 위해
\kfChoice 나 → 몸 → 집: 가장 중요한 것부터 제시하기 위해
\kfChoice 집 → 몸 → 나: 외면적인 것에서 내면적인 것으로 심화하기 위해
\kfChoice 집 → 나 → 몸: 순서를 뒤섞어 독자의 흥미를 유발하기 위해
\kfChoice 나 → 집 → 몸: 결론을 먼저 제시한 뒤 부연하기 위해
\end{kfChoices}
\end{kfQBlock}

\begin{kfQBlock}{03}{객관식}
\kfQStem{윗글에서 글쓴이가 말하는 '나(吾)'의 의미로 가장 적절한 것은?}
\begin{kfChoices}
\kfChoice 육체적으로 건강한 몸을 유지하는 것
\kfChoice 인의예지, 효제충신 등 도덕적 주체로서의 자아
\kfChoice 관직에 복귀하여 사회적 지위를 회복하는 것
\kfChoice 가문의 명예를 후손에게 전하는 것
\kfChoice 유배지에서 학문적 저술을 완성하는 것
\end{kfChoices}
\end{kfQBlock}

\begin{kfQBlock}{04}{객관식}
\kfQStem{윗글의 마지막 문단에서 '집도 잃고 벼슬도 잃었으나, 나를 아직 잃지 않았으니 진실로 다행한 일이다'가 드러내는 글쓴이의 태도로 가장 적절한 것은?}
\begin{kfChoices}
\kfChoice 유배 생활의 고통을 외면하며 현실로부터 도피하려는 태도
\kfChoice 외적 조건의 상실에도 도덕적 자아를 지켰다는 자긍심
\kfChoice 유배를 보낸 권력자들에 대한 깊은 원망과 복수 의지
\kfChoice 집과 벼슬을 머지않아 되찾을 수 있으리라는 낙관적 기대
\kfChoice 세속적 가치를 모두 버리고 산속에 은둔하겠다는 결심
\end{kfChoices}
\end{kfQBlock}

\begin{kfQBlock}{05}{객관식}
\kfQStem{윗글에서 글쓴이가 '나를 잃으면 다시 구할 길이 없다'고 한 까닭으로 가장 적절한 것은?}
\begin{kfChoices}
\kfChoice 집이나 몸과 달리 도덕적 자아는 다른 사람이 대신 지켜 줄 수 없기 때문이다.
\kfChoice 유배지에서는 서적이 매우 부족하여 학문을 다시 시작하기가 어렵기 때문이다.
\kfChoice 한번 잃은 관직은 조정의 허락 없이는 되찾을 수 없기 때문이다.
\kfChoice 나이가 들면 새로운 것을 배우기 어려워지기 때문이다.
\kfChoice 유배된 사람은 사회에서 영원히 격리되기 때문이다.
\end{kfChoices}
\end{kfQBlock}

\kfEndMC

\kfPassageNote{「지지헌기(止止軒記)」 — 이규보}{%
이른바, 지지(止止)라는 것은 능히 그 그칠 곳을 알아서 그치는 것이니, 그 그칠 곳이 아닌 데에 그치면, 그 그침은 그칠 곳에 그친 것이 아니다.

또 호랑이와 표범, 고라니와 사슴, 교룡은 늪과 못이나 굴에 있어야 그 그칠 곳을 알아서 그치는 것인데, 가령 본고장을 떠나서 혼잡한 성시(城市) 가운데에 그친다면 사람들이 재앙으로 여기고 따라서 해칠 것은 필연한 일이다.

나는 세상에 있어서 거만스러워 남과 합하는 일이 적으니, 길들여진 물건이 아니다. 만일 다른 사람들과 함께 나가고 나란히 달려서 명리의 지경에 그치게 된다면, 이는 호랑이와 표범, 고라니와 사슴, 교룡이 성시에 그친 것과 무엇이 다르겠는가? 이것은 내가 그 그칠 곳을 구하여 그치는 것이다. 그렇지 않으면 사람들이 재앙으로 여기고 따라서 해치는 자가 이를 것이다.

하였더니, 어떤 이가 말하기를, "자네의 말과 같이 한다면 산림이나 궁곡에 처하여 다른 사람들과 복잡하게 한곳에 있지 않은 연후에야 그칠 곳에 그쳤다고 할 수가 있다. 이제 자네가 그친 곳은 곧 성시의 가운데인데, 오히려 그칠 곳에 그쳤다고 하여, 호랑이와 표범, 고라니와 사슴, 교룡이 늪과 못이나 굴에 처하는 것에 비유하는 것은 무엇인가?" 하기에, 이렇게 대답하였다.

"벌레와 짐승이 늪과 못이나 굴에 처하는 것과 사람이 성시에 처하는 것은 역시 각각 그 그침의 떳떳한 것이다. 가령 사람이 늪과 못에 엎드리고 굴에 들어간다면, 역시 호랑이와 표범, 고라니와 사슴, 교룡이 성시에 들어간 것과 같으니, 독충이나 맹수가 또한 반드시 재앙으로 여기고 떼를 지어 해칠 것이다. 사람이 사람을 피하여 벌레와 짐승에게 해를 당하는 일을 나는 차마 하지 못한다.

또 사람이 사람을 꺼리어 해치기를 꾀하는 것은, 성시가 좁아서 같이 처하는 것에 인색해서가 아니라, 그 구하는 것과 그 이익을 다투기 때문이다. 진실로 사람들과 다투지 아니하여, 비록 대낮에 내 상자를 훔쳐가는 자가 있더라도 피하고 보지 않는다면, 사람이 성시에 처하는 것이 또한 호랑이와 표범, 고라니와 사슴, 교룡이 늪과 못이나 굴에 처하는 것과 같은데, 어찌 해칠 자가 있겠는가? 내가 거처를 이렇게 이름한 것은 대개 이러한 뜻이다."

정묘년 3월 10일에 기(記)한다.
}\kfResetAreaFlag
\par\kfMaybeNeedspace{50\baselineskip}\noindent{\kfTipMark\,\,}{\small\color{kfInk} 빈칸에 알맞은 말을 채워 넣으시오. (초성 힌트 참고)}\par\nopagebreak[4]\smallskip\nopagebreak[4]
\begin{kfActTable}{|>{\centering\arraybackslash}m{2.6cm}|m{6cm}|m{4cm}|}
\rowcolor{kfPrimaryLight!50}\textbf{\color{kfPrimary}항목} & \textbf{\color{kfPrimary}내용 (빈칸 채우기)} & \textbf{\color{kfPrimary}답안 작성}\\\hline
갈래 & ㄱㅈ ㅅㅍ — ㄱ(記) ① & \kfTBlank \\\hline
글쓴이 & 고려시대 문인 ㅇㄱㅂ ② & \kfTBlank \\\hline
출전 & 『ㄷㄱㅇㅅㄱㅈ』 ③ & \kfTBlank \\\hline
핵심 개념 & ㅈㅈ(止止) — 그칠 곳에 그치기 ④ & \kfTBlank \\\hline
비유 소재 & ㅎㄹㅇ·ㅍㅂ·ㄱㄹㄴ·ㅅㅅ·ㄱㄹ ⑤ & \kfTBlank \\\hline
짐승의 자리 & ㄴ과 ㅁ, ㄱ ⑥ & \kfTBlank \\\hline
글쓴이 본성 & ㄱㅁ스러워 남과 합하지 못함 ⑦ & \kfTBlank \\\hline
그쳐선 안 될 자리 & ㅁㄹ의 지경 ⑧ & \kfTBlank \\\hline
어떤 이의 반박 & ㅅㄹ이나 ㄱㄱ에 들어가야 ⑨ & \kfTBlank \\\hline
다툼의 원인 & ㅇㅇ을 다투기 때문 ⑩ & \kfTBlank \\\hline
결론 & ㅁㅇ이 멈추면 성시도 굴과 같음 ⑪ & \kfTBlank \\\hline
주제 & 본성에 맞는 자리에 ㅁㅁㄹ는 ㄱㄷㅇ ⑫ & \kfTBlank \\\hline
\end{kfActTable}
\kfSecHeader{kfAreaLiterature}{문제}{}

\kfBeginMC
\begin{kfQBlock}{01}{객관식}
\kfQStem{윗글에 대한 설명으로 적절하지 \underline{않은} 것은?}
\begin{kfChoices}
\kfChoice 기(記) 양식으로 자신의 거처에 붙인 이름의 뜻을 풀이하고 있다.
\kfChoice 짐승의 예를 들어 '그칠 곳에 그친다'의 뜻을 구체적으로 설명한다.
\kfChoice 어떤 이의 반박을 제시한 뒤 이에 대해 다시 답하는 문답 형식을 취한다.
\kfChoice 글쓴이는 자신을 거만하여 남과 잘 어울리지 못하는 사람으로 평가한다.
\kfChoice 산림이나 궁곡에 들어가야만 진정한 '지지'라고 결론짓는다.
\end{kfChoices}
\end{kfQBlock}

\begin{kfQBlock}{02}{객관식}
\kfQStem{윗글에서 글쓴이가 '호랑이와 표범, 고라니와 사슴, 교룡'을 예로 든 까닭으로 가장 적절한 것은?}
\begin{kfChoices}
\kfChoice 맹수가 인간 사회에 끼치는 위험성을 부각하기 위해서
\kfChoice 각 존재에 마땅히 있어야 할 자리가 있음을 비유하기 위해서
\kfChoice 야생 동물 보호의 필요성을 일깨우기 위해서
\kfChoice 사냥과 어로의 즐거움을 묘사하기 위해서
\kfChoice 맹수가 인간보다 더 지혜로운 존재임을 주장하기 위해서
\end{kfChoices}
\end{kfQBlock}

\begin{kfQBlock}{03}{객관식}
\kfQStem{윗글에서 '명리(名利)의 지경'에 대한 글쓴이의 인식으로 가장 적절한 것은?}
\begin{kfChoices}
\kfChoice 누구나 적극적으로 추구해야 할 인생의 궁극적 목표이다.
\kfChoice 자기 본성에 맞지 않는, 그쳐서는 안 되는 자리이다.
\kfChoice 가난한 자들에게만 허락된 좁은 영역이다.
\kfChoice 왕과 신하만이 도달할 수 있는 특별한 지위이다.
\kfChoice 자연에 묻혀 사는 은자만이 누리는 평온한 공간이다.
\end{kfChoices}
\end{kfQBlock}

\begin{kfQBlock}{04}{객관식}
\kfQStem{윗글의 마지막 단락에서 '진실로 사람들과 다투지 아니하여 \textasciitilde{} 어찌 해칠 자가 있겠는가?'가 의미하는 바로 가장 적절한 것은?}
\begin{kfChoices}
\kfChoice 도시에 살아도 이익을 다투지 않으면 누구도 해치지 않는다.
\kfChoice 재산이 많을수록 도둑의 위협으로부터 자유로워질 수 있다.
\kfChoice 도시를 떠나 깊은 산속으로 들어가야만 평화가 찾아온다.
\kfChoice 법과 제도로 다툼을 강제로 막아야 사회가 안정된다.
\kfChoice 사람들과의 모든 교류를 끊고 혼자 살아야 안전하다.
\end{kfChoices}
\end{kfQBlock}

\begin{kfQBlock}{05}{객관식}
\kfQStem{윗글을 바탕으로 <보기>를 이해한 내용으로 적절하지 \underline{않은} 것은?}
\begin{kfEvidence}문학 작품 속 공간의 의미는 그 공간을 인식하고 경험하는 인물에 따라 달라질 수 있다. 동일한 공간이라도 인물이 그 공간에 어떻게 머무는가, 또 무엇을 추구하는가에 따라 그 공간은 평화로운 자리가 되기도 하고 위험한 자리가 되기도 한다.\end{kfEvidence}
\begin{kfChoices}
\kfChoice 글쓴이에게 '성시'는 다툼을 멈추기만 하면 평화로운 자리가 될 수 있는 공간이겠군.
\kfChoice '어떤 이'에게 '성시'는 자기 그칠 곳에 그치기에는 부적절한 공간으로 인식되겠군.
\kfChoice '명리의 지경'은 본성에 맞지 않게 머무는 자에게는 위험한 자리가 되겠군.
\kfChoice 동일한 '성시'가 글쓴이와 어떤 이에게 다르게 인식되는 것은 공간의 의미가 인물에 따라 달라진다는 점을 보여 주는군.
\kfChoice 글쓴이는 '성시'에서 벗어나 산림에 들어감으로써 비로소 진정한 '지지'를 이루었겠군.
\end{kfChoices}
\end{kfQBlock}

\kfEndMC

\kfStartNonlit{소크라테스와 자기 인식 외 4편}


\kfPassageNote{소크라테스와 자기 인식 (인문)}{%
고대 그리스의 철학자 소크라테스는 '너 자신을 알라'라는 말로 유명하다. 이 말은 원래 그리스 델포이 신전에 새겨져 있던 격언이었는데, 소크라테스는 이를 자신의 철학적 탐구의 출발점으로 삼았다. 그에 따르면, 인간이 진정으로 지혜로워지기 위해서는 먼저 자기 자신이 무엇을 알고 무엇을 모르는지를 정확히 파악해야 한다.

소크라테스는 당시 아테네에서 지혜롭다고 자처하는 사람들을 찾아다니며 대화를 나누었다. 그는 정치가, 시인, 장인 등을 만나 질문을 던졌는데, 대화를 통해 그들이 실제로는 자신이 안다고 생각하는 것을 제대로 알지 못한다는 사실을 드러냈다. 소크라테스는 이러한 대화법을 통해, 자신은 적어도 '자기가 모른다는 것을 안다'는 점에서 그들보다 나은 위치에 있다고 보았다. 이것이 이른바 '무지(無知)의 지(知)'이다.

소크라테스의 대화법은 오늘날 '산파술'이라고도 불린다. 산파가 산모의 출산을 돕듯이, 소크라테스는 상대방이 스스로 진리를 발견하도록 돕는 역할을 자처했기 때문이다. 그는 직접 답을 가르치지 않고, 질문을 통해 상대방의 생각 속에 숨어 있는 모순을 끌어냄으로써 스스로 깨달음에 이르게 했다. 이때 사용한 방법이 바로 '문답법' 또는 '변증법'이다.

이러한 소크라테스의 접근은 단순한 지식의 습득이 아니라, 반성적 사유를 통한 자기 이해를 강조한다. 그에게 철학이란 외부 세계에 대한 지식을 쌓는 것이 아니라, 자기 자신의 무지를 인정하고 끊임없이 자기 자신을 되돌아보는 행위였다. 소크라테스는 이러한 태도를 '검토되지 않은 삶은 살 가치가 없다'라는 유명한 말로 요약했다. 이는 자기 삶을 비판적으로 돌아보는 것이 인간다운 삶의 본질이라는 뜻이다.

소크라테스의 자기 인식론은 이후 서양 철학의 근본적인 방향을 설정했다. 플라톤은 스승의 가르침을 이어받아 이데아론을 발전시켰고, 아리스토텔레스는 관찰과 분류를 통한 체계적 탐구의 길을 열었다. 그러나 소크라테스가 남긴 가장 중요한 유산은 특정 이론이 아니라, '자기 자신에 대한 끊임없는 물음'이라는 철학적 태도 그 자체라고 할 수 있다.
}\kfResetAreaFlag
\kfSecHeader{kfAreaNonlit}{문제}{}

\kfBeginMC
\begin{kfQBlock}{01}{객관식}
\kfQStem{윗글의 내용과 일치하는 것은?}
\begin{kfChoices}
\kfChoice 소크라테스는 '너 자신을 알라'라는 격언을 직접 만들어 신전에 새겼다.
\kfChoice 소크라테스는 아테네의 정치가와 시인들이 참된 지혜를 갖추고 있다고 인정했다.
\kfChoice 소크라테스의 대화법은 상대방에게 답을 직접 가르치는 방식으로 이루어졌다.
\kfChoice 소크라테스는 자신이 모른다는 사실을 아는 것이 지혜의 출발점이라고 보았다.
\kfChoice 소크라테스의 철학은 외부 세계에 대한 객관적 지식을 축적하는 것을 목표로 했다.
\end{kfChoices}
\end{kfQBlock}

\begin{kfQBlock}{02}{객관식}
\kfQStem{윗글에서 '산파술'이라는 이름이 붙은 까닭으로 가장 적절한 것은?}
\begin{kfChoices}
\kfChoice 소크라테스가 의학에 관한 깊은 지식을 가지고 있었기 때문이다.
\kfChoice 대화를 통해 상대방이 스스로 진리를 깨닫도록 돕는 역할을 했기 때문이다.
\kfChoice 산모에게 출산 방법을 가르쳐 주듯이 지식을 전달했기 때문이다.
\kfChoice 철학적 탐구의 고통을 출산의 산고에 비유하여 따뜻하게 위로했기 때문이다.
\kfChoice 소크라테스의 어머니가 산파였기 때문에 그 직업을 존경한 것이다.
\end{kfChoices}
\end{kfQBlock}

\begin{kfQBlock}{03}{객관식}
\kfQStem{윗글을 바탕으로 <보기>를 이해한 내용으로 적절하지 \underline{않은} 것은?}
\begin{kfEvidence}선생님: 철수야, 정의란 무엇이라고 생각하니? \\\relax
철  수: 정의란 강한 자의 이익입니다. \\\relax
선생님: 그렇다면 강한 자가 잘못된 법을 만들면, 그것도 정의가 되는 것이니? \\\relax
철  수: 음... 그건 아닌 것 같습니다. \\\relax
선생님: 그러면 정의가 단순히 강한 자의 이익이라는 네 처음 생각은 충분한 것이었을까? \\\relax
철  수: 아닙니다. 다시 생각해 봐야 할 것 같습니다.\end{kfEvidence}
\begin{kfChoices}
\kfChoice 선생님은 답을 직접 제시하지 않고 질문만으로 철수의 생각을 이끌고 있다.
\kfChoice 철수는 대화를 통해 자신의 처음 생각에 모순이 있음을 스스로 깨닫고 있다.
\kfChoice 이 대화는 소크라테스의 산파술과 유사한 방식으로 진행되고 있다.
\kfChoice 선생님은 철수에게 정의에 대한 올바른 정의를 최종적으로 알려 주고 있다.
\kfChoice 철수는 처음에 자신이 확신한 생각을 반성적으로 되돌아보게 되었다.
\end{kfChoices}
\end{kfQBlock}

\begin{kfQBlock}{04}{객관식}
\kfQStem{㉠ '검토되지 \underline{않은} 삶은 살 가치가 없다'의 의미로 가장 적절한 것은?}
\begin{kfChoices}
\kfChoice 많은 지식을 축적하지 못한 삶은 후회할 만하다.
\kfChoice 다른 사람의 평가를 받지 않은 삶은 성공할 수 없다.
\kfChoice 자기 삶을 비판적으로 돌아보는 것이 인간다운 삶의 본질이다.
\kfChoice 시험에서 좋은 성적을 받기 위해 끊임없이 공부해야 한다.
\kfChoice 사회 규범을 따르지 않는 삶은 가치가 없다는 경고이다.
\end{kfChoices}
\end{kfQBlock}

\kfEndMC

\kfPassageNote{데카르트의 '나는 생각한다' (인문)}{%
근대 철학의 아버지라 불리는 르네 데카르트는 확실한 지식의 토대를 세우기 위해 '방법적 회의'라는 독특한 사유 방식을 사용했다. 방법적 회의란 조금이라도 의심할 수 있는 것은 모두 거짓이라고 가정하고, 더 이상 의심할 수 없는 확실한 것만을 찾아가는 방법이다. 이때 '방법적'이라는 말은 의심 자체가 목적이 아니라, 확실한 진리에 도달하기 위한 수단임을 뜻한다.

데카르트는 먼저 감각 경험을 의심했다. 우리는 꿈속에서도 생생한 경험을 하므로, 지금 눈앞에 보이는 것이 현실인지 꿈인지 확실히 구별할 수 없다. 또한 감각은 때때로 우리를 속이기도 한다. 예를 들어, 먼 곳에 있는 탑은 둥글게 보이지만 가까이 가면 사각형인 경우가 있다. 이처럼 감각이 한 번이라도 우리를 속인 적이 있다면, 감각에서 얻은 모든 지식은 의심의 대상이 된다.

다음으로 데카르트는 수학적 진리마저 의심했다. '2 더하기 3은 5이다'와 같은 명제는 꿈속에서든 현실에서든 참인 것 같지만, 만약 전능한 악마가 우리를 속이고 있다면 이마저도 확실하다고 할 수 없다. 데카르트는 이 사고 실험을 통해 외부 세계에 대한 거의 모든 믿음을 괄호 안에 넣었다.

그러나 이렇게 모든 것을 의심하는 과정에서, 데카르트는 결코 의심할 수 없는 하나의 사실을 발견했다. 그것은 바로 '지금 이 순간 의심하고 있는 나 자신은 존재한다'는 것이다. 의심한다는 것은 생각하는 행위이므로, 생각하고 있는 한 나는 반드시 존재해야 한다. 이것이 바로 '코기토 에르고 숨(Cogito ergo sum)', 즉 '나는 생각한다, 고로 나는 존재한다'라는 유명한 명제이다.

데카르트는 이 명제를 철학의 제1원리로 삼았다. 여기서 '나'란 신체가 아니라 '생각하는 것', 즉 정신을 가리킨다. 데카르트에 따르면 정신과 신체는 본질적으로 다른 두 실체이다. 정신은 생각하는 것을 본질로 하고, 신체는 공간 속에 펼쳐진 것을 본질로 한다. 이러한 입장을 '심신 이원론'이라 한다. 이 이원론은 이후 근대 철학에서 정신과 물질의 관계를 둘러싼 치열한 논쟁을 불러일으켰으며, 오늘날 인지과학과 뇌과학에서도 여전히 중요한 화두가 되고 있다.
}\kfResetAreaFlag
\kfSecHeader{kfAreaNonlit}{문제}{}

\kfBeginMC
\begin{kfQBlock}{01}{객관식}
\kfQStem{윗글에 대한 설명으로 가장 적절한 것은?}
\begin{kfChoices}
\kfChoice 방법적 회의의 과정을 단계적으로 서술하고 그 결론을 제시하고 있다.
\kfChoice 데카르트의 사상에 대한 다양한 학자들의 비판을 소개하고 있다.
\kfChoice 코기토 명제의 논리적 오류를 분석하고 대안을 제안하고 있다.
\kfChoice 경험론과 합리론을 비교하여 데카르트의 한계를 지적하고 있다.
\kfChoice 심신 이원론의 과학적 근거를 실험 결과를 들어 입증하고 있다.
\end{kfChoices}
\end{kfQBlock}

\begin{kfQBlock}{02}{객관식}
\kfQStem{윗글의 내용과 일치하는 것은?}
\begin{kfChoices}
\kfChoice 데카르트는 감각 경험이 항상 정확하므로 의심할 필요가 없다고 보았다.
\kfChoice 데카르트는 수학적 진리는 어떤 경우에도 의심할 수 없다고 주장했다.
\kfChoice 방법적 회의에서 '방법적'이란 의심 자체를 목적으로 한다는 뜻이다.
\kfChoice 데카르트는 '나'를 신체가 아닌 생각하는 정신으로 규정했다.
\kfChoice 심신 이원론은 정신과 신체가 동일한 실체라는 입장이다.
\end{kfChoices}
\end{kfQBlock}

\begin{kfQBlock}{03}{객관식}
\kfQStem{데카르트가 '전능한 악마' 가설을 사용한 까닭으로 가장 적절한 것은?}
\begin{kfChoices}
\kfChoice 종교적 믿음의 중요성을 강조하기 위해서이다.
\kfChoice 수학적 진리마저 의심할 수 있음을 보이기 위해서이다.
\kfChoice 악마의 실재 가능성을 철학적으로 증명하기 위해서이다.
\kfChoice 감각 경험의 신뢰성을 회복하기 위해서이다.
\kfChoice 코기토 명제가 거짓임을 입증하기 위해서이다.
\end{kfChoices}
\end{kfQBlock}

\begin{kfQBlock}{04}{객관식}
\kfQStem{윗글을 바탕으로 <보기>를 이해한 내용으로 적절하지 \underline{않은} 것은?}
\begin{kfEvidence}영희: 나는 지금 내 눈앞에 사과가 있다고 확신해. \\\relax
철수: 그런데 꿈에서도 사과를 볼 수 있잖아. 지금 네가 꿈을 꾸고 있지 않다는 걸 어떻게 증명할 수 있어? \\\relax
영희: 그건... 글쎄. 하지만 1 더하기 1은 2잖아. 그건 꿈이든 현실이든 변하지 않아. \\\relax
철수: 만약 네 머릿속에서 누군가가 계산까지 조작하고 있다면? \\\relax
영희: 그러면 내가 확실히 아는 건 아무것도 없는 거야?\end{kfEvidence}
\begin{kfChoices}
\kfChoice 철수의 첫 번째 반론은 데카르트가 감각을 의심한 것과 유사하다.
\kfChoice 영희의 '1 더하기 1은 2'라는 주장은 데카르트가 의심한 수학적 진리에 해당한다.
\kfChoice 철수의 두 번째 반론은 전능한 악마 가설과 같은 역할을 한다.
\kfChoice 영희의 마지막 반응은 코기토 명제를 발견한 것에 해당한다.
\kfChoice 철수는 방법적 회의의 단계를 순서대로 적용하고 있다.
\end{kfChoices}
\end{kfQBlock}

\begin{kfQBlock}{05}{객관식}
\kfQStem{㉠ '심신 이원론'에 대한 이해로 가장 적절한 것은?}
\begin{kfChoices}
\kfChoice 정신은 신체의 생리적 기능에서 비롯된 부수적 현상이다.
\kfChoice 정신과 신체는 본질이 같은 하나의 실체이다.
\kfChoice 정신의 본질은 생각이고 신체의 본질은 공간적 연장이다.
\kfChoice 신체가 소멸하면 정신도 함께 소멸한다.
\kfChoice 정신보다 신체가 존재론적으로 우위에 있다.
\end{kfChoices}
\end{kfQBlock}

\kfEndMC

\kfPassageNote{인간은 어떻게 세계를 인식하는가 (인문)}{%
인간은 어떻게 세계를 알 수 있는가? 이 물음에 대해 서양 근대 철학에서는 크게 두 가지 입장이 대립했다. 하나는 감각 경험을 통해 지식이 형성된다고 보는 경험론이고, 다른 하나는 이성의 능력을 통해 확실한 지식에 도달할 수 있다고 보는 합리론이다. 이 두 흐름은 인식의 원천을 어디에 두느냐를 둘러싸고 오랜 논쟁을 벌였다.

경험론의 대표적 철학자인 존 로크는 인간의 마음이 태어날 때 아무것도 적혀 있지 않은 '백지(tabula rasa)' 상태라고 주장했다. 그에 따르면, 모든 지식은 감각 경험을 통해 마음에 새겨진다. 우리가 '빨갛다', '차갑다'와 같은 관념을 갖게 되는 것은 사물을 보고 만지는 경험 덕분이다. 경험 없이는 어떠한 관념도 마음속에 존재할 수 없다는 것이 로크의 핵심 주장이다.

로크의 뒤를 이은 데이비드 흄은 경험론을 더욱 철저하게 밀고 나갔다. 흄은 우리가 당연하게 여기는 인과 관계마저 감각 경험에 의한 습관적 기대일 뿐이라고 주장했다. 예를 들어, 불에 손을 대면 뜨겁다는 것은 반복된 경험에서 온 기대이지, 불과 뜨거움 사이에 논리적으로 필연적인 연결이 있는 것은 아니다. 흄에게 경험으로 확인할 수 없는 모든 관념은 공허한 것이었다.

반면, 합리론자인 데카르트와 라이프니츠는 감각 경험보다 이성이 더 확실한 지식의 원천이라고 보았다. 데카르트는 감각이 우리를 속일 수 있지만, 이성에 의한 명석판명한 인식은 확실하다고 주장했다. 라이프니츠는 로크의 백지설에 반대하며, 인간의 마음에는 경험 이전부터 잠재적으로 갖추어진 관념이 있다고 보았다. 그는 마음을 아무것도 없는 백지가 아니라 '결이 있는 대리석'에 비유했는데, 대리석 안에 이미 어떤 무늬의 가능성이 들어 있듯이, 인간의 정신 속에도 선천적 관념의 씨앗이 내재한다는 것이다.

경험론과 합리론의 대립은 이후 임마누엘 칸트에 의해 종합되었다. 칸트는 경험 없는 인식은 공허하고, 개념 없는 경험은 맹목적이라고 하며 양쪽 모두의 한계를 지적했다. 그에 따르면, 인간은 감각을 통해 경험의 재료를 받아들이되, 이 재료를 시간과 공간이라는 선천적 형식 및 인과 관계 같은 범주를 통해 정리함으로써 비로소 지식을 구성한다. 이러한 칸트의 입장은 경험론과 합리론의 논쟁을 새로운 차원으로 끌어올렸다는 평가를 받는다.
}\kfResetAreaFlag
\kfSecHeader{kfAreaNonlit}{문제}{}

\kfBeginMC
\begin{kfQBlock}{01}{객관식}
\kfQStem{윗글의 서술 방식에 대한 설명으로 가장 적절한 것은?}
\begin{kfChoices}
\kfChoice 두 대립적 입장을 비교한 뒤 제삼의 관점으로 종합하고 있다.
\kfChoice 하나의 이론을 시대 순서에 따라 발전 과정을 추적하고 있다.
\kfChoice 실험 결과를 제시하여 특정 입장의 우위를 논증하고 있다.
\kfChoice 하나의 사례를 다양한 관점에서 반복적으로 분석하고 있다.
\kfChoice 전문 용어의 어원을 풀이하여 개념을 설명하고 있다.
\end{kfChoices}
\end{kfQBlock}

\begin{kfQBlock}{02}{객관식}
\kfQStem{윗글의 내용과 일치하지 않는 것은?}
\begin{kfChoices}
\kfChoice 로크는 인간의 마음이 태어날 때 백지 상태라고 보았다.
\kfChoice 흄은 인과 관계가 논리적으로 필연적인 연결이라고 주장했다.
\kfChoice 데카르트는 이성에 의한 명석판명한 인식이 확실하다고 보았다.
\kfChoice 라이프니츠는 마음에 선천적 관념이 잠재한다고 주장했다.
\kfChoice 칸트는 감각의 재료를 선천적 형식으로 정리해야 지식이 된다고 보았다.
\end{kfChoices}
\end{kfQBlock}

\begin{kfQBlock}{03}{객관식}
\kfQStem{㉠ '결이 있는 대리석'이라는 비유가 의미하는 바로 가장 적절한 것은?}
\begin{kfChoices}
\kfChoice 마음은 경험에 의해 처음부터 새로 만들어지는 것이다.
\kfChoice 마음에는 경험 이전부터 잠재적인 관념이 내재해 있다.
\kfChoice 감각 경험이 쌓이면 마음이 점차 단단해진다.
\kfChoice 마음의 구조는 대리석처럼 변하지 않는다.
\kfChoice 선천적 관념은 경험을 통해 완전히 제거될 수 있다.
\end{kfChoices}
\end{kfQBlock}

\begin{kfQBlock}{04}{객관식}
\kfQStem{윗글을 바탕으로 <보기>의 상황을 분석한 것으로 적절하지 \underline{않은} 것은?}
\begin{kfEvidence}갓 태어난 아기가 뜨거운 물체에 손을 대어 본 적이 없는데도 불빛을 보면 손을 움츠린다. 반면 그 아기가 '빨간색'이라는 이름을 알게 되려면 주위 사람들이 빨간 사물을 가리키며 '빨강'이라고 말해 주는 경험이 필요하다.\end{kfEvidence}
\begin{kfChoices}
\kfChoice 아기가 불빛에 손을 움츠리는 것은 경험 이전의 반응이라 할 수 있다.
\kfChoice '빨강'이라는 이름의 학습에는 감각 경험이 필수적이라는 점에서 로크의 입장을 뒷받침한다.
\kfChoice 경험 없이도 나타나는 반응이 있다는 점은 라이프니츠의 선천적 관념과 관련된다.
\kfChoice 칸트라면 불빛에 대한 반응과 '빨강' 학습 모두에서 선천적 형식과 경험의 결합을 강조할 것이다.
\kfChoice 흄이라면 아기의 움츠림 반응도 반복된 경험에서 형성된 습관이라고 설명할 것이다.
\end{kfChoices}
\end{kfQBlock}

\begin{kfQBlock}{05}{객관식}
\kfQStem{윗글에 등장하는 철학자들의 입장을 비교한 것으로 적절한 것은?}
\begin{kfChoices}
\kfChoice 로크와 데카르트는 모두 감각 경험이 지식의 유일한 원천이라고 보았다.
\kfChoice 흄과 라이프니츠는 모두 선천적 관념의 존재를 인정했다.
\kfChoice 로크와 흄은 모두 경험을 지식의 기반으로 중시했으나, 인과 관계에 대한 견해는 달랐다.
\kfChoice 데카르트와 칸트는 모두 감각 경험을 완전히 배제하고 이성만으로 지식을 구성한다고 보았다.
\kfChoice 라이프니츠와 칸트는 모두 경험 이전의 요소를 부정하고 경험만을 강조했다.
\end{kfChoices}
\end{kfQBlock}

\kfEndMC

\kfPassageNote{플라톤의 동굴 비유 (인문)}{%
플라톤은 《국가》에서 인간의 인식 상태를 설명하기 위해 '동굴의 비유'를 제시했다. 이 비유에 따르면, 어둡고 깊은 동굴 안에 사람들이 태어날 때부터 묶여 있다. 그들은 고개를 돌릴 수 없어 오직 동굴 벽만 바라볼 수 있다. 그들의 뒤편에는 불이 타고 있고, 불과 사람들 사이를 다양한 사물의 모형이 지나간다. 벽에는 이 모형들의 그림자가 비치는데, 동굴 속 사람들은 이 그림자가 세계의 전부라고 믿는다.

플라톤은 이 그림자를 '현상'에, 그림자를 만들어 내는 실제 사물을 '실재'에 비유했다. 동굴 속 사람들이 그림자를 실재라고 착각하듯, 인간은 감각으로 포착하는 현상의 세계를 진짜 세계라고 여기며 살아간다. 그러나 플라톤에 따르면, 감각으로 경험하는 세계는 참된 실재의 불완전한 모사(模寫)에 지나지 않는다.

이 비유에서 한 사람이 사슬에서 풀려나 동굴 밖으로 나간다. 처음에는 강렬한 햇빛에 눈이 부셔 아무것도 볼 수 없지만, 점차 적응하면서 그림자가 아닌 실제 사물들을 보게 된다. 그리고 마침내 태양을 직접 바라볼 수 있게 된다. 플라톤은 이 태양을 '선(善)의 이데아'에 비유했다. 태양이 빛을 비추어 사물을 볼 수 있게 해 주듯이, 선의 이데아는 모든 이데아에 존재와 인식 가능성을 부여하는 궁극적 원리이다.

여기서 '이데아'란 감각으로는 파악할 수 없고 오직 이성으로만 인식할 수 있는, 변하지 않는 완전한 본질을 말한다. 예를 들어, 세상에는 아름다운 사물이 많지만 모두 변하고 사라진다. 그러나 '아름다움 자체', 즉 아름다움의 이데아는 영원불변하게 존재한다. 플라톤은 감각 세계에 존재하는 개별 사물들은 이러한 이데아를 '분유(分有)'한 것, 즉 이데아에 참여함으로써 그 성질을 나누어 가진 것이라고 설명했다.

동굴 밖으로 나간 사람은 진리를 깨달은 뒤 다시 동굴로 돌아온다. 그는 여전히 그림자만 보고 있는 사람들에게 바깥세계의 진실을 알려 주려 하지만, 오랫동안 어둠에 익숙해진 눈으로는 도리어 동굴 속에서 잘 보지 못하게 된다. 동굴 속 사람들은 그가 밖에 나갔다 와서 눈이 나빠졌다고 비웃으며, 그의 말을 믿지 않는다. 플라톤은 이 장면을 통해, 참된 지식을 추구하는 철학자가 현실 세계에서 오해와 배척을 받을 수 있음을 보여 주었다. 이는 스승 소크라테스가 아테네 시민들에 의해 사형 선고를 받은 역사적 사건을 연상시키기도 한다.
}\kfResetAreaFlag
\kfSecHeader{kfAreaNonlit}{문제}{}

\kfBeginMC
\begin{kfQBlock}{01}{객관식}
\kfQStem{윗글의 내용과 일치하는 것은?}
\begin{kfChoices}
\kfChoice 동굴 속 사람들은 자유롭게 고개를 돌려 뒤편의 불을 볼 수 있었다.
\kfChoice 동굴 벽에 비친 그림자는 실제 사물의 불완전한 모사에 해당한다.
\kfChoice 동굴 밖으로 나간 사람은 곧바로 태양을 직접 바라볼 수 있었다.
\kfChoice 이데아는 감각을 통해 쉽게 파악할 수 있는 개별 사물을 가리킨다.
\kfChoice 동굴 속 사람들은 돌아온 사람의 말을 즉시 받아들이고 함께 밖으로 나갔다.
\end{kfChoices}
\end{kfQBlock}

\begin{kfQBlock}{02}{객관식}
\kfQStem{플라톤이 '태양'을 '선의 이데아'에 비유한 까닭으로 가장 적절한 것은?}
\begin{kfChoices}
\kfChoice 태양이 모든 사물에 열을 전달하여 생명을 유지시키기 때문이다.
\kfChoice 태양이 가장 먼 곳에 있어 쉽게 도달할 수 없기 때문이다.
\kfChoice 태양이 빛으로 사물을 볼 수 있게 하듯 선의 이데아가 존재와 인식의 근원이 되기 때문이다.
\kfChoice 태양이 밤낮으로 변하듯 선의 이데아도 시대에 따라 변하기 때문이다.
\kfChoice 태양이 그림자를 만들어 내듯 선의 이데아가 현상 세계의 모든 환영을 만들기 때문이다.
\end{kfChoices}
\end{kfQBlock}

\begin{kfQBlock}{03}{객관식}
\kfQStem{윗글의 문맥상 '분유(分有)'의 의미로 가장 적절한 것은?}
\begin{kfChoices}
\kfChoice 이데아가 개별 사물로 쪼개져 점차 소멸해 가는 것
\kfChoice 개별 사물이 이데아에 참여하여 그 성질을 나누어 갖는 것
\kfChoice 감각 세계의 사물들이 서로 성질을 교환해 가는 것
\kfChoice 이성적 사유를 통해 이데아를 온전히 소유하게 되는 것
\kfChoice 현상 세계가 이데아를 대신하여 참된 실재가 되는 것
\end{kfChoices}
\end{kfQBlock}

\begin{kfQBlock}{04}{객관식}
\kfQStem{윗글을 바탕으로 <보기>를 이해한 내용으로 적절하지 \underline{않은} 것은?}
\begin{kfEvidence}어떤 사람이 평생 스마트폰 화면으로만 풍경 사진을 보았다. 어느 날 그는 처음으로 실제 산과 바다를 보게 되었다. 그는 사진과 실물의 차이에 크게 놀랐고, 돌아와서 친구들에게 '화면 속 풍경은 진짜가 아니다'라고 말했다. 그러나 친구들은 '화면이 곧 진짜'라며 그를 이상한 사람 취급했다.\end{kfEvidence}
\begin{kfChoices}
\kfChoice 스마트폰 화면 속 풍경은 동굴 벽의 그림자에 해당한다.
\kfChoice 실제 산과 바다는 이데아의 세계에 비유될 수 있다.
\kfChoice 친구들이 '화면이 곧 진짜'라고 하는 것은 그림자를 실재로 착각하는 것과 유사하다.
\kfChoice 그가 실물을 보고 놀란 것은 동굴 밖의 햇빛에 눈이 적응하는 과정에 해당한다.
\kfChoice 친구들이 그를 이상한 사람 취급한 것은 동굴 밖에서 돌아온 사람을 환영하는 것에 해당한다.
\end{kfChoices}
\end{kfQBlock}

\begin{kfQBlock}{05}{객관식}
\kfQStem{플라톤이 동굴의 비유를 통해 궁극적으로 전달하고자 한 것으로 가장 적절한 것은?}
\begin{kfChoices}
\kfChoice 감각 세계가 참된 실재이며 이성적 탐구는 불필요하다.
\kfChoice 철학자는 동굴 밖에 머물면서 대중과 소통하지 않아야 한다.
\kfChoice 감각을 넘어 이성으로 참된 실재를 인식해야 한다.
\kfChoice 그림자를 관찰하는 것이 가장 정확한 인식 방법이다.
\kfChoice 동굴 속 생활이 동굴 밖보다 안전하고 바람직하다.
\end{kfChoices}
\end{kfQBlock}

\kfEndMC

\kfPassageNote{현상학이란 무엇인가 (인문)}{%
20세기 초, 에드문트 후설은 '사물 자체로 돌아가라'라는 표어를 내걸며 현상학이라는 새로운 철학 운동을 시작했다. 여기서 '사물 자체'란 우리가 일상적으로 접하는 물리적 대상만을 뜻하는 것이 아니라, 선입견이나 이론적 가정을 걷어 내고 의식에 직접 주어지는 것 그대로를 말한다. 후설은 기존 철학이 과학적 이론이나 형이상학적 전제에 지나치게 의존하여 인간 경험의 본질을 제대로 포착하지 못한다고 비판했다.

후설이 제안한 핵심 방법이 바로 '판단중지', 그리스어로 '에포케(epoché)'이다. 에포케란 외부 세계의 존재 여부에 대한 판단을 잠시 중지하고, 의식에 나타나는 현상 자체에만 주목하는 태도를 뜻한다. 예를 들어, 눈앞에 꽃이 있다고 하자. 보통 우리는 '저 꽃이 실제로 존재한다'고 당연하게 받아들인다. 그러나 에포케를 실행하면 꽃이 실재하는지 여부는 괄호 안에 넣어 두고, '나에게 빨갛고 향기로운 것으로 경험되고 있다'는 의식 체험 자체를 분석 대상으로 삼는다.

에포케를 통해 드러나는 의식의 가장 근본적인 특성은 '지향성(intentionality)'이다. 지향성이란 의식이 항상 '무엇인가에 대한' 의식이라는 뜻이다. 우리는 아무것도 생각하지 않으면서 생각할 수 없고, 아무것도 보지 않으면서 볼 수 없다. 기쁨은 기쁨의 대상이 있어야 하고, 슬픔은 슬픔의 대상이 있어야 한다. 이처럼 의식은 언제나 무엇인가를 향해 열려 있으며, 대상 없는 의식이란 있을 수 없다. 후설은 이러한 지향성을 의식의 본질적 구조라고 규정했다.

지향성 개념은 인간과 세계의 관계를 새롭게 이해하는 열쇠가 된다. 전통적으로 철학에서는 인간의 의식과 외부 세계를 분리하여 '의식이 세계를 어떻게 반영하는가'라는 문제를 다루었다. 그러나 후설의 현상학에서 의식은 세계로부터 고립된 내면의 공간이 아니라, 처음부터 세계를 향해 열려 있는 활동이다. 의식과 대상은 서로 떼어 놓을 수 없는 관계에 있으며, 의식은 대상을 구성하는 능동적 역할을 수행한다.

후설의 현상학은 이후 하이데거, 사르트르, 메를로-퐁티 등에게 깊은 영향을 미쳤다. 하이데거는 현상학의 방법을 받아들이면서도 '존재'의 의미를 묻는 존재론으로 방향을 전환했고, 사르트르는 지향성 개념을 바탕으로 자유와 실존의 철학을 전개했으며, 메를로-퐁티는 신체 경험의 현상학을 탐구했다. 이처럼 후설이 세운 현상학적 방법은 20세기 유럽 철학의 주요 흐름을 형성하는 토대가 되었다.
}\kfResetAreaFlag
\kfSecHeader{kfAreaNonlit}{문제}{}

\kfBeginMC
\begin{kfQBlock}{01}{객관식}
\kfQStem{윗글의 서술 방식에 대한 설명으로 가장 적절한 것은?}
\begin{kfChoices}
\kfChoice 핵심 개념을 순차적으로 설명한 뒤 그 영향 관계를 제시하고 있다.
\kfChoice 두 대립적 이론을 비교하여 하나의 우위를 논증하고 있다.
\kfChoice 역사적 사건의 원인과 결과를 시간 순서로 나열하고 있다.
\kfChoice 실험 결과를 제시하고 이를 통해 가설을 검증하고 있다.
\kfChoice 특정 학자의 이론에 대한 반박과 대안을 중심으로 서술하고 있다.
\end{kfChoices}
\end{kfQBlock}

\begin{kfQBlock}{02}{객관식}
\kfQStem{윗글의 내용과 일치하는 것은?}
\begin{kfChoices}
\kfChoice '사물 자체'란 오직 물리적 대상만을 뜻한다.
\kfChoice 에포케는 외부 세계가 존재하지 않는다고 단정하는 태도이다.
\kfChoice 지향성이란 의식이 항상 무엇인가에 대한 의식이라는 뜻이다.
\kfChoice 후설에 따르면 대상 없는 순수한 의식 상태가 가능하다.
\kfChoice 하이데거는 후설의 현상학을 전면 부정하고 전혀 다른 출발점에서 시작했다.
\end{kfChoices}
\end{kfQBlock}

\begin{kfQBlock}{03}{객관식}
\kfQStem{후설이 '에포케'를 제안한 까닭으로 가장 적절한 것은?}
\begin{kfChoices}
\kfChoice 외부 세계의 존재 자체를 형이상학적으로 부정하기 위해서이다.
\kfChoice 의식에 나타나는 현상에 주목해 경험의 본질을 포착하기 위해서이다.
\kfChoice 당대 자연과학적 이론의 객관적 정당성을 입증하기 위해서이다.
\kfChoice 감각 경험을 의심하고 오직 이성적 사유만을 신뢰하기 위해서이다.
\kfChoice 기존의 다양한 철학 이론을 하나의 체계로 종합하기 위해서이다.
\end{kfChoices}
\end{kfQBlock}

\begin{kfQBlock}{04}{객관식}
\kfQStem{윗글을 바탕으로 <보기>를 이해한 내용으로 적절하지 \underline{않은} 것은?}
\begin{kfEvidence}민수는 밤하늘에 떠 있는 달을 바라보고 있다. 그는 문득 '저 달이 정말로 저기에 있는 걸까, 아니면 내 의식이 만들어 낸 것일까?'라고 의문을 품었다. 그래서 달이 실재하는지 여부는 일단 제쳐 두고, '나에게 둥글고 밝은 것으로 경험되고 있다'는 사실에만 집중하기로 했다.\end{kfEvidence}
\begin{kfChoices}
\kfChoice 민수가 달의 실재 여부를 제쳐 둔 것은 에포케를 실행한 것에 해당한다.
\kfChoice 민수의 의식이 달을 향해 있다는 점은 지향성의 사례로 볼 수 있다.
\kfChoice '나에게 둥글고 밝은 것으로 경험되고 있다'에 집중하는 것은 의식 체험의 분석에 해당한다.
\kfChoice 민수는 달이 실재하지 않는다는 결론을 내린 것이다.
\kfChoice 민수의 태도는 후설이 말한 '사물 자체로 돌아가라'는 표어와 관련된다.
\end{kfChoices}
\end{kfQBlock}

\begin{kfQBlock}{05}{객관식}
\kfQStem{윗글의 '지향성' 개념에 비추어 볼 때, 다음 중 지향성을 가장 잘 보여 주는 사례는?}
\begin{kfChoices}
\kfChoice 아무 생각 없이 멍하니 허공을 바라보는 상태
\kfChoice 뇌의 신경 세포가 화학 물질을 분비하는 과정
\kfChoice 어제 친구에게 받은 선물을 떠올리며 기뻐하는 마음
\kfChoice 심장이 무의식적으로 박동을 유지하는 현상
\kfChoice 컴퓨터가 입력된 데이터를 자동으로 처리하는 과정
\end{kfChoices}
\end{kfQBlock}

\kfEndMC

\kfStartWeekly{패턴 워크북 — 총 18문항}


\kfPassageNote{문항 1 / 섞어치기}{%
생명체의 유전 정보는 아데닌, 구아닌, 사이토신, 티민으로 구성된 디엔에이(DNA) 염기 서열에 저장되어 있다. 동일한 염기 서열을 가진 세포들이라도 유전자 발현 양상은 다를 수 있는데, 이러한 현상을 후성 유전으로 설명할 수 있다. 후성 유전이란 DNA 염기 서열의 변화 없이 화학적 변형을 통해 유전자 발현을 조절하는 과정으로, 대표적인 예로는 DNA 메틸화, 히스톤의 아세틸화 등이 있다.

(중략)

DNA 메틸화는 유전자의 전사 조절 부위인 프로모터 내 사이토신에 메틸기가 부착되어 특정 유전자의 전사가 억제되는 과정이다. 이 과정은 주로 사이토신과 구아닌이 인접한 서열(CpG)에서 일어나며, DNA 메틸화 효소(DNMT)의 촉매 작용으로 촉진된다. 메틸화된 DNA는 전사 인자들의 결합을 방해하는 구조를 형성하여 유전자 발현을 차단하는 일종의 정지 신호 역할을 한다. 반면, 히스톤의 아세틸화는 전사를 촉진하는 화학적 변형이다. 히스톤은 DNA에 감겨 있는 단백질 구조체인데, 그 꼬리 부분의 라

(중략)
}
\begin{kfQBlock}{01}{객관식}
\kfQStem{윗글을 이해한 내용으로 적절하지 \underline{않은} 것은?}
\begin{kfChoices}
\kfChoice 보리노스타트는 HDAC를 저해하는 방식으로 종양 억제 유전자에 부착된 메틸기를 제거하여 종양 억제 유전자의 발현을 복원한다.
\kfChoice CpG의 사이토신에 메틸기가 부착되면 전사 인자들의 결합이 어려운 구조로 DNA가 변화한다.
\kfChoice HDAC의 활성이 높으면 히스톤이 탈아세틸화되어 DNA와 히스톤 간의 상호 작용이 강화된다.
\end{kfChoices}
\end{kfQBlock}

\kfPassageNote{문항 2 / 부정·상대어}{%
생명체의 유전 정보는 아데닌, 구아닌, 사이토신, 티민으로 구성된 디엔에이(DNA) 염기 서열에 저장되어 있다. 동일한 염기 서열을 가진 세포들이라도 유전자 발현 양상은 다를 수 있는데, 이러한 현상을 후성 유전으로 설명할 수 있다. 후성 유전이란 DNA 염기 서열의 변화 없이 화학적 변형을 통해 유전자 발현을 조절하는 과정으로, 대표적인 예로는 DNA 메틸화, 히스톤의 아세틸화 등이 있다.

(중략)

DNA 메틸화는 유전자의 전사 조절 부위인 프로모터 내 사이토신에 메틸기가 부착되어 특정 유전자의 전사가 억제되는 과정이다. 이 과정은 주로 사이토신과 구아닌이 인접한 서열(CpG)에서 일어나며, DNA 메틸화 효소(DNMT)의 촉매 작용으로 촉진된다. 메틸화된 DNA는 전사 인자들의 결합을 방해하는 구조를 형성하여 유전자 발현을 차단하는 일종의 정지 신호 역할을 한다. 반면, 히스톤의 아세틸화는 전사를 촉진하는 화학적 변형이다. 히스톤은 DNA에 감겨 있는 단백질 구조체인데, 그 꼬리 부분의 라

(중략)
}
\begin{kfQBlock}{02}{객관식}
\kfQStem{윗글을 바탕으로 <보기>를 이해한 내용으로 적절하지 \underline{않은} 것은?}
\begin{kfEvidence}악성 종양에는 일반적으로 암 재발과 전이에 관여하는 암 줄기세포가 존재한다. 이 세포들은 강한 자기 재생 능력과 약물 배출 능력을 보유하고 있어 기존 항암제에 저항하는 강력한 내성을 보인다. 그런데 최근 연구에 따르면, 히스톤 메틸화가 암 줄기세포의 특성 유지에 중요한 역할을 담당하는 것으로 밝혀졌다. 히스톤의 특정 아미노산에 메틸기가 결합하는 히스톤 메틸화 현상은 메틸기를 제거하는 효소인 히스톤 탈메틸화 효소(HDM)에 의해 정밀하게 조절된다. 그러나 일부 암 줄기세포에서는 HDM이 비정상적으로 과활성화되어 종양 억제 유전자의 발현이 지속적으로 억제된 상태를 유지한다. 따라서 HDM 저해제는 이렇게 억제된 유전자의 발현을 회복시킴으로써 암 줄기세포를 정상 세포로 분화시키거나 세포 사멸을 유도하는 효과를 보인다.\end{kfEvidence}
\begin{kfChoices}
\kfChoice HDM 저해제가 암 줄기세포에 작용하는 현상은 히스톤의 메틸기를 제거하는 반응을 촉진하여 억제되어 있던 종양 억제 유전자의 발현을 유도하는 방식으로 일어날 것이다.
\kfChoice HDM 저해제는 암 줄기세포의 후성 유전적 오류를 교정하여 세포 자살 프로그램을 재가동시키는 방식으로 항암 효과를 나타낼 것이다.
\kfChoice 기존 항암제에 대한 내성은 히스톤의 변형과 같은 후성 유전적 메커니즘에 의해 암 줄기세포의 약물 배출 관련 유전자 발현이 변화하여 나타날 수 있다.
\end{kfChoices}
\end{kfQBlock}

\kfPassageNote{문항 3 / 주체·객체}{%
최근까지 일반적인 컴퓨터 환경에서는 운영 체제를 제어하거나 사용자 인터페이스를 처리하고, 문서를 작성하거나 웹을 탐색하는 등 다양한 범용 작업이 주요 연산 대상이었다. 이러한 작업은 명령어의 종류가 많고 그 순서가 수시로 바뀌며, 연산 간 종속성이 크다는 특징이 있다. 컴퓨터의 연산은 주로 여러 개의 코어로 구성된 ㉠중앙 처리 장치(CPU)에서 수행된다. 코어는 컴퓨터의 두뇌 역할을 하는 것으로 명령어를 해석하고 각종 연산을 실행한다. CPU는 여러 개의 고성능 코어로 구성되어 복잡한 제어 흐름을 정밀하게 조율할 수 있도록 설계되었다. 그러나 인공 지능 학습, 고해상도 영상 처리, 과학 기술 관련 계산 등 대규모 데이터를 대상으로 동일한 연산을 반복 수행해야 하는 작업이 증가하면서 기존의 연산 환경에 변화가 생겼다. CPU는 코어 수에 한계가 있고 개별 코어가 복잡한 제어 흐름에 맞추어 다양한 연산을 처리하도록 설계되어 있기 때문에 동일한 연산을 대량으로 반복 수행해야 하는 현대의 

(중략)
}
\begin{kfQBlock}{03}{객관식}
\kfQStem{윗글을 바탕으로 <보기>의 [수치 모델]을 이해한 내용으로 가장 적절한 것은?}
\begin{kfEvidence}기상 예측 [수치 모델]은 초기 상태 분석과 변수 조율 후, 대규모 격자 기반 계산을 반복 수행한다. 일부 구간에서는 예외 조건으로 연산 흐름이 갈라지며, 시스템은 상황에 맞게 연산 흐름과 데이터를 분담하도록 설계된다.\end{kfEvidence}
\begin{kfChoices}
\kfChoice 지역 단위 동일 계산 반복 중에는 메모리 접근 방식을 조정해 충돌을 줄이겠군.
\kfChoice 초기 상태 정밀 분석과 변수 조율은 병렬 반복 계산이므로 GPU 중심으로 처리하겠군.
\kfChoice 시간 단위 대기 상태 측정은 종속성이 커 스레드 순차 실행이 필수이겠군.
\end{kfChoices}
\end{kfQBlock}

\kfPassageNote{문항 4 / 순서·방향}{%
원자에는 중심에 핵이 있고, 핵 주위에 동심원처럼 퍼져 있는 특정 궤도들을 도는 전자들이 있다. 이 궤도마다 전자가 갖는 에너지 값이 정해져 있는데, 이 값을 에너지 준위라고 한다. 원자가 가진 특정한 몇 개의 궤도에만 전자가 존재하므로, 전자가 가질 수 있는 에너지 준위도 특정한 값으로 정해져 있다. 전자는 에너지를 얻으면 바깥쪽 궤도로, 에너지를 잃으면 안쪽 궤도로 이동한다. 원자가 하나일 때는 전자의 에너지 준위가 특정한 값이지만, 여러 원자가 모이면 여러 원자들의 궤도가 ⓐ중첩되면서 전자들이 서로 영향을 미쳐 에너지 준위에 변화가 생긴다. 이로 인해 에너지 준위는 특정한 값이 아닌 일정한 범위로 나타나게 된다. 이렇게 에너지 값이 범위로 표현된 것을 에너지 밴드라고 한다. 보통의 물질들은 수많은 원자로 구성돼 있기 때문에 대부분 에너지 밴드를 갖고 있다. 전자는 낮은 에너지 상태가 되려는 성질이 있으므로, 발광 물질의 전자에 에너지가 공급되면 전자가 바깥쪽으로 이동한 뒤 다시

(중략)
}
\begin{kfQBlock}{04}{객관식}
\kfQStem{윗글의 Ⓐ와 <보기>의 Ⓑ를 비교하여 이해한 내용으로 적절하지 \underline{않은} 것은?}
\begin{kfEvidence}태양은 파장이 250\textasciitilde{}2,500nm에 이르는 다양한 빛을 지상으로 보내는데, 실리콘 기반 태양 전지는 이 중 500\textasciitilde{}1,000nm의 빛만 활용할 수 있다. 파장이 1,000nm가 넘는 빛은 이 전지를 통과하고, 500nm 이하의 빛은 흡수되지만 열로 전환돼 날아가 버린다. 이 전지에서 낮은 파장의 빛이 대부분 열로 전환되는 이유는 어떤 물질이 매우 높은 에너지를 받았을 때, 물질이 넓은 에너지 밴드를 갖고 있다면 에너지를 받은 전자가 밴드 내에서 이동하면서 열에너지가 방출되기 때문이다. 하지만 물질이 특정한 에너지 준위를 갖고 있다면 상대적으로 적은 열을 발생하며 전자가 이동하여 전류가 발생한다. 이러한 장점을 가진 양자점은 크기를 조절하면 지상에 도달하는 태양 스펙트럼 중 대부분의 영역을 활용할 수 있기 때문에 Ⓑ양자점 기반 태양 전지가 개발되고 있다. 한편, 양자점의 코어에 셸을 씌우면 빛과 달리 전류는 셸을 뚫고 나오기 힘들므로 양자점 기반 태양 전지에서는 셸을 사용하지 않는 경우가 많다. 따라서 코어 자체를 더 안정한 소재로 만들거나 리간드를 더 붙여 표면을 안정시켜야 한다.\end{kfEvidence}
\begin{kfChoices}
\kfChoice 같은 양의 에너지가 공급될 경우, Ⓐ는 LCD 디스플레이보다, Ⓑ는 실리콘 기반 태양 전지보다 더 많은 열에너지를 방출하겠군.
\kfChoice Ⓐ와 Ⓑ 모두 리간드를 사용하여 입자의 표면을 안정화할 수 있군.
\kfChoice Ⓐ와 Ⓑ 모두 양자점의 전자가 넓은 에너지 밴드가 아닌 매우 좁은 범위의 에너지 준위를 갖고 있다는 점을 활용하는군.
\end{kfChoices}
\end{kfQBlock}

\kfPassageNote{문항 5 / 인과 도치}{%
1908년 조지 하디와 빌헬름 바인베르크는 멘델 유전 법칙에 기초한 수학적 모델을 각각 독립적으로 제시했다. 하디는 집단 내 대립 유전자 A와 a의 빈도를 p와 q로 두고 무작위 교배에서 다음 세대의 유전자형 빈도를 분석했다. 바인베르크도 무작위 결합 조건에서 유전자 빈도가 세대를 거쳐 안정될 수 있음을 통계 자료와 함께 제시했다. 이 두 접근이 결합되어 하디-바인베르크 평형 법칙이 정립되었다.

하디-바인베르크 평형은 ㉠집단이 충분히 크고, ㉡무작위 교배가 일어나며, ㉢새 돌연변이가 없고, ㉣이입·이출 같은 이주가 없고, ㉤자연 선택이 작용하지 않을 때 유지된다. [A] 예를 들어 p=0.7, q=0.3이면 p+q=1이고 유전자형 빈도는 AA=p²=0.49, Aa=2pq=0.42, aa=q²=0.09로 계산되며 p²+2pq+q²=1이 성립한다. 유전자 빈도가 유지되면 다음 세대에서도 유전자형 빈도는 일정하게 유지된다. [A]

(중략)
}
\begin{kfQBlock}{05}{객관식}
\kfQStem{㉠\textasciitilde{}㉤을 바탕으로 <보기>를 이해한 내용으로 적절하지 \underline{않은} 것은?}
\begin{kfEvidence}딱정벌레 집단은 충분히 크고 무작위 교배가 이루어졌으며 새로운 색 돌연변이는 없었다. 그러나 3세대에 외부 갈색 개체가 대량 유입되어 갈색 관련 대립 유전자 빈도가 증가했다.\end{kfEvidence}
\begin{kfChoices}
\kfChoice 외부 개체 유입은 곧 자연선택 작용과 동일하다.
\kfChoice 집단이 충분히 크다는 조건은 ㉠과 연결된다.
\kfChoice 무작위 교배 조건은 ㉡과 연결된다.
\end{kfChoices}
\end{kfQBlock}

\kfPassageNote{문항 6 / 의도·목적}{%
생명체의 유전 정보는 아데닌, 구아닌, 사이토신, 티민으로 구성된 디엔에이(DNA) 염기 서열에 저장되어 있다. 동일한 염기 서열을 가진 세포들이라도 유전자 발현 양상은 다를 수 있는데, 이러한 현상을 후성 유전으로 설명할 수 있다. 후성 유전이란 DNA 염기 서열의 변화 없이 화학적 변형을 통해 유전자 발현을 조절하는 과정으로, 대표적인 예로는 DNA 메틸화, 히스톤의 아세틸화 등이 있다.

DNA 메틸화는 유전자의 전사 조절 부위인 프로모터 내 사이토신에 메틸기가 부착되어 특정 유전자의 전사가 억제되는 과정이다. 이 과정은 주로 사이토신과 구아닌이 인접한 서열(CpG)에서 일어나며, DNA 메틸화 효소(DNMT)의 촉매 작용으로 촉진된다. 메틸화된 DNA는 전사 인자들의 결합을 방해하는 구조를 형성하여 유전자 발현을 차단하는 일종의 정지 신호 역할을 한다. 반면, 히스톤의 아세틸화는 전사를 촉진하는 화학적 변형이다. 히스톤은 DNA에 감겨 있는 단백질 구조체인데, 그 꼬리 부분의…(중략)
}
\begin{kfQBlock}{06}{객관식}
\kfQStem{㉠과 ㉡을 비교하여 이해한 내용으로 가장 적절한 것은?}
\begin{kfChoices}
\kfChoice ㉠은 분열 속도가 비정상적으로 빠른 세포를 직접 파괴하는 방식을 취하지만, ㉡은 유전자 발현을 조절하는 본래의 세포 프로그램을 정상으로 되돌리는 방식을 취한다.
\kfChoice ㉠은 분열 속도가 비정상적으로 빠른 세포를 직접 파괴하는 방식을 취하지만, ㉡은 유전자 발현을 조절하는 본래의 세포 프로그램을 정상으로 되돌리는 방식을 취한다.
\kfChoice ㉠은 세포 공격을 통해 암세포를 파괴하는 반면, ㉡은 종양 억제 유전자의 기능을 억제하는 방식으로 세포 자살을 유도한다.
\kfChoice ㉠은 돌연변이를 교정하여 유전자의 기능을 회복시키지만, ㉡은 유전자 발현 자체에는 영향을 미치지 않고 세포 증식만 억제한다.
\end{kfChoices}
\end{kfQBlock}

\kfPassageNote{문항 7 / 상관관계}{%
원자에는 중심에 핵이 있고, 핵 주위에 동심원처럼 퍼져 있는 특정 궤도들을 도는 전자들이 있다. 이 궤도마다 전자가 갖는 에너지 값이 정해져 있는데, 이 값을 에너지 준위라고 한다. 원자가 가진 특정한 몇 개의 궤도에만 전자가 존재하므로, 전자가 가질 수 있는 에너지 준위도 특정한 값으로 정해져 있다. 전자는 에너지를 얻으면 바깥쪽 궤도로, 에너지를 잃으면 안쪽 궤도로 이동한다. 원자가 하나일 때는 전자의 에너지 준위가 특정한 값이지만, 여러 원자가 모이면 여러 원자들의 궤도가 \uline{ⓐ중첩}되면서 전자들이 서로 영향을 미쳐 에너지 준위에 변화가 생긴다. 이로 인해 에너지 준위는 특정한 값이 아닌 일정한 범위로 나타나게 된다. 이렇게 에너지 값이 범위로 표현된 것을 에너지 밴드라고 한다. 보통의 물질들은 수많은 원자로 구성돼 있기 때문에 대부분 에너지 밴드를 갖고 있다. 전자는 낮은 에너지 상태가 되려는 성질이 있으므로, 발광 물질의 전자에 에너지가 공급되면 전자가 바깥쪽으로 이…(중략)
}
\begin{kfQBlock}{07}{객관식}
\kfQStem{<보기>를 바탕으로 판단할 때, 적절하지 \underline{않은} 것은?}
\begin{kfEvidence}같은 재료로 만든 양자점 C와 D가 있다. 같은 조건에서 C는 녹색, D는 청색을 보였다. 실험자는 '더 작은 입자는 에너지 변동 폭이 더 크다'는 원리를 이용해 두 입자의 상대적 특성을 비교하고자 한다.\end{kfEvidence}
\begin{kfChoices}
\kfChoice D가 청색을 보였다면 D의 에너지 변동 폭이 C보다 반드시 더 작다고 봐야 한다.
\kfChoice D가 청색을 보였다면 D의 에너지 변동 폭이 C보다 반드시 더 작다고 봐야 한다.
\kfChoice D가 청색을 보였다면 D가 C보다 더 작을 가능성이 크다.
\kfChoice C가 D보다 크다면 C의 에너지 변동 폭은 D보다 작을 가능성이 크다.
\end{kfChoices}
\end{kfQBlock}

\kfPassageNote{문항 8 / 필요·충분 조건}{%
생명체의 유전 정보는 아데닌, 구아닌, 사이토신, 티민으로 구성된 디엔에이(DNA) 염기 서열에 저장되어 있다. 동일한 염기 서열을 가진 세포들이라도 유전자 발현 양상은 다를 수 있는데, 이러한 현상을 후성 유전으로 설명할 수 있다. 후성 유전이란 DNA 염기 서열의 변화 없이 화학적 변형을 통해 유전자 발현을 조절하는 과정으로, 대표적인 예로는 DNA 메틸화, 히스톤의 아세틸화 등이 있다.

DNA 메틸화는 유전자의 전사 조절 부위인 프로모터 내 사이토신에 메틸기가 부착되어 특정 유전자의 전사가 억제되는 과정이다. 이 과정은 주로 사이토신과 구아닌이 인접한 서열(CpG)에서 일어나며, DNA 메틸화 효소(DNMT)의 촉매 작용으로 촉진된다. 메틸화된 DNA는 전사 인자들의 결합을 방해하는 구조를 형성하여 유전자 발현을 차단하는 일종의 정지 신호 역할을 한다. 반면, 히스톤의 아세틸화는 전사를 촉진하는 화학적 변형이다. 히스톤은 DNA에 감겨 있는 단백질 구조체인데, 그 꼬리 부분의…(중략)
}
\begin{kfQBlock}{08}{객관식}
\kfQStem{㉮의 전제로 가장 적절한 것은?}
\begin{kfChoices}
\kfChoice 후성 유전적 변화는 DNA 염기 서열의 유전적 변이가 아닌 메틸화나 아세틸화 같은 화학적 변형이 발생하거나 제거되는 과정이다.
\kfChoice 후성 유전적 변화는 DNA 염기 서열의 유전적 변이가 아닌 메틸화나 아세틸화 같은 화학적 변형이 발생하거나 제거되는 과정이다.
\kfChoice 후성 유전적 변화는 DNA 염기 서열이 일시적으로 변형될 뿐 영구히 변형되는 과정이 아니다.
\kfChoice 후성 유전적 변화는 우리 몸의 세포가 본래 가지고 있는 복구 효소에 의해 자연적으로 교정된다.
\end{kfChoices}
\end{kfQBlock}

\kfPassageNote{문항 9 / 결합·분리}{%
생명체의 유전 정보는 아데닌, 구아닌, 사이토신, 티민으로 구성된 디엔에이(DNA) 염기 서열에 저장되어 있다. 동일한 염기 서열을 가진 세포들이라도 유전자 발현 양상은 다를 수 있는데, 이러한 현상을 후성 유전으로 설명할 수 있다. 후성 유전이란 DNA 염기 서열의 변화 없이 화학적 변형을 통해 유전자 발현을 조절하는 과정으로, 대표적인 예로는 DNA 메틸화, 히스톤의 아세틸화 등이 있다.

DNA 메틸화는 유전자의 전사 조절 부위인 프로모터 내 사이토신에 메틸기가 부착되어 특정 유전자의 전사가 억제되는 과정이다. 이 과정은 주로 사이토신과 구아닌이 인접한 서열(CpG)에서 일어나며, DNA 메틸화 효소(DNMT)의 촉매 작용으로 촉진된다. 메틸화된 DNA는 전사 인자들의 결합을 방해하는 구조를 형성하여 유전자 발현을 차단하는 일종의 정지 신호 역할을 한다. 반면, 히스톤의 아세틸화는 전사를 촉진하는 화학적 변형이다. 히스톤은 DNA에 감겨 있는 단백질 구조체인데, 그 꼬리 부분의…(중략)
}
\begin{kfQBlock}{09}{객관식}
\kfQStem{윗글과 <보기>를 바탕으로 이해한 내용으로 적절하지 \underline{않은} 것은?}
\begin{kfEvidence}연구팀 D는 암세포에서 히스톤 메틸화를 조절하는 효소 EZH2가 과활성화되면 종양 억제 유전자 발현이 장기간 억제될 수 있다고 보고, EZH2 저해제를 투여했다. 그 결과 일부 세포에서 억제되어 있던 유전자 발현이 회복되고 세포 증식이 둔화되었다. 그러나 동시 측정에서 같은 세포군의 한 대립 유전자에는 비가역적 염기 서열 돌연변이가 확인되었다. 연구팀은 후성 유전 표지 교정만으로 설명되지 않는 잔여 위험을 고려해 병용 치료를 설계하고 있다.\end{kfEvidence}
\begin{kfChoices}
\kfChoice 한 대립 유전자의 비가역적 돌연변이가 확인되었더라도 후성 유전 표지만 교정하면 염기 서열까지 완전 복구되므로 병용 치료는 불필요하다.
\kfChoice 한 대립 유전자의 비가역적 돌연변이가 확인되었더라도 후성 유전 표지만 교정하면 염기 서열까지 완전 복구되므로 병용 치료는 불필요하다.
\kfChoice EZH2 저해제에 따른 발현 회복과 증식 둔화는 후성 유전 교정 전략의 가능성을 보여 준다.
\kfChoice 대립 유전자의 비가역적 돌연변이가 남아 있으면 후성 유전 약물 단독 효과가 제한될 수 있다.
\end{kfChoices}
\end{kfQBlock}

\kfPassageNote{문항 10 / 조건 왜곡}{%
육상 동물과 수생 동물은 각각의 생활 환경에 따라 호흡 과정에서 어려움을 겪는다. 육상 동물은 폐의 표면에 건조한 공기가 닿으면 폐가 건조해질 수 있는데, 이는 폐를 손상시키고 \uline{ⓐ탈수}를 일으키며 기체를 용해하는 폐의 능력을 떨어뜨린다. 한편 수생 동물은 가용 산소가 적은 수중 환경과 수온 차에 의해 변화되는 산소 가용성에 적응해야 한다. 또한 수생 동물은 물속에서 몸을 움직이는 데 많은 근육 운동이 필요하고, 찬 바닷물의 높은 비열 때문에 아가미의 열을 빼앗기며, 삼투압 현상으로 인해 수분 균형에 문제가 발생하기도 한다.

이러한 환경 속에서 수생 동물은 효과적인 호흡을 위해 아가미라는 특수한 기관을 사용한다. 아가미는 물속에 녹아 있는 산소를 흡수하고 체내에 있는 이산화 탄소를 배출하는 기능을 하는데, 외아가미와 내아가미로 나뉜다. 외아가미는 몸 밖으로 \uline{ⓑ돌출}된 형태로, 표면적이 넓고 형태가 다양하며 유생기 도롱뇽이나 일부 무척추동물에서 나타난다. 외아가미…(중략)
}
\begin{kfQBlock}{10}{객관식}
\kfQStem{윗글의 내용과 일치하는 것은?}
\begin{kfChoices}
\kfChoice 수생 동물은 가용 산소가 적은 수중 환경에 적응해야 한다.
\kfChoice 수생 동물은 가용 산소가 적은 수중 환경에 적응해야 한다.
\kfChoice 어류는 환수에 몸의 총사용 에너지 중 1\textasciitilde{}2\%만을 사용한다.
\kfChoice 외아가미는 체내에 위치해 덮개로 보호받는다.
\end{kfChoices}
\end{kfQBlock}

\kfPassageNote{문항 11 / 긍부정 반전}{%
[16\textasciitilde{}20] 다음 글을 읽고 물음에 답하시오.

(가)

헌 누더기 입은 무리가 남자인지 여자인지    어린 자식 등에 업고 자란 자식 손에 끌고    울면서 눈물 씻고 엎어지며 오는 모양    차마 보지 못할네라 나직이 묻는 말씀    어디로서 쫓아오며 어디로 가려는고    주려들 가는 사람인가 가게 되면 얻어 먹나    아무 데도 한가지라 날 따라 도로 가면    자네 원님 가서 보고 안접(安接)하게 하여줌세    겨우겨우 대답하되 우리 곳은 당진(唐津)이라    여러 해 흉년 들어 살길이 없는 중에    도망한 자 신구환(新舊還)을 있는 자에 물리니    제 것도 못 바치며 남의 곡식 어찌할꼬    못 바치면 매 맞으니 매 맞고 더욱 살까    정처 없이 가게 되면 죽을 줄 알건마는    아니 가고 어찌하리 굶고 맞고 죽을 지경    차라리 구렁*에나 염려 없이 뭇치이면   도리어 편할지라 이런 고로 가노메라    급히 급히 넘어가자 이 백성들 살려보세    둘째 령(嶺)을 …(중략)
}
\begin{kfQBlock}{11}{객관식}
\kfQStem{(나)에 대한 설명으로 적절하지 \underline{않은} 것은?}
\begin{kfChoices}
\kfChoice <2수>: 강산을 즐기느라 임금에게 가지 못하는 상황에 대한 미안함을 드러내고 있다.
\kfChoice <2수>: 강산을 즐기느라 임금에게 가지 못하는 상황에 대한 미안함을 드러내고 있다.
\kfChoice <1수>: 돌아온 고향에서 변해 버린 인사(人事)에 대한 슬픔을 나타내고 있다.
\kfChoice <4수>: 세속의 어지러운 소식을 모른 체하며 살고 싶은 심정을 표현하고 있다.
\end{kfChoices}
\end{kfQBlock}

\kfPassageNote{문항 12 / 비교 대상}{%
\#\#\# [22～27] 다음 글을 읽고 물음에 답하시오.

(가)   흰 벽에는 ――   어련히 해들 적마다 나뭇가지가 그림자 되어 떠오를 뿐이었다.   그러한 정밀*이 천년이나 머물렀다 한다.

단청은 연년(年年)이 빛을 잃어 두리기둥에는 틈이 생기고, 볕과 바람이 쓰라리게 스며들었다. 그러나 험상궂어 가는 것이 서럽지 않았다.

기왓장마다 푸른 이끼가 앉고 세월은 소리없이 쌓였으나 ㉠문은 상기 닫혀진 채 멀리 지나가는 바람 소리에 귀를 기울이는 밤이 있었다.

주춧돌 놓인 자리에 가을풀은 우거졌어도 봄이면 돋아나는 푸른 싹이 살고, 그리고 한 그루 진분홍 꽃이 피는 나무가 자랐다.

유달리도 푸른 높은 하늘을 눈물과 함께 아득히 흘러간 별들이 총총히 돌아오고 사납던 비바람이 걷힌 낡은 처마 끝에 찬란히 빛이 쏟아지는 새벽, 오래 닫혀진 문은 산천을 울리며 열리었다.

―― 그립던 깃발이 눈뿌리에 사무치는 푸른 하늘이었다.

* 김종길, 「문」 -   * 정밀 : 고요하고 편안함. …(중략)
}
\begin{kfQBlock}{12}{객관식}
\kfQStem{(가)～(다)에 대한 설명으로 가장 적절한 것은?}
\begin{kfChoices}
\kfChoice (가)는 동일한 색채어를, (나)는 유사한 문장 구조를 반복적으로 제시하며 시상을 전개한다.
\kfChoice (가)는 동일한 색채어를, (나)는 유사한 문장 구조를 반복적으로 제시하며 시상을 전개한다.
\kfChoice (가)는 명시적 청자에게 말을 건네는 방식으로 화자의 감정을 드러낸다.
\kfChoice (가)와 (나)는 모두, 사라져 가는 대상에 대한 화자의 안타까움을 드러낸다.
\end{kfChoices}
\end{kfQBlock}

\kfPassageNote{문항 13 / 주체·객체}{%
\#\#\# [22 \textasciitilde{} 25] 다음 글을 읽고 물음에 답하시오.

“김기사장의 인격이나 기술을 우리 사에서는 믿고 맡기고 있었소. 한 회사라는 것은 그 회사의 사업을 위주로 해서 사람을 쓴다는 것은 두말도 할 필요 없겠지요. 김기사는 우리 회사가 환도 후 재건에 있어서 가장 큰 공로를 세운 사람이라고 생각하고 있소. 그러나 금년 삼월에 들어 발전기와 제23호 인쇄기의 고장은 우리 회사의 치명적인 타격이었소. 이렇게 되면 회사에서는 그 기계를 다루는 기사장의 책임으로 돌리지 않을 수 없소. 기사라는 것은 기계의 고장을 사전에 발견하는 것이라고, 아니 고장이 나지 않게 하는 것이 제일 큰 직책이라고 회사에서는 생각하고 있소. 한두 번이 아닌 고장의 수리가 이삼일이 멀다 해서 또 고장 또 고장이라면 결국 기사는 고장의 원인을 모르는 것이라고 해석하지 않을 수 없겠지요. 그 책임을 기사장이 져야 한다면 현명한 기사장은 자기 처신을 어떻게 해야 한다는 것은 여기서 민망한 말을 하지 않아도 잘 이해…(중략)
}
\begin{kfQBlock}{13}{객관식}
\kfQStem{㉠과 ㉡에 대한 이해로 가장 적절한 것은?}
\begin{kfChoices}
\kfChoice ㉠은 김명학이 상대의 의중을 수용하는 공간이고, ㉡은 김명학이 상대의 반응에 개의치 않고 행동하는 공간이다.
\kfChoice ㉠은 김명학이 상대의 의중을 수용하는 공간이고, ㉡은 김명학이 상대의 반응에 개의치 않고 행동하는 공간이다.
\kfChoice ㉠은 김명학이 상대를 동정하는 공간이고, ㉡은 김명학이 상대에게 동정을 받는 공간이다.
\kfChoice ㉠은 김명학이 자신의 삶을 성찰하는 공간이고, ㉡은 김명학이 상대의 삶을 비판하는 공간이다.
\end{kfChoices}
\end{kfQBlock}

\kfPassageNote{문항 14 / 내용·방법}{%
\#\#\# [22 \textasciitilde{} 25] 다음 글을 읽고 물음에 답하시오.

“김기사장의 인격이나 기술을 우리 사에서는 믿고 맡기고 있었소. 한 회사라는 것은 그 회사의 사업을 위주로 해서 사람을 쓴다는 것은 두말도 할 필요 없겠지요. 김기사는 우리 회사가 환도 후 재건에 있어서 가장 큰 공로를 세운 사람이라고 생각하고 있소. 그러나 금년 삼월에 들어 발전기와 제23호 인쇄기의 고장은 우리 회사의 치명적인 타격이었소. 이렇게 되면 회사에서는 그 기계를 다루는 기사장의 책임으로 돌리지 않을 수 없소. 기사라는 것은 기계의 고장을 사전에 발견하는 것이라고, 아니 고장이 나지 않게 하는 것이 제일 큰 직책이라고 회사에서는 생각하고 있소. 한두 번이 아닌 고장의 수리가 이삼일이 멀다 해서 또 고장 또 고장이라면 결국 기사는 고장의 원인을 모르는 것이라고 해석하지 않을 수 없겠지요. 그 책임을 기사장이 져야 한다면 현명한 기사장은 자기 처신을 어떻게 해야 한다는 것은 여기서 민망한 말을 하지 않아도 잘 이해…(중략)
}
\begin{kfQBlock}{14}{객관식}
\kfQStem{윗글의 서술상의 특징으로 가장 적절한 것은?}
\begin{kfChoices}
\kfChoice 역순행적 구성을 활용하여 사건이 일어난 이유를 밝히고 있다.
\kfChoice 역순행적 구성을 활용하여 사건이 일어난 이유를 밝히고 있다.
\kfChoice 시대적 배경을 설명하여 갈등 해결의 실마리를 제공하고 있다.
\kfChoice 등장인물이 관찰자 입장에서 인물들의 말과 행동을 전달하고 있다.
\end{kfChoices}
\end{kfQBlock}

\kfPassageNote{문항 15 / 내용 왜곡}{%
\#\#\# [22 \textasciitilde{} 25] 다음 글을 읽고 물음에 답하시오.

“김기사장의 인격이나 기술을 우리 사에서는 믿고 맡기고 있었소. 한 회사라는 것은 그 회사의 사업을 위주로 해서 사람을 쓴다는 것은 두말도 할 필요 없겠지요. 김기사는 우리 회사가 환도 후 재건에 있어서 가장 큰 공로를 세운 사람이라고 생각하고 있소. 그러나 금년 삼월에 들어 발전기와 제23호 인쇄기의 고장은 우리 회사의 치명적인 타격이었소. 이렇게 되면 회사에서는 그 기계를 다루는 기사장의 책임으로 돌리지 않을 수 없소. 기사라는 것은 기계의 고장을 사전에 발견하는 것이라고, 아니 고장이 나지 않게 하는 것이 제일 큰 직책이라고 회사에서는 생각하고 있소. 한두 번이 아닌 고장의 수리가 이삼일이 멀다 해서 또 고장 또 고장이라면 결국 기사는 고장의 원인을 모르는 것이라고 해석하지 않을 수 없겠지요. 그 책임을 기사장이 져야 한다면 현명한 기사장은 자기 처신을 어떻게 해야 한다는 것은 여기서 민망한 말을 하지 않아도 잘 이해…(중략)
}
\begin{kfQBlock}{15}{객관식}
\kfQStem{윗글의 내용을 이해한 것으로 적절하지 \underline{않은} 것은?}
\begin{kfChoices}
\kfChoice 김명학은 사직원을 내고 나서 인쇄 공장에 들어갔을 때, 직공들이 인사하는 것을 보지 못하고 지나갔다.
\kfChoice 김명학은 사직원을 내고 나서 인쇄 공장에 들어갔을 때, 직공들이 인사하는 것을 보지 못하고 지나갔다.
\kfChoice 아내는 벽돌을 나르는 김명학의 모습을 보고 무슨 영문인지 몰랐다.
\kfChoice 김명학은 사장과 공장장에게 기계의 고장은 어쩔 수 없는 일이라고 변명했다.
\end{kfChoices}
\end{kfQBlock}

\kfPassageNote{문항 16 / 과잉 해석}{%
\#\#\# [22～27] 다음 글을 읽고 물음에 답하시오.

(가)   흰 벽에는 ――   어련히 해들 적마다 나뭇가지가 그림자 되어 떠오를 뿐이었다.   그러한 정밀*이 천년이나 머물렀다 한다.

단청은 연년(年年)이 빛을 잃어 두리기둥에는 틈이 생기고, 볕과 바람이 쓰라리게 스며들었다. 그러나 험상궂어 가는 것이 서럽지 않았다.

기왓장마다 푸른 이끼가 앉고 세월은 소리없이 쌓였으나 ㉠문은 상기 닫혀진 채 멀리 지나가는 바람 소리에 귀를 기울이는 밤이 있었다.

주춧돌 놓인 자리에 가을풀은 우거졌어도 봄이면 돋아나는 푸른 싹이 살고, 그리고 한 그루 진분홍 꽃이 피는 나무가 자랐다.

유달리도 푸른 높은 하늘을 눈물과 함께 아득히 흘러간 별들이 총총히 돌아오고 사납던 비바람이 걷힌 낡은 처마 끝에 찬란히 빛이 쏟아지는 새벽, 오래 닫혀진 문은 산천을 울리며 열리었다.

―― 그립던 깃발이 눈뿌리에 사무치는 푸른 하늘이었다.

* 김종길, 「문」 -   * 정밀 : 고요하고 편안함. …(중략)
}
\begin{kfQBlock}{16}{객관식}
\kfQStem{<보기>를 참고하여 (가)를 감상한 내용으로 적절하지 \underline{않은} 것은?}
\begin{kfChoices}
\kfChoice ‘흰 벽’에 나뭇가지가 그림자로 나타나는 것은, 천년을 쇠락해 온 인간의 역사가 자연의 힘을 탐색하는 과정에서 자연의 모습에 영향을 미친 결과를 보여 주는군.
\kfChoice ‘흰 벽’에 나뭇가지가 그림자로 나타나는 것은, 천년을 쇠락해 온 인간의 역사가 자연의 힘을 탐색하는 과정에서 자연의 모습에 영향을 미친 결과를 보여 주는군.
\kfChoice ‘두리기둥’의 틈에 볕과 바람이 쓰라리게 스며드는 것을 서럽지 않다고 한 것은, 쇠락해 가는 인간의 역사가 자연이 가진 변화의 힘을 수용함을 드러내는군.
\kfChoice ‘기왓장마다’ 이끼와 세월이 덮여 감에도 멀리 있는 바람 소리에 귀를 기울이는 것은, 자연의 영향을 받으면서도 자연이 가진 변화의 힘에서 생성의 가능성을 찾는 모습이겠군.
\end{kfChoices}
\end{kfQBlock}

\kfPassageNote{문항 17 / 범위 오류}{%
\#\#\# [25\textasciitilde{}27] 다음 글을 읽고 물음에 답하시오.

[앞부분의 줄거리] 영웅적 면모를 지녔으나 먹고 자기만 하며 허송세월을 보내던 이생은 장인이 죽자 장모에 의해 쫓겨나고, 이생의 부인 양 소저는 시부모의 제사를 지내면서 이생을 기다린다. 한편 집을 나온 이생은 곤경에 빠진 사람에게 자신이 가진 삼백 냥의 돈을 다 내어 준 후 어떤 노인의 도움을 받게 된다.

“그대 오늘 큰 적선을 하였으니 깊이 감격하노라.”   이생이 이렇듯 노인의 신기함을 보고 평범한 인물이 아니리라 생각하며 의아해 마지않았다.   “존옹의 물으심이 무슨 일이니이꼬? 저는 적선한 일이 없소이다.”   노인 말하기를,   “대인은 사람 속이기를 아니하나니라.”   이에 또 말하기를,   “그대 저리 먹는 양에 양식도 없이 어찌 살아가려 하느뇨?”   이생이 답하기를,   “이처럼 얻어먹으면 아니 살아가리이까?”   노인이 크게 웃으며 말하였다.   “㉠ 소년의 말이 사리에 어둡도다. 그러나 나는 그대를 알거…(중략)
}
\begin{kfQBlock}{17}{객관식}
\kfQStem{㉠∼㉣에 대한 설명으로 가장 적절한 것은?}
\begin{kfChoices}
\kfChoice ㉢에서는 과거의 일을 근거로 자책하고 있고, ㉣에서는 미래의 일을 예측하며 상대를 만류하고 있다.
\kfChoice ㉢에서는 과거의 일을 근거로 자책하고 있고, ㉣에서는 미래의 일을 예측하며 상대를 만류하고 있다.
\kfChoice ㉠에서는 당위를 내세워 상대를 설득하고, ㉡에서는 감정에 호소하여 상대의 동의를 이끌어 내고 있다.
\kfChoice ㉠에서는 상대보다 우월한 신분을 통해, ㉣에서는 상대에 대한 배려를 통해 상대의 태도 변화를 촉구하고 있다.
\end{kfChoices}
\end{kfQBlock}

\kfPassageNote{문항 18 / 시점·화자}{%
\#\#\# 2025년 11월 고3 수능

\#\#\# [22\textasciitilde{}26] 다음 글을 읽고 물음에 답하시오.

(가)

두고 온 것들이 빛나는 때가 있다   빛나는 때를 위해 소금을 뿌리며   우리는 이 저녁을 떠돌고 있는가   사방을 둘러보아도   등불 하나 켜 든 이 보이지 않고   등불 뒤에 속삭이며 밤을 지키는   발자국 소리 들리지 않는다   잊혀진 목소리가 살아나는 때가 있다   ㉠ 잊혀진 한 목소리 잊혀진 다른 목소리의 끝을 찾아   목메이게 부르짖다 잦아드는 때가 있다   ㉡ 잦아드는 외마디 소리를 찾아 칼날 세우고   우리는 이 새벽길 숨가쁘게 넘고 있는가   하늘 올려보아도   함께 어둠 지새던 별 하나 눈뜨지 않는다   그래도 두고 온 것들은 빛나는가   빛을 뿜으면서 한 번은 되살아나는가   우리가 뿌린 소금들 반짝반짝 별빛이 되어   오던 길 환히 비춰 주고 있으니

- 이시영, 「그리움」 -

(나)

감나무 잎새를 흔드는 게   어찌 바람뿐이랴.   감나무 잎새를 반짝이는…(중략)
}
\begin{kfQBlock}{18}{객관식}
\kfQStem{ⓐ\textasciitilde{}ⓕ를 중심으로 (나)를 이해한 내용으로 가장 적절한 것은?}
\begin{kfChoices}
\kfChoice 화자는 감나무 열매가 자라는 과정에서 ⓒ를 만나기도 하고 ⓓ를 만나기도 하는 일이 유의미하다고 여기고 있다.
\kfChoice 화자는 감나무 열매가 자라는 과정에서 ⓒ를 만나기도 하고 ⓓ를 만나기도 하는 일이 유의미하다고 여기고 있다.
\kfChoice 화자는 ⓐ가 흔드는 것이 감나무 잎새뿐이라고 여기다가 ⓑ를 보며 그 생각을 바로잡고 있다.
\kfChoice 화자는 ⓑ가 내는 소리와 ⓔ의 움직임을 통해 감나무 열매가 충분히 익은 상태임을 짐작하고 있다.
\end{kfChoices}
\end{kfQBlock}



\clearpage

\kfChapterCover{2}{비트겐슈타인1 ch02}{Chapter 2}{이번 챕터: 문법 → 문학 5작품 → 비문학 5지문 → 패턴 워크북.}
\begin{kfChapterAreaList}
\kfChapterAreaItem{문법}{kfAreaGrammar}{음운 / 음절의 끝소리 규칙}
\kfChapterAreaItem{문학}{kfAreaLiterature}{「슬픔이 기쁨에게」 — 정호승 외 4편}
\kfChapterAreaItem{비문학}{kfAreaNonlit}{「공리주의란 무엇인가」 외 4편}
\kfChapterAreaItem{실력 확인}{kfAreaWeekly}{패턴 워크북 — 총 18문항}
\end{kfChapterAreaList}

\kfStartGrammar{음운 / 음절의 끝소리 규칙}


\kfSecHeader{kfAreaGrammar}{문제}{}

\kfBeginMC
\par\needspace{25\baselineskip}
\kfPassageFull{대표음}{%
국어에서 음절의 끝(받침)에 올 수 있는 자음의 수는 제한되어 있다. 자음 19개 중 음절 끝에서 실제로 발음되는 소리는 'ㄱ, ㄴ, ㄷ, ㄹ, ㅁ, ㅂ, ㅇ'의 7개뿐이다. 이를 대표음이라 한다. 대표음이 아닌 자음이 음절 끝에 오면 7개의 대표음 중 하나로 바뀌어 발음되는데, 이를 '음절의 끝소리 규칙'이라 한다. 예를 들어, 'ㅋ, ㄲ'은 [ㄱ]으로, 'ㅌ, ㅈ, ㅊ, ㅅ, ㅆ, ㅎ'은 [ㄷ]으로, 'ㅍ'은 [ㅂ]으로 발음된다. 이 규칙은 국어의 음절 구조 제약에서 비롯된 것으로, 교체(음운의 개수가 변하지 않고 하나의 음운이 다른 음운으로 바뀌는 현상)에 해당한다.
}\par\nobreak\nopagebreak[4]\kfResetAreaFlag
\begin{kfQBlock}{01}{객관식}
\kfQStem{윗글을 바탕으로 할 때, 밑줄 친 단어의 발음이 옳은 것은?}
\begin{kfChoices}
\kfChoice 부엌 → [부억]
\kfChoice 꽃 → [꼿]
\kfChoice 낮 → [낟]
\kfChoice 잎 → [잎]
\kfChoice 옷 → [옫]
\end{kfChoices}
\end{kfQBlock}

\par\needspace{25\baselineskip}
\kfPassageFull{겹받침}{%
겹받침은 두 개의 자음이 받침에 함께 쓰이는 경우이다. 겹받침이 음절 끝에 올 때는 자음군 단순화가 일어나 두 자음 중 하나만 발음된다. 이때 어떤 자음이 남는지는 겹받침의 종류에 따라 다르다. 'ㄳ, ㄵ, ㄼ, ㄽ, ㄾ, ㅄ'은 앞의 자음이 발음되고, 'ㄺ, ㄻ, ㄿ'은 뒤의 자음이 발음된다. 다만 'ㄺ'의 경우 뒤에 'ㄱ'으로 시작하는 어미가 올 때는 [ㄹ]이 발음되기도 한다. 예를 들어 '읽다'는 [익따]로 발음되지만, '읽고'는 [일꼬]로 발음된다. 겹받침의 발음은 음운 변동의 기초가 되므로, 어떤 자음이 남는지를 정확히 알아야 한다.
}\par\nobreak\nopagebreak[4]\kfResetAreaFlag
\begin{kfQBlock}{02}{객관식}
\kfQStem{윗글을 바탕으로 할 때, 겹받침의 발음이 옳지 \underline{않은} 것은?}
\begin{kfChoices}
\kfChoice 넋 → [넉]
\kfChoice 삶 → [삼]
\kfChoice 값 → [갑]
\kfChoice 닭 → [달]
\kfChoice 앉다 → [안따]
\end{kfChoices}
\end{kfQBlock}

\par\needspace{25\baselineskip}
\kfPassageFull{발음 변화}{%
음절의 끝소리 규칙은 다양한 음운 변동의 출발점이 된다. 받침에 온 자음이 대표음으로 바뀐 뒤, 뒤에 오는 자음이나 모음의 영향으로 추가적인 음운 변동이 일어나기도 한다. 예를 들어 '꽃말'은 먼저 'ㅊ'이 대표음 [ㄷ]으로 바뀌고(음절의 끝소리 규칙), 이어서 [ㄷ]이 뒤의 비음 'ㅁ' 앞에서 [ㄴ]으로 바뀌어(비음화) 최종적으로 [꼰말]이 된다. 이처럼 한 단어에 두 가지 이상의 음운 변동이 순차적으로 적용되는 경우가 많다. 음운 변동의 적용 순서는 보통 '음절의 끝소리 규칙(또는 자음군 단순화) → 자음 동화(비음화, 유음화 등) → 된소리되기' 순이다.
}\par\nobreak\nopagebreak[4]\kfResetAreaFlag
\begin{kfQBlock}{03}{객관식}
\kfQStem{윗글을 바탕으로 할 때, <보기>의 단어에 적용된 음운 변동의 순서로 적절한 것은?}
\begin{kfEvidence}'놓는' → [논는]\end{kfEvidence}
\begin{kfChoices}
\kfChoice 비음화 → 음절의 끝소리 규칙
\kfChoice 음절의 끝소리 규칙 → 비음화
\kfChoice 유음화 → 음절의 끝소리 규칙
\kfChoice 된소리되기 → 비음화
\kfChoice 음절의 끝소리 규칙만 적용된다.
\end{kfChoices}
\end{kfQBlock}

\par\needspace{25\baselineskip}
\kfPassageFull{연음}{%
연음(이어 나는 소리)이란 앞 음절의 받침이 뒤 음절의 첫소리로 이어져 발음되는 현상이다. 연음은 받침 뒤에 모음으로 시작하는 형태소가 올 때 일어나는데, 뒤에 오는 형태소가 형식 형태소(조사, 접사, 어미)인지 실질 형태소(자립적 의미를 가진 형태소)인지에 따라 적용 방식이 달라진다. 형식 형태소가 뒤에 올 때는 연음이 바로 적용되어 받침소리가 그대로 뒤 음절로 넘어간다. 예를 들어 '꽃이'는 [꼬치]로 발음된다. 그러나 실질 형태소가 뒤에 올 때는 음절의 끝소리 규칙이 먼저 적용된 후 연음이 일어난다. 예를 들어 '꽃 아래'는 [꼳 아래 → 꼬다래]로 발음된다.
}\par\nobreak\nopagebreak[4]\kfResetAreaFlag
\begin{kfQBlock}{04}{객관식}
\kfQStem{윗글을 바탕으로 할 때, <보기>의 밑줄 친 부분의 발음으로 적절한 것은?}
\begin{kfEvidence}(가) 옷을 입다 \\\relax
(나) 옷 안에 넣다 \\\relax
(다) 밭이 넓다\end{kfEvidence}
\begin{kfChoices}
\kfChoice (가)의 '옷을'은 [오들]로 발음된다.
\kfChoice (가)의 '옷을'은 [오슬]로 발음된다.
\kfChoice (나)의 '옷 안'은 [오산]으로 발음된다.
\kfChoice (다)의 '밭이'는 [바디]로 발음된다.
\kfChoice (나)의 '옷 안'은 [오싼]으로 발음된다.
\end{kfChoices}
\end{kfQBlock}

\par\needspace{25\baselineskip}
\kfPassageFull{끝소리+조사}{%
음절의 끝소리 규칙과 연음의 관계를 정확히 이해하려면, 받침 뒤에 오는 형태소의 종류를 구별할 수 있어야 한다. 조사는 대표적인 형식 형태소로, 조사가 뒤에 올 때는 연음이 바로 적용된다. 그러나 겹받침 'ㄳ, ㅄ, ㄽ'의 'ㅅ'이 뒤 음절로 연음될 때는 된소리 [ㅆ]으로 발음된다는 점을 주의해야 한다. 예를 들어 '넋이'는 [넉씨], '값을'은 [갑쓸], '외곬으로'는 [외골쓰로]로 발음된다. 이는 겹받침에서 뒤의 'ㅅ'이 연음될 때 된소리로 바뀌는 특수한 규칙이다.
}\par\nobreak\nopagebreak[4]\kfResetAreaFlag
\begin{kfQBlock}{05}{객관식}
\kfQStem{윗글을 바탕으로 할 때, <보기>의 발음으로 옳은 것만을 모두 고른 것은?}
\begin{kfEvidence}ㄱ. 넋이 → [넉씨] \\\relax
ㄴ. 값을 → [가블] \\\relax
ㄷ. 삶을 → [사믈] \\\relax
ㄹ. 없어 → [업써]\end{kfEvidence}
\begin{kfChoices}
\kfChoice ㄱ, ㄴ
\kfChoice ㄱ, ㄷ
\kfChoice ㄱ, ㄹ
\kfChoice ㄱ, ㄷ, ㄹ
\kfChoice ㄴ, ㄷ, ㄹ
\end{kfChoices}
\end{kfQBlock}

\kfEndMC


\kfStartLit{「슬픔이 기쁨에게」 — 정호승 외 4편}


\kfPassageNote{「슬픔이 기쁨에게」 — 정호승}{%
나는 이제 너에게도 슬픔을 주겠다. \\[4pt]
사랑보다 소중한 슬픔을 주겠다. \\[4pt]
겨울밤 거리에서 귤 껍질을 까며 \\[4pt]
기다리던 , 추위 속에서 싹트던 \\[4pt]
내 슬픔의 뿌리를 주겠다. \\[14pt]
나는 이제 너에게도 기다림을 주겠다. \\[4pt]
이 세상에 내리던 함박눈을 주겠다. \\[4pt]
하얀 함박눈의 메시지를 주겠다. \\[14pt]
눈이 오는 먼 길 끝에 \\[4pt]
눈이 오는 먼 마을에 \\[4pt]
귤 껍질 끝에 달린 \\[4pt]
내 사랑의 전부를 주겠다. \\[14pt]
슬픔의 힘에 대해서 \\[4pt]
나는 이제 너에게 이야기하겠다. \\[4pt]
나를 가만히 안아 다오.
}\kfResetAreaFlag
\par\kfMaybeNeedspace{50\baselineskip}\noindent{\kfTipMark\,\,}{\small\color{kfInk} 빈칸에 알맞은 말을 채워 넣으시오. (초성 힌트 참고)}\par\nopagebreak[4]\smallskip\nopagebreak[4]
\begin{kfActTable}{|>{\centering\arraybackslash}m{2.6cm}|m{6cm}|m{4cm}|}
\rowcolor{kfPrimaryLight!50}\textbf{\color{kfPrimary}항목} & \textbf{\color{kfPrimary}내용 (빈칸 채우기)} & \textbf{\color{kfPrimary}답안 작성}\\\hline
갈래 & ㅈㅇㅅ, ㅅㅈㅅ ① & \kfTBlank \\\hline
성격 & ㅇㅅㅈ, ㅅㅊㅈ ② & \kfTBlank \\\hline
화자 & ㅅㅍ 자체를 의인화한 존재 ③ & \kfTBlank \\\hline
청자 & ㄱㅃ ④ & \kfTBlank \\\hline
주요 소재 & ㄱ ㄲㅈ, ㅎㅂㄴ, ㅃㄹ ⑤ & \kfTBlank \\\hline
핵심 표현 & 'ㅅㄹ보다 소중한 ㅅㅍ' ⑥ & \kfTBlank \\\hline
표현 기법 1 & 슬픔에 인격을 부여한 ㅇㅇㅎ ⑦ & \kfTBlank \\\hline
표현 기법 2 & 통념을 뒤집는 ㅇㅅㅂ ⑧ & \kfTBlank \\\hline
표현 기법 3 & '\textasciitilde{}주겠다'의 ㅂㅂㅂ ⑨ & \kfTBlank \\\hline
계절 배경 & ㄱㅇ (겨울밤, 추위, 함박눈) ⑩ & \kfTBlank \\\hline
주제 & ㅅㅍ의 ㄱㅊ와 ㄱㅃ과의 ㄱㅈ ⑪ & \kfTBlank \\\hline
마지막 행 & '나를 가만히 ㅇㅇ 다오' ⑫ & \kfTBlank \\\hline
\end{kfActTable}
\kfSecHeader{kfAreaLiterature}{문제}{}

\kfBeginMC
\begin{kfQBlock}{01}{객관식}
\kfQStem{윗시에 대한 설명으로 적절하지 \underline{않은} 것은?}
\begin{kfChoices}
\kfChoice 슬픔이라는 추상적 감정을 인격화하여 화자로 설정하고 있다.
\kfChoice '\textasciitilde{}주겠다'의 반복을 통해 화자의 의지를 강조하고 있다.
\kfChoice '겨울밤', '추위', '함박눈' 등 계절적 이미지를 활용하고 있다.
\kfChoice 슬픔을 극복하고 벗어나야 할 부정적 대상으로 그리고 있다.
\kfChoice '귤 껍질', '뿌리' 등 구체적 사물을 통해 추상적 감정을 형상화하고 있다.
\end{kfChoices}
\end{kfQBlock}

\begin{kfQBlock}{02}{객관식}
\kfQStem{'사랑보다 소중한 슬픔을 주겠다'에 사용된 표현 기법과 그 효과를 가장 적절하게 설명한 것은?}
\begin{kfChoices}
\kfChoice 점층법을 사용하여 사랑에서 슬픔으로 감정이 고조됨을 보여준다.
\kfChoice 역설을 통해 슬픔이 사랑보다 소중할 수 있다는 의외의 인식을 드러낸다.
\kfChoice 대구법을 사용하여 사랑과 슬픔의 대등한 가치를 병렬적으로 제시한다.
\kfChoice 반어법을 통해 실제로는 슬픔을 주고 싶지 않다는 마음을 표현한다.
\kfChoice 과장법을 사용하여 슬픔의 크기를 실제보다 부풀려 표현한다.
\end{kfChoices}
\end{kfQBlock}

\begin{kfQBlock}{03}{객관식}
\kfQStem{화자가 '너'에게 '슬픔'을 주겠다고 말하는 의도로 가장 적절한 것은?}
\begin{kfChoices}
\kfChoice 상대방에게 고통을 안겨 자신의 아픔에 대한 보복을 하려는 것이다.
\kfChoice 상대방이 슬픔을 통해 삶의 깊이와 성숙에 이르기를 바라는 것이다.
\kfChoice 슬픔을 나눔으로써 자신의 슬픔을 덜어 내려는 것이다.
\kfChoice 상대방과의 관계를 단절하겠다는 의지를 표현하는 것이다.
\kfChoice 슬픔을 느끼지 못하는 상대방의 냉정함을 비판하려는 것이다.
\end{kfChoices}
\end{kfQBlock}

\begin{kfQBlock}{04}{객관식}
\kfQStem{윗시의 마지막 행 '나를 가만히 안아 다오'가 지닌 의미로 가장 적절한 것은?}
\begin{kfChoices}
\kfChoice 슬픔이 기쁨에게 자신을 받아들여 달라고 요청하며 공존을 호소한다.
\kfChoice 화자가 연인에게 따뜻한 물리적 위로를 구하며 외로움을 해소하려 한다.
\kfChoice 슬픔이 기쁨을 압도하겠다는 의지를 완곡한 어조로 드러내고 있다.
\kfChoice 화자가 차가운 세상과의 단절을 선언하며 고독 속으로 숨고자 한다.
\kfChoice 기쁨이 자신의 무관심을 뉘우치며 슬픔에게 용서를 구하고 있다.
\end{kfChoices}
\end{kfQBlock}

\begin{kfQBlock}{05}{객관식}
\kfQStem{윗시를 바탕으로 <보기>를 감상한 내용으로 적절하지 \underline{않은} 것은?}
\begin{kfEvidence}정호승의 시에서 슬픔은 삶의 본질적 조건으로 제시된다. 그의 시적 세계에서 슬픔은 회피해야 할 부정적 감정이 아니라, 인간을 성숙하게 하고 타인에 대한 공감 능력을 키워 주는 긍정적 힘으로 작용한다. 이러한 인식은 한국 현대시의 '한(恨)' 정서와도 맞닿아 있으며, 고통을 삶의 일부로 수용하는 동양적 사유를 반영한다.\end{kfEvidence}
\begin{kfChoices}
\kfChoice '사랑보다 소중한 슬픔'이라는 표현은 슬픔을 삶의 본질적 조건으로 보는 화자의 인식이 반영된 것이겠군.
\kfChoice '슬픔의 뿌리를 주겠다'에서 '뿌리'는 슬픔이 표면적 감정이 아니라 삶의 깊은 곳에서 자라난 것임을 보여주겠군.
\kfChoice '함박눈의 메시지'는 슬픔 속에서도 아름다움과 위로가 있음을 암시하여, 슬픔의 긍정적 힘과 연결되겠군.
\kfChoice '귤 껍질을 까며 기다리던' 행위는 서양 철학의 실존적 불안을 상징하는 것이겠군.
\kfChoice '나를 가만히 안아 다오'는 고통을 삶의 일부로 수용해 달라는 동양적 사유와 맞닿아 있겠군.
\end{kfChoices}
\end{kfQBlock}

\kfEndMC

\kfPassageNote{「허생전」 — 박지원}{%
허생은 묵묵히 앉아서 글을 읽기만 하였다. 그의 아내가 삯바느질로 살림을 꾸려 가다가 하루는 배를 주리며 말하였다. \\[4pt]
"당신은 평생 과거(科擧)를 보지 않으니, 글을 읽어서 무엇에 쓸 거예요?" \\[4pt]
허생이 웃으며 대답하였다. \\[4pt]
"나는 아직 글을 읽고 있는 중이오." \\[4pt]
"그러면 장인바치 일이라도 못 하시나요?" \\[4pt]
"장인바치 일은 배운 적이 없으니 못하지." \\[4pt]
"그러면 장사라도 해 보시지 않겠어요?" \\[4pt]
"장사는 밑천이 없으니 못하지." \\[4pt]
아내가 성이 나서 소리를 질렀다. \\[4pt]
"밤낮 글만 읽더니 기껏 '못' 한다는 소리만 배웠소? 당신은 도무지 어떻게 살자는 거예요?" \\[4pt]
허생은 읽던 책을 덮고 일어나 밖으로 나갔다. \\[14pt]
(중략) \\[14pt]
허생이 변씨를 찾아갔다. \\[4pt]
"내가 집이 가난하여 장사를 좀 하려고 하는데, 만 냥을 꾸어 주시오." \\[4pt]
변씨가 허생의 풍채를 보고 만 냥을 선뜻 내주었다. \\[14pt]
(중략) \\[14pt]
허생은 안성에 이르러 그곳 과일을 몽땅 사들였다. 허생이 과일을 독점하니 온 나라 과일 장수가 허생에게 달려와 비싼 값에 과일을 사 갔다. 허생은 열 배의 이익을 남겼다. \\[14pt]
(중략) \\[14pt]
허생이 이완 대장에게 말하였다. \\[4pt]
"내가 세 가지 계책을 내겠으니 들어 보시오. 첫째, 지금 명나라 장졸의 자손이 중국에서 귀화하여 오는데, 이들을 잘 대접하여 청나라의 허실을 정탐하게 하시오." \\[4pt]
이완이 말하였다. \\[4pt]
"그것은 어렵소." \\[4pt]
"둘째, 조선의 아름다운 여자를 뽑아 청나라 귀족에게 시집보내어 안에서 정보를 얻게 하시오." \\[4pt]
"그것도 어렵소." \\[4pt]
"셋째, 종실과 재상의 자제를 청나라에 볼모로 보내어 신의를 보이고 안으로 대비하시오." \\[4pt]
"그것도 어렵소." \\[4pt]
허생이 크게 꾸짖었다. \\[4pt]
"세 가지 계책을 다 어렵다고 하니, 나라에 아무 인물이 없구나!"
}\kfResetAreaFlag
\par\kfMaybeNeedspace{50\baselineskip}\noindent{\kfTipMark\,\,}{\small\color{kfInk} 빈칸에 알맞은 말을 채워 넣으시오. (초성 힌트 참고)}\par\nopagebreak[4]\smallskip\nopagebreak[4]
\begin{kfActTable}{|>{\centering\arraybackslash}m{2.6cm}|m{6cm}|m{4cm}|}
\rowcolor{kfPrimaryLight!50}\textbf{\color{kfPrimary}항목} & \textbf{\color{kfPrimary}내용 (빈칸 채우기)} & \textbf{\color{kfPrimary}답안 작성}\\\hline
갈래 & ㅎㅁ ㄷㅍ ㅅㅅ ① & \kfTBlank \\\hline
성격 & ㅍㅈㅈ, ㅂㅍㅈ ② & \kfTBlank \\\hline
시대 배경 & ㅈㅅ ㅎㄱ (17세기) ③ & \kfTBlank \\\hline
사상적 배경 & ㅂㅎㅍ 실학 사상 ④ & \kfTBlank \\\hline
주인공 & 남산 목릉동에서 글만 읽던 ㅅㅂ ⑤ & \kfTBlank \\\hline
핵심 사건 1 & ㅂㅆ에게 ㅁ ㄴ을 빌림 ⑥ & \kfTBlank \\\hline
핵심 사건 2 & 안성에서 과일을 ㅁㅈ하여 열 배 이익 ⑦ & \kfTBlank \\\hline
핵심 사건 3 & 이완 대장에게 ㅅ ㄱㅈ ㄱㅊ 제안 → 거절 ⑧ & \kfTBlank \\\hline
비판 대상 1 & ㅇㅌ ㄱㅈ의 취약함 ⑨ & \kfTBlank \\\hline
비판 대상 2 & ㅈㅂㅊ의 무능과 ㅇㅅ ⑩ & \kfTBlank \\\hline
표현 기법 & 기이한 인물을 통한 ㅍㅈ ⑪ & \kfTBlank \\\hline
출전 & 『ㅇㅇㅈ』 「ㅇㄱㅇㅎ」 ⑫ & \kfTBlank \\\hline
\end{kfActTable}
\kfSecHeader{kfAreaLiterature}{문제}{}

\kfBeginMC
\begin{kfQBlock}{01}{객관식}
\kfQStem{윗글에 대한 설명으로 적절하지 \underline{않은} 것은?}
\begin{kfChoices}
\kfChoice 대화를 통해 인물 간의 갈등과 성격을 드러내고 있다.
\kfChoice 허생이 물품을 독점 매입하여 이익을 얻는 과정이 나타나 있다.
\kfChoice 허생이 이완 대장에게 제안한 계책이 모두 수용되어 실행에 옮겨진다.
\kfChoice 현실의 모순을 기이한 인물의 행동을 통해 풍자하고 있다.
\kfChoice 아내와의 대화에서 반복적인 문답 구조가 활용되고 있다.
\end{kfChoices}
\end{kfQBlock}

\begin{kfQBlock}{02}{객관식}
\kfQStem{허생이 안성에서 과일을 '몽땅 사들'이는 행위가 작품에서 지닌 기능으로 가장 적절한 것은?}
\begin{kfChoices}
\kfChoice 허생의 사치스러운 소비 성향을 드러내기 위한 것이다.
\kfChoice 조선 시대 유통 구조의 취약함을 간접적으로 비판하기 위한 것이다.
\kfChoice 과일 재배 농민을 구제하려는 허생의 인도주의적 면모를 부각하기 위한 것이다.
\kfChoice 변씨에게 빌린 돈을 빨리 갚기 위한 단순한 경제 행위이다.
\kfChoice 허생이 상인 계층의 일원으로 편입되는 과정을 보여주기 위한 것이다.
\end{kfChoices}
\end{kfQBlock}

\begin{kfQBlock}{03}{객관식}
\kfQStem{허생의 아내가 '밤낮 글만 읽더니 기껏 못 한다는 소리만 배웠소?'라고 한 말의 의미로 가장 적절한 것은?}
\begin{kfChoices}
\kfChoice 남편의 학문이 실생활에 아무 쓸모가 없다는 비판이다.
\kfChoice 남편이 과거 시험을 거듭 보고도 떨어진 것에 대한 원망의 표현이다.
\kfChoice 남편이 글을 읽는 속도가 다른 선비보다 너무 느리다는 불만이다.
\kfChoice 남편이 장인바치 기술을 익힐 만한 손재주가 없다는 평가이다.
\kfChoice 남편의 학문적 깊이가 다른 선비들보다 한참 뒤떨어진다는 한탄이다.
\end{kfChoices}
\end{kfQBlock}

\begin{kfQBlock}{04}{객관식}
\kfQStem{허생이 이완 대장에게 세 가지 계책을 제안하고 모두 거절당한 뒤 '나라에 아무 인물이 없구나!'라고 꾸짖는 대목에서 드러나는 작가의 의도로 가장 적절한 것은?}
\begin{kfChoices}
\kfChoice 이완 개인의 군사적 무능을 고발하려는 것이다.
\kfChoice 허생의 과대망상적 성격이 현실에 부딪히는 장면을 보여주려는 것이다.
\kfChoice 조선 지배층의 무능과 개혁 의지 부재를 비판하려는 것이다.
\kfChoice 청나라에 대한 사대 정책의 정당성을 옹호하려는 것이다.
\kfChoice 허생에게 정치에서 손을 떼라는 교훈을 전하려는 것이다.
\end{kfChoices}
\end{kfQBlock}

\begin{kfQBlock}{05}{객관식}
\kfQStem{윗글을 바탕으로 <보기>를 감상한 내용으로 적절하지 \underline{않은} 것은?}
\begin{kfEvidence}박지원은 북학파(北學派)의 대표적 실학자로, 청나라의 선진 문물을 적극 수용하여 조선 사회를 개혁해야 한다고 주장하였다. 그의 소설에는 당시 양반 사회의 허위의식과 경제 구조의 모순에 대한 날카로운 비판이 담겨 있으며, 파격적 인물 설정과 풍자적 기법이 특징적이다.\end{kfEvidence}
\begin{kfChoices}
\kfChoice 허생이 선비이면서도 장사에 성공하는 설정은 양반이 상업을 천시하던 당대의 통념을 뒤집는 파격적 인물 설정이겠군.
\kfChoice 허생의 매점매석으로 시장이 쉽게 교란되는 것은 조선 경제 구조의 모순을 드러내는 풍자적 장치이겠군.
\kfChoice 이완이 세 계책을 모두 거절하는 것은 <보기>에서 말한 양반 사회의 개혁 의지 부재를 보여주겠군.
\kfChoice 허생이 변씨에게 만 냥을 빌리는 것은 당시 금융 제도가 완비되어 있었음을 보여주는 근거이겠군.
\kfChoice 아내의 비판에 허생이 바로 밖으로 나가는 것은 실천 없는 학문에 대한 반성이 계기가 된 것이겠군.
\end{kfChoices}
\end{kfQBlock}

\kfEndMC

\kfPassageNote{「아홉 켤레의 구두로 남은 사내」 — 윤흥길}{%
권씨는 새벽부터 밤늦게까지 구두를 수선했다. 쉬는 날이라고는 없었다. 비가 오나 눈이 오나 그는 길모퉁이 좌판 앞에 쪼그리고 앉아 구두를 만지작거렸다. 얼핏 보기에도 그의 손은 가죽과 구분이 안 될 만큼 딱딱하게 굳어 있었다.

(중략)

"권씨, 일요일인데 좀 쉬지 그래요." 내가 가끔 이렇게 말하면 그는 멋쩍게 웃었다. "하루만 놀아도 쌀이 떨어져요. 아홉 식구가 굶게 생겼는데 어떻게 놀겠어요."

(중략)

권씨의 아내가 찾아온 것은 그가 죽고 나서 사흘 뒤였다. 그녀는 권씨가 쓰러진 자리 밑에서 구두 아홉 켤레를 발견했다고 했다. 가족 한 사람에 한 켤레씩이었다. 다 만들지 못한 것도 있었고, 거의 완성된 것도 있었다. 큰아이 것은 거의 다 되어 있었지만, 막내 것은 한쪽만 만들다 만 상태였다.

"이 사람이 언제 이걸 만들었는지 모르겠어요. 밤에 잠도 안 자고 만들었나 봐요."

그녀는 그 구두들을 안고 울었다. 나도 코끝이 찡했다. 권씨는 아홉 켤레의 구두로 남았다. 아니, 그의 삶 전체가 바로 그 구두들이었다.

나는 권씨를 위해 아무것도 해 준 것이 없다는 사실이 문득 부끄러웠다. 가난한 이웃의 고통을 곁에서 지켜보면서도 말 한마디 제대로 건네지 못한 나 자신이 한심스러웠다.
}\kfResetAreaFlag
\par\kfMaybeNeedspace{50\baselineskip}\noindent{\kfTipMark\,\,}{\small\color{kfInk} 빈칸에 알맞은 말을 채워 넣으시오. (초성 힌트 참고)}\par\nopagebreak[4]\smallskip\nopagebreak[4]
\begin{kfActTable}{|>{\centering\arraybackslash}m{2.6cm}|m{6cm}|m{4cm}|}
\rowcolor{kfPrimaryLight!50}\textbf{\color{kfPrimary}항목} & \textbf{\color{kfPrimary}내용 (빈칸 채우기)} & \textbf{\color{kfPrimary}답안 작성}\\\hline
갈래 & 현대 ㄷㅍ ㅅㅅ ① & \kfTBlank \\\hline
성격 & ㅅㅅㅈ, ㅂㅍㅈ ② & \kfTBlank \\\hline
시대 배경 & 1970년대 ㅅㅇㅎ 시기 ③ & \kfTBlank \\\hline
시점 & ㅇㅇㅊ ㄱㅊㅈ 시점 ④ & \kfTBlank \\\hline
주인공 & 구두 수선공 ㄱㅆ ⑤ & \kfTBlank \\\hline
핵심 상징 & ㅇㅎ ㄱㄹ의 ㄱㄷ ⑥ & \kfTBlank \\\hline
구두의 의미 1 & 가족에 대한 말없는 ㅅㄹ의 표현 ⑦ & \kfTBlank \\\hline
구두의 의미 2 & ㄴㄷ으로만 환원된 비극적 삶 ⑧ & \kfTBlank \\\hline
권씨의 죽음 원인 & ㄱㄹ와 ㅇㅇ ㅅㅈ ⑨ & \kfTBlank \\\hline
화자의 역할 & 무력한 ㅈㅅㅇ 관찰자 ⑩ & \kfTBlank \\\hline
주제 & 산업화 시대 ㅂㅈㄱㅈ ㄴㄷㅈ의 ㅂㄱ ⑪ & \kfTBlank \\\hline
부가 주제 & ㅈㅅㅇ의 ㅁㄹㅎ과 ㅈㄱ ㅂㅅ ⑫ & \kfTBlank \\\hline
\end{kfActTable}
\kfSecHeader{kfAreaLiterature}{문제}{}

\kfBeginMC
\begin{kfQBlock}{01}{객관식}
\kfQStem{윗글에 대한 설명으로 적절하지 \underline{않은} 것은?}
\begin{kfChoices}
\kfChoice '나'라는 1인칭 관찰자 시점을 통해 권씨의 삶을 서술하고 있다.
\kfChoice 권씨의 굳어진 손 묘사를 통해 고된 노동의 흔적을 구체적으로 드러내고 있다.
\kfChoice 권씨는 가족에 대한 사랑을 말로 직접 표현하며 감정을 드러내고 있다.
\kfChoice 아홉 켤레의 구두가 핵심 소재로 작품의 제목이자 주제를 함축하고 있다.
\kfChoice 화자 '나'의 부끄러움을 통해 지식인의 무력함이라는 부가 주제가 드러나고 있다.
\end{kfChoices}
\end{kfQBlock}

\begin{kfQBlock}{02}{객관식}
\kfQStem{'아홉 켤레의 구두'가 지닌 상징적 의미로 가장 적절한 것은?}
\begin{kfChoices}
\kfChoice 권씨가 구두 장인으로서 이루고자 했던 평생의 직업적 성취의 결정체
\kfChoice 가족에 대한 말없는 사랑이자, 노동으로만 환원된 한 인간의 삶 전체
\kfChoice 산업화 시대에 대량 생산되는 공산품의 비인간적 속성
\kfChoice 아내에 대한 미안함을 물질로 보상하려는 속죄의 행위
\kfChoice 가난에서 벗어나기 위해 구두를 팔려던 미완의 사업 계획
\end{kfChoices}
\end{kfQBlock}

\begin{kfQBlock}{03}{객관식}
\kfQStem{권씨가 '하루만 놀아도 쌀이 떨어져요'라고 말한 것에서 알 수 있는 것으로 가장 적절한 것은?}
\begin{kfChoices}
\kfChoice 권씨는 자기 일을 즐기며 스스로 쉬지 않으려 하는 것이다.
\kfChoice 권씨의 가난은 하루 벌어 하루 먹는 수준으로 쉴 수 없는 상황이다.
\kfChoice 권씨는 저축을 전혀 하지 않는 비계획적 성격의 소유자이다.
\kfChoice 권씨의 아내가 생활비를 지나치게 헤프게 써서 가계가 어려운 것이다.
\kfChoice 권씨는 바쁜 일을 핑계 삼아 가족과의 관계를 의도적으로 회피하고 있다.
\end{kfChoices}
\end{kfQBlock}

\begin{kfQBlock}{04}{객관식}
\kfQStem{윗글의 마지막 부분에서 화자 '나'가 '부끄러웠다', '한심스러웠다'라고 느끼는 것이 작품 전체에서 지니는 의미로 가장 적절한 것은?}
\begin{kfChoices}
\kfChoice 권씨의 죽음을 막지 못한 직접적 책임을 통감하는 죄의식의 표현이다.
\kfChoice 이웃의 고통을 방관해 온 지식인의 무력함에 대한 자기 반성이다.
\kfChoice 권씨보다 경제적으로 넉넉한 자신의 상대적 풍요에 대한 미안함이다.
\kfChoice 권씨 가족을 더 적극적으로 도와주지 못한 것에 대한 개인적 후회이다.
\kfChoice 구두 수선공이라는 직업을 마음속으로 낮추어 본 데 대한 반성이다.
\end{kfChoices}
\end{kfQBlock}

\begin{kfQBlock}{05}{객관식}
\kfQStem{윗글을 바탕으로 <보기>를 감상한 내용으로 적절하지 \underline{않은} 것은?}
\begin{kfEvidence}1970년대 한국은 급격한 산업화가 진행되면서 경제 성장의 혜택이 고르게 분배되지 않았다. 저임금 노동자와 비정규직 근로자는 열악한 환경에서 장시간 노동을 해야 했으며, 이들의 희생 위에 경제 성장이 이루어졌다. 이 시기 리얼리즘 소설은 이러한 소외 계층의 삶을 사실적으로 그려내어 사회적 관심을 환기하는 역할을 하였다.\end{kfEvidence}
\begin{kfChoices}
\kfChoice 권씨가 비가 오나 눈이 오나 쉬지 않고 일하는 모습은 <보기>에서 말한 저임금 노동자의 열악한 근무 환경을 보여주겠군.
\kfChoice 아홉 식구를 먹여 살리기 위해 밤낮으로 일하다 죽은 권씨의 운명은 경제 성장의 혜택이 고르게 분배되지 않은 현실을 반영하겠군.
\kfChoice 권씨가 구두를 만들며 보여주는 장인 정신은 1970년대 한국 수출 산업의 기술력 향상에 기여한 사례로 볼 수 있겠군.
\kfChoice '나'의 부끄러움은 소외 계층의 고통에 무관심했던 중산층 지식인의 자기 반성으로 읽을 수 있겠군.
\kfChoice 권씨의 죽음과 남겨진 구두를 사실적으로 묘사한 것은 리얼리즘 소설의 특징과 연결되겠군.
\end{kfChoices}
\end{kfQBlock}

\kfEndMC

\kfPassageNote{「이 몸이 죽어 가서」 — 성삼문}{%
이 몸이 죽어 가서 무어이 될꼬 하니 봉래산(蓬萊山) 제일봉(第一峰)에 낙락장송(落落長松) 되었다가 백설(白雪)이 만건곤(滿乾坤)할 제 독야청청(獨也靑靑)하리라
}\kfResetAreaFlag
\par\kfMaybeNeedspace{50\baselineskip}\noindent{\kfTipMark\,\,}{\small\color{kfInk} 빈칸에 알맞은 말을 채워 넣으시오. (초성 힌트 참고)}\par\nopagebreak[4]\smallskip\nopagebreak[4]
\begin{kfActTable}{|>{\centering\arraybackslash}m{2.6cm}|m{6cm}|m{4cm}|}
\rowcolor{kfPrimaryLight!50}\textbf{\color{kfPrimary}항목} & \textbf{\color{kfPrimary}내용 (빈칸 채우기)} & \textbf{\color{kfPrimary}답안 작성}\\\hline
갈래 & ㅍㅅㅈ (고전 시조) ① & \kfTBlank \\\hline
성격 & ㅂㅈㅁ, ㄱㅇㅈ ② & \kfTBlank \\\hline
작가 & ㅅㅇㅅ 중 하나인 ㅅㅅㅁ ③ & \kfTBlank \\\hline
시대 배경 & 세조의 ㅇㅇ ㅊㅌ 이후 ④ & \kfTBlank \\\hline
충성 대상 & ㄷㅈ ⑤ & \kfTBlank \\\hline
핵심 소재 1 & ㄴㄹㅈㅅ (소나무) — 절개의 상징 ⑥ & \kfTBlank \\\hline
핵심 소재 2 & ㅂㅅ — 시련·역경의 상징 ⑦ & \kfTBlank \\\hline
핵심 한자성어 & ㄷㅇㅊㅊ — 홀로 푸르름을 지킴 ⑧ & \kfTBlank \\\hline
공간 상징 & ㅂㄹㅅ ㅈㅇㅂ — 최고의 가치 ⑨ & \kfTBlank \\\hline
표현 기법 & 소나무에 의지를 빗댄 ㅅㅈㅈ ㅂㅇ ⑩ & \kfTBlank \\\hline
주제 & ㅊㅇ와 ㅈㄱ, 죽어서도 변치 않겠다는 ㄱㅇ ⑪ & \kfTBlank \\\hline
역사적 결과 & 단종 복위 실패 후 ㅊㅎ ⑫ & \kfTBlank \\\hline
\end{kfActTable}
\kfSecHeader{kfAreaLiterature}{문제}{}

\kfBeginMC
\begin{kfQBlock}{01}{객관식}
\kfQStem{윗시조에 대한 설명으로 적절하지 \underline{않은} 것은?}
\begin{kfChoices}
\kfChoice 초장에서 화자 자신의 죽음 이후를 가정하여 시상을 전개하고 있다.
\kfChoice 소나무의 속성을 활용하여 화자의 의지를 상징적으로 드러내고 있다.
\kfChoice '봉래산 제일봉'이라는 공간 설정을 통해 높은 지향을 나타내고 있다.
\kfChoice 자연에 은거하여 세속의 번뇌에서 벗어나고자 하는 탈속 의지를 드러내고 있다.
\kfChoice 종장에서 '독야청청'이라는 한자성어를 통해 결연한 자세를 집약하고 있다.
\end{kfChoices}
\end{kfQBlock}

\begin{kfQBlock}{02}{객관식}
\kfQStem{'낙락장송 되었다가'에서 '소나무'가 상징하는 바로 가장 적절한 것은?}
\begin{kfChoices}
\kfChoice 세월의 흐름에 따라 자연스럽게 늙어 가는 인생의 무상함
\kfChoice 어떤 시련에도 변하지 않는 충의와 절개
\kfChoice 자연과 합일하여 마음의 평화를 얻으려는 소망
\kfChoice 권력의 정상에 오르고자 하는 정치적 야심
\kfChoice 고독한 삶을 택하여 세상과 단절하려는 의지
\end{kfChoices}
\end{kfQBlock}

\begin{kfQBlock}{03}{객관식}
\kfQStem{'백설이 만건곤할 제'에서 '백설'이 의미하는 것으로 가장 적절한 것은?}
\begin{kfChoices}
\kfChoice 아름답고 순수한 자연의 경치
\kfChoice 화자가 지향하는 순결한 이상 세계
\kfChoice 화자의 충의를 시험하는 시련이나 역경
\kfChoice 세상 사람들의 순수한 마음
\kfChoice 죽음 이후에 맞이하게 될 평화로운 세계
\end{kfChoices}
\end{kfQBlock}

\begin{kfQBlock}{04}{객관식}
\kfQStem{윗시조의 주제 의식과 가장 거리가 먼 것은?}
\begin{kfChoices}
\kfChoice 죽어서도 변하지 않을 충의와 절개의 다짐
\kfChoice 시련 앞에서 홀로 지조를 지키겠다는 결연한 의지
\kfChoice 자신의 신념을 위해 죽음도 두려워하지 않는 비장한 각오
\kfChoice 세속의 명리를 버리고 자연에서 유유자적하려는 풍류 정신
\kfChoice 부당한 권력에 굴하지 않는 선비의 기개
\end{kfChoices}
\end{kfQBlock}

\begin{kfQBlock}{05}{객관식}
\kfQStem{윗시조를 바탕으로 <보기>를 감상한 내용으로 적절하지 \underline{않은} 것은?}
\begin{kfEvidence}성삼문은 세종 때 집현전 학사로 훈민정음 창제에 참여한 인물이다. 세조가 조카인 단종의 왕위를 빼앗자, 성삼문은 박팽년·하위지 등과 함께 단종 복위를 도모하였으나 실패하여 처형당하였다. 이들을 '사육신(死六臣)'이라 부른다. 성삼문의 시조는 이러한 역사적 맥락에서 충신의 변하지 않는 절의(節義)를 노래한 것으로 높이 평가받고 있다.\end{kfEvidence}
\begin{kfChoices}
\kfChoice '이 몸이 죽어 가서'에서 화자가 자신의 죽음을 전제하는 것은 처형을 앞둔 성삼문의 실제 상황과 연결되겠군.
\kfChoice '봉래산 제일봉'은 자신의 충의가 최고의 가치임을 드러내는 것으로, <보기>에서 말한 충신의 절의와 맞닿겠군.
\kfChoice '백설이 만건곤'하는 상황은 세조 정권의 탄압이라는 역사적 시련에 대응하는 것이겠군.
\kfChoice '독야청청하리라'는 세조의 정당성을 인정하면서도 아쉬움을 표현한 것이겠군.
\kfChoice 죽어서 소나무가 되겠다는 발상은 죽음으로써 절의를 완성하겠다는 사육신의 정신과 통하겠군.
\end{kfChoices}
\end{kfQBlock}

\kfEndMC

\kfPassageNote{「오해」 — 박완서}{%
시어머니는 내가 무엇을 해도 마음에 안 드는 눈치였다. 나는 아침마다 식탁에 정성껏 밥을 차렸지만, 어머니는 한마디 칭찬 없이 젓가락만 들었다. 가끔은 반찬 투정을 하시기도 했다. 나는 속으로 울컥했지만 내색하지 않았다.

(중략)

"요즘 젊은 것들은 시부모 앞에서 다리를 꼬고 앉아." 어머니가 이웃집 아주머니에게 하는 말이 복도까지 들렸다. 나는 문 앞에서 발길을 돌렸다. 그날 저녁 나는 평소보다 반찬을 하나 줄여 내놓았다. 어머니가 눈치를 챘는지 아무 말 없이 밥을 드셨다.

(중략)

겨울이었다. 아이가 열이 나서 꼼짝을 못 하고 있는데, 현관문이 열리더니 어머니가 들어오셨다. 손에는 약봉지와 죽 냄비를 들고 계셨다. "빨리 이거 먹여. 약국에 가니까 이 약이 좋다고 하더라." 어머니는 아이의 이마를 짚어 보시고는 혀를 끌끌 차셨다. 그러고는 부엌으로 들어가 직접 죽을 끓이기 시작하셨다.

나는 부엌 앞에 서서 어머니의 뒷모습을 보았다. 구부정한 등, 후줄근한 손, 바쁘게 움직이는 발. 나는 그제야 어머니가 나를 미워한 적이 없다는 걸 알았다. 단지 사랑하는 방법이 달랐을 뿐이었다.

나는 모르게 눈물이 흘렀다. 어머니가 돌아서며 말씀하셨다. "왜 울어, 이 바보야. 애 아프면 엄마가 오는 거지."

그 순간 나는 어머니를 안았다. 어머니는 처음에는 어색해하시더니 곧 내 등을 토닥이셨다. 오해가 풀리는 데는 말이 필요 없었다.
}\kfResetAreaFlag
\par\kfMaybeNeedspace{50\baselineskip}\noindent{\kfTipMark\,\,}{\small\color{kfInk} 빈칸에 알맞은 말을 채워 넣으시오. (초성 힌트 참고)}\par\nopagebreak[4]\smallskip\nopagebreak[4]
\begin{kfActTable}{|>{\centering\arraybackslash}m{2.6cm}|m{6cm}|m{4cm}|}
\rowcolor{kfPrimaryLight!50}\textbf{\color{kfPrimary}항목} & \textbf{\color{kfPrimary}내용 (빈칸 채우기)} & \textbf{\color{kfPrimary}답안 작성}\\\hline
갈래 & 현대 ㄷㅍ ㅅㅅ ① & \kfTBlank \\\hline
성격 & ㅅㅅㅈ, ㅅㄹㅈ ② & \kfTBlank \\\hline
시점 & ㅇㅇㅊ ㅈㅇㅊ 시점 (며느리) ③ & \kfTBlank \\\hline
갈등 구조 & ㅅㅇㅁㄴ와 ㅁㄴㄹ 사이의 ㅅㅌ ㅂㅈ ④ & \kfTBlank \\\hline
갈등 원인 & 세대 간 ㄱㅊㄱ ㅊㅇ와 ㅇㅎ ⑤ & \kfTBlank \\\hline
핵심 키워드 & 사랑하는 ㅂㅂ의 차이 ⑥ & \kfTBlank \\\hline
전환점 & 아이가 아플 때 어머니가 ㅇ과 ㅈ을 가져옴 ⑦ & \kfTBlank \\\hline
며느리의 깨달음 & '어머니가 나를 ㅁㅇ한 적이 없다' ⑧ & \kfTBlank \\\hline
화해 방식 & 말 없이 ㅍㅇ과 등 ㅌㄷㅇ ⑨ & \kfTBlank \\\hline
서술 특징 & 일상적 ㅅㅈ를 통한 심리 묘사 ⑩ & \kfTBlank \\\hline
주제 & ㅅㄷ ㄱ ㅅㅌ ㅂㅈ와 ㅎㅎ ⑪ & \kfTBlank \\\hline
마지막 문장 & 'ㅇㅎ가 풀리는 데는 ㅁ이 필요 없었다' ⑫ & \kfTBlank \\\hline
\end{kfActTable}
\kfSecHeader{kfAreaLiterature}{문제}{}

\kfBeginMC
\begin{kfQBlock}{01}{객관식}
\kfQStem{윗글에 대한 설명으로 적절하지 \underline{않은} 것은?}
\begin{kfChoices}
\kfChoice 며느리의 시점에서 시어머니와의 관계를 서술하고 있다.
\kfChoice 시어머니의 행동에 대한 며느리의 내면 심리가 직접 드러나고 있다.
\kfChoice 일상적인 소재(식사, 간병)를 통해 인물 간 갈등과 화해를 그리고 있다.
\kfChoice 시어머니는 처음부터 며느리를 미워했으나 아이의 병을 계기로 마음을 바꾼다.
\kfChoice 서사의 전환점에서 며느리가 시어머니의 진심을 깨닫는 장면이 제시되어 있다.
\end{kfChoices}
\end{kfQBlock}

\begin{kfQBlock}{02}{객관식}
\kfQStem{며느리가 '반찬을 하나 줄여 내놓'은 행동이 의미하는 것으로 가장 적절한 것은?}
\begin{kfChoices}
\kfChoice 시어머니에 대한 서운함과 감정적 반응의 간접적 표출
\kfChoice 시어머니의 건강을 걱정하여 음식 양을 조절한 배려
\kfChoice 경제적 어려움으로 인한 절약의 실천
\kfChoice 시어머니의 반찬 투정에 대한 적극적 대응
\kfChoice 요리 실력의 한계를 인정하는 솔직한 태도
\end{kfChoices}
\end{kfQBlock}

\begin{kfQBlock}{03}{객관식}
\kfQStem{며느리가 '어머니가 나를 미워한 적이 없다는 걸 알았다'라고 깨닫게 된 직접적 계기로 가장 적절한 것은?}
\begin{kfChoices}
\kfChoice 시어머니가 그동안의 험담에 대해 며느리에게 직접 사과를 건넸기 때문이다.
\kfChoice 아이가 열이 나자 약과 죽을 들고 와 손수 돌봐 주는 모습을 본 까닭이다.
\kfChoice 이웃 아주머니가 시어머니의 속마음을 며느리에게 살며시 전해 주었기 때문이다.
\kfChoice 남편이 시어머니와 며느리 사이의 갈등을 적극적으로 중재해 주었기 때문이다.
\kfChoice 시어머니가 며느리의 노고를 위로하며 큰돈을 쥐여 주었기 때문이다.
\end{kfChoices}
\end{kfQBlock}

\begin{kfQBlock}{04}{객관식}
\kfQStem{윗글의 마지막 문장 '오해가 풀리는 데는 말이 필요 없었다'가 지닌 의미로 가장 적절한 것은?}
\begin{kfChoices}
\kfChoice 시어머니와 며느리가 서로에 대한 감정을 영원히 말하지 않기로 합의한 것이다.
\kfChoice 말로 표현하지 않아도 행동과 정서적 교감을 통해 진심이 전달될 수 있다는 의미이다.
\kfChoice 두 사람의 갈등이 아직 해결되지 않았음을 암시하는 것이다.
\kfChoice 며느리가 시어머니에게 할 말을 포기한 채 체념한 것이다.
\kfChoice 시어머니가 청각에 문제가 있어 대화가 불가능했다는 뜻이다.
\end{kfChoices}
\end{kfQBlock}

\begin{kfQBlock}{05}{객관식}
\kfQStem{윗글을 바탕으로 <보기>를 감상한 내용으로 적절하지 \underline{않은} 것은?}
\begin{kfEvidence}박완서는 일상의 사소한 경험 속에서 인간관계의 본질을 포착하는 작가로 평가받는다. 그녀의 소설에서 가족 간의 갈등은 대부분 '악의'가 아닌 '무지'에서 비롯된다. 상대방의 입장을 충분히 이해하지 못해서 생기는 오해가 갈등의 핵심이며, 진심이 전달되는 순간 관계는 회복된다. 이러한 서사 구조는 인간에 대한 근본적 신뢰를 바탕으로 한다.\end{kfEvidence}
\begin{kfChoices}
\kfChoice 시어머니의 반찬 투정과 험담은 며느리에 대한 '악의'가 아니라 사랑 표현 방식의 차이에서 비롯된 것이겠군.
\kfChoice 며느리가 반찬을 줄인 것은 상대방의 입장을 이해하지 못한 '무지'에서 나온 행동이겠군.
\kfChoice 아이가 아플 때 달려온 시어머니의 행동은 진심이 전달되는 순간을 만들어 관계 회복의 계기가 되겠군.
\kfChoice 시어머니와 며느리의 갈등은 근본적으로 해소가 불가능한 세대 간 대립의 비극이겠군.
\kfChoice '오해가 풀리는 데는 말이 필요 없었다'는 진심이 전달되면 관계가 회복된다는 <보기>의 관점과 일치하겠군.
\end{kfChoices}
\end{kfQBlock}

\kfEndMC

\kfStartNonlit{공리주의란 무엇인가 외 4편}


\kfPassageNote{공리주의란 무엇인가 (윤리)}{%
공리주의는 18세기 후반 영국에서 체계화된 윤리 이론으로, 행위의 도덕적 가치를 그 행위가 가져오는 결과, 특히 행복과 쾌락의 총량에 의해 판단한다. 공리주의의 창시자로 평가받는 제러미 벤담은 '최대 다수의 최대 행복'이라는 원칙을 제시하며, 도덕적으로 올바른 행위란 관련된 모든 사람의 행복의 총합을 극대화하는 행위라고 주장했다.

벤담은 쾌락과 고통이 인간 행위의 두 가지 근본적인 동기라고 보았다. 그는 쾌락의 가치를 객관적으로 측정하기 위해 '쾌락 계산법'을 고안했는데, 여기에는 쾌락의 강도, 지속성, 확실성, 원근성, 다산성, 순수성, 범위의 일곱 가지 기준이 포함된다. 벤담에 따르면, 어떤 행위가 가져올 쾌락과 고통을 이 기준에 따라 계산한 뒤, 쾌락이 고통보다 큰 행위를 선택하는 것이 합리적이다. 이러한 접근은 모든 쾌락을 양적으로 동일하게 취급한다는 특징이 있다.

그러나 존 스튜어트 밀은 벤담의 양적 공리주의에 중요한 수정을 가했다. 밀은 쾌락에는 질적 차이가 있으며, 지적 쾌락이나 도덕적 감정에서 오는 쾌락은 단순한 신체적 쾌락보다 더 높은 가치를 지닌다고 주장했다. 밀의 유명한 표현에 따르면 '배부른 돼지보다 배고픈 소크라테스가 낫다'. 이는 쾌락의 양이 적더라도 질적으로 우월한 쾌락을 추구하는 삶이 더 바람직하다는 뜻이다. 밀은 두 종류의 쾌락을 모두 경험한 사람이라면 질적으로 높은 쾌락을 선호할 것이라고 보았다.

공리주의는 도덕 판단에 명확한 기준을 제시한다는 장점이 있지만, 비판도 만만치 않다. 첫째, 다수의 행복을 위해 소수의 권리를 희생할 수 있다는 문제가 있다. 예컨대 한 사람을 고문하여 백 명을 구할 수 있다면, 공리주의 논리에 따르면 고문이 정당화될 수 있다. 둘째, 행복의 총량을 정확히 계산하는 것이 현실적으로 가능한지에 대한 의문이 제기된다. 셋째, 행위의 동기나 과정을 무시하고 결과만으로 도덕성을 판단하는 것이 적절한지도 논란이 된다.

이러한 비판에도 불구하고, 공리주의는 현대 사회의 공공 정책 수립에 여전히 큰 영향을 미치고 있다. 비용 편익 분석이나 사회 복지 정책의 설계에서 '최대 다수의 최대 행복'이라는 원칙은 핵심적인 판단 기준으로 활용된다. 공리주의가 제기한 '행위의 결과가 가져오는 전체적인 행복'이라는 관점은, 개인의 이익만을 고려하는 이기주의를 넘어 사회 전체의 이익을 고려하도록 이끈다는 점에서 의의가 있다.
}\kfResetAreaFlag
\kfSecHeader{kfAreaNonlit}{문제}{}

\kfBeginMC
\begin{kfQBlock}{01}{객관식}
\kfQStem{윗글의 내용과 일치하는 것은?}
\begin{kfChoices}
\kfChoice 벤담은 쾌락에 질적 차이가 있으므로 지적 쾌락을 우선해야 한다고 보았다.
\kfChoice 밀은 쾌락의 양만을 기준으로 행위의 도덕성을 판단해야 한다고 주장했다.
\kfChoice 벤담의 쾌락 계산법에는 강도, 지속성, 확실성 등 일곱 가지 기준이 포함된다.
\kfChoice 공리주의는 행위의 동기와 과정을 결과보다 더 중시하는 윤리 이론이다.
\kfChoice 밀은 신체적 쾌락이 지적 쾌락보다 질적으로 더 높은 가치를 지닌다고 보았다.
\end{kfChoices}
\end{kfQBlock}

\begin{kfQBlock}{02}{객관식}
\kfQStem{윗글에서 밀이 '배부른 돼지보다 배고픈 소크라테스가 낫다'고 한 까닭으로 가장 적절한 것은?}
\begin{kfChoices}
\kfChoice 소크라테스가 역사적으로 돼지보다 훨씬 유명한 인물로 평가되기 때문이다.
\kfChoice 쾌락의 양이 적더라도 질적으로 우월한 쾌락의 삶이 더 바람직하기 때문이다.
\kfChoice 배고픔을 묵묵히 견디는 인내심이 도덕적 행위의 핵심 조건이기 때문이다.
\kfChoice 동물은 쾌락을 느낄 수 없으므로 행복 계산에서 반드시 제외해야 하기 때문이다.
\kfChoice 벤담의 쾌락 계산법을 동물에게까지 확장하여 적용한 대표적 사례이기 때문이다.
\end{kfChoices}
\end{kfQBlock}

\begin{kfQBlock}{03}{객관식}
\kfQStem{윗글을 바탕으로 <보기>를 이해한 내용으로 적절하지 \underline{않은} 것은?}
\begin{kfEvidence}정부가 새로운 고속도로를 건설하려 한다. 이 도로가 완공되면 수만 명의 시민이 교통 편의를 누리지만, 건설 과정에서 50가구가 강제로 이주해야 한다. 정부는 비용 편익 분석을 실시하여 전체 사회적 편익이 비용을 초과한다는 결론을 내렸다.\end{kfEvidence}
\begin{kfChoices}
\kfChoice 정부의 비용 편익 분석은 공리주의적 판단 방식을 적용한 사례에 해당한다.
\kfChoice 수만 명의 편의를 위해 50가구를 이주시키는 것은 최대 다수의 최대 행복 원칙과 관련된다.
\kfChoice 이 사례에서는 소수의 권리가 다수의 행복을 위해 희생될 수 있다는 공리주의의 문제점이 드러난다.
\kfChoice 50가구의 이주 의사와 무관하게 결정한 것은 행위의 동기를 중시한 결과이다.
\kfChoice 전체 사회적 편익이 비용을 초과한다는 판단은 행복의 총량 계산에 기초한 것이다.
\end{kfChoices}
\end{kfQBlock}

\begin{kfQBlock}{04}{객관식}
\kfQStem{윗글에 따를 때, 벤담과 밀의 공통점과 차이점에 대한 설명으로 가장 적절한 것은?}
\begin{kfChoices}
\kfChoice 벤담과 밀은 모두 쾌락의 질적 차이를 인정하되, 그 기준에서 차이를 보였다.
\kfChoice 벤담은 행복의 총합을, 밀은 개인의 행복만을 도덕 판단의 기준으로 삼았다.
\kfChoice 벤담과 밀은 모두 결과에 따라 행위를 평가하되, 쾌락의 가치를 보는 관점이 달랐다.
\kfChoice 벤담은 결과주의를, 밀은 의무론을 각각 지지하여 근본적으로 대립하였다.
\kfChoice 벤담과 밀은 모두 행위의 동기를 가장 중요한 도덕 판단 기준으로 보았다.
\end{kfChoices}
\end{kfQBlock}

\begin{kfQBlock}{05}{객관식}
\kfQStem{윗글에서 공리주의에 대한 비판으로 제시된 내용이 \underline{아닌} 것은?}
\begin{kfChoices}
\kfChoice 다수의 행복을 위해 소수의 권리가 희생될 수 있다.
\kfChoice 행복의 총량을 정확히 계산하는 것이 현실적으로 어렵다.
\kfChoice 행위의 동기나 과정을 고려하지 않고 결과만을 중시한다.
\kfChoice 개인의 이익만을 고려하여 사회 전체의 이익을 무시한다.
\kfChoice 한 사람을 고문하여 백 명을 구하는 것이 정당화될 수 있다.
\end{kfChoices}
\end{kfQBlock}

\kfEndMC

\kfPassageNote{칸트의 의무 윤리학 (윤리)}{%
임마누엘 칸트는 18세기 독일의 철학자로, 행위의 결과가 아니라 행위의 동기와 원칙에 도덕적 가치를 부여하는 의무 윤리학을 정립하였다. 칸트에 따르면, 어떤 행위가 도덕적으로 선하기 위해서는 그 행위가 의무감에서 비롯되어야 하며, 단순히 좋은 결과를 가져왔다고 해서 도덕적인 것은 아니다. 예컨대 상인이 고객에게 정직하게 대하는 것이 이익이 되기 때문에 정직한 것과, 정직이 의무이기 때문에 정직한 것은 도덕적으로 전혀 다른 행위이다.

칸트 윤리학의 핵심 개념은 '선의지(善意志)'이다. 칸트는 이 세계에서, 그리고 이 세계 밖에서도 아무런 제한 없이 선하다고 볼 수 있는 것은 오직 선의지뿐이라고 선언했다. 선의지란 도덕 법칙을 따르려는 의지, 즉 옳다고 판단한 것을 옳기 때문에 행하려는 의지를 말한다. 재능이나 용기, 결단력 같은 성질은 선의지 없이는 오히려 해로울 수 있다. 범죄자의 용기나 독재자의 결단력이 그 예이다.

칸트는 도덕 법칙의 형식을 '정언 명령'으로 제시하였다. 정언 명령이란 어떠한 조건이나 목적에도 의존하지 않고 그 자체로 무조건적으로 따라야 하는 명령을 뜻한다. 칸트가 제시한 정언 명령의 첫 번째 형식은 '네 의지의 준칙이 언제나 동시에 보편적 법칙이 될 수 있도록 행위하라'이다. 이는 자신의 행동 원칙이 모든 사람에게 적용되더라도 모순이 없는지 검토하라는 뜻이다. 예를 들어 '거짓말을 해도 된다'는 준칙은 모든 사람이 따를 경우 약속과 신뢰의 체계 자체가 붕괴하므로 보편화될 수 없다.

정언 명령의 두 번째 형식은 '인간을 단지 수단으로만 대하지 말고, 언제나 동시에 목적으로 대하라'이다. 이는 인간이 존엄한 존재이며, 다른 목적을 달성하기 위한 도구로만 이용되어서는 안 된다는 뜻이다. 칸트에 따르면 인간은 이성적 존재로서 스스로 도덕 법칙을 세울 수 있는 자율적 주체이며, 이러한 자율성이 인간 존엄성의 근거가 된다. 따라서 어떤 사람을 속이거나 강제하여 자신의 이익을 위해 이용하는 것은, 설령 좋은 결과를 낳더라도 도덕적으로 허용될 수 없다.

칸트의 의무 윤리학은 도덕의 보편적이고 절대적인 기준을 제시했다는 점에서 높이 평가된다. 그러나 비판도 존재한다. 도덕적 의무들이 서로 충돌하는 상황, 예컨대 살인자에게 쫓기는 친구의 위치를 물었을 때 거짓말을 해야 하는지의 문제에서, 칸트의 윤리학은 유연한 해결을 제시하기 어렵다는 지적이 있다. 또한 결과를 전혀 고려하지 않는 것이 현실적으로 적절한지에 대한 논란도 이어지고 있다.
}\kfResetAreaFlag
\kfSecHeader{kfAreaNonlit}{문제}{}

\kfBeginMC
\begin{kfQBlock}{01}{객관식}
\kfQStem{윗글의 내용과 일치하지 않는 것은?}
\begin{kfChoices}
\kfChoice 칸트는 행위의 결과가 아니라 동기와 원칙에 도덕적 가치를 부여하였다.
\kfChoice 선의지란 도덕 법칙을 따르려는 의지, 즉 옳기 때문에 행하려는 의지이다.
\kfChoice 정언 명령은 특정 목적을 달성하기 위한 조건부 명령의 한 형태이다.
\kfChoice 칸트에 따르면 인간의 자율성이 존엄성의 근거가 된다.
\kfChoice 재능이나 용기는 선의지 없이는 오히려 해로울 수 있다.
\end{kfChoices}
\end{kfQBlock}

\begin{kfQBlock}{02}{객관식}
\kfQStem{윗글에서 칸트가 '거짓말을 해도 된다'는 준칙이 보편화될 수 없다고 본 까닭으로 가장 적절한 것은?}
\begin{kfChoices}
\kfChoice 거짓말이 개인의 양심에 정면으로 반하는 행위이기 때문이다.
\kfChoice 거짓말을 하면 결국 항상 나쁜 결과가 뒤따르게 되기 때문이다.
\kfChoice 모든 사람이 거짓말을 하면 약속과 신뢰 체계가 무너지기 때문이다.
\kfChoice 거짓말은 신이 인간에게 내린 계명에 어긋나는 행위이기 때문이다.
\kfChoice 거짓말을 허용하면 사회의 법적 처벌 체계가 무력화되기 때문이다.
\end{kfChoices}
\end{kfQBlock}

\begin{kfQBlock}{03}{객관식}
\kfQStem{윗글을 바탕으로 <보기>를 이해한 내용으로 적절하지 \underline{않은} 것은?}
\begin{kfEvidence}회사원 A는 거래처로부터 부당한 요구를 받았다. 거래처의 요구를 들어주면 회사에 큰 이익이 생기지만, A는 '부당한 요구에 응하는 것은 옳지 않다'고 판단하여 거절하였다.\end{kfEvidence}
\begin{kfChoices}
\kfChoice A는 행위의 결과보다 행위의 원칙을 기준으로 판단한 것이다.
\kfChoice A의 행동은 의무감에서 비롯된 것으로, 칸트 윤리학에서 도덕적으로 선한 행위에 해당한다.
\kfChoice A는 자신의 행동 원칙이 보편적 법칙이 될 수 있는지를 검토한 것으로 볼 수 있다.
\kfChoice A가 거절한 것은 거래처와의 관계에서 더 큰 이익을 기대했기 때문이다.
\kfChoice A의 판단은 부당한 요구에 응하는 것이 옳지 않다는 도덕 법칙에 따른 것이다.
\end{kfChoices}
\end{kfQBlock}

\begin{kfQBlock}{04}{객관식}
\kfQStem{'인간을 단지 수단으로만 대하지 말고, 언제나 동시에 목적으로 대하라'는 정언 명령의 의미로 가장 적절한 것은?}
\begin{kfChoices}
\kfChoice 인간은 어떤 경우에도 타인의 도움을 받아서는 안 된다.
\kfChoice 인간은 존엄한 존재이므로 다른 목적의 도구로만 이용되어서는 안 된다.
\kfChoice 인간은 도덕적 의무를 다하기 위해 자신의 이익을 완전히 포기해야 한다.
\kfChoice 인간 사이의 모든 거래는 도덕적으로 금지되어야 한다.
\kfChoice 인간은 타인을 위해 자신을 희생하는 것이 최고의 덕이다.
\end{kfChoices}
\end{kfQBlock}

\kfEndMC

\kfPassageNote{덕 윤리학의 부활 (윤리)}{%
덕 윤리학은 '어떤 행위가 옳은가'보다 '어떤 사람이 되어야 하는가'를 묻는 윤리 이론이다. 근대 이후 공리주의와 의무론이 윤리학의 중심을 차지하면서, 행위의 규칙이나 결과에 초점을 맞추는 접근이 주류가 되었다. 그러나 20세기 후반, 행위 중심의 윤리학이 인간의 삶 전체와 품성을 간과한다는 비판이 제기되면서 고대 그리스 철학, 특히 아리스토텔레스의 덕 윤리학이 다시 주목받기 시작했다.

아리스토텔레스에 따르면, 인간의 궁극적 목표는 '에우다이모니아', 즉 행복 혹은 잘 사는 삶이다. 그런데 여기서 행복이란 일시적인 쾌락이 아니라, 덕을 갖추고 그에 따라 행동하는 삶 전체를 통해 실현되는 것이다. 덕이란 감정이나 행동에서 과도함과 부족함의 양 극단을 피하고 중간 상태를 유지하는 성품이다. 아리스토텔레스는 이를 '중용(中庸)'이라고 불렀다. 예를 들어 용기는 무모함과 비겁함의 중간이며, 관용은 낭비와 인색함의 중간이다.

그러나 중용은 단순히 산술적 평균을 뜻하지 않는다. 어떤 상황에서 어떤 정도가 적절한 중간인지를 판단하려면 '실천적 지혜', 그리스어로 '프로네시스'가 필요하다. 실천적 지혜란 구체적인 상황에서 무엇이 옳은지를 파악하고 올바른 행위를 선택하는 능력으로, 이론적 지식과 구별되는 실천적 판단력이다. 아리스토텔레스에 따르면, 이 지혜는 타고나는 것이 아니라 반복적인 실천과 경험을 통해 습득된다. 마치 악기 연주가 반복적인 연습을 통해 숙달되듯이, 덕도 올바른 행위의 반복을 통해 품성으로 자리 잡는다.

20세기 후반, 앨러스데어 매킨타이어는 아리스토텔레스의 전통을 현대에 되살리는 데 핵심적인 역할을 했다. 매킨타이어는 근대 윤리학이 보편적인 도덕 규칙을 찾으려는 프로젝트에 실패했다고 진단했다. 그에 따르면, 도덕적 판단은 추상적 원칙에서 도출되는 것이 아니라, 구체적인 공동체의 전통과 서사 속에서 형성된다. 매킨타이어는 '실천'이라는 개념을 제안했는데, 이는 공동체 안에서 공유되는 활동을 통해 내적 가치를 추구하며 탁월성의 기준이 확장되는 사회적 과정을 뜻한다.

덕 윤리학은 행위의 규칙보다 행위자의 품성과 삶의 맥락을 중시한다는 점에서, 기계적인 규칙 적용의 한계를 보완한다. 그러나 '어떤 덕이 진정한 덕인가'에 대한 합의가 문화마다 다를 수 있다는 점, 구체적인 행위 지침을 제시하기 어렵다는 점에서 비판이 제기되기도 한다. 그럼에도 덕 윤리학은 도덕 교육과 인성 형성의 중요성을 일깨운다는 점에서 현대 윤리학의 중요한 한 축을 이루고 있다.
}\kfResetAreaFlag
\kfSecHeader{kfAreaNonlit}{문제}{}

\kfBeginMC
\begin{kfQBlock}{01}{객관식}
\kfQStem{윗글의 내용과 일치하는 것은?}
\begin{kfChoices}
\kfChoice 아리스토텔레스는 중용을 두 극단의 산술적 평균으로 정의하였다.
\kfChoice 실천적 지혜는 이론적 지식과 동일한 성격의 능력이다.
\kfChoice 매킨타이어는 도덕적 판단이 공동체의 전통과 서사 속에서 형성된다고 보았다.
\kfChoice 덕 윤리학은 행위의 규칙을 행위자의 품성보다 더 중시하는 이론이다.
\kfChoice 아리스토텔레스에 따르면 실천적 지혜는 타고나는 선천적 능력이다.
\end{kfChoices}
\end{kfQBlock}

\begin{kfQBlock}{02}{객관식}
\kfQStem{윗글에서 아리스토텔레스가 덕의 습득을 '악기 연주의 연습'에 비유한 까닭으로 가장 적절한 것은?}
\begin{kfChoices}
\kfChoice 악기 연주처럼 덕도 뛰어난 재능이 있어야만 발휘할 수 있기 때문이다.
\kfChoice 악기 연주가 반복 연습으로 숙달되듯이 덕도 반복적 실천으로 품성에 자리 잡기 때문이다.
\kfChoice 악기 연주와 도덕적 행위가 동일한 뇌 영역에서 이루어지기 때문이다.
\kfChoice 악기 연주가 청중에게 감동을 주듯이 덕도 타인에게 영향을 미치기 때문이다.
\kfChoice 악기 연주 실력이 측정 가능하듯이 덕의 수준도 객관적으로 계량할 수 있기 때문이다.
\end{kfChoices}
\end{kfQBlock}

\begin{kfQBlock}{03}{객관식}
\kfQStem{윗글에 따를 때, 공리주의 및 의무론과 비교하여 덕 윤리학이 갖는 특징으로 가장 적절한 것은?}
\begin{kfChoices}
\kfChoice 덕 윤리학은 행위의 결과를 극대화하는 보편적 규칙을 명확히 제시한다.
\kfChoice 덕 윤리학은 규칙이나 결과보다 행위자의 품성과 삶의 맥락을 중시한다.
\kfChoice 덕 윤리학은 의무론과 마찬가지로 행위의 동기를 최우선의 기준으로 삼는다.
\kfChoice 덕 윤리학은 공리주의와 마찬가지로 행복의 총량을 정량적으로 계산한다.
\kfChoice 덕 윤리학은 상황마다 따라야 할 구체적 행위 지침을 명확하게 제시한다.
\end{kfChoices}
\end{kfQBlock}

\begin{kfQBlock}{04}{객관식}
\kfQStem{윗글을 바탕으로 <보기>를 이해한 내용으로 적절하지 \underline{않은} 것은?}
\begin{kfEvidence}축구 선수 B는 어릴 때부터 팀 동료를 배려하고 공정하게 경기하는 태도를 꾸준히 실천하여, 성인이 된 뒤에도 자연스럽게 동료를 존중하고 페어플레이를 실천하는 선수로 성장하였다.\end{kfEvidence}
\begin{kfChoices}
\kfChoice B가 배려와 공정성을 꾸준히 실천한 것은, 덕이 반복적 실천을 통해 습득된다는 관점과 부합한다.
\kfChoice B가 성인이 된 뒤 자연스럽게 품성을 발휘하는 것은, 덕이 품성으로 자리 잡는 과정을 보여준다.
\kfChoice B의 사례는 '어떤 행위가 옳은가'보다 '어떤 사람이 되어야 하는가'를 묻는 덕 윤리학의 관심과 통한다.
\kfChoice B가 페어플레이를 실천하는 것은, 그가 이론적 지식을 학습하여 규칙을 암기한 결과이다.
\kfChoice B의 태도는 과도함과 부족함의 양 극단을 피한 중용의 실천으로 해석할 수 있다.
\end{kfChoices}
\end{kfQBlock}

\kfEndMC

\kfPassageNote{도덕적 딜레마와 트롤리 문제 (윤리)}{%
도덕적 딜레마란 서로 양립할 수 없는 두 가지 이상의 도덕적 요구가 충돌하여, 어떤 선택을 하더라도 도덕적 잘못을 피할 수 없는 상황을 말한다. 이러한 딜레마는 윤리학의 주요 이론들이 실제 상황에서 어떻게 작동하는지를 시험하는 중요한 사고 실험의 소재가 된다. 그 가운데 가장 널리 알려진 것이 바로 '트롤리 문제'이다.

트롤리 문제는 영국의 철학자 필리파 풋이 1967년에 처음 제안한 사고 실험이다. 기본 시나리오는 다음과 같다. 브레이크가 고장 난 전차(트롤리)가 선로 위의 다섯 명을 향해 돌진하고 있다. 당신은 선로 전환기를 조작하여 전차를 다른 선로로 돌릴 수 있지만, 그 선로에는 한 명이 서 있다. 전환기를 당기면 다섯 명은 살지만 한 명이 죽고, 당기지 않으면 다섯 명이 죽고 한 명은 산다. 당신은 어떻게 해야 하는가?

결과주의, 특히 공리주의의 관점에서 이 문제의 답은 비교적 명확하다. 다섯 명의 생명을 구하는 것이 한 명의 생명을 구하는 것보다 더 큰 행복의 총량을 산출하므로, 전환기를 당기는 것이 도덕적으로 옳다. 그러나 의무론의 관점에서는 문제가 달라진다. 칸트의 정언 명령에 따르면 인간을 단지 수단으로만 대해서는 안 되는데, 전환기를 당기는 행위는 한 명을 다섯 명을 살리기 위한 수단으로 이용하는 것이므로 도덕적으로 정당화되기 어렵다.

이 딜레마를 더욱 첨예하게 만드는 변형이 '다리 위의 뚱뚱한 남자' 시나리오이다. 이 경우 당신은 다리 위에 서 있고, 옆에 큰 체구의 남자가 있다. 이 남자를 다리 아래로 밀어 전차를 멈추면 다섯 명을 구할 수 있다. 흥미로운 점은, 결과만 놓고 보면 전환기 시나리오와 동일하게 한 명이 죽고 다섯 명이 사는 것인데, 대부분의 사람은 전환기를 당기는 것은 허용하면서도 사람을 직접 밀어 떨어뜨리는 것은 허용하지 않는다는 사실이다. 이는 결과가 같더라도 행위의 방식이 도덕적 판단에 영향을 미친다는 것을 보여준다.

이러한 직관의 차이를 설명하기 위해 등장한 것이 '이중 결과 원칙'이다. 이 원칙에 따르면, 선한 목적을 위한 행위의 부수적 결과로 나쁜 결과가 발생하는 것은 허용될 수 있지만, 나쁜 결과 자체를 수단으로 의도하는 것은 허용되지 않는다. 전환기 시나리오에서 한 명의 죽음은 다섯 명을 구하려는 행위의 의도하지 않은 부수적 결과이지만, 다리 시나리오에서는 한 명을 죽이는 것이 다섯 명을 구하기 위한 직접적 수단이므로 도덕적으로 다르게 평가된다. 트롤리 문제는 단순한 사고 실험을 넘어, 인공지능의 윤리적 판단 설계나 의료 자원 배분과 같은 현실 문제에서도 중요한 참조점이 되고 있다.
}\kfResetAreaFlag
\kfSecHeader{kfAreaNonlit}{문제}{}

\kfBeginMC
\begin{kfQBlock}{01}{객관식}
\kfQStem{윗글의 내용과 일치하는 것은?}
\begin{kfChoices}
\kfChoice 트롤리 문제는 독일의 철학자 칸트가 18세기에 처음 제안한 사고 실험이다.
\kfChoice 공리주의 관점에서는 전환기를 당기지 않는 것이 도덕적으로 올바른 선택이다.
\kfChoice 다리 시나리오와 전환기 시나리오는 결과적으로 동일하게 한 명이 죽고 다섯 명이 산다.
\kfChoice 대부분의 사람은 전환기를 당기는 것과 사람을 밀어 떨어뜨리는 것을 동일하게 허용한다.
\kfChoice 이중 결과 원칙에 따르면 나쁜 결과를 수단으로 의도하는 것도 선한 목적이면 허용된다.
\end{kfChoices}
\end{kfQBlock}

\begin{kfQBlock}{02}{객관식}
\kfQStem{윗글에 따를 때, 트롤리 문제의 전환기 시나리오에 대한 결과주의와 의무론의 판단을 비교한 것으로 가장 적절한 것은?}
\begin{kfChoices}
\kfChoice 결과주의는 전환기를 당기는 것을 강하게 반대하지만, 의무론은 이를 적극 지지한다.
\kfChoice 결과주의와 의무론 모두 전환기를 당기는 행위를 도덕적으로 옳다고 동일하게 평가한다.
\kfChoice 결과주의는 행복의 총량을, 의무론은 인간을 수단으로 대하는지를 기준으로 다른 결론을 내린다.
\kfChoice 결과주의는 이중 결과 원칙을, 의무론은 쾌락 계산법을 일관되게 적용하여 판단을 내린다.
\kfChoice 결과주의와 의무론 모두 한 명의 생명도 절대 희생해서는 안 된다는 동일한 결론에 도달한다.
\end{kfChoices}
\end{kfQBlock}

\begin{kfQBlock}{03}{객관식}
\kfQStem{윗글에서 '이중 결과 원칙'이 전환기 시나리오와 다리 시나리오를 다르게 평가하는 근거로 가장 적절한 것은?}
\begin{kfChoices}
\kfChoice 전환기 시나리오에서는 한 명조차 죽지 않지만 다리 시나리오에서는 사람이 죽게 되기 때문이다.
\kfChoice 전환기 시나리오의 죽음은 부수적 결과이지만 다리 시나리오에서는 죽음 자체가 직접적 수단이기 때문이다.
\kfChoice 전환기 시나리오는 기계 장치를 조작하지만 다리 시나리오는 사람의 신체에 직접 접촉해야 하기 때문이다.
\kfChoice 전환기 시나리오에서는 행위자가 위험에 함께 노출되지만 다리 시나리오에서는 안전하게 분리되기 때문이다.
\kfChoice 전환기 시나리오에서 희생되는 사람은 이미 선로 위에서 위험에 처해 있던 인물이기 때문이다.
\end{kfChoices}
\end{kfQBlock}

\begin{kfQBlock}{04}{객관식}
\kfQStem{윗글을 바탕으로 <보기>를 이해한 내용으로 적절하지 \underline{않은} 것은?}
\begin{kfEvidence}전염병이 확산되어 백신이 부족한 상황이다. 한정된 백신을 고위험군 환자 다섯 명에게 투여하면 모두 살 수 있지만, 같은 백신을 희귀 질환 환자 한 명에게 투여하면 그 한 명만 살 수 있다. 보건 당국은 고위험군 다섯 명에게 백신을 배분하기로 결정했다.\end{kfEvidence}
\begin{kfChoices}
\kfChoice 이 결정은 공리주의적 관점에서 더 많은 사람의 생명을 구하는 선택으로 해석할 수 있다.
\kfChoice 희귀 질환 환자의 입장에서는 의무론적 관점에서 자신이 수단으로 취급되었다고 반발할 수 있다.
\kfChoice 이 상황은 서로 양립할 수 없는 도덕적 요구가 충돌하는 도덕적 딜레마에 해당한다.
\kfChoice 이 사례는 트롤리 문제처럼 윤리적 판단 기준에 따라 다른 결론이 도출될 수 있는 상황이다.
\kfChoice 이중 결과 원칙에 따르면 희귀 질환 환자의 죽음은 의도적 수단이므로 이 결정은 허용되지 않는다.
\end{kfChoices}
\end{kfQBlock}

\begin{kfQBlock}{05}{객관식}
\kfQStem{윗글에 따를 때, '도덕적 딜레마'에 대한 이해로 적절하지 \underline{않은} 것은?}
\begin{kfChoices}
\kfChoice 서로 양립할 수 없는 도덕적 요구가 충돌하는 상황이다.
\kfChoice 어떤 선택을 하더라도 도덕적 잘못을 완전히 피하기 어렵다.
\kfChoice 윤리 이론들이 실제 상황에서 어떻게 작동하는지를 시험하는 소재가 된다.
\kfChoice 하나의 올바른 답이 존재하여 이론 간 갈등 없이 해결된다.
\kfChoice 인공지능의 윤리적 판단 설계에도 참조점이 되고 있다.
\end{kfChoices}
\end{kfQBlock}

\kfEndMC

\kfPassageNote{동물 권리와 윤리적 소비 (윤리)}{%
오스트레일리아의 철학자 피터 싱어는 1975년 출간한 《동물 해방》에서 '종차별주의'라는 개념을 제시하며 동물 윤리에 관한 현대적 논의의 포문을 열었다. 종차별주의란 단지 다른 종에 속한다는 이유만으로 동물의 이익을 무시하거나 차별하는 태도를 말한다. 싱어에 따르면 이는 인종이나 성별에 따른 차별과 구조적으로 동일한 편견이다. 그는 고통을 느낄 수 있는 능력, 즉 감각 능력이 도덕적 고려의 핵심 기준이라고 주장했다.

싱어의 논증은 공리주의에 기초한다. 쾌락과 고통이 도덕적 판단의 기본 단위라면, 고통을 느끼는 존재라면 종에 관계없이 그 고통을 동등하게 고려해야 한다. 이를 '이익 평등 고려의 원칙'이라고 한다. 이 원칙은 동물의 이익을 인간의 이익과 완전히 동일하게 취급하라는 뜻이 아니라, 유사한 수준의 고통에 대해서는 유사한 수준의 도덕적 비중을 부여해야 한다는 것이다. 예컨대 개가 느끼는 육체적 고통과 인간이 느끼는 유사한 육체적 고통은 동등하게 고려되어야 한다.

싱어가 특히 비판한 것은 공장식 축산이다. 현대의 공장식 축산에서 동물들은 극도로 좁은 공간에 밀집 사육되며, 자연스러운 행동을 할 수 없는 환경에서 고통받는다. 닭은 날개를 펼 수 없는 배터리 케이지에 갇히고, 돼지는 몸을 돌릴 수 없는 임신 틀에 가둬진다. 싱어에 따르면, 인간이 고기를 먹는 쾌락을 위해 동물에게 이 정도의 고통을 가하는 것은 도덕적으로 정당화될 수 없다. 고기를 먹지 않고도 건강하게 살 수 있는 대안이 존재하는 이상, 동물에게 가해지는 고통의 총량은 인간이 얻는 쾌락의 총량을 크게 초과하기 때문이다.

이러한 문제의식에서 '윤리적 소비' 운동이 확산되고 있다. 윤리적 소비란 상품의 생산 과정에서 발생하는 윤리적 문제를 고려하여 소비를 결정하는 행위를 말한다. 동물 복지 인증 제품의 구매, 공정 무역 상품의 선택, 환경 친화적 제품의 소비 등이 이에 해당한다. 윤리적 소비의 핵심은 소비자가 단순히 가격과 품질만이 아니라, 그 상품이 만들어지는 과정에서 누가 어떤 대가를 치렀는지를 고려한다는 데 있다.

그러나 윤리적 소비에 대한 비판도 있다. 윤리적 제품은 일반 제품보다 가격이 높아 경제적 여유가 있는 소비자만 선택할 수 있으며, 이로 인해 도덕적 소비가 계층적 특권이 될 수 있다는 우려가 있다. 또한 개별 소비자의 선택만으로 구조적인 문제가 해결될 수 있는지에 대한 의문도 제기된다. 그럼에도 불구하고, 싱어가 제기한 종차별주의 비판과 동물 고통에 대한 도덕적 고려는 현대 윤리학에서 인간 중심주의를 넘어서는 중요한 전환점으로 평가받고 있다.
}\kfResetAreaFlag
\kfSecHeader{kfAreaNonlit}{문제}{}

\kfBeginMC
\begin{kfQBlock}{01}{객관식}
\kfQStem{윗글의 내용과 일치하지 않는 것은?}
\begin{kfChoices}
\kfChoice 싱어는 고통을 느끼는 능력이 도덕적 고려의 핵심 기준이라고 주장했다.
\kfChoice 종차별주의란 다른 종에 속한다는 이유만으로 동물의 이익을 차별하는 태도이다.
\kfChoice 이익 평등 고려의 원칙은 동물과 인간의 이익을 똑같은 비중으로 다루어야 한다는 뜻이다.
\kfChoice 싱어의 논증은 쾌락과 고통을 도덕적 판단의 기본 단위로 보는 공리주의에 기초한다.
\kfChoice 싱어에 따르면 종차별주의는 인종이나 성별에 따른 차별과 구조적으로 동일한 편견이다.
\end{kfChoices}
\end{kfQBlock}

\begin{kfQBlock}{02}{객관식}
\kfQStem{윗글에서 싱어가 공장식 축산을 도덕적으로 정당화할 수 없다고 본 까닭으로 가장 적절한 것은?}
\begin{kfChoices}
\kfChoice 공장식 축산이 환경 오염을 유발하여 인간의 건강을 해치기 때문이다.
\kfChoice 동물의 고통이 인간이 고기를 먹어 얻는 쾌락보다 크고, 대안이 존재하기 때문이다.
\kfChoice 공장식 축산 종사자의 노동 환경이 비인간적이기 때문이다.
\kfChoice 동물은 인간과 완전히 동일한 권리를 가지므로 어떤 이용도 불가능하기 때문이다.
\kfChoice 육식이 인체에 해롭다는 과학적 증거가 확립되었기 때문이다.
\end{kfChoices}
\end{kfQBlock}

\begin{kfQBlock}{03}{객관식}
\kfQStem{윗글을 바탕으로 <보기>를 이해한 내용으로 적절하지 \underline{않은} 것은?}
\begin{kfEvidence}C씨는 마트에서 계란을 고를 때, 일반 계란보다 500원 비싼 동물 복지 인증 계란을 선택하였다. C씨는 '닭이 자유롭게 움직일 수 있는 환경에서 생산된 계란을 사는 것이 옳다'고 생각했기 때문이다.\end{kfEvidence}
\begin{kfChoices}
\kfChoice C씨의 행동은 상품의 생산 과정에서 발생하는 윤리적 문제를 고려한 소비에 해당한다.
\kfChoice C씨는 가격과 품질만이 아니라 생산 과정에서의 동물 복지를 소비 기준에 포함시켰다.
\kfChoice C씨의 선택은 배터리 케이지에서 사육되는 닭의 고통을 줄이려는 윤리적 소비의 사례이다.
\kfChoice C씨의 행동은 종차별주의를 강화하여 인간의 이익을 동물의 이익보다 우선시한 것이다.
\kfChoice C씨가 더 비싼 제품을 선택한 것은 윤리적 소비가 경제적 여유와 관련될 수 있음을 보여준다.
\end{kfChoices}
\end{kfQBlock}

\begin{kfQBlock}{04}{객관식}
\kfQStem{'이익 평등 고려의 원칙'에 대한 이해로 가장 적절한 것은?}
\begin{kfChoices}
\kfChoice 동물에게도 인간과 완전히 동일한 법적 권리를 보장해야 한다는 원칙이다.
\kfChoice 유사한 수준의 고통에는 종을 떠나 유사한 도덕적 비중을 두어야 한다는 원칙이다.
\kfChoice 모든 종의 생명 가치가 동등하므로 어떠한 동물도 결코 죽여서는 안 된다는 원칙이다.
\kfChoice 인간의 이익이 어떤 상황에서도 동물의 이익보다 늘 우선한다는 원칙이다.
\kfChoice 고통을 전혀 느끼지 못하는 존재까지도 동등하게 고려해야 한다는 원칙이다.
\end{kfChoices}
\end{kfQBlock}

\begin{kfQBlock}{05}{객관식}
\kfQStem{윗글에서 윤리적 소비에 대해 제기된 비판으로 적절한 것만을 <보기>에서 모두 고른 것은?}
\begin{kfEvidence}ㄱ. 윤리적 제품의 높은 가격이 도덕적 소비를 계층적 특권으로 만들 수 있다. \\\relax
ㄴ. 개별 소비자의 선택만으로 구조적 문제가 해결될 수 있는지 의문이다. \\\relax
ㄷ. 윤리적 소비가 확산되면 기업의 이윤이 감소하여 경제가 위축된다.\end{kfEvidence}
\begin{kfChoices}
\kfChoice ㄱ
\kfChoice ㄴ
\kfChoice ㄱ, ㄴ
\kfChoice ㄱ, ㄷ
\kfChoice ㄴ, ㄷ
\end{kfChoices}
\end{kfQBlock}

\kfEndMC

\kfStartWeekly{패턴 워크북 — 총 18문항}


\kfPassageNote{문항 1 / 섞어치기}{%
원자에는 중심에 핵이 있고, 핵 주위에 동심원처럼 퍼져 있는 특정 궤도들을 도는 전자들이 있다. 이 궤도마다 전자가 갖는 에너지 값이 정해져 있는데, 이 값을 에너지 준위라고 한다. 원자가 가진 특정한 몇 개의 궤도에만 전자가 존재하므로, 전자가 가질 수 있는 에너지 준위도 특정한 값으로 정해져 있다. 전자는 에너지를 얻으면 바깥쪽 궤도로, 에너지를 잃으면 안쪽 궤도로 이동한다. 원자가 하나일 때는 전자의 에너지 준위가 특정한 값이지만, 여러 원자가 모이면 여러 원자들의 궤도가 ⓐ중첩되면서 전자들이 서로 영향을 미쳐 에너지 준위에 변화가 생긴다. 이로 인해 에너지 준위는 특정한 값이 아닌 일정한 범위로 나타나게 된다. 이렇게 에너지 값이 범위로 표현된 것을 에너지 밴드라고 한다. 보통의 물질들은 수많은 원자로 구성돼 있기 때문에 대부분 에너지 밴드를 갖고 있다. 전자는 낮은 에너지 상태가 되려는 성질이 있으므로, 발광 물질의 전자에 에너지가 공급되면 전자가 바깥쪽으로 이동한 뒤 다시

(중략)
}
\begin{kfQBlock}{01}{객관식}
\kfQStem{윗글의 Ⓐ와 <보기>의 Ⓒ를 비교해 이해한 내용으로 적절하지 \underline{않은} 것은?}
\begin{kfEvidence}Ⓒ양자점 기반 광센서는 특정 파장대 신호를 선택적으로 받아 전기 신호로 변환한다. 센서 설계자는 열 잡음을 줄이기 위해 넓은 에너지 밴드 기반 재료 대신 좁은 에너지 준위 특성을 활용하려 한다. 다만 전하 추출이 중요해 두꺼운 셸을 쓰면 출력 전류가 감소할 수 있어 표면 안정화는 리간드 제어를 중심으로 수행한다.\end{kfEvidence}
\begin{kfChoices}
\kfChoice Ⓐ와 Ⓒ 모두 전류 추출을 위해 셸을 반드시 두껍게 설계해야 성능이 향상된다.
\kfChoice Ⓐ와 Ⓒ 모두 좁은 에너지 준위 특성을 이용해 원치 않는 에너지 손실을 줄이려는 공통점이 있다.
\kfChoice Ⓒ에서는 전하 추출을 위해 두꺼운 셸 사용을 제한하고 리간드 중심 안정화를 택할 수 있다.
\end{kfChoices}
\end{kfQBlock}

\kfPassageNote{문항 2 / 부정·상대어}{%
원자에는 중심에 핵이 있고, 핵 주위에 동심원처럼 퍼져 있는 특정 궤도들을 도는 전자들이 있다. 이 궤도마다 전자가 갖는 에너지 값이 정해져 있는데, 이 값을 에너지 준위라고 한다. 원자가 가진 특정한 몇 개의 궤도에만 전자가 존재하므로, 전자가 가질 수 있는 에너지 준위도 특정한 값으로 정해져 있다. 전자는 에너지를 얻으면 바깥쪽 궤도로, 에너지를 잃으면 안쪽 궤도로 이동한다. 원자가 하나일 때는 전자의 에너지 준위가 특정한 값이지만, 여러 원자가 모이면 여러 원자들의 궤도가 ⓐ중첩되면서 전자들이 서로 영향을 미쳐 에너지 준위에 변화가 생긴다. 이로 인해 에너지 준위는 특정한 값이 아닌 일정한 범위로 나타나게 된다. 이렇게 에너지 값이 범위로 표현된 것을 에너지 밴드라고 한다. 보통의 물질들은 수많은 원자로 구성돼 있기 때문에 대부분 에너지 밴드를 갖고 있다. 전자는 낮은 에너지 상태가 되려는 성질이 있으므로, 발광 물질의 전자에 에너지가 공급되면 전자가 바깥쪽으로 이동한 뒤 다시

(중략)
}
\begin{kfQBlock}{02}{객관식}
\kfQStem{문맥상 ⓐ\textasciitilde{}ⓔ와 바꿔 쓰기에 적절하지 \underline{않은} 것은?}
\begin{kfChoices}
\kfChoice ⓒ: 조금이나마
\kfChoice ⓐ: 겹치면서
\kfChoice ⓑ: 내보내면서
\end{kfChoices}
\end{kfQBlock}

\kfPassageNote{문항 3 / 주체·객체}{%
일반적으로 혼인은 남자와 여자가 부부가 되는 일을 말한다. 과거에는 특정한 법적 절차 없이도 결혼식과 같은 사회적 관습에 의한 의식을 치르면 혼인이 성립한 것으로 보았다. 하지만 현재 대부분의 국가는 법에서 규정한 절차를 따라야만 혼인이 성립되는 법률혼 제도를 두고 있으며, 사회적 관습보다는 혼인 신고와 같은 법적 절차가 더 중요하다. 혼인은 사회를 구성하는 기본 단위인 가족을 형성하는 일이기 때문에 국가의 근간을 이루는 중요한 과정이고 이에 따라 법을 통해 관계를 보호하는 것이다. 따라서 현대 사회에서 혼인은 법적 절차에 의해 이루어지게 되고 여기에 참여한 남녀 간에는 법적 효력이 발생하게 된다.

(중략)

혼인이 성립하기 위해서는 법에서 규정하고 있는 실질적 요건과 형식적 요건을 모두 만족하여야 한다. 실질적 요건에서 가장 중요한 것은 혼인 의사의 합치이다. 혼인 의사는 부부 관계를 성립하겠다는 의사로 반드시 혼인 당사자 간의 합의가 필요하다. 또 다른 실질적 요건은 연령과 혼인 자격에 대한

(중략)
}
\begin{kfQBlock}{03}{객관식}
\kfQStem{㉮와 관련한 이해로 가장 적절한 것은?}
\begin{kfChoices}
\kfChoice 협의상 이혼에서 법원이 이혼 사유를 심사하지 않는 이유는 당사자 자율 합의를 중심으로 절차가 설계되어 있기 때문이다.
\kfChoice 협의상 이혼에서도 법원은 재판상 이혼처럼 이혼 사유 존재를 본안 심리한다.
\kfChoice 협의상 이혼은 숙려 기간 없이 즉시 신고 가능하다.
\end{kfChoices}
\end{kfQBlock}

\kfPassageNote{문항 4 / 순서·방향}{%
1908년 조지 하디와 빌헬름 바인베르크는 멘델 유전 법칙에 기초한 수학적 모델을 각각 독립적으로 제시했다. 하디는 집단 내 대립 유전자 A와 a의 빈도를 p와 q로 두고 무작위 교배에서 다음 세대의 유전자형 빈도를 분석했다. 바인베르크도 무작위 결합 조건에서 유전자 빈도가 세대를 거쳐 안정될 수 있음을 통계 자료와 함께 제시했다. 이 두 접근이 결합되어 하디-바인베르크 평형 법칙이 정립되었다.

하디-바인베르크 평형은 ㉠집단이 충분히 크고, ㉡무작위 교배가 일어나며, ㉢새 돌연변이가 없고, ㉣이입·이출 같은 이주가 없고, ㉤자연 선택이 작용하지 않을 때 유지된다. [A] 예를 들어 p=0.7, q=0.3이면 p+q=1이고 유전자형 빈도는 AA=p²=0.49, Aa=2pq=0.42, aa=q²=0.09로 계산되며 p²+2pq+q²=1이 성립한다. 유전자 빈도가 유지되면 다음 세대에서도 유전자형 빈도는 일정하게 유지된다. [A]

(중략)
}
\begin{kfQBlock}{04}{객관식}
\kfQStem{[A]를 바탕으로 <보기>를 이해한 내용으로 적절하지 \underline{않은} 것은?}
\begin{kfEvidence}X-연관 열성 질환에서 남성 환자 빈도가 9\%(q=0.09)일 때, p=0.91이며 여성 보인자 빈도는 2pq, 여성 환자 빈도는 q²로 계산된다.\end{kfEvidence}
\begin{kfChoices}
\kfChoice 여성 환자 빈도가 남성 환자 빈도보다 높다.
\kfChoice 남성 표현형 빈도를 q로 둘 수 있다.
\kfChoice 여성 정상 동형 접합자 빈도는 p²로 계산할 수 있다.
\end{kfChoices}
\end{kfQBlock}

\kfPassageNote{문항 5 / 인과 도치}{%
지구의 과거 자기를 연구하는 고지자기 연구는 1940년대 말 저명한 실험 물리학자인 블래킷의 가설에 의해 학계의 관심을 끌었다. 블래킷은 지구를 포함해서 모든 회전하는 물체는 그것의 각운동량에 비례하는 자기장을 주변에 형성한다고 주장하였다. 이 주장이 맞다면 지구의 자전 방향이 일정할 경우 지구 자기장의 방향은 바뀔 수가 없다. 지구 자기장의 방향이 바뀌었는지 여부는 암석에 잔류하는 자기, 즉 잔류 지자기의 미약한 흔적을 통해 확인할 수 있었다. 왜냐하면 과거에 용암이나 마그마가 굳어서 형성된 바위에는 용암이나 마그마가 식기 전에 그 안의 자성 물질이 지구 자기장의 방향을 따라 정렬된 채 남아 있기 때문이다. 측정 결과, 지자기의 방향이 바뀐다는 것이 드러나자 지구의 자전 방향이 바뀔 가능성이 없다는 판단에 따라 블래킷의 가설은 반박되었지만 잔류 지자기에 대한 연구자들의 관심은 지속되었다.

(중략)

잔류 지자기의 방향이 지질 시대의 다른 시점에서 조금씩 다르게 나타나는 현상을 설명하기 위해 대

(중략)
}
\begin{kfQBlock}{05}{객관식}
\kfQStem{㉠에 대한 이해로 가장 적절한 것은?}
\begin{kfChoices}
\kfChoice 시료가 형성 후 외부 교란이 적을수록 포타슘-아르곤 비율 해석의 신뢰도가 높아진다.
\kfChoice 포타슘 붕괴율은 조건에 따라 크게 변해야 연대 계산이 가능하다.
\kfChoice 초기 아르곤이 완전히 남아 있어야 암석 나이를 정확히 구할 수 있다.
\end{kfChoices}
\end{kfQBlock}

\kfPassageNote{문항 6 / 의도·목적}{%
무선 통신에 대한 수요가 폭발적으로 증가함에 따라 미국 정부는 \uline{ⓐ경매}를 통해 통신 사업자들에게 통신 주파수 대역들을 판매하게 되었다. 경매가 시행되기 전에는 \uline{ⓑ심사}나 \uline{ⓒ제비뽑기}와 같은 방식을 활용하였다. 심사 방식은 여러 통신 사업자들이 사업 계획서를 제출하면 평가자들이 사업 계획서에 근거하여 무선 통신 사업을 가장 잘 운영할 만한 사업자를 선정하고, 그렇게 선정된 사업자가 주파수를 할당받는 것이었다. 그러나 사업자가 고용한 법률가들과 로비스트들이 서로 자기들의 사업 계획이 최선이라고 주장하는 상황에서 사업 계획서에 대한 평가 절차와 이의 제기 과정을 전부 거쳐 주파수를 할당하기까지 수년이 걸리기도 하였다.

1982년 무렵에는 더 이상 심사 방식으로 수많은 주파수를 할당하는 것이 어려운 수준에 이르렀고, 결국 의회는 제비뽑기로 주파수를 판매하는 것에 동의하게 되었다. 제비뽑기 방식은 주파수를 할당하는 데 드는 시간과 노력을 대폭 줄여 주었지만 …(중략)
}
\begin{kfQBlock}{06}{객관식}
\kfQStem{윗글에 대한 이해로 가장 적절한 것은?}
\begin{kfChoices}
\kfChoice 동시 다중 라운드 경매는 라운드별 정보 공개를 통해 참여자의 전략 수정을 가능하게 했다.
\kfChoice 동시 다중 라운드 경매는 라운드별 정보 공개를 통해 참여자의 전략 수정을 가능하게 했다.
\kfChoice 심사 방식은 주파수 대량 할당에 시간이 거의 들지 않았다.
\kfChoice 제비뽑기 방식에서는 투기 목적 참가가 구조적으로 차단되었다.
\end{kfChoices}
\end{kfQBlock}

\kfPassageNote{문항 7 / 상관관계}{%
근대 경제학은 국민 소득, 물가, 고용, 경기 변동 등 경제의 \uline{거시적 현상}을 설명하기 위해 다양한 \uline{모형}을 \uline{ⓐ발전시켜} 왔다. \uline{고전학파}는 시장의 \uline{자율 조정} 능력을 강조하며, 가격과 임금의 \uline{유연성}을 전제로 \uline{완전 고용}이 자연스럽게 달성된다고 보았다. 그러나 20세기 초 \uline{대공황}을 계기로 시장의 자율 조정에 한계가 있다는 비판이 제기되었고, 이는 \uline{케인스 경제학}의 등장으로 이어졌다. 케인스는 \uline{총수요}(한 경제 내에서 주어진 기간과 물가 수준하에 최종 재화와 서비스에 대한 수요의 총합)가 자동으로 \uline{총공급}(주어진 기간과 물가 수준하에 국가 경제에서 재화와 서비스에 대한 공급의 총합)에 맞춰진다고 보는 고전학파의 이론을 비판하며, 낮은 총수요가 소득의 감소 및 실업률 상승을 초래할 수 있다고 보았다. 따라서 그는 정부 지출 확대 등 정책 개입을 통해 총수요를 관리해야…(중략)
}
\begin{kfQBlock}{07}{객관식}
\kfQStem{다음 <보기>의 상황을 IS-MP 모형으로 이해한 내용으로 가장 적절한 것은?}
\begin{kfEvidence}다음 그래프는 현재 경제의 균형 상태(E)를 보여 주는 IS-MP 모형이다. 이후 중앙은행이 경기 부양을 위해 완화적 정책 기조로 전환하였다.

[그래프] 세로축: 이자율, 가로축: 국민 소득. 우하향하는 IS 곡선과 우상향하는 MP 곡선이 교차하며, 교차점이 E이다.\end{kfEvidence}
\begin{kfChoices}
\kfChoice 중앙은행이 완화적 기조로 전환하면 MP 곡선은 아래쪽으로 이동하여 새로운 균형이 형성될 수 있다.
\kfChoice 중앙은행이 완화적 기조로 전환하면 MP 곡선은 아래쪽으로 이동하여 새로운 균형이 형성될 수 있다.
\kfChoice 중앙은행이 완화적 기조로 전환하면 MP 곡선은 위쪽으로 이동할 것이다.
\kfChoice 중앙은행의 정책 기조 변화는 내생적 변화이므로 MP 곡선 위에서 점만 이동할 것이다.
\end{kfChoices}
\end{kfQBlock}

\kfPassageNote{문항 8 / 필요·충분 조건}{%
육상 동물과 수생 동물은 각각의 생활 환경에 따라 호흡 과정에서 어려움을 겪는다. 육상 동물은 폐의 표면에 건조한 공기가 닿으면 폐가 건조해질 수 있는데, 이는 폐를 손상시키고 \uline{ⓐ탈수}를 일으키며 기체를 용해하는 폐의 능력을 떨어뜨린다. 한편 수생 동물은 가용 산소가 적은 수중 환경과 수온 차에 의해 변화되는 산소 가용성에 적응해야 한다. 또한 수생 동물은 물속에서 몸을 움직이는 데 많은 근육 운동이 필요하고, 찬 바닷물의 높은 비열 때문에 아가미의 열을 빼앗기며, 삼투압 현상으로 인해 수분 균형에 문제가 발생하기도 한다.

이러한 환경 속에서 수생 동물은 효과적인 호흡을 위해 아가미라는 특수한 기관을 사용한다. 아가미는 물속에 녹아 있는 산소를 흡수하고 체내에 있는 이산화 탄소를 배출하는 기능을 하는데, 외아가미와 내아가미로 나뉜다. 외아가미는 몸 밖으로 \uline{ⓑ돌출}된 형태로, 표면적이 넓고 형태가 다양하며 유생기 도롱뇽이나 일부 무척추동물에서 나타난다. 외아가미…(중략)
}
\begin{kfQBlock}{08}{객관식}
\kfQStem{윗글의 역류 관계에 대한 이해로 가장 적절한 것은?}
\begin{kfChoices}
\kfChoice 역류 관계에서는 물이 이동하며 계속 산소가 낮은 혈액과 접촉하므로 확산이 계속될 수 있다.
\kfChoice 역류 관계에서는 물이 이동하며 계속 산소가 낮은 혈액과 접촉하므로 확산이 계속될 수 있다.
\kfChoice 물과 혈액이 같은 방향으로 흐를 때 압력 기울기가 오래 유지되어 교환이 최대화된다.
\kfChoice 역류 관계에서는 물의 흐름 때문에 혈액이 구심성에서 원심성으로 흐르지 못하게 된다.
\end{kfChoices}
\end{kfQBlock}

\kfPassageNote{문항 9 / 결합·분리}{%
(가) 가상 공간에서 영생의 삶을 누린다는 흥미로운 이야기가 SF 영화의 소재로만 그치지 않는다고 주장하는 입장이 있다. 커즈와일이나 보스트룀은 기억, 가치관, 태도, 정서적 성향처럼 인간의 자아를 구성하는 것들을 특정 정보 패턴으로 만들어 보존할 수 있다면 신체적인 죽음 이후에도 인간이 여전히 생존할 수 있다고 주장한다. 이 입장은 인격 동일성에 관한 심리적 연속성 이론과 관련이 있다. 이 이론에서는 ‘나’가 누구인지를 규정하는 것은 ‘나’의 기억, 신념, 생각, 감정, 희망, 두려움 등의 묶음이라고 보며, 이러한 심리적인 특성들의 집합이 유지·보존되느냐에 ‘나’의 지속 여부가 달려 있다고 말한다.

커즈와일이나 보스트룀은 마음이나 정신 상태를 의미하는 심적 상태의 본성을 계산 기능주의의 관점에서 이해한다. 계산 기능주의에서는 심적 상태의 독특한 인과적 역할이 두뇌의 물리적 상태에 의해서 이루어진다고 보기 때문에 정신 작용을 두뇌에 구현된 계산 체계의 작동으로 이해하며, 인간의 사…(중략)
}
\begin{kfQBlock}{09}{객관식}
\kfQStem{윗글의 내용과 일치하는 것으로 가장 적절한 것은?}
\begin{kfChoices}
\kfChoice 복수 실현 가능성을 인정해도 물리적 동일성의 자동 보장은 성립하지 않는다고 본다.
\kfChoice 복수 실현 가능성을 인정해도 물리적 동일성의 자동 보장은 성립하지 않는다고 본다.
\kfChoice 계산 기능주의는 정신 상태를 물질 구성과 무관한 순수 도덕 범주로 본다.
\kfChoice 체화된 인지는 신체 종류가 달라도 개념 이해 방식은 항상 같다고 본다.
\end{kfChoices}
\end{kfQBlock}

\kfPassageNote{문항 10 / 조건 왜곡}{%
컴퓨터는 사람의 언어를 직접 이해할 수 없다. 따라서 컴퓨터가 사람의 언어를 이해할 수 있도록 표현해 주어야 한다. 컴퓨터가 사람의 언어를 이해할 수 있도록 표현하는 방법에는 여러 가지가 있을 수 있는데, 그중 \uline{㉠원-핫 인코딩}은 표현하고 싶은 단어의 색인 값에 1을 부여하고 나머지에는 0을 부여하는 방법이다. 원-핫 인코딩을 수행하기 위해서는 먼저 대상이 되는 텍스트의 단어 집합을 만든다. 예를 들어 ‘오늘 날씨 정말 좋다.’라는 한 문장으로 이루어진 텍스트를 대상으로 한다고 가정하면, 단어 집합은 \{오늘, 날씨, 정말, 좋다\}가 되고 집합의 원소가 되는 단어의 개수가 4개이므로 단어 집합의 크기는 4가 된다. 단어 집합의 크기에 따라 각각의 단어에 [1, 0, 0, 0], [0, 1, 0, 0], [0, 0, 1, 0], [0, 0, 0, 1]과 같이 4차원의 고유한 숫자를 부여하는 것이 원-핫 인코딩이다. 원-핫 인코딩은 직관적이고 단순하지만 단어들이 개별적, 독립적…(중략)
}
\begin{kfQBlock}{10}{객관식}
\kfQStem{윗글을 바탕으로 <보기>를 이해한 내용으로 적절하지 \underline{않은} 것은?}
\begin{kfEvidence}대상 텍스트: 오늘 기분 정말 좋다. 오늘 날씨 정말 좋다. 오늘 날씨 매우 나쁘다. \\\relax
단어 집합: 오늘, 기분, 날씨, 정말, 매우, 좋다, 나쁘다 \\\relax
가정: [A] 방식이며 이웃 단어 범위는 한 문장 내\end{kfEvidence}
\begin{kfChoices}
\kfChoice 이웃 범위를 1로 고려하면 ‘오늘’의 이웃 단어는 ‘기분·날씨·정말·매우·좋다·나쁘다’가 될 수 있겠군.
\kfChoice 이웃 범위를 1로 고려하면 ‘오늘’의 이웃 단어는 ‘기분·날씨·정말·매우·좋다·나쁘다’가 될 수 있겠군.
\kfChoice ‘오늘 기분 매우 나쁘다.’라는 문장을 대상 텍스트에 추가하여도 단어 집합은 달라지지 않겠군.
\kfChoice 단어 집합의 크기가 7이므로 [1,0,0,0,0,0,0] 같은 7차원 숫자를 부여할 수 있겠군.
\end{kfChoices}
\end{kfQBlock}

\kfPassageNote{문항 11 / 긍부정 반전}{%
\#\#\# [31 \textasciitilde{} 34] 다음 글을 읽고 물음에 답하시오.

(가)   지팡이 짚고 바람 쐬며 좌우를 돌아보니   누대의 맑은 경치 아마도 깨끗하구나.   ㉠ 물도 하늘 같고 하늘도 물 같으니   푸른 물과 긴 하늘이 한빛이 되었거든   물가에 갈매기는 오는 듯 가는 듯 그칠 줄을 모르네.   ㉡ 바위 위 산꽃은 수놓은 병풍 되었고   시냇가 버들은 초록 장막 되었는데,   좋은 날 좋은 경치 나 혼자 거느리고   ㉢ 꽃피는 시절 허송하지 말리라 하고   아이 불러 하는 말, 이 깊은 산속에서 해산물을 볼쏘냐.   ㉣ 살진 고사리, 향기로운 당귀를 돼지고기, 사슴고기 섞어서   크나큰 바구니에 흡족히 담아두고   붕어회에다 눌어, 꿩 섞어 먹음직하게 구워지거든   술동이의 맑은 술을 술잔에 가득 부어   한잔, 또 한잔 취토록 먹은 후에,   ㉤ 복숭아꽃 붉은 비 되어 취한 낯에 뿌리는데   낚시터 넓은 돌을 높이 베고 누우니   무회씨 때 사람인가, 갈천씨 때 백성*인가.*   *…(중략)
}
\begin{kfQBlock}{11}{객관식}
\kfQStem{(나)의 ‘빛나는 생명의 예술가’가 갖추어야 할 태도로 가장 적절한 것은?}
\begin{kfChoices}
\kfChoice 직관을 통해 자연에 대한 솔직한 감각을 드러낼 수 있어야 한다.
\kfChoice 직관을 통해 자연에 대한 솔직한 감각을 드러낼 수 있어야 한다.
\kfChoice 자연의 모든 것을 알아낼 수 있다는 확신으로 탐구에 임해야 한다.
\kfChoice 여러 기록을 참고하며 자연의 새로운 경지를 소개할 수 있어야 한다.
\end{kfChoices}
\end{kfQBlock}

\kfPassageNote{문항 12 / 비교 대상}{%
[16\textasciitilde{}20] 다음 글을 읽고 물음에 답하시오.

(가)

어버이 낳으시고 임금이 먹이시니    낳은 덕 먹인 은을 다 갚고자 하였더니    숙연히 칠십이 넘으니 할 일 없어 하노라

<제1수>

연하에 깊이 곤 병 약이 효험 없어    강호에 버리언져 십 년 밖이 되었어라    그러나 이제 다 못 죽음도 그 성은인가 하노라

<제3수>

달 밝고 바람 자니 물결이 비단일다    단정*을 비껴 놓아 오락가락 하는 흥을    ⓐ 백구야 하 즐겨 말고려 세상 알까 하노라

<제5수>

피 소주 무우저리* 우습다 어른 대접    남은셔 이른 말이 초초타* 하건마는    두어라 이도 내 분이니 분내사*인가 하노라

<제8수>

- 나위소, 「강호구가」 -

* 단정: 작은 배. * 피 소주 무우저리: 풀로 만든 소주와 무절임. * 초초타: 보잘것없이 초라하다. * 분내사: 분수에 맞는 일.

(나)

들은 지 오래더니 보았구나 백상루    야속하다 강산이 날 기다리고 있었던가    처음으로…(중략)
}
\begin{kfQBlock}{12}{객관식}
\kfQStem{(가)\textasciitilde{}(다)에 대한 설명으로 가장 적절한 것은?}
\begin{kfChoices}
\kfChoice (가)와 (다)는 비유적 표현을 통해 대상의 속성을 드러내고 있다.
\kfChoice (가)와 (다)는 비유적 표현을 통해 대상의 속성을 드러내고 있다.
\kfChoice (가)는 색채어의 대비를 통해 대상을 생생하게 제시하고 있다.
\kfChoice (나)는 명령형 어미를 통해 시적 긴장감을 조성하고 있다.
\end{kfChoices}
\end{kfQBlock}

\kfPassageNote{문항 13 / 주체·객체}{%
[35 \textasciitilde{} 39] 다음 글을 읽고 물음에 답하시오.

(가)   강호에 봄이 드니 미친 흥이 절로 난다   시냇가 막걸리에 쏘가리 안주로다   이 몸이 한가한 것도 역시 임금의 은혜로다

㉠ 강호에 여름이 드니 초당에 일이 없다   미더운 강 물결이 보내는 것은 바람이로다   이 몸이 서늘한 것도 역시 임금의 은혜로다

강호에 가을이 드니 고기마다 살져 있다   조그마한 배에 그물 실어 흐르게 던져두고   이 몸이 소일하는 것도 역시 임금의 은혜로다

강호에 겨울이 드니 눈 깊이 자가 넘다   삿갓 비껴쓰고 도롱이로 옷을 삼아   이 몸이 춥지 않은 것도 역시 임금의 은혜로다

* 맹사성, ｢강호사시가｣ -

(나)   이보게 이웃 사람들아 산수구경 가자꾸나   산책은 오늘하고 목욕은 내일하세   아침에 나물캐고 저녁에 낚시하세   ㉡ 이제 막 익은 술을 갈건으로 걸러놓고   꽃나무 가지 꺾어 잔을 세면서 먹으리라   화풍(和風)이 문득 불어 시내를 건너오니   청향(淸香)은 잔에 …(중략)
}
\begin{kfQBlock}{13}{객관식}
\kfQStem{㉠ \textasciitilde{} ㉤에 대한 설명으로 적절하지 \underline{않은} 것은?}
\begin{kfChoices}
\kfChoice ㉡ : 자연과 동화되고 싶은 화자의 바람이 드러나 있다.
\kfChoice ㉡ : 자연과 동화되고 싶은 화자의 바람이 드러나 있다.
\kfChoice ㉠ : 여름날 한가한 초당의 모습이 드러나 있다.
\kfChoice ㉢ : 변화된 들판을 보며 감탄하는 화자의 모습이 드러나 있다.
\end{kfChoices}
\end{kfQBlock}

\kfPassageNote{문항 14 / 내용·방법}{%
\#\#\# [18～21] 다음 글을 읽고 물음에 답하시오.

황상과 만조백관이 어찌할 줄 모르더니 좌장군 서경태가 급히 입직군을 동원하여 칼을 들고 내달아 크게 꾸짖길,   “이 몹쓸 흉악한 놈아, 어찌 이런 변을 짓느냐?”   하고 칼을 들어 치니 아귀가 몸을 기울여 피하고 입을 벌려 숨을 들이쉬니 서경태가 날리어 아귀 입으로 들어갔다. 상이 보시다가 크게 놀라,   “짐이 여러 번 전장을 지내었으되 이런 일은 보도 듣도 못하였으니 제신 중에 뉘 이 짐승을 잡아 짐의 한을 씻으리오.”   정서장군 한세충이 나와 아뢰길,   “소장이 비록 재주 없으나 저것을 베어 황상께 바치리이다.”   하고 황금 투구에 엄신갑을 입고 팔 척 장창을 들고 청룡마를 내달아 외쳐 말하길,   “흉적은 목을 늘여 내 칼을 받으라.”   아귀가 크게 웃고 말하길,   “아까는 내 숨을 들이쉬니 모기 같은 것도 삼켰으니 지금은 숨을 내쉴 것이니 네 눈을 부릅뜨고 자세히 보라.”   하고 입을 벌려 숨을 내부니 …(중략)
}
\begin{kfQBlock}{14}{객관식}
\kfQStem{[A]의 서술상 특징에 대한 설명으로 가장 적절한 것은?}
\begin{kfChoices}
\kfChoice 대화를 통해 인물 간의 위계나 관계를 보여 주고 있다.
\kfChoice 대화를 통해 인물 간의 위계나 관계를 보여 주고 있다.
\kfChoice 서술자가 개입하여 인물에 대한 평가를 제시하고 있다.
\kfChoice 현재와 과거를 교차하여 장면의 전환을 보여 주고 있다.
\end{kfChoices}
\end{kfQBlock}

\kfPassageNote{문항 15 / 내용 왜곡}{%
[24 \textasciitilde{} 26] 다음 글을 읽고 물음에 답하시오.

(가)   바닷가 햇빛 바른 바위 위에   습한 간을 펴서 말리우자,

코카서스 산중에서 도망해 온 토끼처럼   둘러리를 ㉠ 빙빙 돌며 간을 지키자,

내가 오래 기르던 여윈 독수리야!   와서 뜯어 먹어라, 시름없이

너는 살찌고   나는 ㉡ 여위어야지, 그러나,

거북이야!   다시는 용궁의 유혹에 안 떨어진다.

프로메테우스 불쌍한 프로메테우스   불 도적한 죄로 목에 맷돌을 달고   끝없이 침전하는 프로메테우스.

* 윤동주, ｢간｣ -

(나)   큰일났다. 가만히 있어도 목구멍으로   시가 술술 쏟아져 나오니.

천기누설이다.

머리에 이가 있고   거북 등처럼 손이 튼 계집애가   제 짝이라는 것을   누군 모르랴.

그런데 감히 여왕을 사모함은   전생에 지은 이 무슨 아름다운 업보인가.

세상에 못 맺을 사랑이란 없다는 것을   떠꺼머리, 너는 ㉢ 무엄하게도 알아 버렸구나.

길 비켜라.   사랑이 사랑을 찾아간다…(중략)
}
\begin{kfQBlock}{15}{객관식}
\kfQStem{㉠ \textasciitilde{} ㉤에 대한 이해로 가장 적절한 것은?}
\begin{kfChoices}
\kfChoice ㉠을 활용하여 소중한 대상을 지키려는 의지를 드러낸다.
\kfChoice ㉠을 활용하여 소중한 대상을 지키려는 의지를 드러낸다.
\kfChoice ㉡을 활용하여 현재 상황에서 벗어날 수 없는 절망감을 드러낸다.
\kfChoice ㉢을 활용하여 사랑의 진리를 깨닫지 못한 이에 대한 질책을 드러낸다.
\end{kfChoices}
\end{kfQBlock}

\kfPassageNote{문항 16 / 과잉 해석}{%
\#\#\# 2025년 11월 고3 수능

\#\#\# [22\textasciitilde{}26] 다음 글을 읽고 물음에 답하시오.

(가)

두고 온 것들이 빛나는 때가 있다   빛나는 때를 위해 소금을 뿌리며   우리는 이 저녁을 떠돌고 있는가   사방을 둘러보아도   등불 하나 켜 든 이 보이지 않고   등불 뒤에 속삭이며 밤을 지키는   발자국 소리 들리지 않는다   잊혀진 목소리가 살아나는 때가 있다   ㉠ 잊혀진 한 목소리 잊혀진 다른 목소리의 끝을 찾아   목메이게 부르짖다 잦아드는 때가 있다   ㉡ 잦아드는 외마디 소리를 찾아 칼날 세우고   우리는 이 새벽길 숨가쁘게 넘고 있는가   하늘 올려보아도   함께 어둠 지새던 별 하나 눈뜨지 않는다   그래도 두고 온 것들은 빛나는가   빛을 뿜으면서 한 번은 되살아나는가   우리가 뿌린 소금들 반짝반짝 별빛이 되어   오던 길 환히 비춰 주고 있으니

- 이시영, 「그리움」 -

(나)

감나무 잎새를 흔드는 게   어찌 바람뿐이랴.   감나무 잎새를 반짝이는…(중략)
}
\begin{kfQBlock}{16}{객관식}
\kfQStem{<보기>를 참고하여 (다)를 감상한 내용으로 적절하지 \underline{않은} 것은?}
\begin{kfChoices}
\kfChoice '소리'가 지닌 상반된 특성들이 서로 균형을 이루어야 '좋은 소리'임을 제시하여, 문장이 궁극적으로 도달해야 할 바를 드러내고 있군.
\kfChoice '소리'가 지닌 상반된 특성들이 서로 균형을 이루어야 '좋은 소리'임을 제시하여, 문장이 궁극적으로 도달해야 할 바를 드러내고 있군.
\kfChoice '만물'이 소리 나는 이치에서 시작하여 '사람'이 소리를 내는 이치를 밝히며, 소리를 화두로 삼아 문장에 대해 말하고 있군.
\kfChoice 최립의 문장이 완성된 것은 아니지만 '참으로 어려운 일'에 가까움을 언급하며, 그의 문장에 대한 평가를 드러내고 있군.
\end{kfChoices}
\end{kfQBlock}

\kfPassageNote{문항 17 / 범위 오류}{%
\#\#\# 2024년 11월 고3 수능

[22-27] 다음 글을 읽고 물음에 답하시오.

(가)   배를 민다   배를 밀어보는 것은 아주 드문 경험   희번덕이는 잔잔한 가을 바닷물 위에   배를 밀어넣고는   온몸이 아주 추락하지 않을 순간의 한 허공에서   밀던 힘을 한껏 더해 밀어주고는   아슬아슬히 배에서 떨어진 손, 순간 환해진 손을   허공으로부터 거둔다

사랑은 참 부드럽게도 떠나지   뵈지도 않는 길을 부드럽게도

배를 한껏 세게 밀어내듯이 슬픔도   그렇게 밀어내는 것이지

배가 나가고 남은 빈 물 위의 흉터   잠시 머물다 가라앉고

그런데 오, 내 안으로 들어오는 배여   아무 소리 없이 밀려들어오는 배여   - 장석남, <배를 밀며>

(나)   당신……, 당신이라는 말 참 좋지요, 그래서 불러봅니다 킥킥거리며 한때 적요로움의 울음이 있었던 때, 한 슬픔이 문을 닫으면 또 한 슬픔이 문을 여는 것을 이만큼 살아옴의 상처에 기대, 나 킥킥……, 당신을 부릅니다 단…(중략)
}
\begin{kfQBlock}{17}{객관식}
\kfQStem{문항별}
\begin{kfChoices}
\kfChoice 번은 '통제할 수 없는 익명의 욕구'를 상대를 향한 사랑이 운명적이어서 멈출 수 없다는 뜻으로 해석했습니다.
\kfChoice 번은 '통제할 수 없는 익명의 욕구'를 상대를 향한 사랑이 운명적이어서 멈출 수 없다는 뜻으로 해석했습니다.
\kfChoice 번은 배가 들어오는 것을 '대상과의 재회가 예상대로 이루어짐'이라고 해석했습니다.
\kfChoice 번은 '당신'을 화자의 내면에 사는 '병자'로서 연민의 대상이라고 설명했습니다.
\end{kfChoices}
\end{kfQBlock}

\kfPassageNote{문항 18 / 시점·화자}{%
[16\textasciitilde{}20] 다음 글을 읽고 물음에 답하시오.

(가)

어버이 낳으시고 임금이 먹이시니    낳은 덕 먹인 은을 다 갚고자 하였더니    숙연히 칠십이 넘으니 할 일 없어 하노라

<제1수>

연하에 깊이 곤 병 약이 효험 없어    강호에 버리언져 십 년 밖이 되었어라    그러나 이제 다 못 죽음도 그 성은인가 하노라

<제3수>

달 밝고 바람 자니 물결이 비단일다    단정*을 비껴 놓아 오락가락 하는 흥을    ⓐ 백구야 하 즐겨 말고려 세상 알까 하노라

<제5수>

피 소주 무우저리* 우습다 어른 대접    남은셔 이른 말이 초초타* 하건마는    두어라 이도 내 분이니 분내사*인가 하노라

<제8수>

- 나위소, 「강호구가」 -

* 단정: 작은 배. * 피 소주 무우저리: 풀로 만든 소주와 무절임. * 초초타: 보잘것없이 초라하다. * 분내사: 분수에 맞는 일.

(나)

들은 지 오래더니 보았구나 백상루    야속하다 강산이 날 기다리고 있었던가    처음으로…(중략)
}
\begin{kfQBlock}{18}{객관식}
\kfQStem{ⓐ와 ⓑ에 대한 설명으로 가장 적절한 것은?}
\begin{kfChoices}
\kfChoice ⓐ는 화자의 즐거움이 투영된, ⓑ는 글쓴이의 반가움을 유발하는 대상이다.
\kfChoice ⓐ는 화자의 즐거움이 투영된, ⓑ는 글쓴이의 반가움을 유발하는 대상이다.
\kfChoice ⓐ는 화자의 내적 성찰을 일으키는, ⓑ는 글쓴이의 경험을 환기하는 대상이다.
\kfChoice ⓐ는 화자에게 세월의 흐름을, ⓑ는 글쓴이에게 계절의 변화를 인식하게 하는 대상이다.
\end{kfChoices}
\end{kfQBlock}



\clearpage

\kfChapterCover{3}{비트겐슈타인1 ch03}{Chapter 3}{이번 챕터: 문법 → 문학 5작품 → 비문학 5지문 → 패턴 워크북.}
\begin{kfChapterAreaList}
\kfChapterAreaItem{문법}{kfAreaGrammar}{음운 / 비음화·유음화}
\kfChapterAreaItem{문학}{kfAreaLiterature}{「상춘곡」 — 정극인 외 4편}
\kfChapterAreaItem{비문학}{kfAreaNonlit}{「공자의 인(仁)과 예(禮)」 외 4편}
\kfChapterAreaItem{실력 확인}{kfAreaWeekly}{패턴 워크북 — 총 18문항}
\end{kfChapterAreaList}

\kfStartGrammar{음운 / 비음화·유음화}


\kfSecHeader{kfAreaGrammar}{문제}{}

\kfBeginMC
\par\needspace{25\baselineskip}
\kfPassageFull{비음화 조건}{%
비음화는 비음이 아닌 자음이 비음으로 바뀌는 현상이다. 비음은 공기를 코로 내보내면서 내는 소리로, 'ㅁ, ㄴ, ㅇ'이 해당한다. 비음화가 일어나는 환경은 크게 세 가지이다. 첫째, 받침 'ㄱ, ㄷ, ㅂ'이 비음 'ㄴ, ㅁ' 앞에서 각각 'ㅇ, ㄴ, ㅁ'으로 바뀐다. 예를 들어 '국물'은 [궁물], '닫는다'는 [단는다], '밥물'은 [밤물]로 발음된다. 둘째, 받침 'ㄱ, ㄷ, ㅂ' 뒤에 'ㄹ'이 오면 'ㄹ'이 먼저 'ㄴ'으로 바뀌고, 이어서 앞의 'ㄱ, ㄷ, ㅂ'도 비음으로 바뀐다. 예를 들어 '독립'은 [동닙]으로 발음된다. 셋째, 받침 'ㅁ, ㅇ' 뒤에 'ㄹ'이 오면 'ㄹ'이 'ㄴ'으로 바뀐다.
}\par\nobreak\nopagebreak[4]\kfResetAreaFlag
\begin{kfQBlock}{01}{객관식}
\kfQStem{윗글을 바탕으로 할 때, 비음화가 일어나는 예로 적절하지 \underline{않은} 것은?}
\begin{kfChoices}
\kfChoice 먹는 → [멍는]
\kfChoice 합리 → [함니]
\kfChoice 종로 → [종노]
\kfChoice 칼날 → [칼랄]
\kfChoice 십리 → [심니]
\end{kfChoices}
\end{kfQBlock}

\par\needspace{25\baselineskip}
\kfPassageFull{유음화 조건}{%
유음화는 유음이 아닌 자음 'ㄴ'이 유음 'ㄹ' 앞이나 뒤에서 유음 'ㄹ'로 바뀌는 현상이다. 유음은 혀끝을 잇몸에 가볍게 대었다가 떼거나, 잇몸에 댄 채 공기를 흘려 내보내면서 내는 소리로, 국어에서는 'ㄹ' 하나뿐이다. 유음화의 환경은 두 가지이다. 첫째, 'ㄴ' 뒤에 'ㄹ'이 올 때 'ㄴ'이 'ㄹ'로 바뀐다. 예를 들어 '신라'는 [실라], '논리'는 [놀리]로 발음된다. 둘째, 'ㄹ' 뒤에 'ㄴ'이 올 때에도 'ㄴ'이 'ㄹ'로 바뀐다. 예를 들어 '칼날'은 [칼랄], '설날'은 [설랄]로 발음된다.
}\par\nobreak\nopagebreak[4]\kfResetAreaFlag
\begin{kfQBlock}{02}{객관식}
\kfQStem{윗글을 바탕으로 할 때, 유음화가 일어나는 단어를 <보기>에서 모두 고른 것은?}
\begin{kfEvidence}ㄱ. 천리 → [철리] \\\relax
ㄴ. 담력 → [담녁] \\\relax
ㄷ. 물난리 → [물랄리] \\\relax
ㄹ. 줄넘기 → [줄럼끼]\end{kfEvidence}
\begin{kfChoices}
\kfChoice ㄱ, ㄴ
\kfChoice ㄱ, ㄷ
\kfChoice ㄱ, ㄷ, ㄹ
\kfChoice ㄱ, ㄴ, ㄷ
\kfChoice ㄱ, ㄴ, ㄷ, ㄹ
\end{kfChoices}
\end{kfQBlock}

\par\needspace{25\baselineskip}
\kfPassageFull{비음화+유음화 복합}{%
하나의 단어에 비음화와 유음화가 복합적으로 적용되는 경우도 있다. 특히 겹받침이 포함된 단어에서 자음군 단순화 후 추가적인 음운 변동이 일어나는 경우를 주의해야 한다. 예를 들어 '뚫는'은 먼저 겹받침 'ㄾ'에서 자음군 단순화가 일어나 'ㅌ'이 탈락하고 'ㄹ'만 남은 뒤, 'ㄹ+ㄴ'에서 'ㄴ'이 'ㄹ'로 바뀌는 유음화가 적용되어 최종적으로 [뚤른]이 된다. 반면 '독립'은 'ㄱ+ㄹ'에서 'ㄹ'이 먼저 'ㄴ'으로 바뀌고(비음화), 'ㄱ'도 'ㅇ'으로 바뀌어(비음화) [동닙]이 된다. 이처럼 같은 'ㄹ'이 관여하더라도, 앞의 자음이 무엇이냐에 따라 비음화가 일어날 수도, 유음화가 일어날 수도 있다.
}\par\nobreak\nopagebreak[4]\kfResetAreaFlag
\begin{kfQBlock}{03}{객관식}
\kfQStem{윗글을 바탕으로 할 때, <보기>의 ㉠\textasciitilde{}㉢에 적용된 음운 변동으로 옳은 것은?}
\begin{kfEvidence}㉠ 앓는 → [알른] \\\relax
㉡ 백리 → [뱅니] \\\relax
㉢ 광한루 → [광할루]\end{kfEvidence}
\begin{kfChoices}
\kfChoice ㉠: 자음군 단순화 → 비음화
\kfChoice ㉠: 자음군 단순화 → 유음화
\kfChoice ㉡: 유음화만 적용
\kfChoice ㉢: 비음화만 적용
\kfChoice ㉠, ㉡, ㉢ 모두 비음화가 적용된다.
\end{kfChoices}
\end{kfQBlock}

\par\needspace{25\baselineskip}
\kfPassageFull{발음 표기}{%
비음화와 유음화는 모두 자음 동화에 해당하는 음운 변동이다. 동화란 인접한 음운의 영향을 받아 그와 비슷한 성질의 소리로 바뀌는 현상이다. 비음화에서는 비음이 아닌 자음이 비음의 영향을 받아 비음으로 바뀌고, 유음화에서는 유음이 아닌 'ㄴ'이 유음 'ㄹ'의 영향을 받아 유음 'ㄹ'로 바뀐다. 동화의 방향에 따라 순행 동화(앞 소리가 뒤 소리에 영향)와 역행 동화(뒤 소리가 앞 소리에 영향)로 구분할 수 있다. 예를 들어 '국물'[궁물]은 뒤의 'ㅁ'이 앞의 'ㄱ'에 영향을 주었으므로 역행 동화이고, '종로'[종노]는 앞의 'ㅇ'이 뒤의 'ㄹ'에 영향을 주었으므로 순행 동화이다.
}\par\nobreak\nopagebreak[4]\kfResetAreaFlag
\begin{kfQBlock}{04}{객관식}
\kfQStem{윗글을 바탕으로 할 때, <보기>의 밑줄 친 단어의 발음이 옳고, 동화의 방향도 바르게 짝지은 것은?}
\begin{kfEvidence}(가) 설날 [설랄] \\\relax
(나) 입는 [임는] \\\relax
(다) 신라 [실라]\end{kfEvidence}
\begin{kfChoices}
\kfChoice (가) — 역행 동화, (나) — 순행 동화
\kfChoice (가) — 순행 동화, (나) — 순행 동화
\kfChoice (가) — 순행 동화, (나) — 역행 동화
\kfChoice (나) — 역행 동화, (다) — 순행 동화
\kfChoice (가) — 역행 동화, (다) — 역행 동화
\end{kfChoices}
\end{kfQBlock}

\par\needspace{25\baselineskip}
\kfPassageFull{음운 변동 횟수}{%
음운 변동의 횟수를 세는 문제는 수능에서 자주 출제된다. 음운 변동 횟수를 셀 때는, 하나의 음운이 다른 음운으로 바뀌거나 탈락하거나 첨가되는 것을 각각 1회로 센다. 예를 들어 '꽃말'이 [꼰말]로 발음되는 과정에서는 ①'ㅊ'→[ㄷ](음절의 끝소리 규칙: 1회), ②[ㄷ]→[ㄴ](비음화: 1회)로 총 2회의 음운 변동이 일어난다. '독립'이 [동닙]으로 발음되는 과정에서는 ①'ㄹ'→[ㄴ](비음화: 1회), ②'ㄱ'→[ㅇ](비음화: 1회)로 총 2회의 음운 변동이 일어난다. 이때 하나의 규칙이 적용되더라도 바뀌는 음운의 수만큼 횟수를 세야 한다.
}\par\nobreak\nopagebreak[4]\kfResetAreaFlag
\begin{kfQBlock}{05}{객관식}
\kfQStem{윗글을 바탕으로 할 때, <보기>의 단어에서 일어나는 음운 변동의 총 횟수가 같은 것끼리 묶은 것은?}
\begin{kfEvidence}ㄱ. 꽃망울 [꼰망울] \\\relax
ㄴ. 놓는 [논는] \\\relax
ㄷ. 협력 [혐녁] \\\relax
ㄹ. 앞마당 [암마당]\end{kfEvidence}
\begin{kfChoices}
\kfChoice ㄱ과 ㄴ
\kfChoice ㄱ과 ㄷ
\kfChoice ㄴ과 ㄹ
\kfChoice ㄱ, ㄴ, ㄷ
\kfChoice ㄱ, ㄴ, ㄷ, ㄹ 모두 같다.
\end{kfChoices}
\end{kfQBlock}

\kfEndMC


\kfStartLit{「상춘곡」 — 정극인 외 4편}


\kfPassageNote{「상춘곡」 — 정극인}{%
紅塵에 묻힌 분네 이 내 生涯 어떠한고 \\[4pt]
녯 사람 風流를 미칠가 못 미칠가 \\[4pt]
天地間 男子 몸이 날만한 이 하건마는 \\[4pt]
山林에 묻혀 이셔 지럴 줄을 모르더니 \\[14pt]
수간모옥(數間茅屋)을 벽계수(碧溪水) 앞에 두고 \\[4pt]
松竹 울울리(鬱鬱裏)에 풍월주인(風月主人) 되어셔라 \\[4pt]
엇그제 겨울 지나 새 봄이 도라오니 \\[4pt]
桃花杏花는 석양리(夕陽裏)에 퓌여 있고 \\[4pt]
綠楊芳草는 細雨 中에 프르도다 \\[14pt]
칼로 물을 치니 소리 다시 흘러가고 \\[4pt]
술잔을 들어 시름 다 잊는다 \\[4pt]
功名도 날 꺼리고 富貴도 날 꺼리니 \\[4pt]
淸風에 月을 싣고 碧波에 배를 띄워 \\[4pt]
이 中에 시름 업스니 漁父의 生涯이라 \\[14pt]
닛태 百年行樂이 이만한들 어떠하리 \\[4pt]
점동새 울 제면 거러 보고 \\[4pt]
꽃 지고 새 울음에 봄 뜻이 적어 가니 \\[4pt]
아마도 林泉 한 興을 비길 곧이 업세라
}\kfResetAreaFlag
\par\kfMaybeNeedspace{50\baselineskip}\noindent{\kfTipMark\,\,}{\small\color{kfInk} 빈칸에 알맞은 말을 채워 넣으시오. (초성 힌트 참고)}\par\nopagebreak[4]\smallskip\nopagebreak[4]
\begin{kfActTable}{|>{\centering\arraybackslash}m{2.6cm}|m{6cm}|m{4cm}|}
\rowcolor{kfPrimaryLight!50}\textbf{\color{kfPrimary}항목} & \textbf{\color{kfPrimary}내용 (빈칸 채우기)} & \textbf{\color{kfPrimary}답안 작성}\\\hline
갈래 & ㄱㅅ (4음보 연속체) ① & \kfTBlank \\\hline
시대 & ㅈㅅ 전기 ② & \kfTBlank \\\hline
성격 & ㅈㅇㅊㅁㅈ, ㅍㄹㅈ ③ & \kfTBlank \\\hline
홍진의 의미 & 번잡한 ㅅㅅ 세계 ④ & \kfTBlank \\\hline
풍월주인 & 바람과 달의 ㅈㅇ, 자연의 주인 ⑤ & \kfTBlank \\\hline
주요 배경 & ㅂ 계절의 자연 풍경 ⑥ & \kfTBlank \\\hline
사상적 기반 & ㅇㅂㄴㄷ (가난해도 도를 즐김) ⑦ & \kfTBlank \\\hline
표현 기법 & ㄷㄱㅂ, ㅇㄱㅂ ⑧ & \kfTBlank \\\hline
어부의 생애 & 자연 속 시름 없는 삶의 ㅂㅇ ⑨ & \kfTBlank \\\hline
임천한흥 & ㅅ과 ㅅ의 한가로운 흥취 ⑩ & \kfTBlank \\\hline
주제 & ㅈㅇ ㅊㅁ과 ㅇㅂㄴㄷ의 삶 ⑪ & \kfTBlank \\\hline
문학사적 의의 & 우리말로 쓰인 ㅊㅊ의 ㄱㅅ 작품 ⑫ & \kfTBlank \\\hline
\end{kfActTable}
\kfSecHeader{kfAreaLiterature}{문제}{}

\kfBeginMC
\begin{kfQBlock}{01}{객관식}
\kfQStem{윗글에 대한 설명으로 적절하지 \underline{않은} 것은?}
\begin{kfChoices}
\kfChoice 4음보 연속체의 가사 형식으로 이루어져 있다.
\kfChoice 자연 속 삶의 만족감과 풍류를 노래하고 있다.
\kfChoice 계절의 변화를 통해 시간의 흐름을 나타내고 있다.
\kfChoice 화자는 공명과 부귀를 쟁취하려는 의지를 드러내고 있다.
\kfChoice 감각적 이미지를 활용하여 봄 풍경을 생동감 있게 그려 내고 있다.
\end{kfChoices}
\end{kfQBlock}

\begin{kfQBlock}{02}{객관식}
\kfQStem{'홍진에 묻힌 분네'에서 '홍진(紅塵)'이 상징하는 것으로 가장 적절한 것은?}
\begin{kfChoices}
\kfChoice 붉은 꽃이 흩날리는 아름다운 자연 풍경
\kfChoice 관직과 명예를 둘러싼 번잡한 세속 세계
\kfChoice 전쟁으로 황폐해진 민중의 삶의 터전
\kfChoice 사계절이 순환하는 자연의 이치
\kfChoice 화자가 동경하지만 도달할 수 없는 이상 세계
\end{kfChoices}
\end{kfQBlock}

\begin{kfQBlock}{03}{객관식}
\kfQStem{화자가 '어부의 생애'를 언급한 의도로 가장 적절한 것은?}
\begin{kfChoices}
\kfChoice 고기잡이를 생업으로 삼아 경제적 자립을 이루겠다는 뜻
\kfChoice 어부처럼 강호에 묻혀 근심 없이 자연을 즐기겠다는 뜻
\kfChoice 어부의 고된 노동을 통해 삶의 의미를 깨닫겠다는 뜻
\kfChoice 은퇴 후 낚시터에서 여생을 보내겠다는 현실적 계획
\kfChoice 바다에 나가 미지의 세계를 탐험하겠다는 모험 정신
\end{kfChoices}
\end{kfQBlock}

\begin{kfQBlock}{04}{객관식}
\kfQStem{윗글에 사용된 표현 기법에 대한 설명으로 적절하지 \underline{않은} 것은?}
\begin{kfChoices}
\kfChoice '도화행화'와 '녹양방초'의 나열을 통해 봄 풍경을 풍부하게 묘사하고 있다.
\kfChoice '석양리에 퓌여 있고'와 '세우 중에 프르도다'는 대구의 형식을 취하고 있다.
\kfChoice '청풍에 월을 싣고'에서 바람에 달을 싣는다는 비유적 표현을 사용하고 있다.
\kfChoice '칼로 물을 치니'에서 시각과 청각의 감각적 전이(공감각)를 활용하고 있다.
\kfChoice 한문 어구와 우리말을 혼용하여 격조 있는 문체를 형성하고 있다.
\end{kfChoices}
\end{kfQBlock}

\begin{kfQBlock}{05}{객관식}
\kfQStem{윗글을 바탕으로 <보기>를 감상한 내용으로 적절하지 \underline{않은} 것은?}
\begin{kfEvidence}조선 전기 사림파 문인들은 관직에서 물러나거나 처음부터 벼슬에 나아가지 않고 자연에 은거하는 삶을 이상적으로 여겼다. 이들은 자연을 단순한 풍경이 아니라 도(道)가 구현된 공간으로 인식했으며, 자연 속에서의 삶을 통해 유교적 덕목인 안빈낙도를 실천하고자 했다. 이러한 자연관은 가사 문학의 중요한 주제 의식을 형성하였다.\end{kfEvidence}
\begin{kfChoices}
\kfChoice '산림에 묻혀'라는 표현은 <보기>에서 말한 자연 은거의 삶을 실천하는 모습으로 볼 수 있겠군.
\kfChoice '풍월주인 되어셔라'는 자연 속에서 풍류의 주체가 된 화자의 자부심을 드러낸 것이겠군.
\kfChoice '공명도 날 꺼리고 부귀도 날 꺼리니'는 안빈낙도를 실천하기 위해 세속적 가치를 멀리하는 태도이겠군.
\kfChoice '임천한흥(林泉閑興)'은 자연을 도가 구현된 공간으로 인식하며 그 속에서 느끼는 흥취를 가리키겠군.
\kfChoice 화자가 '녯 사람 풍류'를 언급한 것은 과거 불교 승려들의 수행 전통을 계승하겠다는 의지이겠군.
\end{kfChoices}
\end{kfQBlock}

\kfEndMC

\kfPassageNote{「오우가」 — 윤선도}{%
내 벗이 몇인가 하니 수석(水石)과 송죽(松竹)이라 \\[4pt]
동산(東山)에 달 오르니 그 더욱 반갑고야 \\[4pt]
두어라 이 다섯밖에 또 더하여 무엇하리 \\[14pt]
구름 빛이 좋다 하나 검기를 자로 한다 \\[4pt]
바람이 맑다 하나 그칠 적이 하노매라 \\[4pt]
좋고도 그칠 뜻이 없기는 물뿐인가 하노라 \\[14pt]
꽃은 무슨 일로 피면서 쉬이 지고 \\[4pt]
풀은 어이하여 푸르는 듯 누르나니 \\[4pt]
아마도 변치 아니할은 바위뿐인가 하노라 \\[14pt]
더우면 꽃 피고 추우면 잎 지거늘 \\[4pt]
솔아 너는 어찌 눈서리를 모르는다 \\[4pt]
구천(九泉)에 뿌리 곧은 줄을 글로 하여 아노라 \\[14pt]
나무도 아닌 것이 풀도 아닌 것이 \\[4pt]
곧기는 뉘 시키며 속은 어이 비었는다 \\[4pt]
저러고도 사시(四時)에 푸르니 그를 좋아하노라 \\[14pt]
작은 것이 높이 떠서 만물(萬物)을 다 비추니 \\[4pt]
밤중의 광명(光明)이 너만한 이 또 있느냐 \\[4pt]
보고도 말 아니하니 내 벗인가 하노라
}\kfResetAreaFlag
\par\kfMaybeNeedspace{50\baselineskip}\noindent{\kfTipMark\,\,}{\small\color{kfInk} 빈칸에 알맞은 말을 채워 넣으시오. (초성 힌트 참고)}\par\nopagebreak[4]\smallskip\nopagebreak[4]
\begin{kfActTable}{|>{\centering\arraybackslash}m{2.6cm}|m{6cm}|m{4cm}|}
\rowcolor{kfPrimaryLight!50}\textbf{\color{kfPrimary}항목} & \textbf{\color{kfPrimary}내용 (빈칸 채우기)} & \textbf{\color{kfPrimary}답안 작성}\\\hline
갈래 & ㅇㅅㅈ (6수 연작) ① & \kfTBlank \\\hline
시대 & ㅈㅅ 중기, 보길도 은거 시기 ② & \kfTBlank \\\hline
다섯 벗 & ㅁ·ㅂㅇ·ㅅㄴㅁ·ㄷㄴㅁ·ㄷ ③ & \kfTBlank \\\hline
물의 덕 & 쉬지 않고 흐르는 ㅅㅅ ④ & \kfTBlank \\\hline
바위의 덕 & 변치 않는 ㅎㄱㅅ ⑤ & \kfTBlank \\\hline
소나무의 덕 & 눈서리에도 푸른 ㅈㅈ ⑥ & \kfTBlank \\\hline
대나무의 덕 & 속 비고 곧은 ㄱㅈ ⑦ & \kfTBlank \\\hline
달의 덕 & 말없이 비추는 ㄱㅁ ⑧ & \kfTBlank \\\hline
표현 기법 & 자연물에 인격 부여 = ㅇㅎ ⑨ & \kfTBlank \\\hline
대비 구조 & 변하는 것(꽃·풀·구름 등) vs 변치 않는 ㅂ ⑩ & \kfTBlank \\\hline
사상적 기반 & ㅇㄱ적 ㄱㅈ론 ⑪ & \kfTBlank \\\hline
주제 & ㅈㅇ ㅇㅊ과 ㄱㅈ의 ㄷ ⑫ & \kfTBlank \\\hline
\end{kfActTable}
\kfSecHeader{kfAreaLiterature}{문제}{}

\kfBeginMC
\begin{kfQBlock}{01}{객관식}
\kfQStem{윗글에 대한 설명으로 적절한 것은?}
\begin{kfChoices}
\kfChoice 화자는 인간관계의 소중함을 다섯 친구의 우정을 통해 노래하고 있다.
\kfChoice 화자는 다섯 자연물을 벗으로 삼아 각각의 덕성을 예찬하고 있다.
\kfChoice 화자는 자연의 변화무쌍한 모습에서 인생의 무상함을 느끼고 있다.
\kfChoice 화자는 자연과 문명의 조화를 통해 이상 사회를 구현하고자 한다.
\kfChoice 화자는 자연물의 결함을 지적하며 인간의 우월함을 강조하고 있다.
\end{kfChoices}
\end{kfQBlock}

\begin{kfQBlock}{02}{객관식}
\kfQStem{제3수에서 '꽃'과 '풀'을 먼저 언급한 뒤 '바위'를 제시한 표현 방식의 효과로 가장 적절한 것은?}
\begin{kfChoices}
\kfChoice 꽃과 풀의 아름다움을 통해 바위의 초라함을 부각하려는 것이다.
\kfChoice 변하는 것(꽃·풀)과 대비하여 변하지 않는 바위의 항구성을 강조하려는 것이다.
\kfChoice 꽃과 풀이 바위보다 우월한 존재임을 역설적으로 표현하려는 것이다.
\kfChoice 세 자연물을 동등하게 나열하여 자연 전체의 조화로움을 보여 주려는 것이다.
\kfChoice 계절의 순환을 시간순으로 보여 주기 위한 서사적 구성이다.
\end{kfChoices}
\end{kfQBlock}

\begin{kfQBlock}{03}{객관식}
\kfQStem{제5수의 '나무도 \underline{아닌} 것이 풀도 \underline{아닌} 것이'에서 화자가 대나무를 이렇게 표현한 까닭으로 가장 적절한 것은?}
\begin{kfChoices}
\kfChoice 대나무가 나무와 풀 어디에도 속하지 못하는 존재라서 안타까워하는 것이다.
\kfChoice 대나무를 식물학적으로 정확히 분류하기 위한 학술적 서술이다.
\kfChoice 나무도 풀도 아니면서 사시에 푸른 독특한 속성을 부각하여 경탄하는 것이다.
\kfChoice 대나무가 자연에서 소외된 존재임을 강조하여 동정심을 유발하는 것이다.
\kfChoice 나무와 풀의 결합으로 새로운 생명체가 탄생했음을 노래하는 것이다.
\end{kfChoices}
\end{kfQBlock}

\begin{kfQBlock}{04}{객관식}
\kfQStem{제6수의 '보고도 말 아니하니 내 벗인가 하노라'에 담긴 화자의 태도로 가장 적절한 것은?}
\begin{kfChoices}
\kfChoice 달이 말을 하지 못하는 것에 대한 서운함과 아쉬움
\kfChoice 말없이 만물을 비추는 달의 과묵한 덕성에 대한 공감
\kfChoice 달이 자신의 비밀을 지켜 줄 것이라는 실용적 기대
\kfChoice 인간의 언어보다 침묵이 더 아름답다는 노장 사상의 표현
\kfChoice 달에게 말을 건네 봤으나 응답이 없어 체념하는 모습
\end{kfChoices}
\end{kfQBlock}

\begin{kfQBlock}{05}{객관식}
\kfQStem{윗글을 바탕으로 <보기>를 감상한 내용으로 적절하지 \underline{않은} 것은?}
\begin{kfEvidence}윤선도는 병자호란 이후 정치적 좌절을 겪고 보길도에 은거하며 자연을 벗 삼아 살았다. 그의 시조에 등장하는 자연물은 단순한 풍경이 아니라 유교적 '군자'의 덕목을 상징한다. 조선 시대 문인들은 자연물의 속성에서 인간이 본받아야 할 품성을 발견하는 물아일체(物我一體)의 자연관을 지녔다.\end{kfEvidence}
\begin{kfChoices}
\kfChoice 화자가 '수석과 송죽'을 벗으로 삼은 것은 정치적 좌절 이후 자연에서 위안을 찾는 모습이겠군.
\kfChoice '변치 아니할은 바위뿐'이라는 표현에는 불변하는 절개라는 군자의 덕목이 투영되어 있겠군.
\kfChoice '구천에 뿌리 곧은 줄'이라는 구절에서 소나무의 깊은 뿌리는 확고한 신념을 상징하겠군.
\kfChoice 구름과 바람을 먼저 언급하고 물을 대비시킨 것은 변하는 것과 변치 않는 것의 대비를 통한 강조이겠군.
\kfChoice 다섯 자연물이 각각 다른 덕목을 지닌 것은 자연이 인간보다 도덕적으로 열등함을 보여 주는 것이겠군.
\end{kfChoices}
\end{kfQBlock}

\kfEndMC

\kfPassageNote{「동짓달 기나긴 밤을」 — 황진이}{%
동짓달 기나긴 밤을 한 허리를 베어 내여 \\[4pt]
춘풍(春風) 니불 아래 서리서리 넣었다가 \\[4pt]
어론 님 오신 날 밤이여든 구뷔구뷔 펴리라
}\kfResetAreaFlag
\par\kfMaybeNeedspace{50\baselineskip}\noindent{\kfTipMark\,\,}{\small\color{kfInk} 빈칸에 알맞은 말을 채워 넣으시오. (초성 힌트 참고)}\par\nopagebreak[4]\smallskip\nopagebreak[4]
\begin{kfActTable}{|>{\centering\arraybackslash}m{2.6cm}|m{6cm}|m{4cm}|}
\rowcolor{kfPrimaryLight!50}\textbf{\color{kfPrimary}항목} & \textbf{\color{kfPrimary}내용 (빈칸 채우기)} & \textbf{\color{kfPrimary}답안 작성}\\\hline
갈래 & ㅍㅅㅈ ① & \kfTBlank \\\hline
시대 & ㅈㅅ 중기 (16세기) ② & \kfTBlank \\\hline
성격 & ㄱㅈㅈ, ㄴㅁㅈ ③ & \kfTBlank \\\hline
핵심 발상 & 추상적 ㅅㄱ을 구체적 ㅅㅁ로 변환 ④ & \kfTBlank \\\hline
초장 행위 & 밤을 한 허리 ㅂㅇ ㄴㅇ ⑤ & \kfTBlank \\\hline
보관 방식 & 춘풍 니불 아래 ㅅㄹㅅㄹ 넣음 ⑥ & \kfTBlank \\\hline
사용 시점 & ㄴ 오신 날 밤 ⑦ & \kfTBlank \\\hline
종장 행위 & ㄱㅂㄱㅂ 펴리라 ⑧ & \kfTBlank \\\hline
표현 기법 & ㄱㅊㅎ + ㄷㅂ적 구조 ⑨ & \kfTBlank \\\hline
정서 & 님에 대한 ㄱㄹ과 ㅈㅎ의 소망 ⑩ & \kfTBlank \\\hline
주제 & ㅅㄹ의 ㅎㅅㅎ와 ㄴ에 대한 ㄱㄹ ⑪ & \kfTBlank \\\hline
출전 & 『ㅊㄱㅇㅇ』 수록 ⑫ & \kfTBlank \\\hline
\end{kfActTable}
\kfSecHeader{kfAreaLiterature}{문제}{}

\kfBeginMC
\begin{kfQBlock}{01}{객관식}
\kfQStem{윗시조에 대한 설명으로 적절하지 \underline{않은} 것은?}
\begin{kfChoices}
\kfChoice 추상적인 시간을 구체적인 물체처럼 다루는 발상이 나타나 있다.
\kfChoice 초장·중장·종장의 3장 형식을 갖춘 평시조이다.
\kfChoice 님의 부재에서 오는 그리움과 재회의 소망이 담겨 있다.
\kfChoice 화자는 동짓달의 추위에서 벗어나기 위해 봄바람을 기다리고 있다.
\kfChoice '베어 내여'와 '구뷔구뷔 펴리라'에서 대비적 구조가 드러난다.
\end{kfChoices}
\end{kfQBlock}

\begin{kfQBlock}{02}{객관식}
\kfQStem{'춘풍 니불 아래 서리서리 넣었다가'에서 '서리서리'가 환기하는 이미지로 가장 적절한 것은?}
\begin{kfChoices}
\kfChoice 실타래나 천을 또아리 지어 감아 두는 모양
\kfChoice 밤하늘에 서리가 하얗게 내려앉는 풍경
\kfChoice 바람이 거세게 불어 이불이 날리는 장면
\kfChoice 긴 밤 동안 잠을 이루지 못하고 뒤척이는 행위
\kfChoice 서러운 감정이 눈물이 되어 흐르는 모습
\end{kfChoices}
\end{kfQBlock}

\begin{kfQBlock}{03}{객관식}
\kfQStem{화자가 동짓달의 긴 밤을 '베어 내'는 행위에 담긴 의미로 가장 적절한 것은?}
\begin{kfChoices}
\kfChoice 님이 없는 밤의 고독을 견디기 어려워 그 시간을 줄이고 싶은 마음
\kfChoice 계절의 흐름을 인위적으로 바꾸어 빨리 봄을 맞이하고 싶은 소망
\kfChoice 밤 시간을 잘 활용하여 생산적인 일에 쓰려는 실용적 의도
\kfChoice 긴 겨울밤의 추위를 이겨 내기 위한 주술적 행위
\kfChoice 흘러가는 시간을 붙잡아 영원히 멈추게 하려는 욕망
\end{kfChoices}
\end{kfQBlock}

\begin{kfQBlock}{04}{객관식}
\kfQStem{이 시조에서 '동짓달 기나긴 밤'과 '님 오신 날 밤'의 관계를 가장 적절하게 설명한 것은?}
\begin{kfChoices}
\kfChoice 동일한 밤이 반복되어 화자의 일상이 단조로움을 보여 준다.
\kfChoice 님의 부재와 존재에 따라 같은 시간이 정반대의 정서적 가치를 지닌다.
\kfChoice 겨울밤이 봄밤으로 계절이 바뀌면서 화자의 처지가 개선된다.
\kfChoice 과거의 밤과 미래의 밤을 서로 대비하여 흐르는 세월의 무상함을 노래한다.
\kfChoice 두 밤 모두 화자에게 고통스러운 시간으로 인식되고 있다.
\end{kfChoices}
\end{kfQBlock}

\begin{kfQBlock}{05}{객관식}
\kfQStem{윗시조를 바탕으로 <보기>를 감상한 내용으로 적절하지 \underline{않은} 것은?}
\begin{kfEvidence}황진이는 조선 시대 대표적인 여류 시인으로, 그녀의 시조에는 대담한 발상과 절묘한 비유가 자주 등장한다. 특히 추상적 개념을 구체적 사물로 변환하는 기법이 뛰어나며, 이를 통해 사랑과 이별의 정서를 독창적으로 표현하였다. 이러한 발상은 단순한 수사 기교를 넘어 시적 상상력의 정수를 보여 준다.\end{kfEvidence}
\begin{kfChoices}
\kfChoice '밤을 베어 내여'에서 시간을 천(布)처럼 자를 수 있는 것으로 다루는 대담한 발상이 드러나겠군.
\kfChoice '서리서리 넣었다가'에서 잘라 낸 시간을 실처럼 감아 보관하는 구체적 변환이 이루어지고 있겠군.
\kfChoice '구뷔구뷔 펴리라'에서 보관해 둔 시간을 다시 풀어 사용하겠다는 발상은 시적 상상력의 독창성을 보여 주겠군.
\kfChoice '춘풍 니불'은 봄바람처럼 따뜻한 이불이라는 비유로, 님과의 재회에 대한 따뜻한 기대를 환기하겠군.
\kfChoice 이 시조의 발상은 시간의 물리적 법칙을 과학적으로 탐구하려는 근대적 사고의 선구적 표현이겠군.
\end{kfChoices}
\end{kfQBlock}

\kfEndMC

\kfPassageNote{「관동별곡」 — 정철}{%
금강(金剛)이 몃 첩(疊)이완대 동해(東海)로 담아 두고 \\[4pt]
어와 조화옹(造化翁)이 물건(物件) 마다 헌사롭다 \\[4pt]
날거든 뛰거든 우뚝우뚝 솟거든 \\[4pt]
은(銀)같이 풍류(風流)같이 벽공(碧空)에 솟아 있으니 \\[4pt]
금강대(金剛臺) 디나거든 옥(玉)같은 산(山)이로다 \\[4pt]
선인(仙人)을 못 보거든 만폭동(萬瀑洞)에 가보리라 \\[14pt]
만폭동 십리(十里)허리 은(銀) 같은 무지개 \\[4pt]
놋 같은 부연(浮鉛)에 사선(四仙)이 나리온 듯 \\[4pt]
고성(高城)을 자로 지나 삼일포(三日浦)를 바라보니 \\[4pt]
단발령(斷髮嶺) 높은 봉(峰)에 동해(東海)를 겨우 건너 \\[4pt]
총석정(叢石亭) 거즈런 돌 하늘에 치솟아 있네 \\[4pt]
공수(恐水)도 돌도 아닌 것이 눈 아래 구뷔치니 \\[4pt]
은하(銀河)에 학(鶴)을 타고 광한전(廣寒殿)에 올라가니 \\[4pt]
옥(玉)잔에 가득 부은 술을 그 뉘 와 권(勸)하려는고 \\[4pt]
님을 그려 울고져 하니 눈물이 절로 나니 \\[4pt]
인간(人間) 어디메오 사자(使者)를 못 보실까 \\[4pt]
만리(萬里) 창천(蒼天)에 외로운 학(鶴)만 떠 있네 \\[4pt]
하늘을 원망하니 구름도 많고 많다 \\[4pt]
인간(人間)의 꿈을 깨니 달빛만 삼경(三更)이라
}\kfResetAreaFlag
\par\kfMaybeNeedspace{50\baselineskip}\noindent{\kfTipMark\,\,}{\small\color{kfInk} 빈칸에 알맞은 말을 채워 넣으시오. (초성 힌트 참고)}\par\nopagebreak[4]\smallskip\nopagebreak[4]
\begin{kfActTable}{|>{\centering\arraybackslash}m{2.6cm}|m{6cm}|m{4cm}|}
\rowcolor{kfPrimaryLight!50}\textbf{\color{kfPrimary}항목} & \textbf{\color{kfPrimary}내용 (빈칸 채우기)} & \textbf{\color{kfPrimary}답안 작성}\\\hline
갈래 & ㄱㅅ (기행 가사) ① & \kfTBlank \\\hline
시대/작자 & ㅈㅅ 중기, ㅈㅊ(호: 송강) ② & \kfTBlank \\\hline
창작 계기 & ㄱㅇ ㄱㅊㅅ 부임 (1580년) ③ & \kfTBlank \\\hline
주요 배경 1 & ㄱㄱㅅ의 일만이천 봉 ④ & \kfTBlank \\\hline
주요 배경 2 & 동해안 ㄱㄷ ㅍㄱ ⑤ & \kfTBlank \\\hline
선계 상상 & 은하에 ㅎ을 타고 ㄱㅎㅈ에 오름 ⑥ & \kfTBlank \\\hline
'님'의 정체 & 화자를 부임시킨 ㅇㄱ ⑦ & \kfTBlank \\\hline
표현 기법 & 한문과 우리말 ㅎㅇ, ㄱㄱㅈ 이미지 ⑧ & \kfTBlank \\\hline
비유 대상 & 봉우리를 ㅇ·ㅇ에 비유 ⑨ & \kfTBlank \\\hline
공존하는 정서 & ㅈㅇ ㅇㅊ + ㅇㄱ ㅊㅈ ⑩ & \kfTBlank \\\hline
주제 & 관동 ㅈㅇ의 풍광과 ㅇㄱ에 대한 ㅊㅅ ⑪ & \kfTBlank \\\hline
문학사적 의의 & ㅅㅁㅇㄱ과 함께 가사 문학의 쌍벽 ⑫ & \kfTBlank \\\hline
\end{kfActTable}
\kfSecHeader{kfAreaLiterature}{문제}{}

\kfBeginMC
\begin{kfQBlock}{01}{객관식}
\kfQStem{윗글에 대한 설명으로 적절하지 \underline{않은} 것은?}
\begin{kfChoices}
\kfChoice 금강산과 동해안의 자연 풍광을 기행하며 묘사한 기행 가사이다.
\kfChoice 한문 어구와 우리말을 혼용하여 격조 있는 문체를 형성하고 있다.
\kfChoice 시각적·청각적 이미지를 결합하여 자연의 웅장함을 전달하고 있다.
\kfChoice 화자는 자연 풍광에 대한 감탄만을 노래하며 다른 정서는 보이지 않는다.
\kfChoice 금강산의 봉우리를 은·옥 등에 비유하여 그 아름다움을 형상화하고 있다.
\end{kfChoices}
\end{kfQBlock}

\begin{kfQBlock}{02}{객관식}
\kfQStem{'님을 그려 울고져 하니 눈물이 절로 나니'에서 '님'이 가리키는 대상으로 가장 적절한 것은?}
\begin{kfChoices}
\kfChoice 금강산에서 만난 도사(道士)
\kfChoice 화자가 고향에 두고 온 부인
\kfChoice 화자를 임지에 보낸 임금
\kfChoice 화자의 어린 시절 스승
\kfChoice 자연 속에서 느끼는 이상적 자아
\end{kfChoices}
\end{kfQBlock}

\begin{kfQBlock}{03}{객관식}
\kfQStem{'총석정 거즈런 돌 하늘에 치솟아 있네'에서 화자의 태도로 가장 적절한 것은?}
\begin{kfChoices}
\kfChoice 총석정의 위험한 지형에 대한 두려움과 경계심
\kfChoice 하늘을 찌를 듯한 기암의 장엄한 형상에 대한 감탄
\kfChoice 인공적으로 쌓아 올린 석탑의 정교함에 대한 찬사
\kfChoice 자연이 쇠퇴하여 돌만 남은 황량한 풍경에 대한 서글픔
\kfChoice 돌기둥이 무너질 것을 걱정하는 현실적 우려
\end{kfChoices}
\end{kfQBlock}

\begin{kfQBlock}{04}{객관식}
\kfQStem{윗글에서 '은하에 학을 타고 광한전에 올라가니'라는 표현의 기능으로 가장 적절한 것은?}
\begin{kfChoices}
\kfChoice 화자가 실제로 신선이 되어 하늘을 나는 체험을 서술한 것이다.
\kfChoice 금강산의 절경이 주는 감동을 선계(仙界) 체험으로 상상하여 극대화한 것이다.
\kfChoice 임금을 만나기 위해 하늘 궁전으로 이동하는 실제 여정을 기록한 것이다.
\kfChoice 불교적 깨달음을 얻어 해탈의 경지에 이른 것을 비유한 것이다.
\kfChoice 밤하늘의 천체를 과학적으로 관찰한 결과를 문학적으로 기록한 것이다.
\end{kfChoices}
\end{kfQBlock}

\begin{kfQBlock}{05}{객관식}
\kfQStem{윗글을 바탕으로 <보기>를 감상한 내용으로 적절하지 \underline{않은} 것은?}
\begin{kfEvidence}정철(송강)은 조선 시대 가사 문학을 대표하는 작가로, 우리말의 아름다움을 극대화한 시인으로 평가받는다. 「관동별곡」은 강원 관찰사로 부임한 정철이 관동 지방을 기행하며 쓴 작품으로, 자연에 대한 감탄과 임금에 대한 충정이 조화를 이루고 있다. 조선 시대 사대부 문인에게 자연은 심미적 대상인 동시에 임금을 그리는 정치적 공간이기도 했다.\end{kfEvidence}
\begin{kfChoices}
\kfChoice '금강이 몃 첩이완대 동해로 담아 두고'에서 금강산을 동해와 결합하여 조화로운 자연미를 드러내고 있겠군.
\kfChoice '어와 조화옹이 물건 마다 헌사롭다'에서 자연을 조물주의 작품으로 인식하며 감탄하고 있겠군.
\kfChoice '님을 그려 울고져 하니'에서 관동의 자연을 즐기면서도 임금에 대한 충정을 놓지 않는 사대부의 이중적 태도가 드러나겠군.
\kfChoice '인간의 꿈을 깨니 달빛만 삼경이라'에서 선계 상상에서 깨어나는 것은 자연 감상을 포기하고 현실 정치로 복귀하겠다는 선언이겠군.
\kfChoice 한문 어구와 우리말의 혼용은 <보기>에서 말한 '우리말의 아름다움을 극대화'한 정철의 문체적 특징과 연관되겠군.
\end{kfChoices}
\end{kfQBlock}

\kfEndMC

\kfPassageNote{「정읍사」 — 작자 미상}{%
[원문] \\[4pt]
달하 노피곰 도다샤 \\[4pt]
어긔야 머리곰 비취오시라 \\[4pt]
저재 녀러신고요 \\[4pt]
어긔야 어강됴리 \\[4pt]
아으 다롱디리 \\[14pt]
전디를 드디울셰라 \\[4pt]
어긔야 즌 대를 드디울셰라 \\[4pt]
어긔야 어강됴리 \\[4pt]
아으 다롱디리 \\[14pt]
[현대어 해석] \\[4pt]
달님이시여, 높이 돋으소서 \\[4pt]
아아, 멀리멀리 비추어 주소서 \\[4pt]
(님이) 저잣거리에 들르셨는가요 \\[4pt]
아아 어강됴리(여음) \\[4pt]
아으 다롱디리(여음) \\[14pt]
진(질퍽한) 데를 밟으실까 두렵습니다 \\[4pt]
아아, 진 데를 디딜세라(밟으실까 두렵습니다) \\[4pt]
아아 어강됴리(여음) \\[4pt]
아으 다롱디리(여음)
}\kfResetAreaFlag
\par\kfMaybeNeedspace{50\baselineskip}\noindent{\kfTipMark\,\,}{\small\color{kfInk} 빈칸에 알맞은 말을 채워 넣으시오. (초성 힌트 참고)}\par\nopagebreak[4]\smallskip\nopagebreak[4]
\begin{kfActTable}{|>{\centering\arraybackslash}m{2.6cm}|m{6cm}|m{4cm}|}
\rowcolor{kfPrimaryLight!50}\textbf{\color{kfPrimary}항목} & \textbf{\color{kfPrimary}내용 (빈칸 채우기)} & \textbf{\color{kfPrimary}답안 작성}\\\hline
갈래 & ㄱㄷㄱㅇ ① & \kfTBlank \\\hline
시대 & ㅂㅈ ② & \kfTBlank \\\hline
작자 & ㅁㅅ (행상 나간 남편을 기다리는 아내) ③ & \kfTBlank \\\hline
핵심 소재 & ㄷ (밤길을 비추는 자연물) ④ & \kfTBlank \\\hline
화자의 행위 & 달에게 높이 돋아 멀리 ㅂㅊ달라고 ㄱㅇ ⑤ & \kfTBlank \\\hline
화자의 걱정 & 님이 ㅈ 데를 밟으실까 두려움 ⑥ & \kfTBlank \\\hline
여음 & '어긔야 ㅇㄱㄷㄹ', 'ㅇㅇ 다롱디리' ⑦ & \kfTBlank \\\hline
정서 & 님에 대한 ㄱㄹ과 ㅂㅇ, ㅅㄹ ⑧ & \kfTBlank \\\hline
사고 방식 & 자연물에 기대어 비는 ㅇㄴㅁㅈ적 사고 ⑨ & \kfTBlank \\\hline
수록 문헌 & 『ㅇㅎㄱㅂ』(1493) ⑩ & \kfTBlank \\\hline
주제 & 님의 안전 귀환을 비는 ㅁㅂ의 정서 ⑪ & \kfTBlank \\\hline
문학사적 의의 & ㅂㅈ의 유일한 ㅎㅈ 가요 ⑫ & \kfTBlank \\\hline
\end{kfActTable}
\kfSecHeader{kfAreaLiterature}{문제}{}

\kfBeginMC
\begin{kfQBlock}{01}{객관식}
\kfQStem{윗글에 대한 설명으로 적절한 것은?}
\begin{kfChoices}
\kfChoice 백제의 노래 중 가사가 전하는 유일한 현전 고대가요로 알려져 있다.
\kfChoice 고려 시대 궁중 연회에서 불린 속요 갈래의 축제 노래에 해당한다.
\kfChoice 조선 시대에 창작된 가사 작품으로 4음보 연속체 형식을 취한다.
\kfChoice 신라 시대 향가의 형식 중 10구체로 정제되어 전해지는 노래이다.
\kfChoice 작자가 정읍이라는 문인으로 자신의 고향 풍경을 노래한 작품이다.
\end{kfChoices}
\end{kfQBlock}

\begin{kfQBlock}{02}{객관식}
\kfQStem{'달하 노피곰 도다샤'에서 화자가 달에게 '높이 돋으소서'라고 한 의도로 가장 적절한 것은?}
\begin{kfChoices}
\kfChoice 달빛 아래에서 축제를 즐기기 위해 밝은 조명이 필요했기 때문이다.
\kfChoice 달이 높이 떠서 밤길을 밝혀 주면 님이 안전하게 돌아올 수 있기 때문이다.
\kfChoice 달의 위치로 시간을 측정하여 님의 도착 시각을 예측하려 했기 때문이다.
\kfChoice 달이 높이 뜨면 조수간만의 차가 커져 배가 빨리 돌아올 수 있기 때문이다.
\kfChoice 달에게 자신의 존재를 알려 님이 찾아오게 하려는 신호의 목적이다.
\end{kfChoices}
\end{kfQBlock}

\begin{kfQBlock}{03}{객관식}
\kfQStem{'전디를 드디울셰라'(진 데를 밟으실까 두렵습니다)에서 드러나는 화자의 심리로 가장 적절한 것은?}
\begin{kfChoices}
\kfChoice 님이 험한 길을 가다가 위험에 처할까 걱정하는 마음
\kfChoice 님이 다른 여인에게 빠져 돌아오지 않을까 의심하는 마음
\kfChoice 님이 장사에 실패하여 빈손으로 돌아올까 근심하는 마음
\kfChoice 님이 너무 빨리 돌아와 준비가 되지 않을까 초조한 마음
\kfChoice 님이 돌아와도 옛 정을 잊었을까 서운해하는 마음
\end{kfChoices}
\end{kfQBlock}

\begin{kfQBlock}{04}{객관식}
\kfQStem{윗글에서 '어긔야 어강됴리', '아으 다롱디리' 같은 구절의 기능으로 가장 적절한 것은?}
\begin{kfChoices}
\kfChoice 화자의 구체적 감정을 언어로 정확히 전달하는 서술적 기능을 한다.
\kfChoice 작품의 내용을 요약 제시하여 독자의 이해를 돕는 해설적 기능이다.
\kfChoice 뜻 없는 소리의 반복으로 노래의 리듬과 음악성을 형성하는 여음 기능이다.
\kfChoice 다음 연으로 넘어가기 전 시간적 전환을 알리는 구조적 표지의 기능이다.
\kfChoice 달에게 주문을 외워 소원을 성취하려는 주술적 의례의 기능을 한다.
\end{kfChoices}
\end{kfQBlock}

\begin{kfQBlock}{05}{객관식}
\kfQStem{윗글을 바탕으로 <보기>를 감상한 내용으로 적절하지 \underline{않은} 것은?}
\begin{kfEvidence}고대가요는 문자로 정착되기 이전에 구전(口傳)으로 전해진 노래이다. 「정읍사」는 백제 시대에 행상을 나간 남편을 기다리는 아내가 부른 것으로 전하며, 『악학궤범』에 가사가 수록되어 조선 시대까지 궁중 음악으로 전승되었다. 고대가요에는 당시 민중의 소박한 정서가 담겨 있으며, 자연물에 기대어 소원을 비는 애니미즘적 사고가 반영되어 있기도 하다.\end{kfEvidence}
\begin{kfChoices}
\kfChoice 달에게 높이 돋아 비추어 달라고 비는 것은 자연물에 기대어 소원을 비는 애니미즘적 사고가 반영된 것이겠군.
\kfChoice '어긔야 어강됴리' 같은 여음은 이 작품이 본래 노래로 불렸음을 보여 주는 근거이겠군.
\kfChoice 님이 질퍽한 길을 밟을까 걱정하는 마음에는 민중의 소박한 정서가 담겨 있겠군.
\kfChoice 이 작품이 『악학궤범』에 수록된 것은 백제의 문자 기록 체계가 고도로 발달했음을 증명하는 것이겠군.
\kfChoice 원문이 고대 한국어로 되어 있어 현대어 해석을 병기해야 내용을 이해할 수 있겠군.
\end{kfChoices}
\end{kfQBlock}

\kfEndMC

\kfStartNonlit{공자의 인(仁)과 예(禮) 외 4편}


\kfPassageNote{공자의 인(仁)과 예(禮) (인문)}{%
유교의 창시자 공자는 인간이 마땅히 지녀야 할 덕목으로 '인(仁)'을 제시했다. 인이란 사람을 사랑하고 배려하는 마음, 곧 타인에 대한 진심 어린 관심과 존중을 뜻한다. 공자는 제자 번지(樊遲)가 인에 대해 물었을 때 '사람을 사랑하는 것(愛人)'이라고 답한 바 있으며, 이는 인이 추상적 원리가 아니라 구체적인 인간관계 속에서 실현되는 실천적 덕목임을 보여 준다.

인의 실천에서 핵심적 원리는 '서(恕)'이다. 서란 '자기가 원하지 않는 것을 남에게 하지 말라(己所不欲 勿施於人)'는 것으로, 상대방의 입장에서 생각하는 역지사지(易地思之)의 태도를 가리킨다. 공자는 이 서의 원리를 인을 실현하는 가장 기본적인 방법으로 보았다. 나아가 인은 단순히 해를 끼치지 않는 소극적 태도에 그치지 않고, 자기가 서고 싶으면 남도 세워 주고 자기가 이루고 싶으면 남도 이루게 해 주는(己欲立而立人 己欲達而達人) 적극적 실천까지 포함한다.

그런데 인이 내면의 도덕적 마음이라면, '예(禮)'는 그것이 외부로 드러나는 형식이자 사회적 규범이다. 예란 사회의 질서를 유지하기 위한 행동 규범과 의례 절차를 아우르는 개념으로, 일상의 언행에서부터 국가의 정치 제도에 이르기까지 모든 영역에 걸쳐 있다. 공자는 인이 예를 통해 비로소 구체적으로 실현되고, 예는 인이라는 정신이 담길 때 비로소 형식적 껍데기를 넘어 참된 의미를 갖는다고 보았다.

인과 예의 관계를 가장 잘 보여 주는 개념이 '극기복례(克己復禮)'이다. 극기복례란 자기의 사적인 욕심을 이기고 예로 돌아가는 것을 뜻한다. 공자는 제자 안연(顏淵)에게 '극기복례가 곧 인이다(克己復禮爲仁)'라고 가르쳤는데, 이는 개인의 이기적 욕구를 절제하고 사회적 규범에 맞게 행동하는 것이야말로 인의 실현이라는 의미이다. 여기서 중요한 것은 극기복례가 단순히 외적 규범에 복종하라는 뜻이 아니라, 내면의 도덕적 자각을 바탕으로 자발적으로 예를 따르는 것이라는 점이다.

이처럼 공자의 윤리학은 인이라는 내면적 덕과 예라는 외면적 규범이 유기적으로 결합된 체계이다. 인 없는 예는 공허한 형식에 불과하고, 예 없는 인은 사회적으로 실현될 수 없다. 이 양자의 균형과 조화를 통해 개인의 도덕적 완성과 사회의 질서 유지가 동시에 달성된다는 것이 유교 윤리의 핵심이다.
}\kfResetAreaFlag
\kfSecHeader{kfAreaNonlit}{문제}{}

\kfBeginMC
\begin{kfQBlock}{01}{객관식}
\kfQStem{윗글의 내용과 일치하는 것은?}
\begin{kfChoices}
\kfChoice 공자는 인을 외부의 행동 규범이자 사회 제도로 규정했다.
\kfChoice 서(恕)의 원리는 자기가 원하는 것을 남에게도 적극적으로 행하라는 뜻이다.
\kfChoice 극기복례는 외적 규범에 대한 무조건적 복종을 강조한다.
\kfChoice 공자는 인을 타인에 대한 사랑과 배려의 마음이라고 보았다.
\kfChoice 예는 인과 무관하게 독립적으로 의미를 갖는 형식이다.
\end{kfChoices}
\end{kfQBlock}

\begin{kfQBlock}{02}{객관식}
\kfQStem{윗글에서 공자가 '극기복례'를 인의 실현이라고 본 까닭으로 가장 적절한 것은?}
\begin{kfChoices}
\kfChoice 사적 욕심을 이기고 사회적 규범에 자발적으로 따르는 것이 인의 본질이기 때문이다.
\kfChoice 외적 규범에 복종하면 자연스럽게 내면의 도덕심이 형성되기 때문이다.
\kfChoice 극기복례를 실천하면 사회적 지위가 점점 높아져서 남을 도울 수 있기 때문이다.
\kfChoice 자기 욕심을 억제하는 고통 자체가 도덕적 수양의 목적이기 때문이다.
\kfChoice 예를 회복하면 국가의 법률 체계가 정비되기 때문이다.
\end{kfChoices}
\end{kfQBlock}

\begin{kfQBlock}{03}{객관식}
\kfQStem{윗글을 바탕으로 <보기>를 이해한 내용으로 적절하지 \underline{않은} 것은?}
\begin{kfEvidence}학생 A는 학교 축제에서 자신이 원하는 행사만 하려 했으나, 친구들의 의견을 듣고 모두가 만족할 수 있는 프로그램을 함께 기획했다. 또한 학급 회의에서 정해진 발표 순서와 토론 규칙을 스스로 지키며 회의를 이끌었다.\end{kfEvidence}
\begin{kfChoices}
\kfChoice A가 자신의 욕구를 절제하고 규칙을 지킨 것은 극기복례의 태도에 해당한다.
\kfChoice A가 친구들의 의견을 경청한 것은 역지사지의 태도와 관련된다.
\kfChoice 발표 순서와 토론 규칙은 예(禮)에 해당하는 사회적 규범이다.
\kfChoice A는 규칙을 따르되 자발적으로 지켰으므로 인과 예가 결합된 모습이다.
\kfChoice A가 친구들의 프로그램을 대신 기획해 준 것은 서(恕)의 소극적 실천이다.
\end{kfChoices}
\end{kfQBlock}

\begin{kfQBlock}{04}{객관식}
\kfQStem{윗글에 대한 설명으로 가장 적절한 것은?}
\begin{kfChoices}
\kfChoice 인(仁)의 개념을 먼저 설명한 뒤 예(禮)와의 관계를 밝히고 있다.
\kfChoice 공자와 다른 사상가의 견해를 비교하여 인의 독창성을 논증하고 있다.
\kfChoice 유교 윤리의 역사적 변천 과정을 시대순으로 서술하고 있다.
\kfChoice 예(禮)의 구체적 의례 절차를 사례를 들어 상세히 기술하고 있다.
\kfChoice 현대 사회에서 유교 윤리의 한계를 비판적으로 분석하고 있다.
\end{kfChoices}
\end{kfQBlock}

\begin{kfQBlock}{05}{객관식}
\kfQStem{㉠ '인 \underline{없는} 예는 공허한 형식에 불과하고, 예 \underline{없는} 인은 사회적으로 실현될 수 없다'의 의미로 가장 적절한 것은?}
\begin{kfChoices}
\kfChoice 인과 예는 서로 무관한 독립적 개념이므로 각각 따로 실천해야 한다.
\kfChoice 예만 갖추면 인은 자연스럽게 형성되므로 예의 실천이 우선이다.
\kfChoice 내면의 도덕적 마음과 외면의 사회적 규범이 서로를 필요로 한다.
\kfChoice 인을 실현하기 위해서는 모든 사회적 규범을 폐지해야 한다.
\kfChoice 예는 강제적 규범이므로 인보다 상위의 가치를 지닌다.
\end{kfChoices}
\end{kfQBlock}

\kfEndMC

\kfPassageNote{노자의 무위자연 (인문)}{%
도가(道家) 사상의 창시자 노자는 만물의 근원이자 우주 운행의 원리를 '도(道)'라고 불렀다. 도는 천지가 생겨나기 전부터 존재하던 것으로, 눈에 보이지도 않고 귀에 들리지도 않으며 손에 잡히지도 않는다. 그러나 도는 만물을 낳고 기르면서도 자기 것이라 내세우지 않고, 만물을 이루게 하면서도 그 공을 자랑하지 않는다. 노자는 이러한 도의 성질을 설명하면서 '도가도 비상도(道可道 非常道)', 즉 '말로 표현할 수 있는 도는 참된 도가 아니다'라고 선언했다.

도의 핵심 속성은 '무위(無爲)'이다. 무위란 아무것도 하지 않는다는 뜻이 아니라, 인위적으로 꾸미거나 억지로 작용하지 않는다는 의미이다. 자연은 누가 시키지 않아도 봄에 꽃이 피고, 물은 억지로 흐르는 것이 아니라 낮은 곳으로 자연스럽게 흘러간다. 노자는 이처럼 만물이 자기 본성에 따라 저절로 그러한 상태를 '자연(自然)', 곧 '스스로 그러함'이라고 표현했다. 무위자연이란 바로 이 자연의 순리에 따르는 삶의 태도를 가리킨다.

노자는 무위자연의 이상을 물의 비유로 설명했다. '상선약수(上善若水)', 즉 '최고의 선은 물과 같다'는 것이다. 물은 만물을 이롭게 하면서도 다투지 않고, 모든 사람이 싫어하는 낮은 곳에 머문다. 또한 물은 그릇에 따라 형태를 바꾸되 자기 본질을 잃지 않으며, 부드러워 보이지만 오랜 시간 바위를 깎아 낼 만큼 강한 힘을 지닌다. 노자는 물의 이러한 성질에서 도에 가장 가까운 삶의 모습을 보았다.

이러한 무위자연의 사상은 정치 영역에도 적용되었다. 노자에 따르면, 최선의 통치는 백성이 통치자가 있는지조차 모르는 상태이다. 통치자가 법령과 규제를 남발하면 할수록 백성의 삶은 오히려 궁핍해지고 사회는 혼란해진다. 이에 노자는 '무위이무불위(無爲而無不爲)', 즉 '하지 않음으로써 하지 않는 것이 없게 된다'는 역설적 원리를 제시했다. 이는 인위적 간섭을 최소화할 때 만물이 제 본성대로 잘 이루어진다는 뜻이다.

노자의 무위자연 사상은 유교의 적극적 실천 윤리와 대비된다. 유교가 인의예지(仁義禮智)라는 덕목을 세우고 교육과 제도를 통해 인간을 교화하려 했다면, 노자는 오히려 이러한 인위적 덕목과 제도가 인간 본연의 순수함을 해친다고 보았다. 노자에게 이상적 삶이란 인위적 욕망과 지식을 버리고 자연 그대로의 소박한 상태로 돌아가는 것이었다.
}\kfResetAreaFlag
\kfSecHeader{kfAreaNonlit}{문제}{}

\kfBeginMC
\begin{kfQBlock}{01}{객관식}
\kfQStem{윗글의 내용과 일치하지 않는 것은?}
\begin{kfChoices}
\kfChoice 도(道)는 천지가 생겨나기 전부터 존재하던 것으로, 감각으로 파악할 수 없다.
\kfChoice 무위(無爲)란 인위적으로 꾸미거나 억지로 작용하지 않는다는 의미이다.
\kfChoice 노자는 법령과 규제를 강화해야 백성의 삶이 안정된다고 보았다.
\kfChoice 상선약수는 최고의 선이 물과 같다는 뜻이다.
\kfChoice 노자는 인위적 덕목과 제도가 인간의 순수함을 해친다고 보았다.
\end{kfChoices}
\end{kfQBlock}

\begin{kfQBlock}{02}{객관식}
\kfQStem{윗글에서 노자가 물(水)을 '최고의 선'에 비유한 까닭으로 가장 적절한 것은?}
\begin{kfChoices}
\kfChoice 물이 모든 오염 물질을 정화하는 힘을 가지고 있기 때문이다.
\kfChoice 물이 만물을 이롭게 하면서도 다투지 않고 낮은 곳에 머물기 때문이다.
\kfChoice 물이 항상 같은 형태를 유지하여 변하지 않는 진리를 상징하기 때문이다.
\kfChoice 물이 높은 곳에서 낮은 곳으로 흘러 권위를 보여 주기 때문이다.
\kfChoice 물이 다른 원소들과 결합하여 새로운 물질을 만들어 내기 때문이다.
\end{kfChoices}
\end{kfQBlock}

\begin{kfQBlock}{03}{객관식}
\kfQStem{㉠ '무위이무불위(無爲而無不爲)'의 의미로 가장 적절한 것은?}
\begin{kfChoices}
\kfChoice 아무런 행위도 하지 않으면 결과도 전혀 없다.
\kfChoice 인위적 간섭을 하지 않을 때 오히려 모든 것이 저절로 이루어진다.
\kfChoice 하고 싶은 일을 모두 하면 이루지 못할 것이 없다.
\kfChoice 어떤 일이든 열심히 노력하면 반드시 큰 성과를 거둘 수 있다.
\kfChoice 통치자가 강력한 의지로 모든 일을 직접 처리해야 한다.
\end{kfChoices}
\end{kfQBlock}

\begin{kfQBlock}{04}{객관식}
\kfQStem{윗글을 바탕으로 <보기>를 이해한 내용으로 적절하지 \underline{않은} 것은?}
\begin{kfEvidence}어느 마을의 촌장은 주민들에게 세세한 규칙을 정해 주지 않았다. 그는 주민들이 스스로 농사 시기를 정하고, 서로 도우며 분쟁을 해결하도록 맡겨 두었다. 마을은 큰 갈등 없이 오랫동안 평화롭게 유지되었고, 주민들은 촌장의 존재를 의식하지 못했다.\end{kfEvidence}
\begin{kfChoices}
\kfChoice 촌장이 세세한 규칙을 정하지 않은 것은 무위의 태도에 해당한다.
\kfChoice 주민들이 촌장의 존재를 의식하지 못한 것은 노자가 말한 최선의 통치와 부합한다.
\kfChoice 주민들이 스스로 농사 시기를 정한 것은 자연(自然)의 원리에 해당한다.
\kfChoice 마을이 평화롭게 유지된 것은 무위이무불위의 결과로 볼 수 있다.
\kfChoice 촌장은 주민들에게 덕목을 가르쳐 교화한 결과 마을이 안정된 것이다.
\end{kfChoices}
\end{kfQBlock}

\begin{kfQBlock}{05}{객관식}
\kfQStem{윗글에서 언급된 유교와 도가의 입장 차이로 가장 적절한 것은?}
\begin{kfChoices}
\kfChoice 유교는 자연 상태로의 회귀를, 도가는 인위적 제도를 통한 적극적 교화를 강조한다.
\kfChoice 유교는 인의예지를 세워 교화하려 하고, 도가는 인위적 덕목이 순수함을 해친다고 본다.
\kfChoice 유교는 개인의 자유를, 도가는 사회적 질서를 중시한다.
\kfChoice 유교와 도가 모두 법률을 통한 통치를 이상으로 삼았다.
\kfChoice 유교는 도(道)의 개념을 중시하고, 도가는 인(仁)의 실천을 강조한다.
\end{kfChoices}
\end{kfQBlock}

\kfEndMC

\kfPassageNote{맹자와 순자의 인성론 (인문)}{%
인간의 본성이 선한가, 악한가 하는 물음은 동양 철학에서 가장 오래되고 치열한 논쟁 중 하나이다. 이 논쟁의 대표적 두 축은 맹자의 성선설(性善說)과 순자의 성악설(性惡說)이다. 두 사상가는 모두 공자의 유학을 계승했지만, 인간 본성에 대한 견해에서 정반대의 입장을 취했다.

맹자는 인간이 태어날 때부터 선한 본성을 지니고 있다고 주장했다. 그 근거로 맹자는 '사단(四端)', 즉 네 가지 선한 마음의 싹을 제시했다. 첫째는 측은지심(惻隱之心)으로, 남의 고통을 차마 보지 못하는 불쌍히 여기는 마음이다. 맹자는 어린아이가 우물에 빠지려 하면 누구나 깜짝 놀라며 구하려 드는 것이 바로 이 마음의 증거라고 보았다. 둘째는 수오지심(羞惡之心)으로, 옳지 못한 것을 부끄러워하고 미워하는 마음이다. 셋째는 사양지심(辭讓之心)으로, 남에게 양보하고 공경하는 마음이다. 넷째는 시비지심(是非之心)으로, 옳고 그름을 분별하는 마음이다. 맹자는 이 네 가지 마음이 각각 인(仁), 의(義), 예(禮), 지(智)라는 덕목의 단서라고 보았다.

맹자에게 도덕적 수양이란 이미 내면에 갖추어진 선한 싹을 잘 길러 키우는 것이었다. 마치 콩나물에 물을 주듯이, 타고난 선한 본성을 정성껏 보존하고 확충하면 누구나 성인(聖人)이 될 수 있다는 것이다. 인간이 악행을 저지르는 것은 본성 자체가 악하기 때문이 아니라, 외부 환경의 영향으로 선한 본성을 잃어버렸기 때문이라고 맹자는 설명했다.

이에 반해 순자는 인간의 본성이 악하다고 주장했다. 순자가 말하는 '성(性)'이란 인간이 태어나면서 갖는 자연적 욕구와 감정을 가리킨다. 사람은 본래 이익을 좋아하고, 시기하고 미워하며, 감각적 쾌락을 추구하는 경향을 가지고 있다. 이러한 자연적 성향을 그대로 따르면 다툼과 혼란이 생기므로, 순자는 인간의 본성이 악하다고 본 것이다. 다만 순자의 성악설에서 '악'이란 도덕적으로 사악하다는 뜻이라기보다, 자연 상태 그대로 두면 사회적 질서가 무너진다는 의미에 가깝다.

따라서 순자에게 도덕적 수양이란 인위적 노력을 통해 본성을 교정하는 것이었다. 순자는 이 인위적 노력을 '위(僞)'라고 불렀는데, 이는 '사람이 만든 것'이라는 뜻이다. 예(禮)와 의(義)는 성인이 인위적으로 만든 것으로, 인간은 이러한 예의 교육을 받고 끊임없이 자기를 갈고닦음으로써 비로소 선해질 수 있다. 맹자가 내면의 선을 보존하는 데 초점을 두었다면, 순자는 외부의 규범과 교육을 통한 변화에 초점을 두었다는 점에서 두 사람의 교육관은 뚜렷이 갈린다.
}\kfResetAreaFlag
\kfSecHeader{kfAreaNonlit}{문제}{}

\kfBeginMC
\begin{kfQBlock}{01}{객관식}
\kfQStem{윗글의 내용과 일치하는 것은?}
\begin{kfChoices}
\kfChoice 맹자는 인간이 악행을 저지르는 원인을 본성 자체의 악함에서 찾았다.
\kfChoice 순자는 인간의 자연적 욕구를 그대로 따르면 사회적 질서가 유지된다고 보았다.
\kfChoice 맹자의 사단 중 수오지심은 옳고 그름을 분별하는 마음이다.
\kfChoice 순자는 예(禮)와 의(義)가 성인이 인위적으로 만든 것이라고 보았다.
\kfChoice 맹자와 순자는 인간 본성에 대해 동일한 견해를 가졌다.
\end{kfChoices}
\end{kfQBlock}

\begin{kfQBlock}{02}{객관식}
\kfQStem{맹자가 '어린아이가 우물에 빠지려 하는 상황'을 예로 든 까닭으로 가장 적절한 것은?}
\begin{kfChoices}
\kfChoice 어린아이가 위험 상황 자체를 인식하지 못한다는 점을 분명히 지적하기 위해서이다.
\kfChoice 인간에게 남의 고통을 차마 보지 못하는 선한 마음이 본래 있음을 보이기 위해서이다.
\kfChoice 교육을 충분히 받지 못한 아이의 무모한 행동이 위험함을 경고하기 위해서이다.
\kfChoice 어린아이를 구하는 행위가 사회 규범에 따른 마땅한 의무임을 강조하기 위해서이다.
\kfChoice 오직 도덕적으로 완성된 성인(聖人)만이 아이를 구할 수 있음을 드러내기 위해서이다.
\end{kfChoices}
\end{kfQBlock}

\begin{kfQBlock}{03}{객관식}
\kfQStem{맹자와 순자의 '도덕적 수양' 방법의 차이로 가장 적절한 것은?}
\begin{kfChoices}
\kfChoice 맹자는 외부 교육을, 순자는 내면 성찰을 수양의 핵심으로 보았다.
\kfChoice 맹자는 내면의 선한 본성을 보존·확충하고, 순자는 예의 교육으로 본성을 교정하려 했다.
\kfChoice 맹자는 자연적 욕구를 긍정하고, 순자는 도덕적 감정을 부정했다.
\kfChoice 맹자와 순자는 모두 예(禮)를 통한 교화만이 유일하고 옳은 수양 방법이라고 보았다.
\kfChoice 맹자는 성인만이 수양할 수 있다고 보고, 순자는 누구나 가능하다고 보았다.
\end{kfChoices}
\end{kfQBlock}

\begin{kfQBlock}{04}{객관식}
\kfQStem{윗글을 바탕으로 <보기>를 이해한 내용으로 적절하지 \underline{않은} 것은?}
\begin{kfEvidence}영수는 길에서 할머니가 무거운 짐을 들고 힘들어하는 모습을 보고 자신도 모르게 다가가 도움을 드렸다. 한편 민지는 처음에는 남의 물건을 탐냈지만, 부모님의 가르침과 학교 교육을 통해 점차 남의 것을 존중하는 습관을 갖게 되었다.\end{kfEvidence}
\begin{kfChoices}
\kfChoice 영수의 행동은 맹자가 말한 측은지심의 발현으로 볼 수 있다.
\kfChoice 영수가 '자신도 모르게' 도운 것은 선한 본성이 자발적으로 드러난 것이다.
\kfChoice 민지의 변화 과정은 순자가 말한 인위적 노력을 통한 교정에 해당한다.
\kfChoice 민지가 처음 남의 물건을 탐낸 것은 맹자가 말한 사양지심에 해당한다.
\kfChoice 영수는 맹자의, 민지는 순자의 관점에서 각각 설명될 수 있다.
\end{kfChoices}
\end{kfQBlock}

\begin{kfQBlock}{05}{객관식}
\kfQStem{순자가 말한 '위(僞)'의 의미로 가장 적절한 것은?}
\begin{kfChoices}
\kfChoice 거짓된 행동으로 남을 속이는 것이다.
\kfChoice 타고난 본성을 그대로 유지하는 것이다.
\kfChoice 사람이 인위적 노력으로 만들어 낸 것이다.
\kfChoice 자연적 욕구에 따라 행동하는 것이다.
\kfChoice 성인만이 가질 수 있는 타고난 능력이다.
\end{kfChoices}
\end{kfQBlock}

\kfEndMC

\kfPassageNote{불교의 연기설과 공(空) (인문)}{%
불교의 핵심 교리 가운데 하나인 연기설(緣起說)은 '모든 존재와 현상은 홀로 존재하는 것이 아니라 여러 조건이 서로 의존하여 생겨난다'는 가르침이다. '연(緣)'은 조건, '기(起)'는 일어남을 뜻하므로, 연기란 '조건에 의해 일어남'이라는 의미이다. 석가모니는 이를 '이것이 있으므로 저것이 있고, 이것이 생기므로 저것이 생긴다. 이것이 없으면 저것이 없고, 이것이 사라지면 저것도 사라진다'라는 말로 표현했다.

연기설에 따르면, 세상의 모든 존재는 무수한 원인과 조건의 그물 속에서 서로 연결되어 있다. 예를 들어, 하나의 꽃이 피기 위해서는 씨앗, 흙, 물, 햇빛, 적절한 온도 등 수많은 조건이 함께 갖추어져야 한다. 이 조건 가운데 하나라도 빠지면 꽃은 피지 못한다. 이처럼 어떤 존재도 다른 조건과 무관하게 독립적으로 존재할 수 없으며, 모든 것은 관계의 그물 속에서 비로소 성립한다.

연기설에서 자연스럽게 도출되는 핵심 개념이 '공(空)'이다. 공이란 '비어 있다'는 뜻이지만, 아무것도 존재하지 않는다는 의미가 아니다. 공은 모든 존재에 고정된 실체, 즉 '자성(自性)'이 없다는 것을 뜻한다. 모든 것이 조건에 의해 생겨나고 조건이 바뀌면 변하거나 사라지므로, 영원히 변하지 않는 고정된 본질이란 존재하지 않는다. 나무는 씨앗에서 자라 꽃을 피우고 열매를 맺다가 결국 쓰러지는데, 이 과정 어디에도 영원히 변하지 않는 '나무 자체'라는 실체는 없다.

이러한 연기와 공의 이해는 인간의 고통을 해결하기 위한 가르침으로 이어진다. 석가모니는 인간이 겪는 고통의 원인을 '사성제(四聖諦)'로 체계화했다. 첫째, 삶은 고통이다(고제, 苦諦). 둘째, 고통의 원인은 갈애(渴愛), 즉 집착이다(집제, 集諦). 셋째, 집착을 소멸하면 고통도 소멸한다(멸제, 滅諦). 넷째, 고통의 소멸에 이르는 길이 있다(도제, 道諦). 여기서 고통의 근본 원인인 집착은 사물에 고정된 실체가 있다고 착각하는 데서 비롯된다. 연기와 공을 깨달으면, 영원한 것에 대한 집착에서 벗어나 고통을 극복할 수 있는 것이다.

사성제의 네 번째인 도제의 구체적 내용이 '팔정도(八正道)'이다. 팔정도란 바른 견해(정견), 바른 사유(정사유), 바른 말(정어), 바른 행위(정업), 바른 생활(정명), 바른 노력(정정진), 바른 마음챙김(정념), 바른 집중(정정)의 여덟 가지 바른 길을 말한다. 이 팔정도는 지혜(정견·정사유), 계율(정어·정업·정명), 선정(정정진·정념·정정)의 세 영역으로 구성되어, 앎과 행동과 마음의 수련을 동시에 추구한다는 점에서 불교 수행의 총체적 성격을 보여 준다.
}\kfResetAreaFlag
\kfSecHeader{kfAreaNonlit}{문제}{}

\kfBeginMC
\begin{kfQBlock}{01}{객관식}
\kfQStem{윗글의 내용과 일치하지 않는 것은?}
\begin{kfChoices}
\kfChoice 연기란 모든 존재가 조건에 의해 서로 의존하며 생겨남을 뜻한다.
\kfChoice 공(空)이란 아무것도 존재하지 않는다는 뜻이다.
\kfChoice 사성제에 따르면 고통의 원인은 집착이다.
\kfChoice 팔정도는 지혜, 계율, 선정의 세 영역으로 구성된다.
\kfChoice 석가모니는 연기를 '이것이 있으므로 저것이 있다'는 말로 표현했다.
\end{kfChoices}
\end{kfQBlock}

\begin{kfQBlock}{02}{객관식}
\kfQStem{윗글에서 '꽃이 피는 과정'을 예로 든 까닭으로 가장 적절한 것은?}
\begin{kfChoices}
\kfChoice 꽃이 아름답기 때문에 불교의 이상을 상징할 수 있어서이다.
\kfChoice 어떤 존재도 여러 조건 없이는 독립적으로 성립할 수 없음을 보여 주기 위해서이다.
\kfChoice 식물의 생장이 인간의 수행 과정과 동일하다는 점을 강조하기 위해서이다.
\kfChoice 자연 현상이 인간의 의지와 무관하게 일어난다는 점을 보여 주기 위해서이다.
\kfChoice 꽃이 피고 지는 순환이 윤회의 원리를 설명하기 위해서이다.
\end{kfChoices}
\end{kfQBlock}

\begin{kfQBlock}{03}{객관식}
\kfQStem{'공(空)'의 개념에 대한 이해로 가장 적절한 것은?}
\begin{kfChoices}
\kfChoice 세상의 모든 존재는 실제로 존재하지 않는 환상이다.
\kfChoice 우주의 시작에는 텅 빈 공간만 있었다는 뜻이다.
\kfChoice 모든 존재에 영원히 변하지 않는 고정된 실체가 없다는 뜻이다.
\kfChoice 인간의 마음을 비워 아무 생각도 하지 않는 상태를 가리킨다.
\kfChoice 물질적 존재만 비어 있고 정신적 존재는 영원하다는 뜻이다.
\end{kfChoices}
\end{kfQBlock}

\begin{kfQBlock}{04}{객관식}
\kfQStem{윗글을 바탕으로 <보기>를 이해한 내용으로 적절하지 \underline{않은} 것은?}
\begin{kfEvidence}선호는 자신이 가장 아끼는 스마트폰이 고장 나자 큰 괴로움을 느꼈다. 그는 '이 폰은 영원히 내 것이어야 해'라는 생각에서 벗어나지 못했다. 이후 선호는 불교 강좌를 듣고 '모든 것은 조건에 의해 생겨나고 조건이 바뀌면 변한다'는 가르침을 접한 뒤, 점차 괴로움이 줄어드는 것을 경험했다.\end{kfEvidence}
\begin{kfChoices}
\kfChoice 선호의 괴로움은 사성제에서 말하는 '고(苦)'에 해당한다.
\kfChoice 선호가 폰이 영원히 자기 것이어야 한다고 생각한 것은 '집착'에 해당한다.
\kfChoice 선호의 집착은 스마트폰에 고정된 실체가 있다고 착각한 데서 비롯된다.
\kfChoice 선호가 연기의 가르침을 이해한 후 괴로움이 줄어든 것은 집착의 소멸 과정이다.
\kfChoice 선호는 팔정도의 수행을 완성하여 완전한 깨달음에 도달한 것이다.
\end{kfChoices}
\end{kfQBlock}

\begin{kfQBlock}{05}{객관식}
\kfQStem{윗글에 대한 설명으로 가장 적절한 것은?}
\begin{kfChoices}
\kfChoice 연기설의 원리에서 공 개념을 끌어낸 뒤 수행 체계로 이어가는 순차적 서술이다.
\kfChoice 불교와 유교의 교리를 항목별로 대조하여 불교의 우위를 논증하고 있다.
\kfChoice 사성제와 팔정도의 역사적 변천 과정을 시대순으로 정리하여 보여 주고 있다.
\kfChoice 연기설에 대한 여러 학파의 해석을 소개하며 그 차이의 의미를 분석하고 있다.
\kfChoice 과학적 실험과 통계 자료를 근거로 연기설의 타당성을 입증하고 있다.
\end{kfChoices}
\end{kfQBlock}

\kfEndMC

\kfPassageNote{양명학과 지행합일 (인문)}{%
중국 명나라의 사상가 왕양명(王陽明)은 당시 지배적이던 주자학(朱子學)에 도전하며 양명학(陽明學)이라는 독자적 학문 체계를 수립했다. 주자학의 창시자 주희(朱熹)는 사물의 이치를 탐구한 뒤에야 비로소 올바른 행동이 가능하다고 보았는데, 이를 '선지후행(先知後行)', 즉 '앎이 먼저이고 행함이 뒤따른다'는 입장으로 요약할 수 있다. 주희에 따르면 인간은 격물치지(格物致知), 곧 사물에 나아가 그 이치를 궁구함으로써 앎에 이르러야 하며, 이렇게 얻은 앎이 축적되어야 올바른 실천이 가능해진다.

왕양명은 이러한 주자학의 입장에 근본적 의문을 제기했다. 그는 젊은 시절 주희의 가르침대로 대나무 앞에 앉아 이치를 궁구하려 했으나, 며칠이 지나도 깨달음을 얻지 못하고 오히려 병이 났다고 전해진다. 이 경험을 통해 왕양명은 이치가 사물 밖에 있는 것이 아니라 인간의 마음속에 이미 갖추어져 있다고 확신하게 되었다.

왕양명 사상의 핵심은 '치양지(致良知)'이다. 양지(良知)란 인간이 태어날 때부터 갖추고 있는 도덕적 앎, 즉 옳고 그름을 가리는 선천적 도덕 판단 능력을 뜻한다. 예를 들어, 부모를 보면 자연스럽게 효도해야 한다는 것을 알고, 어려운 사람을 보면 도와야 한다는 것을 안다. 이러한 앎은 학습을 통해 얻는 것이 아니라 본래 마음속에 갖추어져 있는 것이다. 치양지란 이 양지를 가리고 있는 사적 욕심을 제거하여 양지가 온전히 발현되도록 하는 것이다.

왕양명 철학에서 가장 독창적인 주장은 '지행합일(知行合一)'이다. 이는 앎과 행함이 본래 하나라는 것이다. 왕양명에 따르면, 참된 앎은 반드시 행함을 수반하며, 행하지 않는 앎은 아직 참으로 안 것이 아니다. 예컨대 효도해야 한다는 것을 알면서도 효도하지 않는 사람은, 왕양명의 기준에서는 아직 효도를 참으로 아는 것이 아니다. 반대로 참된 앎이 있으면 그에 따른 행동은 자연스럽게 이루어진다.

이러한 지행합일론은 주자학의 선지후행과 뚜렷이 대비된다. 주자학이 앎과 행함을 시간적으로 분리하여 앎이 먼저 완성된 뒤에 행함이 따른다고 본 반면, 왕양명은 앎과 행함이 동시에 일어나며 둘을 나눌 수 없다고 주장했다. 이 차이는 단순한 순서의 문제가 아니라, 앎의 본질에 대한 근본적 시각 차이에서 비롯된다. 주자학에서 앎은 사물을 탐구하여 얻는 객관적 지식인 반면, 양명학에서 앎은 마음속에 이미 있는 도덕적 자각이기 때문이다.
}\kfResetAreaFlag
\kfSecHeader{kfAreaNonlit}{문제}{}

\kfBeginMC
\begin{kfQBlock}{01}{객관식}
\kfQStem{윗글의 내용과 일치하는 것은?}
\begin{kfChoices}
\kfChoice 주희는 이치가 인간의 마음속에 이미 갖추어져 있다고 보았다.
\kfChoice 왕양명은 사물에 나아가 이치를 궁구해야 참된 앎에 이른다고 주장했다.
\kfChoice 양지(良知)란 학습을 통해 후천적으로 얻는 도덕 지식이다.
\kfChoice 왕양명에 따르면 참된 앎은 반드시 행함을 수반한다.
\kfChoice 주자학과 양명학은 앎의 본질에 대해 동일한 시각을 가지고 있다.
\end{kfChoices}
\end{kfQBlock}

\begin{kfQBlock}{02}{객관식}
\kfQStem{왕양명이 주자학에 의문을 제기하게 된 계기로 가장 적절한 것은?}
\begin{kfChoices}
\kfChoice 주희의 저서를 읽고 논리적 모순을 발견했기 때문이다.
\kfChoice 대나무 앞에서 이치를 궁구했으나 깨달음을 얻지 못했기 때문이다.
\kfChoice 스승으로부터 양지의 개념을 전수받았기 때문이다.
\kfChoice 거듭된 정치적 실패를 통해 실천의 중요성을 깊이 깨달았기 때문이다.
\kfChoice 불교의 공(空) 사상에서 영감을 받았기 때문이다.
\end{kfChoices}
\end{kfQBlock}

\begin{kfQBlock}{03}{객관식}
\kfQStem{주자학과 양명학의 차이에 대한 설명으로 적절하지 \underline{않은} 것은?}
\begin{kfChoices}
\kfChoice 주자학은 앎이 먼저이고 행함이 뒤따른다고 보았다.
\kfChoice 양명학은 앎과 행함이 동시에 일어나며 결코 나눌 수 없다고 보았다.
\kfChoice 주자학에서 앎은 사물 탐구를 통해 얻는 객관적 지식이다.
\kfChoice 양명학에서 앎은 마음속에 이미 있는 도덕적 자각이다.
\kfChoice 주자학과 양명학의 차이는 앎의 순서에 국한되며 앎의 본질과는 무관하다.
\end{kfChoices}
\end{kfQBlock}

\begin{kfQBlock}{04}{객관식}
\kfQStem{윗글을 바탕으로 <보기>를 이해한 내용으로 적절하지 \underline{않은} 것은?}
\begin{kfEvidence}수진이는 '환경 보호가 중요하다'는 것을 알고 있었지만, 일회용 컵을 습관적으로 사용하고 분리수거도 하지 않았다. 한편 지호는 환경 보호에 대해 깊이 공부하지 않았지만, 길에 버려진 쓰레기를 보면 자연스럽게 주워서 쓰레기통에 넣었다.\end{kfEvidence}
\begin{kfChoices}
\kfChoice 왕양명의 관점에서, 수진이는 환경 보호를 참으로 알고 있는 것이 아니다.
\kfChoice 주자학의 관점에서, 수진이는 앎은 갖추었으나 행함이 부족한 상태이다.
\kfChoice 왕양명의 관점에서, 지호의 행동은 양지의 자연스러운 발현으로 볼 수 있다.
\kfChoice 주자학의 관점에서, 지호는 격물치지를 통해 환경의 이치를 깊이 깨달은 사람이다.
\kfChoice 수진이와 지호의 사례는 주자학과 양명학의 시각 차이를 보여 준다.
\end{kfChoices}
\end{kfQBlock}

\begin{kfQBlock}{05}{객관식}
\kfQStem{'치양지(致良知)'의 의미로 가장 적절한 것은?}
\begin{kfChoices}
\kfChoice 외부 사물을 두루 탐구하여 새로운 지식을 차근차근 축적하는 일이다.
\kfChoice 스승으로부터 도덕적 가르침을 전수받아 그대로 실천하는 일이다.
\kfChoice 양지를 가리는 사적 욕심을 제거해 양지가 온전히 드러나게 하는 일이다.
\kfChoice 선한 행동을 꾸준히 반복하여 도덕적 습관을 형성해 나가는 일이다.
\kfChoice 타인의 양지를 깨워 주기 위해 도덕 교육을 베푸는 일이다.
\end{kfChoices}
\end{kfQBlock}

\kfEndMC

\kfStartWeekly{패턴 워크북 — 총 18문항}


\kfPassageNote{문항 1 / 섞어치기}{%
육상 동물과 수생 동물은 각각의 생활 환경에 따라 호흡 과정에서 어려움을 겪는다. 육상 동물은 폐의 표면에 건조한 공기가 닿으면 폐가 건조해질 수 있는데, 이는 폐를 손상시키고 ⓐ탈수를 일으키며 기체를 용해하는 폐의 능력을 떨어뜨린다. 한편 수생 동물은 가용 산소가 적은 수중 환경과 수온 차에 의해 변화되는 산소 가용성에 적응해야 한다. 또한 수생 동물은 물속에서 몸을 움직이는 데 많은 근육 운동이 필요하고, 찬 바닷물의 높은 비열 때문에 아가미의 열을 빼앗기며, 삼투압 현상으로 인해 수분 균형에 문제가 발생하기도 한다.

(중략)

이러한 환경 속에서 수생 동물은 효과적인 호흡을 위해 아가미라는 특수한 기관을 사용한다. 아가미는 물속에 녹아 있는 산소를 흡수하고 체내에 있는 이산화 탄소를 배출하는 기능을 하는데, 외아가미와 내아가미로 나뉜다. 외아가미는 몸 밖으로 ⓑ돌출된 형태로, 표면적이 넓고 형태가 다양하며 유생기 도롱뇽이나 일부 무척추동물에서 나타난다. 외아가미는 주변 물과의 접촉을 극대화

(중략)
}
\begin{kfQBlock}{01}{객관식}
\kfQStem{윗글의 내용과 일치하는 것은?}
\begin{kfChoices}
\kfChoice 수생 동물은 가용 산소가 적은 수중 환경에 적응해야 한다.
\kfChoice 어류는 환수에 몸의 총사용 에너지 중 1\textasciitilde{}2\%만을 사용한다.
\kfChoice 외아가미는 체내에 위치해 덮개로 보호받는다.
\end{kfChoices}
\end{kfQBlock}

\kfPassageNote{문항 2 / 부정·상대어}{%
원자에는 중심에 핵이 있고, 핵 주위에 동심원처럼 퍼져 있는 특정 궤도들을 도는 전자들이 있다. 이 궤도마다 전자가 갖는 에너지 값이 정해져 있는데, 이 값을 에너지 준위라고 한다. 원자가 가진 특정한 몇 개의 궤도에만 전자가 존재하므로, 전자가 가질 수 있는 에너지 준위도 특정한 값으로 정해져 있다. 전자는 에너지를 얻으면 바깥쪽 궤도로, 에너지를 잃으면 안쪽 궤도로 이동한다. 원자가 하나일 때는 전자의 에너지 준위가 특정한 값이지만, 여러 원자가 모이면 여러 원자들의 궤도가 ⓐ중첩되면서 전자들이 서로 영향을 미쳐 에너지 준위에 변화가 생긴다. 이로 인해 에너지 준위는 특정한 값이 아닌 일정한 범위로 나타나게 된다. 이렇게 에너지 값이 범위로 표현된 것을 에너지 밴드라고 한다. 보통의 물질들은 수많은 원자로 구성돼 있기 때문에 대부분 에너지 밴드를 갖고 있다. 전자는 낮은 에너지 상태가 되려는 성질이 있으므로, 발광 물질의 전자에 에너지가 공급되면 전자가 바깥쪽으로 이동한 뒤 다시

(중략)
}
\begin{kfQBlock}{02}{객관식}
\kfQStem{<보기>를 바탕으로 판단할 때, 적절하지 \underline{않은} 것은?}
\begin{kfEvidence}같은 재료로 만든 양자점 C와 D가 있다. 같은 조건에서 C는 녹색, D는 청색을 보였다. 실험자는 '더 작은 입자는 에너지 변동 폭이 더 크다'는 원리를 이용해 두 입자의 상대적 특성을 비교하고자 한다.\end{kfEvidence}
\begin{kfChoices}
\kfChoice D가 청색을 보였다면 D의 에너지 변동 폭이 C보다 반드시 더 작다고 봐야 한다.
\kfChoice D가 청색을 보였다면 D가 C보다 더 작을 가능성이 크다.
\kfChoice C가 D보다 크다면 C의 에너지 변동 폭은 D보다 작을 가능성이 크다.
\end{kfChoices}
\end{kfQBlock}

\kfPassageNote{문항 3 / 주체·객체}{%
우리는 일반적으로 마음속 상태를 겉으로 분명하게 드러내거나 나타낼 때 ‘표현’이라는 말을 사용한다. 그러나 표현이라는 말이 단지 자신의 심리적 상태를 드러내는 데에만 쓰이는 것은 아니다. 때로는 자연을 대하는 우리의 태도나 반응에도 표현이라는 말이 적용된다. 예를 들어 축 늘어진 나뭇가지를 보고 “애처롭다.”라고 말하는 경우가 있다. 이는 나뭇가지 자체가 감정을 지닌 것은 아니지만, 우리가 그것을 바라보며 내면의 정서를 투영하고 외부로 드러내는 행위를 통해 표현이라는 말을 사용하게 되는 것이다. 하지만 표현이라는 말이 가장 널리 사용되는 분야는 다름 아닌 예술이며, 예술과 표현의 관계에 대해서는 오래전부터 다양한 견해가 논의되어 왔다.

(중략)

첫째로, ㉠표현을 예술가가 창작 과정 속에서 행하는 내적 활동으로 보는 견해가 있다. 이 입장에서 표현은 예술가의 마음속에서 일어나는 정신적 과정, 즉 내적인 작용으로 이해된다. 이러한 입장을 대표하는 사상가가 크로체와 콜링우드이다. 크로체는 예술의 본

(중략)
}
\begin{kfQBlock}{03}{객관식}
\kfQStem{윗글의 전개 방식으로 가장 적절한 것은?}
\begin{kfChoices}
\kfChoice 예술과 표현의 관계에 대해 각각의 견해와 그 견해의 한계를 제시하고 있다.
\kfChoice 예술가가 자신의 감정을 표현하는 방식에 대한 견해를 역사적으로 살펴보고 있다.
\kfChoice 예술의 아름다움에 대한 전통적인 관점과 현대의 관점을 비교한 후 이를 절충하고 있다.
\end{kfChoices}
\end{kfQBlock}

\kfPassageNote{문항 4 / 순서·방향}{%
컴퓨터는 사람의 언어를 직접 이해할 수 없다. 따라서 컴퓨터가 사람의 언어를 이해할 수 있도록 표현해 주어야 한다. 컴퓨터가 사람의 언어를 이해할 수 있도록 표현하는 방법에는 여러 가지가 있을 수 있는데, 그중 ㉠원-핫 인코딩은 표현하고 싶은 단어의 색인 값에 1을 부여하고 나머지에는 0을 부여하는 방법이다. 원-핫 인코딩을 수행하기 위해서는 먼저 대상이 되는 텍스트의 단어 집합을 만든다. 예를 들어 ‘오늘 날씨 정말 좋다.’라는 한 문장으로 이루어진 텍스트를 대상으로 한다고 가정하면, 단어 집합은 \{오늘, 날씨, 정말, 좋다\}가 되고 집합의 원소가 되는 단어의 개수가 4개이므로 단어 집합의 크기는 4가 된다. 단어 집합의 크기에 따라 각각의 단어에 [1, 0, 0, 0], [0, 1, 0, 0], [0, 0, 1, 0], [0, 0, 0, 1]과 같이 4차원의 고유한 숫자를 부여하는 것이 원-핫 인코딩이다. 원-핫 인코딩은 직관적이고 단순하지만 단어들이 개별적, 독립적으로 표현되기

(중략)
}
\begin{kfQBlock}{04}{객관식}
\kfQStem{윗글을 이해한 내용으로 가장 적절한 것은?}
\begin{kfChoices}
\kfChoice 코사인 유사도 0은 두 벡터가 직각임을 뜻한다.
\kfChoice 원-핫 인코딩은 실수값으로 의미 관계를 직접 반영한다.
\kfChoice 워드투벡은 단어 집합 없이 문장 하나를 한 벡터로만 만든다.
\end{kfChoices}
\end{kfQBlock}

\kfPassageNote{문항 5 / 인과 도치}{%
리더는 목표 달성을 위해 구성원들을 이끌고 그들에게 영향력을 미치는 조직의 핵심 인물이다. 오래전부터 리더로서의 능력을 나타내는 리더십은 무엇이며, 어떠한 리더십이 효과적인지에 관한 연구가 행해졌다. 피들러는 리더를 업무의 목표 달성을 중시하는 업무 지향적 스타일을 가진 리더와 구성원들과의 관계, 소통 등을 중시하는 관계 지향적 스타일을 가진 리더로 구분하였다. 그리고 리더가 자신의 집단에 영향을 미칠 수 있는 정도를 나타내는 상황의 호의성이 ‘리더와 구성원들의 관계, 업무 구조, 리더의 권한’이라는 세 가지 변수로 결정된다고 보았다. 즉 리더와 구성원들이 서로 믿고 존중하는 관계일수록, 업무가 명확하게 구조화되어 있을수록, 채용 및 승진 등에 대한 리더의 공식적 권한이 클수록 상황의 호의성이 높고, 반대의 경우는 호의성이 낮다고 본 것이다. 그는 상황의 호의성에 따라 적합한 리더십의 종류가 다르다고 보았는데, 상황의 호의성이 매우 낮거나 매우 높을 때는 업무 지향적 리더가, 중간 

(중략)
}
\begin{kfQBlock}{05}{객관식}
\kfQStem{윗글을 바탕으로 <보기>를 이해한 반응으로 적절하지 \underline{않은} 것은?}
\begin{kfEvidence}\textbackslash{}n갑이 경영하고 있는 기업 K는 얼마 전 회사에 오랫동안 근무하던 구성원들이 대거 퇴직하여 대부분이 입사한 지 오래되지 않은 이들로 구성되어 있다. 이들은 업무에 임하는 태도는 매우 적극적이지만, 아직 업무에 숙달되지 않은 상태이다. 한편 을이 경영하고 있는 기업 P는 업무에 숙달된 구성원들이 매우 많지만 전반적으로 구성원들이 자신감이 부족하고 근로 의욕이 저하된 상태이다.\textbackslash{}n갑은 각 구성원이 수행해야 할 업무와 성과에 따른 보상과 처벌을 명확히 정하여 제시하였다. 을은 구성원에게 기업이 추구하는 목표를 제시하고 구성원들이 회사의 의사 결정에 적극적으로 참여할 수 있게 하였다. 그리고 목표를 달성한 구성원에게 추가 수당이라는 보상을 제공하였다. 기업 K의 구성원들은 대부분 목표를 달성하지 못했고, 기업 P의 구성원들은 대부분 목표를 쉽게 달성하였다.\end{kfEvidence}
\begin{kfChoices}
\kfChoice 데시와 라이언은 기업 P의 구성원들이 목표를 달성한 이유를 을이 구성원에게 제공한 보상 때문이라고 여기겠군.
\kfChoice 허시와 블랜차드는 갑이 구성원들과 인간적 관계를 형성하면서도 명확한 지시를 내려야 한다고 보겠군.
\kfChoice 로크는 을이 구성원들의 동기를 자극하기 위해서는 현재보다 더 어려운 목표를 제시해야 한다고 생각하겠군.
\end{kfChoices}
\end{kfQBlock}

\kfPassageNote{문항 6 / 의도·목적}{%
미니멀리즘 음악은 1960년대 이후 미국에서 태동한 현대 음악의 한 흐름으로, 라 몬테 영, 테리 라일리, 스티브 라이히, 필립 글래스 등이 대표적 작곡가로 꼽힌다. 이들은 소수의 \uline{음형}(몇 개의 음이 연속하여 어떤 가락이나 악곡의 요소가 되는 음의 모양)이나 단순한 \uline{패턴}을 반복하는 방식으로 기존 음악의 복잡한 서사나 화성 전개를 과감히 축소하여, 감상자가 소리 그 자체를 경험하도록 유도한다.

전통적인 서양 음악에서 반복은 주제를 강조하여 이를 서사적으로 강화하기 위해 사용되었다. 예컨대 제시부-발전부-재현부로 구성된 소나타 형식은 제시부에서 등장한 주제가 발전부에서 여러 차례 변주되다가 재현부에서 제시부에 등장한 형태로 회귀함으로써, 작곡가가 설정한 주제를 확인 및 강화하는 구조를 이룬다. 철학자 \uline{들뢰즈}의 표현을 빌리면, 이와 같은 ㉠ \uline{‘같은 것의 반복’}에서는 특정 주제가 중심에 놓이고 음악의 각 부분은 그 주제를 재확인하는 역할만…(중략)
}
\begin{kfQBlock}{06}{객관식}
\kfQStem{㉠과 ㉡을 이해한 내용으로 가장 적절한 것은?}
\begin{kfChoices}
\kfChoice ㉠은 주제를 강조함으로써 통일성을 강화하는 반면, ㉡은 표면적 동일성 속에서 미세한 음향적 차이를 지속적으로 생성한다.
\kfChoice ㉠은 주제를 강조함으로써 통일성을 강화하는 반면, ㉡은 표면적 동일성 속에서 미세한 음향적 차이를 지속적으로 생성한다.
\kfChoice ㉠은 작곡가가 의도한 주제를 강화하는 반복이지만, ㉡은 작곡가의 의도를 감상자가 수동적으로 재확인하게 만드는 반복이다.
\kfChoice ㉠은 반복 과정에서 주제의 변화가 일어나지 않는 반면, ㉡은 반복 속에서 주제가 점진적으로 변화되어 새로운 주제로 발전한다.
\end{kfChoices}
\end{kfQBlock}

\kfPassageNote{문항 7 / 상관관계}{%
(가) 생수병이나 찻주전자의 내용물을 따를 때 액체가 용기 바깥 면을 타고 흘러 당황한 경험이 있다. 왈칵 따를 때에는 잘 쏟아지던 액체가 속력이 줄면 바깥 면을 타고 흘러내리는데, 이런 현상을 \uline{드리블링} 또는 찻주전자 효과라고 부른다. 찻주전자 효과를 체계적으로 연구한 대표 사례는 2010년 뒤에즈 연구 팀의 실험이다. 이들은 원인으로 주둥이의 젖음성에 주목했고, 이 현상은 점성과 무관하다고 \uline{ⓐ밝혔다}. 젖음성은 고체 표면에서 액체가 퍼지는 정도를 뜻하며, 젖음성이 클수록 액체는 더 얇게 퍼진다. 물에 대한 젖음성이 큰 경우를 친수성, 작은 경우를 소수성이라고 한다. 연구 팀은 유속, 주둥이 젖음성, 가장자리 곡률, 점성에 따라 이탈 각이 어떻게 달라지는지 실험으로 확인했다. 주둥이에서 이탈하는 액체가 연직선과 이루는 각을 이탈 각이라 하며, 이탈 각이 작을수록 찻주전자 효과가 커진다. 유속이 충분히 크면 이탈 각은 90°에 가깝게 유지되지만 유속이 줄어들면 …(중략)
}
\begin{kfQBlock}{07}{객관식}
\kfQStem{(가)의 조건 관계를 바르게 적용한 것은?}
\begin{kfChoices}
\kfChoice 유속을 높이고 곡률 반경을 줄이면 이탈에 유리한 조건을 만들 수 있다.
\kfChoice 유속을 높이고 곡률 반경을 줄이면 이탈에 유리한 조건을 만들 수 있다.
\kfChoice 유속을 낮추고 젖음성을 키우면 이탈 각이 커지기 쉽다.
\kfChoice 점성만 바꾸면 항상 큰 차이가 난다.
\end{kfChoices}
\end{kfQBlock}

\kfPassageNote{문항 8 / 필요·충분 조건}{%
최근까지 일반적인 컴퓨터 환경에서는 운영 체제를 제어하거나 사용자 인터페이스를 처리하고, 문서를 작성하거나 웹을 탐색하는 등 다양한 범용 작업이 주요 연산 대상이었다. 이러한 작업은 명령어의 종류가 많고 그 순서가 수시로 바뀌며, 연산 간 종속성이 크다는 특징이 있다. 컴퓨터의 연산은 주로 여러 개의 코어로 구성된 \uline{㉠중앙 처리 장치(CPU)}에서 수행된다. 코어는 컴퓨터의 두뇌 역할을 하는 것으로 명령어를 해석하고 각종 연산을 실행한다. CPU는 여러 개의 고성능 코어로 구성되어 복잡한 제어 흐름을 정밀하게 조율할 수 있도록 설계되었다. 그러나 인공 지능 학습, 고해상도 영상 처리, 과학 기술 관련 계산 등 대규모 데이터를 대상으로 동일한 연산을 반복 수행해야 하는 작업이 증가하면서 기존의 연산 환경에 변화가 생겼다. CPU는 코어 수에 한계가 있고 개별 코어가 복잡한 제어 흐름에 맞추어 다양한 연산을 처리하도록 설계되어 있기 때문에 동일한 연산을 대량으로 반복 수행해야 …(중략)
}
\begin{kfQBlock}{08}{객관식}
\kfQStem{<보기>를 바탕으로 GPU 작동을 이해한 내용으로 적절하지 \underline{않은} 것은?}
\begin{kfEvidence}연산 유닛 A에는 스레드 블록 4개, B와 C에는 각각 1개가 배치되어 있다. 블록 내부 스레드는 동일 명령어를 실행하며, 같은 메모리 주소를 동시에 접근하는 작업이 일부 포함되어 있다.\end{kfEvidence}
\begin{kfChoices}
\kfChoice 전역 메모리는 연산 중 접근 자체가 불가능하므로 결과 저장 외에는 절대 사용할 수 없다.
\kfChoice 전역 메모리는 연산 중 접근 자체가 불가능하므로 결과 저장 외에는 절대 사용할 수 없다.
\kfChoice A에 작업이 과도하게 몰리면 B·C 유닛보다 대기 가능성이 커질 수 있다.
\kfChoice 블록 내부에서 동일 주소 동시 접근이 많으면 메모리 충돌로 속도가 떨어질 수 있다.
\end{kfChoices}
\end{kfQBlock}

\kfPassageNote{문항 9 / 결합·분리}{%
현대 경제는 과거와는 다른 양상으로 움직인다. 기술은 시시각각 변하고 세계 금융 시장은 얽히고설켜 예측하기 어려운 변동성을 보인다. 이러한 복잡한 경제 상황을 이해하기 위해 미국의 산타페 연구소를 중심으로 등장한 복잡계 경제학은 경제 주체가 언제나 가장 합리적 선택을 한다는 명제는 진리가 아니라고 주장한다. 그러면서 경제에 대한 근본적인 생각의 전환을 요구한다.

복잡계 경제학에서는 현대 경제를 수학적 계산이나 도식적 모델의 결과물이 아니라, 동적으로 변화하는 복잡한 시스템이라고 본다. 또한 현대 경제의 특징은 다양한 주체들이 제한된 정보 속에서 상호 작용하며 예측 불가능한 질서를 만들어 내는 비정형성이라고 설명한다. 어떤 기술이 표준이 될지, 어떤 기업이 시장을 주도할지, 어떤 정책이 장기적으로 성공할지는 누구도 확신할 수 없다는 것이다. 예를 들어 기술의 선택에서는 효율성보다는 경로 의존성이 더 크게 작용하기 때문에 새로운 기술 개발의 효과를 예측하는 것은 매우 어렵다. 여기서 …(중략)
}
\begin{kfQBlock}{09}{객관식}
\kfQStem{[01] 현대 경제에 대한 '복잡계 경제학'의 관점에 부합하지 않는 것은?}
\begin{kfChoices}
\kfChoice 자발적으로 형성된 시장의 질서는 개별 주체들에서 보편적으로 발견되는 특성이다.
\kfChoice 자발적으로 형성된 시장의 질서는 개별 주체들에서 보편적으로 발견되는 특성이다.
\kfChoice 경로 고착은 정책이나 제도에서도 나타난다.
\kfChoice 경제 주체들의 상호 작용에 의해 창발이 나타난다.
\end{kfChoices}
\end{kfQBlock}

\kfPassageNote{문항 10 / 조건 왜곡}{%
에너지 대사는 생명 유지를 위해 필요한 에너지를 얻고 사용하는 과정으로 인체에 필요한 에너지는 대부분 음식물 섭취를 통해 이루어진다. 섭취를 통해 체내에 \uline{ⓐ유입된} 음식물은 그대로 에너지로 소비되지 않고 소화 과정을 거쳐 체내 각 조직에 공급된다. 음식물에 함유된 탄수화물, 단백질, 지방은 각각 일련의 과정을 거쳐 세포 활동에 직접 사용되는 에너지원인 ATP로 전환된 후 필요한 조직에 사용된다. 에너지로 바로 사용되지 않은 탄수화물 일부는 포도당의 형태로 간이나 근육에 저장되었다가 필요할 때 다시 ATP로 전환된다. 섭취된 에너지와 인체가 소비하는 에너지는 균형을 유지해야 하지만, 섭취한 에너지의 양이 소비되는 양보다 많아지는 상태가 지속되면 남는 에너지는 지방산으로 전환되어 지방 세포에 저장된다. 이러한 상태가 \uline{ⓑ심화되면} 지방 세포의 크기가 커지게 된다. 지방산이 과도하게 축적되면 체내의 에너지 조절 체계가 정상적으로 작동하지 않게 되고 비만을 비롯한 다양한…(중략)
}
\begin{kfQBlock}{10}{객관식}
\kfQStem{윗글을 읽고 알 수 있는 내용으로 적절하지 \underline{않은} 것은?}
\begin{kfChoices}
\kfChoice 간·근육 저장 포도당은 필요 시 그대로 즉시 에너지원으로 활용된다.
\kfChoice 간·근육 저장 포도당은 필요 시 그대로 즉시 에너지원으로 활용된다.
\kfChoice 신체 에너지 요구 충족에는 ATP가 필요하다.
\kfChoice 갈색 지방 세포 양이 줄면 추위 적응력이 떨어질 수 있다.
\end{kfChoices}
\end{kfQBlock}

\kfPassageNote{문항 11 / 긍부정 반전}{%
\#\#\# [22～27] 다음 글을 읽고 물음에 답하시오.

(가)   첩첩산중에도 없는 마을이 여긴 있습니다. 잎 진 사잇길 저 모랫둑, 그 너머 강기슭에서도 보이진 않습니다. 허방다리* 들어내면 보이는 마을.   갱 속 같은 마을. ㉠꼴깍, 해가, 노루꼬리 해가 지면 집집마다 봉당에 불을 켜지요. 콩깍지, 콩깍지처럼 후미진 외딴집, 외딴집에도 불빛은 앉아 이슥토록 창문은 모과빛입니다.   기인 밤입니다. 외딴집 노인은 홀로 잠이 깨어 출출한 나머지 무우를 깎기도 하고 고구마를 깎다, 문득 바람도 없는데 시나브로 풀려 풀려 내리는 짚단, 짚오라기의 설레임을 듣습니다. 귀를 모으고 듣지요. ㉡후루룩 후루룩 처마 깃에 나래 묻는 이름 모를 새, 새들의 온기를 생각합니다. 숨을 죽이고 생각하지요.   참 오래오래, 노인의 자리맡에 밭은기침 소리도 없을 양이면 벽 속에서 겨울 귀뚜라미는 울지요. 떼를 지어 웁니다, 벽이 무너지라고 웁니다.   어느덧 밖에는 눈발이라도 치는지, 펄펄 함박눈이라…(중략)
}
\begin{kfQBlock}{11}{객관식}
\kfQStem{(나)에 대한 설명으로 적절하지 \underline{않은} 것은?}
\begin{kfChoices}
\kfChoice 4연에서 ‘외로이 자랐다’와 이어진 ‘하얀 넋’은 ‘붉은 발자욱’에 함축된 정서와 상반되는 의미를 이끌어 내고 있다.
\kfChoice 4연에서 ‘외로이 자랐다’와 이어진 ‘하얀 넋’은 ‘붉은 발자욱’에 함축된 정서와 상반되는 의미를 이끌어 내고 있다.
\kfChoice 1연에서 ‘연’과 ‘연실’의 모습에 빗대어 ‘내 어린 날’의 기억을 ‘아슴풀하다’라고 표현하고 있다.
\kfChoice 2연에서 ‘조매롭고’로 표현된 ‘연실’의 긴장은 3연에서 연실이 ‘바람 일어 끊어지던 날’의 정서를 고조하고 있다.
\end{kfChoices}
\end{kfQBlock}

\kfPassageNote{문항 12 / 비교 대상}{%
\#\#\# [35 \textasciitilde{} 38] 다음 글을 읽고 물음에 답하시오.

(가)   청춘에 병이 들어 공산에 누웠더니   한 조각 남은 꿈에 호랑나비의 날개 빌려   강한 바람에 고삐 잡혀 남포(南浦)로 내려가니   초선도(招仙島)가 어디메뇨 개암정(皆巖亭)이 여기로다   어주(漁舟)를 흘러 타고 백구(白鷗)에게 길을 물어   굽이굽이 돌아드니 수석(水石)도 맑고 곱다   시냇가에 누은 돌은 석국(石局)\textit{처럼 벌려 잇고}   \textit{돌 틈에 솟은 물에 표주박 술잔이 띄워 있다}   \textit{등나무를 후려잡고 석국을 디뎌 밟아}   \textit{높은 창을 바삐 열고 주인 영감께 인사하니}   \textit{건장하도다 뛰어난 풍채 거룩하도다 고령의 건강}   \textit{의좋은 삼 형제는 아버님과 함께 놀고}   \textit{색동옷 입은 아이 재롱 노래자(老萊子)가 부러울소냐}   \textit{거문고 비파와 서책은 책상 위에 둘러 있고}   \textit{뛰어난 자식들은 뜰 앞에 벌려 있다}   \textit{인사를 마친 후에 난간에 빗겨 앉아}   *원근(遠近)의 산천(山川)을 한눈에 굽…(중략)
}
\begin{kfQBlock}{12}{객관식}
\kfQStem{(가)와 (나)의 공통점으로 가장 적절한 것은?}
\begin{kfChoices}
\kfChoice 비유적 표현을 통해 대상의 속성을 드러내고 있다.
\kfChoice 비유적 표현을 통해 대상의 속성을 드러내고 있다.
\kfChoice 색채의 대비를 통해 주제 의식을 강조하고 있다.
\kfChoice 문답의 방식을 통해 인식의 전환을 드러내고 있다.
\end{kfChoices}
\end{kfQBlock}

\kfPassageNote{문항 13 / 주체·객체}{%
[22 \textasciitilde{} 27] 다음 글을 읽고 물음에 답하시오.

(가)   몰아라 어서 보자 총석정 어서 보자   총석정 좋단 말을 일찍이 들었거니   바람 불면 못 보려니 몰아라 어서 보자   벽해 위의 높은 집이 저것이 총석정인가   올라 보니 후면이라 전면으로 보오리라   배 대어라 사공들아 풍랑이 일지 않아   층파로 돌아 저어 총석 전면 보게 하라   배 띄워라 굽이마다 따라 저어 볼 양이면   영소전 태을궁*을 지으려고 경영턴가*   \textit{돌기둥 천백 개를 육모로 깎아 내어}   \textit{개개이 묶어 세워 몇 만 년이 되었던지}   \textit{황량한 데 벌였으니 배 없어 못 실린가}   \textit{(중략)}   \textit{하우씨 도끼뿔이 용문을 뚫었으나}   *이 돌*을 만났으면 이같이 깎을세며   영장*이 신묘하여 코끝의 것 찍었으나*   \textit{이 돌을 다듬는다고 이같이 곧을쏘냐}   \textit{어떠한 도끼로 용이히 깎았으며}   *어떠한 승묵*으로 천연히 골랐는고   끈 없이 묶었으되 틈 없이 묶었으며   풀 없이 붙였으되 흔적…(중략)
}
\begin{kfQBlock}{13}{객관식}
\kfQStem{㉠, ㉡에 대한 설명으로 가장 적절한 것은?}
\begin{kfChoices}
\kfChoice ㉠은 화자가 위치한 공간적 배경을 제시하고 있다.
\kfChoice ㉠은 화자가 위치한 공간적 배경을 제시하고 있다.
\kfChoice ㉡은 세상과 거리를 두려는 글쓴이의 태도와 관련이 있다.
\kfChoice ㉡은 아이들이 파도를 피해 움직이는 모습을 나타내고 있다.
\end{kfChoices}
\end{kfQBlock}

\kfPassageNote{문항 14 / 내용·방법}{%
[18 \textasciitilde{} 23] 다음 글을 읽고 물음에 답하시오.

(가)   삼 년을 임을 떠나 해도(海島)에 유배되니   ㉠ 내 언제 무심하여 임에게 득죄했나   임이 언제 박정(薄情)하여 날 대접 소홀히 했나   내 얼굴 고왔던지 질투하는 건 뭇 여자로다   유한한* 이내 몸을 음란하다 이르로세   (중략)   긴 소매 들고 앉아 옛 잘못을 헤아리니   우직하기 본성이오 망령됨도 내 죄로되   근본을 생각하면 임 위한 정성일세   일월 같은 우리 임이 거의 아니 굽어볼까   날 살리신 이 은혜를 결초(結草)하기 생각하나   광주리의 가을 부채 어느 날 다시 날꼬   황금을 못 얻으니 장문부*를 어이 사리*   \textit{마름과 연(蓮)으로 옷을 짓고 부용(芙蓉)으로 치마 지어}   \textit{상자 안에 두어신들 눌 위하여 단장할꼬}   \textit{고향에 돌아갈 꿈 벽해(碧海)를 밟아 건너}   \textit{옥루(玉樓) 높은 곳에 밤마다 임을 모셔}   \textit{일당우불에 수답이 여향하니}   가까이 다가앉아 귀신을 묻던 가태부 이러한가…(중략)
}
\begin{kfQBlock}{14}{객관식}
\kfQStem{㉠ \textasciitilde{} ㉤의 표현상의 특징으로 적절하지 \underline{않은} 것은?}
\begin{kfChoices}
\kfChoice ㉢ : 과장법을 사용하여 임을 향한 사랑을 포기해야 하는 것에 대한 화자의 절망감을 강조하고 있다.
\kfChoice ㉢ : 과장법을 사용하여 임을 향한 사랑을 포기해야 하는 것에 대한 화자의 절망감을 강조하고 있다.
\kfChoice ㉠ : 대구적 표현을 사용하여 운율감을 조성하고 있다.
\kfChoice ㉡ : 감각적 심상을 활용하여 화자의 그리움을 부각하고 있다.
\end{kfChoices}
\end{kfQBlock}

\kfPassageNote{문항 15 / 내용 왜곡}{%
\#\#\# [18 \textasciitilde{} 23] 다음 글을 읽고 물음에 답하시오.

(가)   녹양방초 언덕에 소 먹이는 저 아이야   인간 영욕을 아는가 모르는가   인생 백년이 풀 끝에 이슬이라   삼만 육천 일이 그 아니 보잘것없는가   하물며 장수 단명이 운명이어니 사생(死生)을 정할쏘냐   여관 같은 세상에 하루살이같이 나왔다가   ㉠ 공명도 못 이루고 초목같이 썩어지면   ⓐ 공산 백골이 그 아니 느꺼우냐   하늘의 뜻을 이어 법칙을 세움은 옛 성인의 사업이요   아름다운 이름을 후세에 전함은 대장부의 할 일이라   생애는 유한하고 사일(死日)은 무궁하니   유한한 생애로 썩지 않을 이름을   영구히 전하여 ⓑ 천지와 함께 무궁하려고   (중략)   어와 그 뉘신고 어떠한 사람인고   형용이 초췌하니 초나라 대부 굴원이신가   잔혼이 영락하니 학사 유자후이신가*   눈썹을 찡그리시니 근심이 많으신가   발끝으로 서시니 어디를 바라는고   아름다운 기약을 바라는가 이별의 슬픔이 중하신가   ⓒ 해…(중략)
}
\begin{kfQBlock}{15}{객관식}
\kfQStem{(나)의 ‘나’의 관점에서 ㉠, ㉡을 이해한 내용으로 가장 적절한 것은?}
\begin{kfChoices}
\kfChoice ㉠을 이루는 일로부터 물러나 ㉡을 즐겨야 하는 때도 있다.
\kfChoice ㉠을 이루는 일로부터 물러나 ㉡을 즐겨야 하는 때도 있다.
\kfChoice ㉠과 ㉡을 모두 멀리해야 하는 때도 있다.
\kfChoice ㉠에서 생겨나는 ㉡의 즐거움에 만족하며 살아야 한다.
\end{kfChoices}
\end{kfQBlock}

\kfPassageNote{문항 16 / 과잉 해석}{%
\#\#\# [35 \textasciitilde{} 38] 다음 글을 읽고 물음에 답하시오.

(가)   청춘에 병이 들어 공산에 누웠더니   한 조각 남은 꿈에 호랑나비의 날개 빌려   강한 바람에 고삐 잡혀 남포(南浦)로 내려가니   초선도(招仙島)가 어디메뇨 개암정(皆巖亭)이 여기로다   어주(漁舟)를 흘러 타고 백구(白鷗)에게 길을 물어   굽이굽이 돌아드니 수석(水石)도 맑고 곱다   시냇가에 누은 돌은 석국(石局)\textit{처럼 벌려 잇고}   \textit{돌 틈에 솟은 물에 표주박 술잔이 띄워 있다}   \textit{등나무를 후려잡고 석국을 디뎌 밟아}   \textit{높은 창을 바삐 열고 주인 영감께 인사하니}   \textit{건장하도다 뛰어난 풍채 거룩하도다 고령의 건강}   \textit{의좋은 삼 형제는 아버님과 함께 놀고}   \textit{색동옷 입은 아이 재롱 노래자(老萊子)가 부러울소냐}   \textit{거문고 비파와 서책은 책상 위에 둘러 있고}   \textit{뛰어난 자식들은 뜰 앞에 벌려 있다}   \textit{인사를 마친 후에 난간에 빗겨 앉아}   *원근(遠近)의 산천(山川)을 한눈에 굽…(중략)
}
\begin{kfQBlock}{16}{객관식}
\kfQStem{<보기>를 참고하여 (가), (나)를 감상한 내용으로 적절하지 \underline{않은} 것은? [3점]}
\begin{kfChoices}
\kfChoice (가)의 ‘우리 일문’에 ‘장수한 어른이 많기도 많’다며 ‘거룩’하다고 감탄하는 것에서 가문 의식을 실현하는, (나)의 ‘내면은 잊’고 ‘텅 비고 밝은 본체’를 기르는 ‘이군’의 모습에서 인격 수양을 실천하는 공간으로서의 자연을 엿볼 수 있군.
\kfChoice (가)의 ‘우리 일문’에 ‘장수한 어른이 많기도 많’다며 ‘거룩’하다고 감탄하는 것에서 가문 의식을 실현하는, (나)의 ‘내면은 잊’고 ‘텅 비고 밝은 본체’를 기르는 ‘이군’의 모습에서 인격 수양을 실천하는 공간으로서의 자연을 엿볼 수 있군.
\kfChoice (가)의 ‘돌 틈에 솟은 물에 표주박 술잔이 띄워 있’는 모습에서 풍류를 즐기는 공간으로서의 자연을 엿볼 수 있군.
\kfChoice (가)의 ‘봉우리’들이 ‘수려하고 기암괴석’이 ‘절경이로구나’라고 감탄하는 것에서 미적 가치를 인식하는 공간으로서의 자연을 엿볼 수 있군.
\end{kfChoices}
\end{kfQBlock}

\kfPassageNote{문항 17 / 범위 오류}{%
\#\#\# 2024년 11월 고3 수능

[22-27] 다음 글을 읽고 물음에 답하시오.

(가)   배를 민다   배를 밀어보는 것은 아주 드문 경험   희번덕이는 잔잔한 가을 바닷물 위에   배를 밀어넣고는   온몸이 아주 추락하지 않을 순간의 한 허공에서   밀던 힘을 한껏 더해 밀어주고는   아슬아슬히 배에서 떨어진 손, 순간 환해진 손을   허공으로부터 거둔다

사랑은 참 부드럽게도 떠나지   뵈지도 않는 길을 부드럽게도

배를 한껏 세게 밀어내듯이 슬픔도   그렇게 밀어내는 것이지

배가 나가고 남은 빈 물 위의 흉터   잠시 머물다 가라앉고

그런데 오, 내 안으로 들어오는 배여   아무 소리 없이 밀려들어오는 배여   - 장석남, <배를 밀며>

(나)   당신……, 당신이라는 말 참 좋지요, 그래서 불러봅니다 킥킥거리며 한때 적요로움의 울음이 있었던 때, 한 슬픔이 문을 닫으면 또 한 슬픔이 문을 여는 것을 이만큼 살아옴의 상처에 기대, 나 킥킥……, 당신을 부릅니다 단…(중략)
}
\begin{kfQBlock}{17}{객관식}
\kfQStem{문항별}
\begin{kfChoices}
\kfChoice 번은 '통제할 수 없는 익명의 욕구'를 상대를 향한 사랑이 운명적이어서 멈출 수 없다는 뜻으로 해석했습니다.
\kfChoice 번은 '통제할 수 없는 익명의 욕구'를 상대를 향한 사랑이 운명적이어서 멈출 수 없다는 뜻으로 해석했습니다.
\kfChoice 번은 배가 들어오는 것을 '대상과의 재회가 예상대로 이루어짐'이라고 해석했습니다.
\kfChoice 번은 '당신'을 화자의 내면에 사는 '병자'로서 연민의 대상이라고 설명했습니다.
\end{kfChoices}
\end{kfQBlock}

\kfPassageNote{문항 18 / 시점·화자}{%
[16\textasciitilde{}20] 다음 글을 읽고 물음에 답하시오.

(가)

어버이 낳으시고 임금이 먹이시니    낳은 덕 먹인 은을 다 갚고자 하였더니    숙연히 칠십이 넘으니 할 일 없어 하노라

<제1수>

연하에 깊이 곤 병 약이 효험 없어    강호에 버리언져 십 년 밖이 되었어라    그러나 이제 다 못 죽음도 그 성은인가 하노라

<제3수>

달 밝고 바람 자니 물결이 비단일다    단정*을 비껴 놓아 오락가락 하는 흥을    ⓐ 백구야 하 즐겨 말고려 세상 알까 하노라

<제5수>

피 소주 무우저리* 우습다 어른 대접    남은셔 이른 말이 초초타* 하건마는    두어라 이도 내 분이니 분내사*인가 하노라

<제8수>

- 나위소, 「강호구가」 -

* 단정: 작은 배. * 피 소주 무우저리: 풀로 만든 소주와 무절임. * 초초타: 보잘것없이 초라하다. * 분내사: 분수에 맞는 일.

(나)

들은 지 오래더니 보았구나 백상루    야속하다 강산이 날 기다리고 있었던가    처음으로…(중략)
}
\begin{kfQBlock}{18}{객관식}
\kfQStem{<보기>를 참고하여 (다)를 감상한 내용으로 적절하지 \underline{않은} 것은?}
\begin{kfChoices}
\kfChoice ‘논밭 일에 소를 부릴 때면 으레 그 노래를 부르는 것에서 현실적 삶을 배제한 산골의 모습이 드러나는군.
\kfChoice ‘논밭 일에 소를 부릴 때면 으레 그 노래를 부르는 것에서 현실적 삶을 배제한 산골의 모습이 드러나는군.
\kfChoice 집 대부분이 ‘쓰러질 듯한 헌 초가’라는 것에서 산골 사람들의 궁핍한 삶의 모습이 짐작되는군.
\kfChoice ‘산천의 풍경’을 ‘하나 흠잡을 데 없는 귀여운 전원’이라고 한 것은 산골을 낭만적인 전원으로 제시한 것이군.
\end{kfChoices}
\end{kfQBlock}



\clearpage

\kfChapterCover{4}{비트겐슈타인1 ch04}{Chapter 4}{이번 챕터: 문법 → 문학 5작품 → 비문학 5지문 → 패턴 워크북.}
\begin{kfChapterAreaList}
\kfChapterAreaItem{문법}{kfAreaGrammar}{음운 / 구개음화·된소리되기}
\kfChapterAreaItem{문학}{kfAreaLiterature}{「완장」 — 윤흥길 외 4편}
\kfChapterAreaItem{비문학}{kfAreaNonlit}{「법이란 무엇인가」 외 4편}
\kfChapterAreaItem{실력 확인}{kfAreaWeekly}{패턴 워크북 — 총 18문항}
\end{kfChapterAreaList}

\kfStartGrammar{음운 / 구개음화·된소리되기}


\kfSecHeader{kfAreaGrammar}{문제}{}

\kfBeginMC
\par\needspace{25\baselineskip}
\kfPassageFull{구개음화 조건}{%
구개음화란 끝소리가 'ㄷ, ㅌ'인 형태소가 모음 'ㅣ'나 반모음 'ㅣ'로 시작하는 형식 형태소(조사, 접미사, 어미)와 만날 때, 'ㄷ'이 'ㅈ'으로, 'ㅌ'이 'ㅊ'으로 바뀌는 현상이다. 예를 들어 '맏이'는 [마지], '같이'는 [가치], '해돋이'는 [해도지]로 발음된다. 구개음화는 자음 'ㄷ, ㅌ'이 모음 'ㅣ'와 조음 위치가 비슷한 센입천장소리(경구개음) 'ㅈ, ㅊ'으로 바뀌는 것이므로 동화 현상에 해당한다. 다만, 실질 형태소 앞에서는 구개음화가 일어나지 않으며, 하나의 형태소 내부에서도 일어나지 않는다.
}\par\nobreak\nopagebreak[4]\kfResetAreaFlag
\begin{kfQBlock}{01}{객관식}
\kfQStem{윗글을 바탕으로 할 때, 구개음화가 일어나는 것만을 <보기>에서 모두 고른 것은?}
\begin{kfEvidence}ㄱ. 밭이 → [바치] \\\relax
ㄴ. 굳이 → [구지] \\\relax
ㄷ. 잔디 → [잔지] \\\relax
ㄹ. 붙이다 → [부치다]\end{kfEvidence}
\begin{kfChoices}
\kfChoice ㄱ, ㄴ
\kfChoice ㄱ, ㄴ, ㄹ
\kfChoice ㄱ, ㄷ, ㄹ
\kfChoice ㄴ, ㄷ
\kfChoice ㄱ, ㄴ, ㄷ, ㄹ
\end{kfChoices}
\end{kfQBlock}

\par\needspace{25\baselineskip}
\kfPassageFull{된소리되기 규칙}{%
된소리되기란 예사소리(ㄱ, ㄷ, ㅂ, ㅅ, ㅈ)가 된소리(ㄲ, ㄸ, ㅃ, ㅆ, ㅉ)로 바뀌는 현상이다. 된소리되기가 일어나는 환경은 여러 가지가 있다. 첫째, 받침 'ㄱ, ㄷ, ㅂ' 뒤에 예사소리 'ㄱ, ㄷ, ㅂ, ㅅ, ㅈ'이 오면 된소리로 발음된다. 예를 들어 '국가'는 [국까], '닫지'는 [닫찌]로 발음된다. 둘째, 용언 어간의 받침 'ㄴ, ㅁ' 뒤에 어미의 첫소리 'ㄱ, ㄷ, ㅅ, ㅈ'이 오면 된소리로 발음된다. 예를 들어 '신다'는 [신따], '삼고'는 [삼꼬]로 발음된다. 셋째, 한자어에서 'ㄹ' 받침 뒤에 'ㄷ, ㅅ, ㅈ'이 오면 된소리로 발음된다.
}\par\nobreak\nopagebreak[4]\kfResetAreaFlag
\begin{kfQBlock}{02}{객관식}
\kfQStem{윗글을 바탕으로 할 때, <보기>의 밑줄 친 부분에서 된소리되기가 일어나지 않는 것은?}
\begin{kfEvidence}(가) 책방 [책빵] \\\relax
(나) 신고 [신꼬] \\\relax
(다) 발전 [발쩐] \\\relax
(라) 바람도 [바람또] \\\relax
(마) 국수 [국쑤]\end{kfEvidence}
\begin{kfChoices}
\kfChoice (가)
\kfChoice (나)
\kfChoice (다)
\kfChoice (라)
\kfChoice (마)
\end{kfChoices}
\end{kfQBlock}

\par\needspace{25\baselineskip}
\kfPassageFull{음운 변동 유형 분류}{%
음운 변동은 음운의 개수 변화 여부에 따라 네 가지 유형으로 나뉜다. 교체(바꿈)는 한 음운이 다른 음운으로 바뀌는 것으로, 음운의 개수가 변하지 않는다. 탈락(빠짐)은 음운 하나가 없어지는 것으로, 음운의 개수가 줄어든다. 첨가(더함)는 없던 음운이 새로 생기는 것으로, 음운의 개수가 늘어난다. 축약(줄임)은 두 음운이 합쳐져 하나의 음운이 되는 것으로, 음운의 개수가 줄어든다. 음절의 끝소리 규칙, 비음화, 유음화, 구개음화, 된소리되기는 모두 교체에 해당하고, 자음군 단순화·'ㅎ' 탈락·'ㄹ' 탈락·모음 탈락 등은 탈락에 해당한다.
}\par\nobreak\nopagebreak[4]\kfResetAreaFlag
\begin{kfQBlock}{03}{객관식}
\kfQStem{윗글을 바탕으로 할 때, <보기>의 ㉠\textasciitilde{}㉣에 해당하는 음운 변동 유형을 바르게 짝지은 것은?}
\begin{kfEvidence}㉠ 꽃 → [꼳] (ㅊ→ㄷ) \\\relax
㉡ 놀- + -는 → 노는 (ㄹ 탈락) \\\relax
㉢ 좋- + -아 → 좋아[조아] (ㅎ 탈락) \\\relax
㉣ 국밥 → [국빱] (ㅂ→ㅃ)\end{kfEvidence}
\begin{kfChoices}
\kfChoice ㉠ 교체, ㉡ 교체, ㉢ 탈락, ㉣ 교체
\kfChoice ㉠ 탈락, ㉡ 탈락, ㉢ 탈락, ㉣ 교체
\kfChoice ㉠ 교체, ㉡ 탈락, ㉢ 탈락, ㉣ 교체
\kfChoice ㉠ 교체, ㉡ 탈락, ㉢ 교체, ㉣ 첨가
\kfChoice ㉠ 축약, ㉡ 탈락, ㉢ 탈락, ㉣ 첨가
\end{kfChoices}
\end{kfQBlock}

\par\needspace{25\baselineskip}
\kfPassageFull{교체·탈락·첨가·축약 구분}{%
음운 변동의 네 유형을 정확히 구분하기 위해서는 변동 전후의 음운 개수를 비교하는 것이 효과적이다. 교체는 음운 A가 음운 B로 바뀌는 것이므로 음운의 총 개수가 변하지 않는다. 탈락은 음운 하나가 없어지는 것이므로 음운의 총 개수가 1개 줄어든다. 첨가는 없던 음운이 새로 생기는 것이므로 음운의 총 개수가 1개 늘어난다. 축약은 두 음운이 합쳐져 하나의 새로운 음운이 되는 것이므로 음운의 총 개수가 1개 줄어든다. 예를 들어 '놓다'가 [노타]로 발음될 때, 'ㅎ+ㄷ'이 'ㅌ'으로 합쳐지는 자음 축약(거센소리되기)이 일어나 음운이 1개 줄어든다.
}\par\nobreak\nopagebreak[4]\kfResetAreaFlag
\begin{kfQBlock}{04}{객관식}
\kfQStem{윗글을 바탕으로 할 때, 음운 변동 전후의 음운 개수 변화가 나머지와 \underline{다른 하나}는?}
\begin{kfChoices}
\kfChoice 입학 → [이팍] (ㅂ+ㅎ→ㅍ, 자음 축약)
\kfChoice 놓다 → [노타] (ㅎ+ㄷ→ㅌ, 자음 축약)
\kfChoice 닦다 → [닥따] (ㄲ→ㄱ, ㄷ→ㄸ)
\kfChoice 좋다 → [조타] (ㅎ+ㄷ→ㅌ, 자음 축약)
\kfChoice 축하 → [추카] (ㄱ+ㅎ→ㅋ, 자음 축약)
\end{kfChoices}
\end{kfQBlock}

\par\needspace{25\baselineskip}
\kfPassageFull{구개음화+된소리되기 종합}{%
구개음화와 된소리되기는 모두 음운의 교체에 해당한다. 그러나 두 현상은 발생 환경이 다르다. 구개음화는 자음과 모음 사이에서 일어나는 동화 현상이고, 된소리되기는 자음과 자음 사이에서 일어나는 현상이다. 또한 구개음화는 'ㄷ, ㅌ'이 센입천장소리 'ㅈ, ㅊ'으로 바뀌는 것이고, 된소리되기는 예사소리가 된소리로 바뀌는 것이다. 한 단어에 구개음화와 된소리되기가 동시에 적용될 수도 있다. 예를 들어 '밭전'에서는 먼저 음절의 끝소리 규칙으로 'ㅌ'이 [ㄷ]으로 바뀐 후, 된소리되기가 적용되어 [받쩐]→[밧쩐]으로 발음된다. 반면 '밭이'에서는 'ㅌ'이 모음 'ㅣ'(형식 형태소) 앞에서 구개음화가 적용되어 [바치]로 발음된다.
}\par\nobreak\nopagebreak[4]\kfResetAreaFlag
\begin{kfQBlock}{05}{객관식}
\kfQStem{윗글을 바탕으로 할 때, <보기>에 대한 설명으로 적절하지 \underline{않은} 것은?}
\begin{kfEvidence}(가) 굳히다 → [구치다] \\\relax
(나) 옆집 → [엽찝] \\\relax
(다) 같이 → [가치] \\\relax
(라) 곧이듣다 → [고지듣따]\end{kfEvidence}
\begin{kfChoices}
\kfChoice (가)에서는 자음 축약(ㄷ+ㅎ→ㅌ) 후 구개음화(ㅌ→ㅊ)가 일어났다.
\kfChoice (나)에서는 음절의 끝소리 규칙과 된소리되기가 일어났다.
\kfChoice (다)에서는 구개음화만 일어났다.
\kfChoice (라)에서는 구개음화와 된소리되기가 모두 일어났다.
\kfChoice (나)에서 일어난 음운 변동은 모두 동화 현상이다.
\end{kfChoices}
\end{kfQBlock}

\kfEndMC


\kfStartLit{「완장」 — 윤흥길 외 4편}


\kfPassageNote{「완장」 — 윤흥길}{%
그가 수해 복구 대책 위원회에서 완장을 타 가지고 온 날부터 달라지기 시작했다. 흰 바탕에 빨간 글씨로 '수해 복구 대책 위원'이라 쓴 그 완장은 별것 아닌 것 같았지만 사실은 대단한 물건이었다. 수해민의 구호 물자를 배급하고 복구 작업의 인원을 배치하는 권한이 그 완장 하나에 달려 있었으니까.

사람이란 원래부터 그렇게 속이 좁고 고약한 사람은 아니었다. 완장 차기 전만 해도 이웃에게 제법 상냥한 편이었고, 말도 늘 공손했다. 그런데 그 얇다란 완장 하나를 팔뚝에 감은 뒤로 얼굴이 달라지고 말투가 달라지고 걸음걸이가 달라졌다. 처음에는 약간 뻣뻣해지는가 싶더니 급기야 뻔뻔해지고, 뻔뻔하다 못해 난폭해지기까지 했다.

"야, 이놈아, 줄을 서야지 줄을! 금 밖으로 나와 봐, 쥐어박을 놈!"

구호 물자 배급 줄에서 조금만 새치기해도 그의 호통이 떨어졌다. 이전에는 동네 사람 누구에게도 반말 한마디 못 하던 사람이 이제는 나이 지긋한 어른에게조차 "이놈" "저놈" 하며 손가락질을 했다. 완장 앞에서 사람들은 할 말을 잃었다. 분명 같은 마을 사람이고, 분명 어제까지 "형님" "아저씨" 하던 사이인데, 그 얇은 완장 한 조각이 사이에 끼자 사람과 사람 사이가 까마득히 벌어져 버렸다.

그 뒤로 그는 구호 물자 가운데 좋은 것은 먼저 자기 집으로 빼돌리기 시작했고, 복구 작업에서도 자기 집 주변을 먼저 하도록 인원을 배치했다. 아무도 거기에 이의를 달지 못했다. 아니, 이의를 달 엄두를 내지 못했다고 하는 편이 옳다.
}\kfResetAreaFlag
\par\kfMaybeNeedspace{50\baselineskip}\noindent{\kfTipMark\,\,}{\small\color{kfInk} 빈칸에 알맞은 말을 채워 넣으시오. (초성 힌트 참고)}\par\nopagebreak[4]\smallskip\nopagebreak[4]
\begin{kfActTable}{|>{\centering\arraybackslash}m{2.6cm}|m{6cm}|m{4cm}|}
\rowcolor{kfPrimaryLight!50}\textbf{\color{kfPrimary}항목} & \textbf{\color{kfPrimary}내용 (빈칸 채우기)} & \textbf{\color{kfPrimary}답안 작성}\\\hline
갈래 & ㅎㄷ ㅅㅅ, ㄷㅍ ㅅㅅ ① & \kfTBlank \\\hline
배경 & ㅅㅎ ㅂㄱ 현장 ② & \kfTBlank \\\hline
주인공 이전 성격 & ㅅㄴ하고 ㄱㅅ한 소시민 ③ & \kfTBlank \\\hline
완장의 상징 & 사소하지만 인간을 변질시키는 ㄱㄹ의 상징물 ④ & \kfTBlank \\\hline
변화 과정 & 공손 → ㅃㅃ → ㅃㅃ → ㄴㅍ ⑤ & \kfTBlank \\\hline
권력 남용 사례 1 & ㄱㅎ ㅁㅈ를 자기 집으로 빼돌림 ⑥ & \kfTBlank \\\hline
권력 남용 사례 2 & 복구 작업 ㅇㅇ을 자기 집 주변에 우선 배치 ⑦ & \kfTBlank \\\hline
마을 사람들의 반응 & 이의를 달 ㅇㄷ를 내지 못함 ⑧ & \kfTBlank \\\hline
서술 시점 & ㅈㅈㅈ ㅈㄱ 시점 (3인칭) ⑨ & \kfTBlank \\\hline
주제 1 & ㄱㄹ의 ㅅㅅ과 인간 ㅂㅈ ⑩ & \kfTBlank \\\hline
주제 2 & 사소한 ㄱㅇ도 인간을 ㅌㄹ시킬 수 있음 ⑪ & \kfTBlank \\\hline
핵심 대비 & 완장 차기 ㅈ과 ㅎ의 인물 변화 ⑫ & \kfTBlank \\\hline
\end{kfActTable}
\kfSecHeader{kfAreaLiterature}{문제}{}

\kfBeginMC
\begin{kfQBlock}{01}{객관식}
\kfQStem{윗글에 대한 설명으로 적절하지 \underline{않은} 것은?}
\begin{kfChoices}
\kfChoice 인물의 변화 과정을 시간 순서에 따라 서술하고 있다.
\kfChoice 인물의 말투와 행동 묘사를 통해 성격 변화를 드러내고 있다.
\kfChoice 서술자가 작품 속 인물로 등장하여 자신의 체험을 서술하고 있다.
\kfChoice 구체적인 대화를 삽입하여 인물의 달라진 태도를 보여 주고 있다.
\kfChoice 권력 획득 이전과 이후의 대비를 통해 주제 의식을 강화하고 있다.
\end{kfChoices}
\end{kfQBlock}

\begin{kfQBlock}{02}{객관식}
\kfQStem{윗글에서 '완장'이 지닌 상징적 의미로 가장 적절한 것은?}
\begin{kfChoices}
\kfChoice 수해 복구라는 공적 임무에 대한 책임감과 사명감
\kfChoice 마을 공동체를 하나로 결속시키는 연대의 표지
\kfChoice 재난 상황에서 질서를 유지하기 위한 합리적 제도
\kfChoice 사소하지만 인간을 변질시킬 수 있는 권력의 상징물
\kfChoice 국가 권력에 의해 강제된 억압적 통치 도구
\end{kfChoices}
\end{kfQBlock}

\begin{kfQBlock}{03}{객관식}
\kfQStem{윗글에서 주인공의 변화 양상을 가장 적절하게 정리한 것은?}
\begin{kfChoices}
\kfChoice 공손함 → 뻣뻣함 → 뻔뻔함 → 난폭함의 단계적 타락 과정
\kfChoice 순종 → 반항 → 타협 → 체념의 심리 변화 과정
\kfChoice 고립 → 연대 → 갈등 → 화해의 관계 변화 과정
\kfChoice 무관심 → 참여 → 회의 → 이탈로 이어지는 태도 변화 과정
\kfChoice 겸손 → 자각 → 성찰 → 반성의 내면 성장 과정
\end{kfChoices}
\end{kfQBlock}

\begin{kfQBlock}{04}{객관식}
\kfQStem{윗글의 '분명 같은 마을 사람이고, 분명 어제까지 "형님" "아저씨" 하던 사이인데, 그 얇은 완장 한 조각이 사이에 끼자 사람과 사람 사이가 까마득히 벌어져 버렸다.'에 대한 이해로 가장 적절한 것은?}
\begin{kfChoices}
\kfChoice 수해로 인해 마을 사람들이 뿔뿔이 흩어진 물리적 거리감을 표현한 것이다.
\kfChoice 완장이라는 사소한 권력 표지가 공동체의 평등한 관계를 파괴했음을 보여 준다.
\kfChoice 마을 사람들이 주인공을 시기하여 스스로 거리를 둔 것을 나타낸 것이다.
\kfChoice 주인공이 출세하여 마을을 떠나게 됨으로써 생긴 아쉬움을 표현한 것이다.
\kfChoice 재난 상황에서 개인주의가 팽배해지는 현상을 비판한 것이다.
\end{kfChoices}
\end{kfQBlock}

\begin{kfQBlock}{05}{객관식}
\kfQStem{윗글을 바탕으로 <보기>를 감상한 내용으로 적절하지 \underline{않은} 것은?}
\begin{kfEvidence}윤흥길의 「완장」은 권력의 속성을 보여 주는 대표적 작품이다. 작가는 거대한 국가 권력이 아니라 마을 수준의 사소한 권위를 소재로 삼아, 권력이 크기와 무관하게 인간을 변질시킬 수 있음을 보여 주었다. 이러한 접근은 권력의 문제가 특정 체제나 시대에 국한되지 않는 인간 보편의 문제임을 시사한다.\end{kfEvidence}
\begin{kfChoices}
\kfChoice '수해 복구 대책 위원'이라는 직함은 마을 수준의 사소한 권위에 해당하므로, <보기>에서 말한 '사소한 권위'의 구체적 사례로 볼 수 있겠군.
\kfChoice 주인공이 구호 물자를 빼돌리는 행위는 권력이 인간을 변질시킨 결과이므로, <보기>의 '인간을 변질시킬 수 있음'을 뒷받침하겠군.
\kfChoice 마을 사람들이 이의를 달지 못하는 모습은 권력에 대한 복종을 보여 주므로, 권력이 관계를 왜곡하는 양상을 드러내겠군.
\kfChoice 주인공이 완장을 차기 전에도 내면에 폭력적 성향을 갖고 있었으므로, 권력은 기존의 악을 표출시키는 계기일 뿐이겠군.
\kfChoice 이 작품이 특정 정치 체제가 아닌 마을 공동체를 배경으로 한 것은, <보기>에서 말한 '인간 보편의 문제'를 부각하기 위한 설정이겠군.
\end{kfChoices}
\end{kfQBlock}

\kfEndMC

\kfPassageNote{「춘향전」 — 작자 미상}{%
변 사또 분부하되, "네 비록 천기(賤妓)의 딸이로되 관장의 수청을 어이 거역하느냐? 바삐 수청 들라!" 춘향이 여짜오되, "소녀 비록 천기의 딸이오나 일부종사(一夫從事)는 여자의 대절이라. 한 번 몸을 허(許)한 후에 다시 고칠 수 없사오니, 차라리 죽여 주옵소서." 변 사또 크게 노하여, "이년이 관장을 능멸하니 당장 매를 쳐 다스리라!"

(중략)

이때 어사또 분부하시되, "금부도사 급히 나졸을 분부하여 변 사또를 결박하라!" 일시에 나졸이 달려들어 변 사또를 끌어내니, 좌우의 수령 방백이 벌벌 떨어 제 살기를 도모하더라. 어사또 상좌에 높이 앉아 분부하되, "남원 부사 변학도는 탐학무도하여 백성의 재물을 착취하고, 죄 없는 양민을 가두어 형벌하였으니 그 죄 만사무석(萬死無釋)이라. 즉시 봉고파직(封庫罷職)하라!"

이때 옥중의 춘향이 구사일생으로 놓여나니, 어사또 춘향의 손을 붙잡고 왈, "이제야 맺힌 한을 풀었으니, 우리 백년가약(百年佳約)을 이루리라." 춘향이 기쁜 눈물을 흘리며 어사또를 우러러보더라.
}\kfResetAreaFlag
\par\kfMaybeNeedspace{50\baselineskip}\noindent{\kfTipMark\,\,}{\small\color{kfInk} 빈칸에 알맞은 말을 채워 넣으시오. (초성 힌트 참고)}\par\nopagebreak[4]\smallskip\nopagebreak[4]
\begin{kfActTable}{|>{\centering\arraybackslash}m{2.6cm}|m{6cm}|m{4cm}|}
\rowcolor{kfPrimaryLight!50}\textbf{\color{kfPrimary}항목} & \textbf{\color{kfPrimary}내용 (빈칸 채우기)} & \textbf{\color{kfPrimary}답안 작성}\\\hline
갈래 & ㄱㅈ ㅅㅅ, ㅍㅅㄹㄱ 소설 ① & \kfTBlank \\\hline
작자 & ㅈㅈ ㅁㅅ ② & \kfTBlank \\\hline
배경 & 전라도 ㄴㅇ ③ & \kfTBlank \\\hline
주요 인물 & ㅊㅎ, ㅇㄷㄹ(어사또), ㅂㅎㄷ ④ & \kfTBlank \\\hline
핵심 갈등 & 변학도의 ㅅㅊ 강요에 맞서는 춘향의 저항 ⑤ & \kfTBlank \\\hline
춘향의 거부 논리 & ㅇㅂㅈㅅ는 여자의 대절 ⑥ & \kfTBlank \\\hline
변학도의 처분 & ㅂㄱㅍㅈ (관직 박탈) ⑦ & \kfTBlank \\\hline
반전 장치 & ㅇㅎㅇㅅ 출도 ⑧ & \kfTBlank \\\hline
결말 구조 & ㄱㅅ ㅈㅇ — 선한 자 보상, 악한 자 징벌 ⑨ & \kfTBlank \\\hline
주제 1 & ㅌㄱㅇㄹ에 대한 저항과 ㅈㅇ 실현 ⑩ & \kfTBlank \\\hline
주제 2 & ㅅㄹ의 ㅈㅈ와 ㅅㅂ ㄱㄷ 극복 ⑪ & \kfTBlank \\\hline
시대적 의미 & 조선 후기 ㅅㅁ ㅇㅅ의 성장 반영 ⑫ & \kfTBlank \\\hline
\end{kfActTable}
\kfSecHeader{kfAreaLiterature}{문제}{}

\kfBeginMC
\begin{kfQBlock}{01}{객관식}
\kfQStem{윗글에 대한 설명으로 적절하지 \underline{않은} 것은?}
\begin{kfChoices}
\kfChoice 대화를 통해 인물 사이의 갈등을 직접적으로 드러내고 있다.
\kfChoice 한문 투의 표현과 고전적인 어투가 곳곳에 사용되고 있다.
\kfChoice 서술자의 심리 분석을 통해 인물의 내면 갈등을 세밀하게 묘사한다.
\kfChoice 선과 악의 대립 구도가 인물 배치를 통해 뚜렷하게 나타난다.
\kfChoice 권선징악의 전형적인 결말 구조를 분명하게 보여 주고 있다.
\end{kfChoices}
\end{kfQBlock}

\begin{kfQBlock}{02}{객관식}
\kfQStem{윗글에서 춘향이 수청을 거부하는 근거로 내세운 것은?}
\begin{kfChoices}
\kfChoice 양반 신분이므로 기생 노릇을 할 수 없다는 신분적 자부심
\kfChoice 한 번 허락한 사랑을 바꿀 수 없다는 여성의 절의
\kfChoice 변학도의 통치가 부당하므로 명령에 따를 의무가 없다는 법리적 주장
\kfChoice 이도령이 돌아와 자신을 구해 줄 것이라는 굳건한 믿음
\kfChoice 수청이 법률상 강제할 수 없는 행위라는 제도적 근거
\end{kfChoices}
\end{kfQBlock}

\begin{kfQBlock}{03}{객관식}
\kfQStem{윗글에서 어사또가 변학도에게 내린 죄목으로 언급된 것을 모두 고른 것은?}
\begin{kfChoices}
\kfChoice 백성의 재물 착취
\kfChoice 백성의 재물 착취, 죄 없는 양민 가두어 형벌
\kfChoice 백성의 재물 착취, 죄 없는 양민 가두어 형벌, 관청 재물 횡령
\kfChoice 죄 없는 양민 가두어 형벌, 관청 재물 횡령
\kfChoice 백성의 재물 착취, 관청 재물 횡령, 역모
\end{kfChoices}
\end{kfQBlock}

\begin{kfQBlock}{04}{객관식}
\kfQStem{윗글에서 '좌우의 수령 방백이 벌벌 떨어 제 살기를 도모하더라'라는 서술의 효과로 가장 적절한 것은?}
\begin{kfChoices}
\kfChoice 변학도 한 사람의 비리가 아니라 관리 전체의 부패가 만연했음을 암시한다.
\kfChoice 수령 방백들이 변학도에 대한 우정을 보여 주는 의리의 표현이다.
\kfChoice 암행어사의 권력이 지나치게 강하여 부당하다는 비판적 시각을 드러낸다.
\kfChoice 전란의 위험이 다가왔음을 감지한 관리들의 불안을 묘사한 것이다.
\kfChoice 춘향의 고난에 동정하는 관리들의 양심적 반응을 보여 준다.
\end{kfChoices}
\end{kfQBlock}

\begin{kfQBlock}{05}{객관식}
\kfQStem{윗글을 바탕으로 <보기>를 감상한 내용으로 적절하지 \underline{않은} 것은?}
\begin{kfEvidence}「춘향전」은 조선 후기 사회에서 성장한 서민 의식을 반영한 작품이다. 탐관오리의 횡포에 맞서는 춘향의 저항은 단순한 개인의 사랑 수호를 넘어, 부당한 권력에 대한 민중의 저항 의식을 대변한다. 또한 암행어사 출도라는 장치를 통해 백성들이 꿈꾸던 정의로운 질서의 회복을 형상화하고 있다.\end{kfEvidence}
\begin{kfChoices}
\kfChoice 춘향이 '차라리 죽여 주옵소서'라고 한 것은 부당한 권력에 굴복하기보다 죽음을 택하겠다는 저항 의식의 표현이겠군.
\kfChoice 변학도가 춘향을 '천기의 딸'이라 부르며 수청을 강요하는 것은 신분 질서를 이용한 부당한 권력 행사이겠군.
\kfChoice 어사또의 '봉고파직' 명령은 <보기>에서 말한 정의로운 질서의 회복이 실현되는 장면이겠군.
\kfChoice 춘향의 저항이 성공하는 결말은 조선 시대에 신분제가 실제로 철폐되었음을 반영한 것이겠군.
\kfChoice 변학도의 처벌 장면에서 민중이 느꼈을 통쾌함은 서민 의식의 성장과 관련이 있겠군.
\end{kfChoices}
\end{kfQBlock}

\kfEndMC

\kfPassageNote{「성난 기계」 — 차범석}{%
공장장  (작업대를 두드리며) 이봐, 여기가 어딘 줄 알아? 공장이야, 공장! 공장에서는 사람이 기계에 맞춰야 해. 기계가 사람에 맞추는 게 아니라. 기계가 돌아가면 사람도 돌아가고, 기계가 멈추면 사람도 멈추는 거야. 알겠어?

영호  (주먹을 쥐며) 우리가 기계입니까? 기계는 기름 치면 되지만 사람은 그렇지 않습니다. 사람한테는 밥도 줘야 하고, 잠도 재워야 하고, 아프면 쉬게도 해 줘야 합니다.

공장장  쉬고 싶으면 나가. 밖에 일하고 싶어 줄 선 사람이 얼마나 많은 줄 알아? 너 하나 나간다고 이 공장이 멈추는 줄 알아?

영호  (다른 노동자들을 돌아보며) 들었어? 우리가 뭔지 알겠지? 부속품이야, 부속품! 고장 나면 갈아 끼우면 그만인 부속품!

(기계 소음이 점점 커진다. 노동자들의 대화가 소음에 묻힌다.)

순이  (기계 앞에서 손가락을 다치며) 아! 영호  (달려가며) 순이! 괜찮아? 공장장  (시계를 보며) 작업 중단 없이 처리해. 라인 세우면 손해가 얼만데. 영호  (분노하며) 사람이 다쳤는데 라인이 먼저입니까! 공장장  여기선 라인이 먼저야. 싫으면 나가!
}\kfResetAreaFlag
\par\kfMaybeNeedspace{50\baselineskip}\noindent{\kfTipMark\,\,}{\small\color{kfInk} 빈칸에 알맞은 말을 채워 넣으시오. (초성 힌트 참고)}\par\nopagebreak[4]\smallskip\nopagebreak[4]
\begin{kfActTable}{|>{\centering\arraybackslash}m{2.6cm}|m{6cm}|m{4cm}|}
\rowcolor{kfPrimaryLight!50}\textbf{\color{kfPrimary}항목} & \textbf{\color{kfPrimary}내용 (빈칸 채우기)} & \textbf{\color{kfPrimary}답안 작성}\\\hline
갈래 & ㅎㄱ (연극 대본) ① & \kfTBlank \\\hline
작가 & ㅊㅂㅅ ② & \kfTBlank \\\hline
배경 & ㅅㅇㅎ 시대의 공장 ③ & \kfTBlank \\\hline
제목의 이중 의미 & 실제 ㄱㄱ + 분노하는 ㄴㄷㅈ ④ & \kfTBlank \\\hline
핵심 갈등 & ㅎㅇ(생산성) vs ㅈㅇ(인간 존엄) ⑤ & \kfTBlank \\\hline
공장장의 가치관 & 사람이 ㄱㄱ에 맞춰야 한다 ⑥ & \kfTBlank \\\hline
영호의 비유 & 노동자는 교체 가능한 ㅂㅅㅍ ⑦ & \kfTBlank \\\hline
무대 장치 상징 & ㄱㄱ ㅅㅇ — 노동자 목소리를 삼키는 ㅇㅇ ⑧ & \kfTBlank \\\hline
핵심 사건 & ㅅㅇ가 기계에 다쳐도 ㄹㅇ 중단 불가 ⑨ & \kfTBlank \\\hline
주제 1 & ㅅㅇㅎ 시대 ㅂㅇㄱㅈ 제도 비판 ⑩ & \kfTBlank \\\hline
주제 2 & 인간 ㅈㅇㅅ의 회복 요구 ⑪ & \kfTBlank \\\hline
표현 기법 & ㅅㅅㅈㅇ 기법으로 현실 고발 ⑫ & \kfTBlank \\\hline
\end{kfActTable}
\kfSecHeader{kfAreaLiterature}{문제}{}

\kfBeginMC
\begin{kfQBlock}{01}{객관식}
\kfQStem{윗글에 대한 설명으로 적절하지 \underline{않은} 것은?}
\begin{kfChoices}
\kfChoice 인물 간의 대화를 통해 갈등을 직접적으로 드러내고 있다.
\kfChoice 무대 지시문을 통해 인물의 행동과 분위기를 구체화하고 있다.
\kfChoice 기계 소음이라는 음향 효과가 극적 긴장감을 높이고 있다.
\kfChoice 공장장과 영호의 관계를 통해 경영자와 노동자의 대립을 보여 주고 있다.
\kfChoice 서술자의 논평을 통해 작가의 주제 의식을 직접 전달하고 있다.
\end{kfChoices}
\end{kfQBlock}

\begin{kfQBlock}{02}{객관식}
\kfQStem{윗글에서 '성난 기계'라는 제목이 지닌 이중적 의미로 가장 적절한 것은?}
\begin{kfChoices}
\kfChoice 고장 난 공장 기계와 공장장의 분노
\kfChoice 공장의 실제 기계와 비인간적 체제에 분노하는 노동자
\kfChoice 기계의 소음과 순이의 사고에 대한 공장장의 경악
\kfChoice 파업을 일으키는 노동자와 파업을 진압하는 경찰
\kfChoice 기계 문명에 대한 동경과 그에 따른 환멸
\end{kfChoices}
\end{kfQBlock}

\begin{kfQBlock}{03}{객관식}
\kfQStem{윗글에서 공장장이 '밖에 일하고 싶어 줄 선 사람이 얼마나 많은 줄 알아?'라고 말하는 의도로 가장 적절한 것은?}
\begin{kfChoices}
\kfChoice 공장의 인기가 높아 자부심을 느끼고 있음을 과시하려는 것
\kfChoice 노동자의 대체 가능성을 내세워 불만 제기를 억압하려는 것
\kfChoice 더 많은 인원을 고용하여 노동 환경을 개선하겠다는 약속
\kfChoice 노동 시장의 구조적 문제를 영호와 함께 고민하겠다는 제안
\kfChoice 영호의 능력을 인정하며 다른 사람들과 비교하여 칭찬하려는 것
\end{kfChoices}
\end{kfQBlock}

\begin{kfQBlock}{04}{객관식}
\kfQStem{윗글에서 영호가 '부속품이야, 부속품! 고장 나면 갈아 끼우면 그만인 부속품!'이라고 한 말의 효과로 가장 적절한 것은?}
\begin{kfChoices}
\kfChoice 기계에 비유함으로써 노동자가 처한 비인간적 현실을 동료에게 각성시킨다.
\kfChoice 자신이 이 공장에서 결코 빠질 수 없는 필수 불가결한 존재임을 과시한다.
\kfChoice 공장 기계의 노후화로 발생하는 안전 문제를 지적하여 점검을 요구한다.
\kfChoice 자조적 유머를 통해 긴장한 동료들의 분위기를 부드럽게 풀어 주려는 의도이다.
\kfChoice 기계를 능숙하게 다루는 자신의 전문적인 작업 능력을 비유적으로 표현한 것이다.
\end{kfChoices}
\end{kfQBlock}

\begin{kfQBlock}{05}{객관식}
\kfQStem{윗글을 바탕으로 <보기>를 감상한 내용으로 적절하지 \underline{않은} 것은?}
\begin{kfEvidence}차범석의 「성난 기계」는 1960\textasciitilde{}70년대 한국 산업화 시기를 배경으로, 급속한 경제 성장의 이면에 놓인 노동 소외의 문제를 다루고 있다. 작가는 공장이라는 공간을 통해 인간이 생산 도구로 전락하는 현실을 고발하며, 효율과 이윤 추구가 인간의 존엄성을 어떻게 파괴하는지를 보여 주고 있다.\end{kfEvidence}
\begin{kfChoices}
\kfChoice 공장장이 '기계에 맞춰야 해'라고 한 것은 <보기>에서 말한 '인간이 생산 도구로 전락하는 현실'을 보여 주는 발언이겠군.
\kfChoice 순이가 다쳤는데 '라인 세우면 손해'라고 한 것은 효율과 이윤이 인간의 존엄성보다 우선시되는 현실을 드러내겠군.
\kfChoice 영호가 '사람한테는 밥도 줘야 하고, 잠도 재워야 하고'라고 한 것은 인간으로서의 기본적 권리 요구이겠군.
\kfChoice 기계 소음이 점점 커져 대화가 묻히는 것은 산업화로 인해 노동자의 목소리가 외면당하는 현실의 상징이겠군.
\kfChoice 공장장이 영호에게 퇴직금과 보상을 제시하는 장면은 경영자의 양심적 각성을 보여 주겠군.
\end{kfChoices}
\end{kfQBlock}

\kfEndMC

\kfPassageNote{「파수꾼」 — 이강백}{%
파수꾼 다  적이라……. 난 한 번도 적을 본 적이 없어요. 파수꾼 나  나도 그래. 하지만 파수꾼 가가 알려 주면 나는 너에게 알려 주고, 넌 마을에 알려 주는 거야. 그게 우리가 해야 할 일이야. 파수꾼 다  그런데 파수꾼 가는 정말 적을 본 걸까요? 파수꾼 나  글쎄……. 나도 가끔 그런 생각을 해. 파수꾼 다  만약 적이 없다면요? 만약 이 모든 게 거짓이라면요? 파수꾼 나  (잠시 침묵) 그런 말 하면 안 돼. 촌장님이 들으면……. 파수꾼 다  촌장님이 뭘 어떻게 하는데요? 파수꾼 나  촌장님은 적이 있다고 하셨어. 적이 있으니까 우리가 파수를 서는 거라고.

(촌장 등장)

촌장  무슨 이야기들을 하고 있느냐? 파수꾼 나  아, 아무것도 아닙니다. 촌장  (날카롭게) 적은 반드시 있다. 적이 없다고 말하는 자가 바로 적이다. 알겠느냐? 파수꾼 다  (작은 목소리로) 하지만 저는 한 번도……. 촌장  (끊으며) 의심하는 자가 적의 편이다! 네가 파수꾼으로서 해야 할 일은 의심하는 게 아니라 전하는 것이다!

(파수꾼 다, 고개를 숙인다. 파수꾼 나, 불안한 눈으로 파수꾼 다를 바라본다.)
}\kfResetAreaFlag
\par\kfMaybeNeedspace{50\baselineskip}\noindent{\kfTipMark\,\,}{\small\color{kfInk} 빈칸에 알맞은 말을 채워 넣으시오. (초성 힌트 참고)}\par\nopagebreak[4]\smallskip\nopagebreak[4]
\begin{kfActTable}{|>{\centering\arraybackslash}m{2.6cm}|m{6cm}|m{4cm}|}
\rowcolor{kfPrimaryLight!50}\textbf{\color{kfPrimary}항목} & \textbf{\color{kfPrimary}내용 (빈칸 채우기)} & \textbf{\color{kfPrimary}답안 작성}\\\hline
갈래 & ㅎㄱ, ㅇㅎㄱ ① & \kfTBlank \\\hline
작가 & ㅇㄱㅂ ② & \kfTBlank \\\hline
배경 & 추상적 ㅁㅇ (특정 시대·장소 없음) ③ & \kfTBlank \\\hline
주요 인물 & 파수꾼 ㄱ/ㄴ/ㄷ, ㅊㅈ ④ & \kfTBlank \\\hline
핵심 갈등 & ㅈ의 실재 여부를 둘러싼 의문과 억압 ⑤ & \kfTBlank \\\hline
촌장의 논리 & 적이 없다고 말하는 자가 바로 ㅈ이다 ⑥ & \kfTBlank \\\hline
파수꾼 나의 태도 & 의문이 있으나 ㅈㄱ ㄱㅇ로 침묵 ⑦ & \kfTBlank \\\hline
핵심 기법 & ㅇㅎ (우화)를 통한 보편적 주제 전달 ⑧ & \kfTBlank \\\hline
적의 모호성 & 적의 실재 여부가 끝까지 ㅎㅇ되지 않음 ⑨ & \kfTBlank \\\hline
주제 1 & ㄱㄹ이 ㅎㄱ적 위협으로 구성원을 통제 ⑩ & \kfTBlank \\\hline
주제 2 & ㅈㅅ 추구와 그에 따르는 대가 ⑪ & \kfTBlank \\\hline
창작 시기 & 1970년대 한국 사회, 그러나 ㅂㅍㅈ 의미 ⑫ & \kfTBlank \\\hline
\end{kfActTable}
\kfSecHeader{kfAreaLiterature}{문제}{}

\kfBeginMC
\begin{kfQBlock}{01}{객관식}
\kfQStem{윗글에 대한 설명으로 적절하지 \underline{않은} 것은?}
\begin{kfChoices}
\kfChoice 인물의 대화를 중심으로 갈등이 전개되고 있다.
\kfChoice 무대 지시문을 통해 인물의 심리와 분위기를 보충하고 있다.
\kfChoice 구체적인 역사적 시간과 장소를 배경으로 설정하여 사실성을 높이고 있다.
\kfChoice 우화적 상황을 통해 권력과 진실의 문제를 다루고 있다.
\kfChoice '적'의 실재 여부가 모호하게 처리되어 극적 긴장을 유발하고 있다.
\end{kfChoices}
\end{kfQBlock}

\begin{kfQBlock}{02}{객관식}
\kfQStem{윗글에서 촌장이 '적이 없다고 말하는 자가 바로 적이다'라고 한 발언의 의도로 가장 적절한 것은?}
\begin{kfChoices}
\kfChoice 적의 정체를 실제로 알고 있기 때문에 경고하려는 것
\kfChoice 파수꾼들의 안전을 걱정하여 경각심을 높이려는 것
\kfChoice 적의 존재에 대한 의문 자체를 차단하여 기존 질서를 유지하려는 것
\kfChoice 마을 사람들 사이의 단결심을 자연스럽게 고취하려는 것
\kfChoice 파수꾼 다의 용기를 시험하여 더 강한 파수꾼으로 키우려는 것
\end{kfChoices}
\end{kfQBlock}

\begin{kfQBlock}{03}{객관식}
\kfQStem{윗글에서 파수꾼 나가 '그런 말 하면 안 돼. 촌장님이 들으면…….'이라고 말한 데에서 드러나는 파수꾼 나의 심리로 가장 적절한 것은?}
\begin{kfChoices}
\kfChoice 적의 존재를 확신하며 파수꾼 다의 무지를 답답해하는 심리
\kfChoice 촌장의 말이 옳다고 믿기에 충성심으로 경고하는 심리
\kfChoice 자신도 의문을 품고 있으나 권력의 보복이 두려워 자기 검열하는 심리
\kfChoice 파수꾼 다를 촌장에게 고발하려는 기회를 엿보는 심리
\kfChoice 촌장의 리더십을 깊이 존경하여 그의 판단을 전적으로 신뢰하는 심리
\end{kfChoices}
\end{kfQBlock}

\begin{kfQBlock}{04}{객관식}
\kfQStem{윗글에서 '적'의 존재가 끝까지 확인되지 않는 것이 작품에서 갖는 효과로 가장 적절한 것은?}
\begin{kfChoices}
\kfChoice 독자에게 적이 분명히 존재한다는 단정적 결론을 자연스럽게 이끌어 낸다.
\kfChoice 적의 실재 여부보다 '적이 있다'는 담론이 권력 유지에 동원되는 구조를 부각한다.
\kfChoice 추리소설적 긴장감을 강하게 유발하여 범인을 쫓는 독자의 흥미를 한층 높여 준다.
\kfChoice 작가가 결말을 끝내 쓰지 못해 미완성 상태로 남겨 둔 우연적 결과를 보여 준다.
\kfChoice 적의 정체가 알 수 없는 외계인임을 은밀히 암시하여 SF적 상상력을 자극한다.
\end{kfChoices}
\end{kfQBlock}

\begin{kfQBlock}{05}{객관식}
\kfQStem{윗글을 바탕으로 <보기>를 감상한 내용으로 적절하지 \underline{않은} 것은?}
\begin{kfEvidence}이강백의 「파수꾼」은 1970년대 한국 사회를 배경으로 창작되었으나, 우화적 형식을 취함으로써 특정 시대에 국한되지 않는 보편적 의미를 획득하고 있다. 이 작품은 권력이 '적'이라는 허구적 위협을 만들어 내고, 이를 통해 구성원의 자유로운 사고를 억압하는 메커니즘을 보여 준다. 진실을 추구하는 개인은 체제에 의해 위험 인물로 낙인찍히게 된다.\end{kfEvidence}
\begin{kfChoices}
\kfChoice 파수꾼 다가 '적을 본 적이 없다'고 말하는 것은 <보기>에서 말한 진실을 추구하는 개인의 모습이겠군.
\kfChoice 촌장이 '의심하는 자가 적의 편'이라고 한 것은 자유로운 사고를 억압하는 메커니즘의 구체적 사례이겠군.
\kfChoice 파수꾼 나가 촌장 앞에서 대화 내용을 부인하는 것은 억압적 체제 아래 자기 검열이 작동하는 모습이겠군.
\kfChoice 마을, 파수꾼, 촌장이라는 추상적 설정은 <보기>에서 말한 보편적 의미 획득과 관련이 있겠군.
\kfChoice 이 작품은 실제 전쟁 상황을 기록한 르포르타주이므로, 적의 존재는 역사적 사실에 기반한 것이겠군.
\end{kfChoices}
\end{kfQBlock}

\kfEndMC

\kfPassageNote{「레디메이드 인생」 — 채만식}{%
P는 졸업장이란 것을 몇 장이나 가지고 있다. 보통학교 졸업장, 고등보통학교 졸업장, 전문학교 졸업장…… 졸업장마다 우수한 성적이라고 씌어 있다. 그러나 그 화려한 졸업장들은 지금 P에게 아무 쓸모가 없다. 취직이 안 되는 것이다.

삼천리 방방곡곡을 돌아다녀 봐도 조선 사람에게 줄 자리가 없다. 은행에 들어가자니 일본 사람이 자리를 차지하고 있고, 회사에 들어가자니 일본 사람이 또 자리를 차지하고 있고, 관청에 들어가자니 역시 일본 사람이다. 남는 것은 초등학교 교원 자리인데, 그마저 줄이 길다.

P는 혼자 중얼거렸다. "대학을 나와도 소용이 없어. 전문학교를 나와도 소용이 없어. 졸업장이 밥을 먹여 주는 것도 아니고, 학문이 돈이 되는 것도 아니야. 나는 레디메이드야, 레디메이드! 기성품이란 말이야. 공장에서 찍어 낸 물건처럼 그저 쓸 데 없이 나뒹구는 거야."

어린 아들 창수가 물었다. "아버지, 나도 크면 학교 다녀요?" P는 쓸쓸히 웃으며 대답했다. "학교? 학교를 다녀서 뭘 하게? 나처럼 졸업장 서너 장 가지고 다니면서 실업자 노릇 하게? 차라리 공장에 가서 기술이나 배우렴. 그래야 밥이라도 먹지."

P는 고개를 떨구었다. 자신이 배운 학문이, 자신이 읽은 책이, 이 세상에서 아무 소용도 없다는 사실이 그를 짓누르고 있었다.
}\kfResetAreaFlag
\par\kfMaybeNeedspace{50\baselineskip}\noindent{\kfTipMark\,\,}{\small\color{kfInk} 빈칸에 알맞은 말을 채워 넣으시오. (초성 힌트 참고)}\par\nopagebreak[4]\smallskip\nopagebreak[4]
\begin{kfActTable}{|>{\centering\arraybackslash}m{2.6cm}|m{6cm}|m{4cm}|}
\rowcolor{kfPrimaryLight!50}\textbf{\color{kfPrimary}항목} & \textbf{\color{kfPrimary}내용 (빈칸 채우기)} & \textbf{\color{kfPrimary}답안 작성}\\\hline
갈래 & ㅎㄷ ㅅㅅ, ㄷㅍ ㅅㅅ ① & \kfTBlank \\\hline
작가 & ㅊㅁㅅ ② & \kfTBlank \\\hline
시대 배경 & 1930년대 ㅅㅁㅈ 조선 ③ & \kfTBlank \\\hline
주인공 & 교육받은 ㅅㅇ ㅈㅅㅇ P ④ & \kfTBlank \\\hline
제목의 의미 & ㄱㅅㅍ — 쓸모없이 방치된 존재 ⑤ & \kfTBlank \\\hline
실업 원인 & ㅇㅂ ㅅㄹ이 주요 직위를 독점 ⑥ & \kfTBlank \\\hline
아들에게 한 말 & 학교 대신 ㄱㅈ에서 ㄱㅅ을 배우라 ⑦ & \kfTBlank \\\hline
표현 기법 1 & ㅍㅈ와 ㅂㅇ를 통한 현실 비판 ⑧ & \kfTBlank \\\hline
표현 기법 2 & 반복적 ㅇㄱ로 절망적 상황 강조 ⑨ & \kfTBlank \\\hline
핵심 갈등 & 개인의 ㄴㄹ vs 사회의 ㄱㅈ ⑩ & \kfTBlank \\\hline
주제 1 & ㅅㅁㅈ ㅅㅎ ㄱㅈ가 만든 지식인의 무력감 ⑪ & \kfTBlank \\\hline
주제 2 & ㄱㅇ의 무용화와 ㅅㅈㅈ 좌절 ⑫ & \kfTBlank \\\hline
\end{kfActTable}
\kfSecHeader{kfAreaLiterature}{문제}{}

\kfBeginMC
\begin{kfQBlock}{01}{객관식}
\kfQStem{윗글에 대한 설명으로 적절하지 \underline{않은} 것은?}
\begin{kfChoices}
\kfChoice 주인공의 독백과 대화를 통해 내면을 드러내고 있다.
\kfChoice 반복적 열거를 통해 주인공이 처한 상황의 막막함을 강조하고 있다.
\kfChoice 자조적 어조를 통해 식민지 현실에 대한 비판 의식을 드러내고 있다.
\kfChoice 아들과의 대화를 통해 세대 간 단절의 원인을 밝히고 있다.
\kfChoice 비유적 표현('레디메이드', '기성품')을 통해 주인공의 처지를 효과적으로 형상화하고 있다.
\end{kfChoices}
\end{kfQBlock}

\begin{kfQBlock}{02}{객관식}
\kfQStem{윗글에서 '레디메이드(기성품)'라는 비유가 P의 상황에 대해 의미하는 바로 가장 적절한 것은?}
\begin{kfChoices}
\kfChoice 공장에서 대량 생산된 규격품처럼 개성 없이 살아가는 현대인의 모습
\kfChoice 교육을 받았으나 사회가 필요로 하지 않아 쓸모없이 방치된 존재
\kfChoice 스스로 노력하지 않고 기성세대의 유산에 기대어 사는 무위도식의 삶
\kfChoice 서양 문물을 무비판적으로 수용한 지식인의 자기 비하
\kfChoice 기성 사회의 관습을 따르지 않는 반항적 지식인의 자부심
\end{kfChoices}
\end{kfQBlock}

\begin{kfQBlock}{03}{객관식}
\kfQStem{윗글에서 P가 아들에게 '차라리 공장에 가서 기술이나 배우렴'이라고 말한 이유로 가장 적절한 것은?}
\begin{kfChoices}
\kfChoice 공장 기술자가 지식인보다 사회적 지위가 높았기 때문이다.
\kfChoice 자신의 경험에 비추어 교육이 생존에 도움이 되지 않았기 때문이다.
\kfChoice 아들이 학업에 재능이 없다고 판단했기 때문이다.
\kfChoice 일본인 기술자가 부족하여 조선인 기술자의 수요가 높았기 때문이다.
\kfChoice 교육보다 기술 습득이 민족 독립에 더 기여한다고 믿었기 때문이다.
\end{kfChoices}
\end{kfQBlock}

\begin{kfQBlock}{04}{객관식}
\kfQStem{윗글의 '은행에 들어가자니 일본 사람이 자리를 차지하고 있고, 회사에 들어가자니 일본 사람이 또 자리를 차지하고 있고'라는 서술에 사용된 표현 기법과 그 효과로 가장 적절한 것은?}
\begin{kfChoices}
\kfChoice 점층법 — P의 분노가 단계적으로 고조되는 효과
\kfChoice 열거법과 반복 — 어디에도 자리가 없는 절망적 상황을 강조하는 효과
\kfChoice 대구법 — P와 일본인 사이의 우정이 깊어지는 과정을 보여 주는 효과
\kfChoice 역설법 — 일본인이 자리를 차지한 것이 오히려 좋은 결과라는 역설적 효과
\kfChoice 과장법 — 실제로는 취직 기회가 많았으나 P가 과장하여 불평하는 효과
\end{kfChoices}
\end{kfQBlock}

\begin{kfQBlock}{05}{객관식}
\kfQStem{윗글을 바탕으로 <보기>를 감상한 내용으로 적절하지 \underline{않은} 것은?}
\begin{kfEvidence}채만식은 1930년대 식민지 조선의 사회적 모순을 풍자적 필치로 그려 낸 작가이다. 「레디메이드 인생」에서 주인공 P의 실업은 개인의 무능이 아니라 식민지라는 구조적 상황이 만들어 낸 것이다. 채만식은 이 작품을 통해 교육받은 지식인마저 사회에서 쓸모없는 존재로 전락시키는 식민 체제의 모순을 고발하고 있다.\end{kfEvidence}
\begin{kfChoices}
\kfChoice P의 '졸업장이란 것을 몇 장이나 가지고 있다'라는 서술은 교육이 무용해진 현실을 역설적으로 보여 주겠군.
\kfChoice '일본 사람이 자리를 차지하고 있고'라는 반복은 <보기>에서 말한 식민지 구조적 모순을 구체화하겠군.
\kfChoice P가 스스로를 '레디메이드'라고 부르는 것은 자조를 통한 현실 비판이겠군.
\kfChoice P가 아들에게 기술을 배우라고 한 것은 식민지 상황과 무관한 순수한 교육관의 표현이겠군.
\kfChoice P의 실업이 개인의 무능이 아닌 구조적 문제임을 보여 주는 것은 <보기>의 관점과 일치하겠군.
\end{kfChoices}
\end{kfQBlock}

\kfEndMC

\kfStartNonlit{법이란 무엇인가 외 4편}


\kfPassageNote{법이란 무엇인가 (사회)}{%
법이란 사회 구성원의 행위를 규율하는 강제적 규범 체계이다. 인간 사회에는 도덕, 관습, 종교 등 다양한 규범이 존재하지만, 법은 국가 권력에 의해 그 이행이 강제된다는 점에서 다른 규범과 구별된다. 법을 위반하면 벌금, 징역 등 구체적인 제재가 뒤따르지만, 도덕을 어겼을 때는 사회적 비난은 있을 수 있어도 국가가 직접 처벌하지는 않는다.

법에 대한 이해는 크게 두 가지 흐름으로 나뉜다. 하나는 자연법론이고, 다른 하나는 법실증주의이다. 자연법론은 인간의 이성이나 자연의 질서에 근거한 보편적 도덕 원리가 존재하며, 실정법은 이 원리에 부합해야만 참된 법으로서의 효력을 가진다고 본다. 예컨대 '인간의 생명을 부당하게 빼앗는 행위는 악이다'라는 명제는 어떤 국가가 법으로 허용하더라도 자연법의 관점에서는 부정의한 것이다. 이러한 자연법 사상은 고대 그리스의 스토아 학파에서 시작되어, 중세의 토마스 아퀴나스, 근대의 로크와 루소에 이르기까지 오랜 전통을 가지고 있다.

반면 법실증주의는 법의 효력을 도덕과 분리하여 파악한다. 법실증주의에 따르면 법이란 정당한 절차에 의해 제정된 실정법 그 자체이며, 그 내용이 도덕적으로 옳은지와 법적 효력은 별개의 문제이다. 대표적 법실증주의자인 한스 켈젠은 법을 순수한 규범의 체계로 보고, 법학이 도덕이나 정치로부터 독립해야 한다고 주장했다. 법실증주의는 법적 안정성과 예측 가능성을 중시한다는 장점이 있지만, 부당한 법에 대한 저항의 근거를 제시하기 어렵다는 비판을 받기도 한다.

이 두 입장은 법의 목적에 대한 이해에서도 차이를 보인다. 자연법론은 법의 궁극적 목적을 정의의 실현에 둔다. 법이 사회의 구성원에게 각자의 몫을 공정하게 배분하고, 부당한 차별을 시정하는 역할을 해야 한다는 것이다. 반면 법실증주의는 법의 목적을 사회 질서의 유지와 구성원의 안전 보장에서 찾는다. 법이 명확한 기준을 제시하여 분쟁을 예방하고, 위반에 대한 제재를 통해 사회의 안정을 도모하는 것이 핵심이라고 본다.

오늘날 대부분의 법 체계는 자연법론과 법실증주의의 요소를 모두 포함하고 있다. 헌법에 기본권을 명시하여 자연법적 가치를 실정법 안에 수용하는 한편, 법률의 제정과 적용에서는 절차적 정당성을 엄격히 요구한다. 이처럼 현대 법 체계는 정의, 질서, 안전이라는 세 가지 목적을 조화롭게 추구하는 방향으로 발전해 왔다.
}\kfResetAreaFlag
\kfSecHeader{kfAreaNonlit}{문제}{}

\kfBeginMC
\begin{kfQBlock}{01}{객관식}
\kfQStem{윗글의 내용과 일치하는 것은?}
\begin{kfChoices}
\kfChoice 법실증주의는 법의 효력이 도덕적 정당성에 의해 결정된다고 본다.
\kfChoice 자연법론은 정당한 절차에 의해 제정된 법이면 모두 참된 법이라고 본다.
\kfChoice 한스 켈젠은 법학이 도덕이나 정치로부터 독립해야 한다고 주장했다.
\kfChoice 도덕 규범은 법 규범과 마찬가지로 국가 권력에 의해 강제된다.
\kfChoice 자연법 사상은 근대 이후에 처음 등장한 비교적 새로운 법 이론이다.
\end{kfChoices}
\end{kfQBlock}

\begin{kfQBlock}{02}{객관식}
\kfQStem{윗글에서 자연법론과 법실증주의를 비교한 내용으로 적절하지 \underline{않은} 것은?}
\begin{kfChoices}
\kfChoice 자연법론은 보편적 도덕 원리의 존재를 전제하지만, 법실증주의는 법의 효력을 도덕과 분리한다.
\kfChoice 자연법론은 법의 궁극적 목적을 정의의 실현에 두지만, 법실증주의는 사회 질서의 유지에서 찾는다.
\kfChoice 자연법론은 부당한 법에 대한 저항의 근거를 제시할 수 있지만, 법실증주의는 그러한 근거를 제시하기 어렵다.
\kfChoice 자연법론은 법적 안정성과 예측 가능성을 중시하지만, 법실증주의는 법의 내용적 정당성을 중시한다.
\kfChoice 자연법론은 실정법이 도덕 원리에 부합해야 한다고 보지만, 법실증주의는 실정법 자체의 효력을 인정한다.
\end{kfChoices}
\end{kfQBlock}

\begin{kfQBlock}{03}{객관식}
\kfQStem{윗글을 바탕으로 <보기>를 이해한 내용으로 적절하지 \underline{않은} 것은?}
\begin{kfEvidence}A국에서는 특정 소수 민족의 재산을 국가가 몰수할 수 있도록 하는 법률을 정당한 입법 절차를 거쳐 제정하였다. 이에 대해 갑은 '이 법률은 절차적으로 정당하므로 법적 효력을 가진다'고 주장하였고, 을은 '어떤 절차를 거쳤든 인간의 기본적 권리를 침해하는 법률은 참된 법이 아니다'라고 주장하였다.\end{kfEvidence}
\begin{kfChoices}
\kfChoice 갑의 입장은 법실증주의적 관점에서 법의 효력을 판단한 것이다.
\kfChoice 을의 입장은 자연법론의 관점에서 실정법의 정당성을 비판한 것이다.
\kfChoice 갑의 입장에서는 이 법률의 도덕적 부당성을 이유로 효력을 부정하기 어렵다.
\kfChoice 을의 입장에서는 보편적 도덕 원리에 어긋나는 법을 거부할 수 있는 근거가 마련된다.
\kfChoice 갑과 을 모두 법의 내용적 정당성보다 절차적 정당성을 우선시하고 있다.
\end{kfChoices}
\end{kfQBlock}

\begin{kfQBlock}{04}{객관식}
\kfQStem{'법은 국가 권력에 의해 그 이행이 강제된다'는 것의 의미로 가장 적절한 것은?}
\begin{kfChoices}
\kfChoice 법을 위반하면 국가가 구체적인 제재를 가할 수 있다는 뜻이다.
\kfChoice 국가가 모든 국민의 행위를 상시 감시한다는 뜻이다.
\kfChoice 법은 국가의 이익만을 위해 존재한다는 뜻이다.
\kfChoice 도덕과 법이 동일한 강제력을 가진다는 뜻이다.
\kfChoice 법을 제정할 수 있는 모든 권한이 오직 국민에게 있다는 뜻이다.
\end{kfChoices}
\end{kfQBlock}

\kfEndMC

\kfPassageNote{헌법의 기본 원리 (사회)}{%
헌법은 한 국가의 최고 법규범으로서, 국가의 조직과 운영의 기본 틀을 정하고 국민의 기본적 권리를 보장하는 역할을 한다. 대한민국 헌법은 여러 기본 원리 위에 세워져 있는데, 그 핵심은 국민 주권의 원리, 권력 분립의 원리, 기본권 보장, 그리고 법치주의이다.

국민 주권의 원리란 국가의 최고 의사 결정 권한인 주권이 국민에게 있다는 것을 뜻한다. 대한민국 헌법 제1조 제2항은 '대한민국의 주권은 국민에게 있고, 모든 권력은 국민으로부터 나온다'고 선언하고 있다. 이 원리에 따라 국민은 선거를 통해 대표를 선출하고, 국민 투표를 통해 중요한 국가적 사안에 직접 참여할 수 있다.

권력 분립의 원리는 국가 권력을 입법권, 행정권, 사법권으로 나누어 서로 다른 기관에 배분함으로써 권력의 집중과 남용을 방지하려는 것이다. 입법권은 법률을 제정하는 국회에, 행정권은 법률을 집행하는 대통령과 정부에, 사법권은 법률을 적용하여 분쟁을 해결하는 법원에 속한다. 이 세 기관은 서로를 견제하고 균형을 유지함으로써 어느 한쪽에 권력이 과도하게 집중되는 것을 막는다. 예컨대 국회는 탄핵 소추권을 통해 대통령을 견제하고, 대통령은 법률안 거부권을 행사하여 국회를 견제하며, 법원은 위헌 법률 심판을 통해 국회의 입법을 통제할 수 있다.

기본권 보장은 헌법이 국민의 기본적 권리를 명시하고 이를 국가가 침해하지 못하도록 보호하는 것을 말한다. 기본권에는 크게 자유권, 평등권, 사회권이 있다. 자유권은 국가의 간섭을 받지 않고 개인의 자유를 누릴 수 있는 권리로, 신체의 자유, 표현의 자유, 종교의 자유 등이 이에 해당한다. 평등권은 합리적 이유 없이 차별받지 않을 권리이며, 사회권은 인간다운 생활을 위해 국가에 적극적으로 급부를 요구할 수 있는 권리로 교육을 받을 권리, 근로의 권리 등이 포함된다.

법치주의는 국가 권력의 행사가 법에 의해, 법에 따라 이루어져야 한다는 원리이다. 이는 단순히 법률에 근거만 있으면 된다는 형식적 법치주의를 넘어, 법의 내용도 정의와 인권에 부합해야 한다는 실질적 법치주의를 포함한다. 실질적 법치주의는 헌법재판소를 통한 위헌 법률 심사, 적법 절차의 원칙, 신뢰 보호의 원칙 등을 통해 구현된다. 이처럼 헌법의 기본 원리들은 서로 유기적으로 연결되어 민주적이고 정의로운 국가 운영의 토대를 이루고 있다.
}\kfResetAreaFlag
\kfSecHeader{kfAreaNonlit}{문제}{}

\kfBeginMC
\begin{kfQBlock}{01}{객관식}
\kfQStem{윗글의 내용과 일치하지 않는 것은?}
\begin{kfChoices}
\kfChoice 국민 주권의 원리에 따라 국민은 선거와 국민 투표를 통해 주권을 행사할 수 있다.
\kfChoice 권력 분립의 원리는 국가 권력의 집중과 남용을 방지하기 위한 것이다.
\kfChoice 자유권은 국가에 적극적으로 급부를 요구할 수 있는 권리에 해당한다.
\kfChoice 실질적 법치주의는 법의 내용도 정의와 인권에 부합해야 한다는 것을 의미한다.
\kfChoice 사법권은 법률을 적용하여 분쟁을 해결하는 법원에 속한다.
\end{kfChoices}
\end{kfQBlock}

\begin{kfQBlock}{02}{객관식}
\kfQStem{윗글에서 '권력 분립의 원리'가 필요한 까닭으로 가장 적절한 것은?}
\begin{kfChoices}
\kfChoice 국가의 업무 처리 속도를 높이기 위해 권한을 효율적으로 분배하기 위함이다.
\kfChoice 입법, 행정, 사법 기관이 각각 독립적으로 활동하여 상호 간섭을 완전히 배제하기 위함이다.
\kfChoice 어느 한 기관에 권력이 과도하게 집중되는 것을 막고 상호 견제와 균형을 유지하기 위함이다.
\kfChoice 국민이 세 기관 중 하나를 선택하여 직접 권력을 행사할 수 있도록 하기 위함이다.
\kfChoice 사법부가 입법부와 행정부보다 우월한 지위를 갖도록 보장하기 위함이다.
\end{kfChoices}
\end{kfQBlock}

\begin{kfQBlock}{03}{객관식}
\kfQStem{윗글을 바탕으로 <보기>를 이해한 내용으로 적절한 것은?}
\begin{kfEvidence}국회가 특정 종교를 국교로 지정하고, 다른 종교의 활동을 전면 금지하는 법률을 제정하였다.\end{kfEvidence}
\begin{kfChoices}
\kfChoice 이 법률은 국회가 제정하였으므로 권력 분립의 원리에 위배되지 않는다.
\kfChoice 이 법률은 평등권을 보장하기 위해 합리적 차별을 시행하는 것이다.
\kfChoice 이 법률은 형식적 법치주의를 충족하므로 실질적 법치주의에도 부합한다.
\kfChoice 이 법률은 종교의 자유라는 자유권을 침해하여 기본권 보장의 원리에 어긋난다.
\kfChoice 이 법률은 사회권을 실현하기 위한 국가의 적극적 조치에 해당한다.
\end{kfChoices}
\end{kfQBlock}

\begin{kfQBlock}{04}{객관식}
\kfQStem{윗글에서 '형식적 법치주의'와 '실질적 법치주의'에 대한 설명으로 적절한 것은?}
\begin{kfChoices}
\kfChoice 형식적 법치주의는 법의 내용이 정의에 부합하는지를 핵심 기준으로 삼는다.
\kfChoice 실질적 법치주의는 법률에 근거만 있으면 국가 권력 행사가 정당하다고 본다.
\kfChoice 형식적 법치주의와 실질적 법치주의는 동일한 의미를 가진 용어이다.
\kfChoice 실질적 법치주의는 법의 절차적 근거뿐 아니라 내용적 정당성까지 요구한다.
\kfChoice 형식적 법치주의만으로도 위헌 법률 심사가 가능하다.
\end{kfChoices}
\end{kfQBlock}

\begin{kfQBlock}{05}{객관식}
\kfQStem{윗글에서 '대통령은 법률안 거부권을 행사하여 국회를 견제'한다는 서술은 권력 분립 원리의 어떤 측면을 보여 주는가?}
\begin{kfChoices}
\kfChoice 입법부가 행정부의 법 집행을 감독하는 측면
\kfChoice 사법부가 행정부의 위법 행위를 심판하는 측면
\kfChoice 행정부가 입법부의 입법 활동에 대해 견제하는 측면
\kfChoice 입법부와 사법부가 공동으로 행정부를 통제하는 측면
\kfChoice 세 기관이 동일한 권한을 공유하여 균형을 이루는 측면
\end{kfChoices}
\end{kfQBlock}

\kfEndMC

\kfPassageNote{형법의 죄형법정주의 (사회)}{%
형법은 어떤 행위가 범죄이며, 그에 대해 어떤 형벌을 부과할 것인지를 규정하는 법이다. 형법의 가장 중요한 원칙은 죄형법정주의인데, 이는 '법률 없으면 범죄 없고, 법률 없으면 형벌도 없다'라는 라틴어 격언에서 유래한 것이다. 즉, 어떤 행위를 처벌하려면 그 행위가 범죄라는 것과 그에 대한 형벌이 미리 법률로 정해져 있어야 한다는 뜻이다. 이 원칙은 국가의 형벌권 행사로부터 국민의 자유와 권리를 보호하기 위한 것이다.

죄형법정주의는 여러 가지 파생 원칙을 포함한다. 첫째, 소급입법 금지의 원칙이다. 이는 행위 시점에 범죄로 규정되지 않았던 행위를 나중에 제정된 법률로 소급하여 처벌할 수 없다는 것을 뜻한다. 예를 들어, 갑이 2024년에 한 행위가 당시 법률에 의하면 합법이었는데 2025년에 그 행위를 범죄로 규정하는 법률이 새로 만들어졌다고 하자. 이 경우 갑의 2024년 행위를 2025년의 새 법률로 처벌하는 것은 소급입법 금지의 원칙에 위배된다.

둘째, 유추해석 금지의 원칙이다. 형법 조문을 해석할 때, 법률에 명시되지 않은 사항을 유사한 규정을 근거로 확장하여 적용하는 것을 금지하는 원칙이다. 예컨대 법률이 '자동차 절도'만을 규정하고 있을 때, '자전거 절도'를 자동차 절도와 유사하다는 이유로 같은 조항을 적용하여 처벌할 수는 없다. 다만, 피고인에게 유리한 방향으로의 유추해석은 허용될 수 있다는 것이 일반적 견해이다.

셋째, 명확성의 원칙이다. 범죄와 형벌을 규정하는 법률은 그 내용이 명확해야 한다. 법률의 내용이 지나치게 모호하거나 추상적이면 국민이 어떤 행위가 허용되고 어떤 행위가 금지되는지를 예측할 수 없게 되기 때문이다. 만약 '사회에 해로운 행위를 한 자는 처벌한다'라는 식으로 법률을 규정한다면, '사회에 해로운 행위'가 구체적으로 무엇인지 명확하지 않으므로 명확성의 원칙에 위배될 수 있다.

이처럼 죄형법정주의의 파생 원칙들은 국가가 자의적으로 형벌권을 행사하는 것을 방지하고, 국민이 자신의 행위에 대한 법적 결과를 미리 예측할 수 있도록 보장한다. 소급입법 금지는 시간적 측면에서, 유추해석 금지는 해석적 측면에서, 명확성의 원칙은 입법적 측면에서 각각 국민의 자유를 보호하는 역할을 한다. 대한민국 헌법 제12조 제1항과 제13조 제1항에서도 이러한 원칙들을 명시적으로 보장하고 있다.
}\kfResetAreaFlag
\kfSecHeader{kfAreaNonlit}{문제}{}

\kfBeginMC
\begin{kfQBlock}{01}{객관식}
\kfQStem{윗글의 내용과 일치하는 것은?}
\begin{kfChoices}
\kfChoice 유추해석은 피고인에게 불리한 방향이든 유리한 방향이든 모두 금지된다.
\kfChoice 죄형법정주의는 형벌의 강도를 법관의 재량에 전적으로 맡기는 원칙이다.
\kfChoice 명확성의 원칙에 따르면 범죄와 형벌을 규정하는 법률의 내용은 명확해야 한다.
\kfChoice 소급입법 금지의 원칙은 미래에 제정될 법률의 효력만을 규율하는 것이다.
\kfChoice 형법은 민사 분쟁의 해결 절차를 규정하는 법률이다.
\end{kfChoices}
\end{kfQBlock}

\begin{kfQBlock}{02}{객관식}
\kfQStem{윗글에서 죄형법정주의가 필요한 근본적 까닭으로 가장 적절한 것은?}
\begin{kfChoices}
\kfChoice 범죄자에 대한 형벌을 최대한 엄격하게 적용하기 위해서이다.
\kfChoice 국가의 형벌권 행사로부터 국민의 자유와 권리를 보호하기 위해서이다.
\kfChoice 법관이 법률 해석에 소모하는 시간을 줄이기 위해서이다.
\kfChoice 모든 범죄에 동일한 형벌을 부과하여 평등을 실현하기 위해서이다.
\kfChoice 입법 기관의 법률 제정 절차를 간소화하기 위해서이다.
\end{kfChoices}
\end{kfQBlock}

\begin{kfQBlock}{03}{객관식}
\kfQStem{윗글을 바탕으로 <보기>의 각 사례를 판단한 내용으로 적절하지 \underline{않은} 것은?}
\begin{kfEvidence}(가) 2025년에 합법이던 행위를 2026년에 범죄로 규정하는 법률을 만들어 2025년의 행위를 처벌하였다. \\\relax
(나) 법률에 '칼로 사람을 해친 경우'를 처벌한다고 되어 있는데, 망치를 사용한 사건에 이 조항을 적용하여 처벌하였다. \\\relax
(다) '공공 질서를 문란하게 하는 일체의 행위를 처벌한다'라는 법률 조항을 제정하였다.\end{kfEvidence}
\begin{kfChoices}
\kfChoice (가)는 행위 당시 합법이던 행위를 나중 법률로 처벌하였으므로 소급입법 금지의 원칙에 위배된다.
\kfChoice (나)는 법률에 명시되지 않은 대상에 유사한 조항을 적용하였으므로 유추해석 금지의 원칙에 위배된다.
\kfChoice (다)는 '일체의 행위'라는 표현이 지나치게 모호하므로 명확성의 원칙에 위배될 수 있다.
\kfChoice (가)는 유추해석 금지의 원칙에, (나)는 소급입법 금지의 원칙에 각각 위배된다.
\kfChoice (가), (나), (다)는 모두 죄형법정주의의 파생 원칙과 관련된 사례이다.
\end{kfChoices}
\end{kfQBlock}

\begin{kfQBlock}{04}{객관식}
\kfQStem{윗글에서 소급입법 금지, 유추해석 금지, 명확성의 원칙이 각각 국민의 자유를 보호하는 측면으로 짝지은 것 중 적절한 것은?}
\begin{kfChoices}
\kfChoice 소급입법 금지 -- 해석적 측면, 유추해석 금지 -- 시간적 측면, 명확성의 원칙 -- 입법적 측면
\kfChoice 소급입법 금지 -- 시간적 측면, 유추해석 금지 -- 해석적 측면, 명확성의 원칙 -- 입법적 측면
\kfChoice 소급입법 금지 -- 입법적 측면, 유추해석 금지 -- 시간적 측면, 명확성의 원칙 -- 해석적 측면
\kfChoice 소급입법 금지 -- 시간적 측면, 유추해석 금지 -- 입법적 측면, 명확성의 원칙 -- 해석적 측면
\kfChoice 소급입법 금지 -- 해석적 측면, 유추해석 금지 -- 입법적 측면, 명확성의 원칙 -- 시간적 측면
\end{kfChoices}
\end{kfQBlock}

\kfEndMC

\kfPassageNote{미성년자의 법적 지위 (사회)}{%
민법은 만 19세 미만인 사람을 미성년자로 규정하고 있다. 미성년자는 아직 판단 능력이 충분히 성숙하지 않았다고 보아, 법률은 이들을 보호하기 위한 여러 장치를 마련하고 있다. 미성년자의 법적 지위를 이해하기 위해서는 행위능력, 법정대리인의 동의, 계약 취소권, 불법행위 책임능력이라는 네 가지 개념을 살펴볼 필요가 있다.

행위능력이란 법률 행위를 단독으로 유효하게 할 수 있는 능력을 말한다. 민법은 미성년자의 행위능력을 제한하여, 미성년자가 법률 행위를 할 때에는 원칙적으로 법정대리인의 동의를 받도록 하고 있다. 법정대리인이란 법률에 의해 미성년자를 대리할 권한이 주어진 사람으로, 일반적으로 부모가 이에 해당한다. 다만, 단순히 권리만을 얻거나 의무만을 면하는 행위, 그리고 법정대리인이 범위를 정하여 처분을 허락한 재산의 처분 등은 법정대리인의 동의 없이도 미성년자가 단독으로 할 수 있다.

미성년자가 법정대리인의 동의 없이 한 법률 행위는 취소할 수 있다. 이 취소권은 미성년자 본인뿐 아니라 법정대리인도 행사할 수 있으며, 취소된 법률 행위는 처음부터 무효인 것으로 간주된다. 이러한 취소권은 미성년자를 경솔한 거래로부터 보호하기 위한 것이다. 그러나 미성년자가 속임수를 써서 자신이 성년자인 것처럼 상대방을 믿게 한 경우에는 취소권이 제한된다. 예컨대 미성년자가 위조된 신분증을 제시하여 성년자로 행세하며 계약을 체결한 경우, 나중에 미성년자임을 이유로 그 계약을 취소할 수 없다.

한편 불법행위 책임능력은 행위능력과는 별개의 개념이다. 불법행위란 고의 또는 과실로 타인에게 손해를 끼치는 행위를 말하는데, 불법행위의 책임을 지기 위해서는 자기 행위의 결과를 인식하고 판단할 수 있는 능력, 즉 책임능력이 있어야 한다. 민법은 책임능력의 유무를 판단하는 구체적인 나이를 명시하지 않고 있으나, 판례는 대체로 만 12세에서 만 14세 사이를 기준으로 삼고 있다. 책임능력이 없는 미성년자가 불법행위를 한 경우에는 그 미성년자 대신 감독 의무자인 부모 등이 손해배상 책임을 진다.

이처럼 민법은 미성년자의 행위능력을 제한하면서도, 법정대리인의 동의와 취소권이라는 보호 장치를 통해 미성년자의 법률적 이익을 지키고 있다. 동시에 불법행위 책임능력 제도를 통해, 아직 충분한 판단력이 갖춰지지 않은 미성년자에게 과도한 법적 책임을 지우지 않도록 하면서도 피해자 구제를 위해 감독 의무자의 책임을 인정하는 균형을 취하고 있다.
}\kfResetAreaFlag
\kfSecHeader{kfAreaNonlit}{문제}{}

\kfBeginMC
\begin{kfQBlock}{01}{객관식}
\kfQStem{윗글의 내용과 일치하는 것은?}
\begin{kfChoices}
\kfChoice 미성년자는 어떠한 경우에도 법정대리인의 동의 없이 법률 행위를 할 수 없다.
\kfChoice 미성년자가 속임수를 써서 성년자로 행세한 경우에도 취소권은 그대로 유지된다.
\kfChoice 불법행위 책임능력이 없는 미성년자의 행위에 대해서는 감독 의무자가 책임을 진다.
\kfChoice 민법은 불법행위 책임능력의 기준 나이를 만 14세로 명시하고 있다.
\kfChoice 법정대리인은 법원에 의해 선임된 변호사만이 될 수 있다.
\end{kfChoices}
\end{kfQBlock}

\begin{kfQBlock}{02}{객관식}
\kfQStem{윗글에서 '행위능력'과 '불법행위 책임능력'에 대한 설명으로 적절한 것은?}
\begin{kfChoices}
\kfChoice 행위능력은 법률 행위의 유효성과 관련되고, 불법행위 책임능력은 손해배상 책임과 관련된다.
\kfChoice 행위능력과 불법행위 책임능력은 동일한 개념이므로 같은 나이 기준이 적용된다.
\kfChoice 행위능력이 제한되면 불법행위 책임능력도 자동으로 제한된다.
\kfChoice 불법행위 책임능력은 만 19세 미만에게는 일률적으로 인정되지 않는다.
\kfChoice 행위능력은 판례에 의해 기준이 결정되고, 불법행위 책임능력은 법률로 나이가 명시되어 있다.
\end{kfChoices}
\end{kfQBlock}

\begin{kfQBlock}{03}{객관식}
\kfQStem{윗글을 바탕으로 <보기>를 이해한 내용으로 적절한 것은?}
\begin{kfEvidence}만 16세인 A는 법정대리인의 동의 없이 고가의 노트북 컴퓨터를 구매하는 계약을 체결하였다. 이후 A의 부모가 이 사실을 알게 되었다.\end{kfEvidence}
\begin{kfChoices}
\kfChoice A의 부모는 법정대리인으로서 이 계약을 취소할 수 있다.
\kfChoice A는 만 16세이므로 이 계약은 동의 없이도 처음부터 유효하다.
\kfChoice 이 계약은 법정대리인의 동의가 없으므로 자동으로 무효가 된다.
\kfChoice A가 노트북을 이미 수령하였다면 부모는 취소권을 행사할 수 없다.
\kfChoice 노트북 구매는 단순히 권리만을 얻는 행위이므로 동의가 필요 없다.
\end{kfChoices}
\end{kfQBlock}

\begin{kfQBlock}{04}{객관식}
\kfQStem{윗글에서 미성년자가 속임수를 쓴 경우 취소권이 제한되는 까닭으로 가장 적절한 것은?}
\begin{kfChoices}
\kfChoice 미성년자가 범죄를 저질렀으므로 처벌의 일환으로 취소권을 박탈하는 것이다.
\kfChoice 속임수를 쓴 미성년자까지 보호하면 거래 상대방의 신뢰를 침해하게 되기 때문이다.
\kfChoice 속임수를 쓸 수 있을 만큼 지능이 높으면 성년자와 동일하게 취급해야 하기 때문이다.
\kfChoice 법정대리인이 속임수를 사전에 허락한 것으로 간주되기 때문이다.
\kfChoice 미성년자의 모든 법률 행위는 성년이 된 이후에 취소해야 하기 때문이다.
\end{kfChoices}
\end{kfQBlock}

\kfEndMC

\kfPassageNote{민법의 기본 원칙 (사회)}{%
민법은 개인 간의 재산 관계와 가족 관계를 규율하는 법률로서, 근대 이후 몇 가지 기본 원칙 위에 세워져 왔다. 대표적인 것이 사적 자치의 원칙, 과실 책임의 원칙, 소유권 절대의 원칙인데, 이 세 원칙은 근대 시민 사회의 자유와 평등이라는 이념을 법적으로 구현한 것이다.

사적 자치의 원칙은 개인이 자신의 법률관계를 자유로운 의사에 따라 형성할 수 있다는 원칙이다. 이를 구체적으로 실현하는 것이 계약 자유의 원칙으로, 계약을 체결할지 여부, 누구와 계약할지, 어떤 내용으로 계약할지를 당사자가 자유롭게 결정할 수 있다. 국가는 원칙적으로 이러한 사적 영역에 개입하지 않으며, 개인의 자율적 결정을 존중한다.

과실 책임의 원칙은 타인에게 손해를 끼친 경우, 그 행위에 고의나 과실이 있을 때에만 배상 책임을 진다는 원칙이다. 즉, 아무런 잘못 없이 타인에게 손해가 발생하였더라도 고의나 과실이 없다면 배상할 의무가 없다는 것이다. 이 원칙은 개인의 자유로운 활동을 보장하기 위한 것으로, 고의나 과실 없이 발생한 손해까지 책임을 지게 하면 개인의 활동이 지나치게 위축될 수 있기 때문이다.

소유권 절대의 원칙은 개인의 소유권은 국가나 다른 개인이 함부로 침해할 수 없으며, 소유자가 자신의 재산을 자유롭게 사용, 수익, 처분할 수 있다는 원칙이다. 근대 시민 혁명 이전에는 국왕이나 영주가 개인의 재산을 임의로 몰수하거나 제한할 수 있었는데, 소유권 절대의 원칙은 이러한 자의적 침해로부터 개인의 재산권을 보호하기 위해 확립되었다.

그러나 자본주의가 발전하면서 이 세 원칙의 무제한적 적용이 사회적 불평등과 약자에 대한 착취를 초래할 수 있다는 문제가 드러났다. 이에 따라 현대 민법은 이 원칙들을 수정하여 적용하고 있다. 사적 자치의 원칙은 공공복리에 반하는 계약을 무효로 하거나 경제적 약자를 보호하는 방향으로 수정되었고, 과실 책임의 원칙은 제조물 책임이나 환경 오염 등의 분야에서 과실 없이도 책임을 지는 무과실 책임의 법리가 도입되었다. 소유권 절대의 원칙 역시 공공복리를 위해 소유권의 행사가 제한될 수 있다는 방향으로 수정되었는데, 대한민국 헌법 제23조 제2항은 '재산권의 행사는 공공복리에 적합하도록 하여야 한다'고 규정하여 이를 명시적으로 뒷받침하고 있다.
}\kfResetAreaFlag
\kfSecHeader{kfAreaNonlit}{문제}{}

\kfBeginMC
\begin{kfQBlock}{01}{객관식}
\kfQStem{윗글의 내용과 일치하지 않는 것은?}
\begin{kfChoices}
\kfChoice 사적 자치의 원칙은 계약 자유의 원칙을 통해 구체적으로 실현된다.
\kfChoice 과실 책임의 원칙에 따르면 고의나 과실이 없으면 배상 의무가 없다.
\kfChoice 소유권 절대의 원칙은 근대 시민 혁명 이전부터 확립되어 있었다.
\kfChoice 현대 민법에서는 무과실 책임의 법리가 일부 분야에 도입되었다.
\kfChoice 헌법 제23조 제2항은 재산권 행사가 공공복리에 적합해야 한다고 규정한다.
\end{kfChoices}
\end{kfQBlock}

\begin{kfQBlock}{02}{객관식}
\kfQStem{윗글에서 '과실 책임의 원칙'이 근대 민법에 자리 잡은 까닭으로 가장 적절한 것은?}
\begin{kfChoices}
\kfChoice 피해자의 손해를 빠짐없이 배상하기 위해서이다.
\kfChoice 개인의 자유로운 활동을 보장하기 위해서이다.
\kfChoice 경제적 약자를 착취로부터 보호하기 위해서이다.
\kfChoice 국가의 재산 몰수를 방지하기 위해서이다.
\kfChoice 모든 손해에 대해 공동 책임을 지도록 하기 위해서이다.
\end{kfChoices}
\end{kfQBlock}

\begin{kfQBlock}{03}{객관식}
\kfQStem{윗글을 바탕으로 <보기>를 이해한 내용으로 적절하지 \underline{않은} 것은?}
\begin{kfEvidence}한 제조 회사가 생산한 가전제품의 결함으로 인해 소비자 B가 화상을 입었다. 이 회사는 제품 제조 과정에서 과실이 없었음을 주장하였으나, 법원은 제조물 책임법에 의거하여 회사에 손해배상 책임을 인정하였다.\end{kfEvidence}
\begin{kfChoices}
\kfChoice 이 판결은 과실 책임의 원칙이 수정된 사례에 해당한다.
\kfChoice 이 사례에는 무과실 책임의 법리가 적용되었다고 볼 수 있다.
\kfChoice 이러한 수정은 자본주의 발전에 따른 사회적 불평등 문제와 관련이 있다.
\kfChoice 제조물 책임은 과실 책임의 원칙을 더욱 강화하여 적용한 결과이다.
\kfChoice 피해자인 소비자 B의 구제를 위해 제조사에 배상 책임이 인정된 사례이다.
\end{kfChoices}
\end{kfQBlock}

\begin{kfQBlock}{04}{객관식}
\kfQStem{윗글에서 민법의 세 가지 기본 원칙과 그 현대적 수정 내용을 짝지은 것으로 적절한 것은?}
\begin{kfChoices}
\kfChoice 사적 자치 -- 무과실 책임 도입, 과실 책임 -- 공공복리에 따른 제한, 소유권 절대 -- 사회적 약자 보호 강화
\kfChoice 사적 자치 -- 약자 보호와 공공복리 위반 계약 무효, 과실 책임 -- 무과실 책임 도입, 소유권 절대 -- 공공복리 제한
\kfChoice 사적 자치 -- 공공복리에 의한 강한 제한, 과실 책임 -- 사회적 약자 보호 강화, 소유권 절대 -- 무과실 책임의 도입
\kfChoice 사적 자치 -- 소유권 행사의 전면적 제한, 과실 책임 -- 계약 자유 축소, 소유권 절대 -- 손해 배상 책임 확대
\kfChoice 사적 자치 -- 무과실 책임의 도입, 과실 책임 -- 사회적 약자 보호 강화, 소유권 절대 -- 위반 계약 무효 처리
\end{kfChoices}
\end{kfQBlock}

\kfEndMC

\kfStartWeekly{패턴 워크북 — 총 18문항}


\kfPassageNote{문항 1 / 섞어치기}{%
육상 동물과 수생 동물은 각각의 생활 환경에 따라 호흡 과정에서 어려움을 겪는다. 육상 동물은 폐의 표면에 건조한 공기가 닿으면 폐가 건조해질 수 있는데, 이는 폐를 손상시키고 ⓐ탈수를 일으키며 기체를 용해하는 폐의 능력을 떨어뜨린다. 한편 수생 동물은 가용 산소가 적은 수중 환경과 수온 차에 의해 변화되는 산소 가용성에 적응해야 한다. 또한 수생 동물은 물속에서 몸을 움직이는 데 많은 근육 운동이 필요하고, 찬 바닷물의 높은 비열 때문에 아가미의 열을 빼앗기며, 삼투압 현상으로 인해 수분 균형에 문제가 발생하기도 한다.

(중략)

이러한 환경 속에서 수생 동물은 효과적인 호흡을 위해 아가미라는 특수한 기관을 사용한다. 아가미는 물속에 녹아 있는 산소를 흡수하고 체내에 있는 이산화 탄소를 배출하는 기능을 하는데, 외아가미와 내아가미로 나뉜다. 외아가미는 몸 밖으로 ⓑ돌출된 형태로, 표면적이 넓고 형태가 다양하며 유생기 도롱뇽이나 일부 무척추동물에서 나타난다. 외아가미는 주변 물과의 접촉을 극대화

(중략)
}
\begin{kfQBlock}{01}{객관식}
\kfQStem{ⓐ\textasciitilde{}ⓔ의 사전적 의미로 적절하지 \underline{않은} 것은?}
\begin{kfChoices}
\kfChoice ⓒ: 아껴서 모아 둠.
\kfChoice ⓐ: 몸속의 수분이 모자라서 일어나는 증상.
\kfChoice ⓑ: 쑥 내밀거나 불거져 있음.
\end{kfChoices}
\end{kfQBlock}

\kfPassageNote{문항 2 / 부정·상대어}{%
육상 동물과 수생 동물은 각각의 생활 환경에 따라 호흡 과정에서 어려움을 겪는다. 육상 동물은 폐의 표면에 건조한 공기가 닿으면 폐가 건조해질 수 있는데, 이는 폐를 손상시키고 ⓐ탈수를 일으키며 기체를 용해하는 폐의 능력을 떨어뜨린다. 한편 수생 동물은 가용 산소가 적은 수중 환경과 수온 차에 의해 변화되는 산소 가용성에 적응해야 한다. 또한 수생 동물은 물속에서 몸을 움직이는 데 많은 근육 운동이 필요하고, 찬 바닷물의 높은 비열 때문에 아가미의 열을 빼앗기며, 삼투압 현상으로 인해 수분 균형에 문제가 발생하기도 한다.

(중략)

이러한 환경 속에서 수생 동물은 효과적인 호흡을 위해 아가미라는 특수한 기관을 사용한다. 아가미는 물속에 녹아 있는 산소를 흡수하고 체내에 있는 이산화 탄소를 배출하는 기능을 하는데, 외아가미와 내아가미로 나뉜다. 외아가미는 몸 밖으로 ⓑ돌출된 형태로, 표면적이 넓고 형태가 다양하며 유생기 도롱뇽이나 일부 무척추동물에서 나타난다. 외아가미는 주변 물과의 접촉을 극대화

(중략)
}
\begin{kfQBlock}{02}{객관식}
\kfQStem{[A]를 바탕으로 <보기>를 이해한 내용으로 적절하지 \underline{않은} 것은?}
\begin{kfEvidence}(아가미의 구조를 나타낸 모식도: 아가미궁(A), 라멜라(B), 필라멘트(C), 원심성 혈관(D), 구심성 혈관(E), 덮개 공간, 모세 혈관, 물 흐름의 방향 표시)\end{kfEvidence}
\begin{kfChoices}
\kfChoice D의 혈액은 E의 혈액에 비해 상대적으로 적은 산소를 함유하고 있겠군.
\kfChoice C가 붙어 있는 A는 아가미를 지탱하는 역할을 하겠군.
\kfChoice B의 모세 혈관에서 흐르는 혈액은 아가미 내부로 들어온 물의 흐름과 반대 방향으로 흐르겠군.
\end{kfChoices}
\end{kfQBlock}

\kfPassageNote{문항 3 / 주체·객체}{%
(가) 인과 관계에 관한 ㉮조건설에 의하면 어떤 범죄적 행위는 그것이 특정한 범죄 결과를 필연적으로 야기하는 전체 조건에 포함된 하나의 조건일 때, 그 결과에 대한 원인이 된다. 조건설은 이와 같이 조건의 전체적 복합을 이루는 개개의 조건들을 모두 동등하게 결과 발생의 원인으로 취급하므로 등가설이라고도 불리는데, 이러한 조건들 중에 어느 하나라도 결여된다면 문제의 결과는 발생하지 않게 된다.

조건설에서는 이른바 조건 관계의 발견을 위해 선행 사실이 없었다고 가정할 경우 후행 사실이 발생했을 것인가를 묻는 조건 공식을 사용하고 있다. 당해 행위를 제거하고 생각해 봤을 때 문제의 결과가 발생하지 않을 것으로 판단된다면, 행위와 결과 사이에는 조건 관계가 성립한다고 보는 것이다. 이러한 ⓐ반사실적 사고 실험의 절차는 일종의 가설적 소거법을 따르는 것이라 할 수 있다.

(중략)
}
\begin{kfQBlock}{03}{객관식}
\kfQStem{(가)와 (나)를 바탕으로 <보기>를 이해한 내용으로 가장 적절한 것은?}
\begin{kfEvidence}ㄱ. X의 물건 위에 그것이 더 이상 버티지 못하고 비가역적인 손상을 입기 시작하게 만드는 무게의 95\%에 해당하는 하중을 A가 가하는 순간, 그러한 무게의 나머지 5\%에 해당하는 하중을 B가 가했다. 이후 X의 물건은 곧 손상되었다. \\\relax
ㄴ. A는 X가 평소 무게 50kg의 추를 여러 개 보관하던 창고에 그것과 모양이 비슷하지만 무게가 90kg에 달하는 추 하나를 놓아두었는데, 마침 추를 사용하려던 X는 60kg까지 버틸 수 있는 기계로 A가 놓아둔 추를 들어 올렸다. 당시 X가 고를 수 있었을 나머지 모든 추들은 B가 A와 마찬가지로 모양이 비슷한 무게 90kg의 추들로 바꿔 둔 상태였다. X가 추를 들어 올리자 곧 기계가 추의 무게를 견디지 못하고 고장을 일으켰다. \\\relax
ㄷ. A는 X가 평소 무게 50kg의 추를 보관하던 창고에 그것과 모양이 비슷하지만 무게가 90kg에 달하는 추를 놓아두었다. 마침 추를 사용하려던 X가 60kg까지 버틸 수 있는 기계로 A가 놓아둔 추를 집으려는 순간, 전날 B가 창고 지붕을 수리하며 허술하게 쌓아 두었던 자재가 기계 위로 쏟아지면서 기계가 고장을 일으켰다. (모든 경우에 A와 B는 서로 알지 못하며, 협력하거나 공모한 적도 없다.)\end{kfEvidence}
\begin{kfChoices}
\kfChoice ㄴ: 조건설에 따르면 A의 행위를 X의 기계가 고장 나게 한 원인으로 보지 않을 것이다.
\kfChoice ㄱ: 조건설에 따르면 A의 행위를 X의 물건이 손상되게 한 원인으로 보지 않을 것이다.
\kfChoice ㄱ: 조건설에 따르면 B의 행위를 X의 물건이 손상되게 한 원인으로 보지 않을 것이다.
\end{kfChoices}
\end{kfQBlock}

\kfPassageNote{문항 4 / 순서·방향}{%
에너지 대사는 생명 유지를 위해 필요한 에너지를 얻고 사용하는 과정으로 인체에 필요한 에너지는 대부분 음식물 섭취를 통해 이루어진다. 섭취를 통해 체내에 ⓐ유입된 음식물은 그대로 에너지로 소비되지 않고 소화 과정을 거쳐 체내 각 조직에 공급된다. 음식물에 함유된 탄수화물, 단백질, 지방은 각각 일련의 과정을 거쳐 세포 활동에 직접 사용되는 에너지원인 ATP로 전환된 후 필요한 조직에 사용된다. 에너지로 바로 사용되지 않은 탄수화물 일부는 포도당의 형태로 간이나 근육에 저장되었다가 필요할 때 다시 ATP로 전환된다. 섭취된 에너지와 인체가 소비하는 에너지는 균형을 유지해야 하지만, 섭취한 에너지의 양이 소비되는 양보다 많아지는 상태가 지속되면 남는 에너지는 지방산으로 전환되어 지방 세포에 저장된다. 이러한 상태가 ⓑ심화되면 지방 세포의 크기가 커지게 된다. 지방산이 과도하게 축적되면 체내의 에너지 조절 체계가 정상적으로 작동하지 않게 되고 비만을 비롯한 다양한 대사 질환의 발생 가능성

(중략)
}
\begin{kfQBlock}{04}{객관식}
\kfQStem{[에너지 항상성]에 대한 설명으로 적절하지 \underline{않은} 것은?}
\begin{kfChoices}
\kfChoice 지방 세포가 충분히 커진 뒤에야 남는 에너지가 지방산으로 전환되기 시작한다.
\kfChoice 시상 하부 문제가 생기면 식욕 조절이 어려워질 수 있다.
\kfChoice 에너지 소비가 원활하지 않을 때 섭취를 줄이고 소비를 늘리려는 방향이 나타날 수 있다.
\end{kfChoices}
\end{kfQBlock}

\kfPassageNote{문항 5 / 인과 도치}{%
대중적 정부의 형태를 채택하는 국가는 항상 시민 사회의 분열과 파벌 형성이라는 위험에 노출되어 있다. 여기서 파벌이란 공동체 전체의 이익과 일치하지 않는, 일부 구성원들만을 위한 이해관계에 기초하여 결합하고 행동하는 시민 집단을 뜻한다. 이들이 비교적 소수의 인원들로 구성되어 있을 경우에는 다수의 시민들이 선거에 의하여 그들의 준동을 좌절시킬 수 있을 것이나, 그와 반대의 상황이라면 파벌이 공공선 및 파벌에 속하지 않은 시민들의 권리를 자신들의 정념(情念)과 이익을 위해 희생시키는 결과가 초래될 수 있다.

(중략)

미국 건국의 아버지 중 한 사람으로 알려진 제임스 매디슨은 바로 이와 같은 문제가 사회적 다수의 지배를 그 핵심으로 하는 ㉮민주주의에 내재해 있는 치명적인 약점이라고 생각했다. 누구도 자기 자신의 송사에서 재판관이 될 수 없다는 것은 정의의 관점에서 지극히 당연한데, 이 말이 의미하는 바를 곱씹어 볼 때 사회적 다수가 자신들의 이익이 걸려 있는 상황에서 정치적 권력을 중립적으로 행사

(중략)
}
\begin{kfQBlock}{05}{객관식}
\kfQStem{㉮에 대한 설명으로 가장 적절한 것은?}
\begin{kfChoices}
\kfChoice 당파적 이익을 우선시하는 세력이 정치적 다수를 형성하면 그 폐단을 막을 수 없다.
\kfChoice 학습과 수양이 부족한 시민들도 정치적 권력의 행사에 직접 참여한다.
\kfChoice 시민들이 자신의 정치적 힘을 적절한 한도 내에서 절제하며 발휘할 수 없다.
\end{kfChoices}
\end{kfQBlock}

\kfPassageNote{문항 6 / 의도·목적}{%
미니멀리즘 음악은 1960년대 이후 미국에서 태동한 현대 음악의 한 흐름으로, 라 몬테 영, 테리 라일리, 스티브 라이히, 필립 글래스 등이 대표적 작곡가로 꼽힌다. 이들은 소수의 \uline{음형}(몇 개의 음이 연속하여 어떤 가락이나 악곡의 요소가 되는 음의 모양)이나 단순한 \uline{패턴}을 반복하는 방식으로 기존 음악의 복잡한 서사나 화성 전개를 과감히 축소하여, 감상자가 소리 그 자체를 경험하도록 유도한다.

전통적인 서양 음악에서 반복은 주제를 강조하여 이를 서사적으로 강화하기 위해 사용되었다. 예컨대 제시부-발전부-재현부로 구성된 소나타 형식은 제시부에서 등장한 주제가 발전부에서 여러 차례 변주되다가 재현부에서 제시부에 등장한 형태로 회귀함으로써, 작곡가가 설정한 주제를 확인 및 강화하는 구조를 이룬다. 철학자 \uline{들뢰즈}의 표현을 빌리면, 이와 같은 ㉠ \uline{‘같은 것의 반복’}에서는 특정 주제가 중심에 놓이고 음악의 각 부분은 그 주제를 재확인하는 역할만…(중략)
}
\begin{kfQBlock}{06}{객관식}
\kfQStem{다음 중 ㉠ ‘같은 것의 반복’과 ㉡ ‘차이를 생성하는 반복’을 이해한 내용으로 적절하지 \underline{않은} 것은?}
\begin{kfChoices}
\kfChoice ㉡에서 반복은 특정 주제를 더 선명하게 전달하기 위한 장치이므로, 감상자는 작곡가가 설정한 의미를 재확인하는 데 초점을 둔다.
\kfChoice ㉡에서 반복은 특정 주제를 더 선명하게 전달하기 위한 장치이므로, 감상자는 작곡가가 설정한 의미를 재확인하는 데 초점을 둔다.
\kfChoice ㉠에서는 특정 주제가 중심에 놓이고, 반복은 그 주제를 확인·강화하는 역할을 한다.
\kfChoice ㉡에서는 표면적으로 같은 음형이 반복되더라도 음색·강세 등 미세한 차이가 드러나 새로운 사건이 생길 수 있다.
\end{kfChoices}
\end{kfQBlock}

\kfPassageNote{문항 7 / 상관관계}{%
(가) \uline{건축 설계}는 공간을 계획하거나 건물의 외형을 아름답게 만드는 것만을 의미하지 않는다. \uline{건축 설계}는 \uline{구조적 안정성}, \uline{기능적 효율성}, \uline{미적 가치}를 위한 과학이자 예술이다. \uline{건축 설계}가 과학이자 예술로 기능하는 데 있어 가장 중요한 두 가지는 \uline{건축 구조}와 \uline{건축 재료}이다.

\uline{건축 구조}는 건물에 가해지는 \uline{하중}을 \uline{㉠지탱}하며 형체를 유지하도록 설계된다. \uline{하중}이란 구조물에 작용하는 모든 힘으로, \uline{정하중}과 \uline{활하중}으로 구분된다. \uline{정하중}은 건축물 자체의 무게, 즉 기둥, 보, 슬래브, 벽체 등 영구적으로 작용하는 무게를 뜻한다. \uline{활하중}은 사람이 건축물 안에서 이동하거나 가구와 같은 \uline{가변적}인 물체가 추가되면서 발생하는 하중을 의미한다. 건축 구조는 이 하중을 지탱할 수 있도록 설계되어야 …(중략)
}
\begin{kfQBlock}{07}{객관식}
\kfQStem{(가)를 바탕으로 추론한 내용으로 가장 적절한 것은?}
\begin{kfChoices}
\kfChoice 철근 콘크리트는 콘크리트의 인장 약점을 보완하므로, 동일 조건에서 콘크리트 단독보다 균열 위험을 낮출 수 있다.
\kfChoice 철근 콘크리트는 콘크리트의 인장 약점을 보완하므로, 동일 조건에서 콘크리트 단독보다 균열 위험을 낮출 수 있다.
\kfChoice 응력은 단위 면적당 힘이므로, 같은 힘에서 단면적이 커지면 응력도 커진다.
\kfChoice 좌굴은 주로 인장력에 의해 부재가 늘어나며 생기는 현상이다.
\end{kfChoices}
\end{kfQBlock}

\kfPassageNote{문항 8 / 필요·충분 조건}{%
(가) 가상 공간에서 영생의 삶을 누린다는 흥미로운 이야기가 SF 영화의 소재로만 그치지 않는다고 주장하는 입장이 있다. 커즈와일이나 보스트룀은 기억, 가치관, 태도, 정서적 성향처럼 인간의 자아를 구성하는 것들을 특정 정보 패턴으로 만들어 보존할 수 있다면 신체적인 죽음 이후에도 인간이 여전히 생존할 수 있다고 주장한다. 이 입장은 인격 동일성에 관한 심리적 연속성 이론과 관련이 있다. 이 이론에서는 ‘나’가 누구인지를 규정하는 것은 ‘나’의 기억, 신념, 생각, 감정, 희망, 두려움 등의 묶음이라고 보며, 이러한 심리적인 특성들의 집합이 유지·보존되느냐에 ‘나’의 지속 여부가 달려 있다고 말한다.

커즈와일이나 보스트룀은 마음이나 정신 상태를 의미하는 심적 상태의 본성을 계산 기능주의의 관점에서 이해한다. 계산 기능주의에서는 심적 상태의 독특한 인과적 역할이 두뇌의 물리적 상태에 의해서 이루어진다고 보기 때문에 정신 작용을 두뇌에 구현된 계산 체계의 작동으로 이해하며, 인간의 사…(중략)
}
\begin{kfQBlock}{08}{객관식}
\kfQStem{[05] 윗글의 내용과 일치하지 않는 것은?}
\begin{kfChoices}
\kfChoice 인간의 뇌를 생물학적 컴퓨터로 보면 물리적 기반이 없더라도 다양한 방식으로 정신의 구현이 가능하다.
\kfChoice 인간의 뇌를 생물학적 컴퓨터로 보면 물리적 기반이 없더라도 다양한 방식으로 정신의 구현이 가능하다.
\kfChoice 커즈와일이나 보스트룀은 심리적인 특성들의 집합을 인간의 자아를 구성하는 요소로 여긴다.
\kfChoice 계산 기능주의의 관점에서는 두뇌가 수행하는 모종의 기능적 역할이 심적 상태를 만든다고 말한다.
\end{kfChoices}
\end{kfQBlock}

\kfPassageNote{문항 9 / 결합·분리}{%
현대 사회에서 디지털 금융 기술의 발달은 금융 거래의 속도와 편의성을 극적으로 향상시켰다. 모바일 뱅킹 등을 통해 몇 번의 터치만으로 송금이 이루어지는 시대가 도래했지만, 이러한 기술 발전은 예상치 못한 문제도 야기했다. 대표적인 사례가 바로 착오 송금이다. 착오 송금이란 송금인이 실수로 의도하지 않은 계좌에 돈을 보내는 경우를 말하며, 입력 오류, 수취인 착오 선택 등 다양한 원인에 의해 발생한다.

착오 송금이 발생했을 때, 송금인은 즉시 해당 은행에 이를 신고하고 반환을 요청하며, 은행은 중개자로서 수취인에게 연락하여 송금 실수임을 알리고 자발적 반환을 유도한다. 이때 수취인이 동의하면 간단하게 환급 절차가 이루어지지만, 반환을 거부하거나 연락이 닿지 않는 경우에는 상황이 복잡해진다. 송금인이 직접 법적 절차를 진행해야 하기 때문이다. 착오 송금인, 즉 실수로 돈을 잘못 보낸 사람은 그 실수를 바로잡기 위해 직접 회수를 시도하지만, 법적 대응에 들어가는 비용과 시간, 절차적 복…(중략)
}
\begin{kfQBlock}{09}{객관식}
\kfQStem{윗글에 대한 설명으로 적절하지 \underline{않은} 것은?}
\begin{kfChoices}
\kfChoice 글의 화제에 대해 여러 주체들이 갈등하는 모습을 구체적 사례를 들어 제시하고 있다.
\kfChoice 글의 화제에 대해 여러 주체들이 갈등하는 모습을 구체적 사례를 들어 제시하고 있다.
\kfChoice 글의 화제에 해당하는 용어의 개념을 설명하고 있다.
\kfChoice 글의 내용과 관련된 구체적인 법률 조항을 제시하고 있다.
\end{kfChoices}
\end{kfQBlock}

\kfPassageNote{문항 10 / 조건 왜곡}{%
개별적이고 특수한 작업을 수행하는 여러 대의 클라이언트와 연결되어 정보나 서비스를 제공하는 컴퓨터 시스템을 서버라고 한다. 서버는 클라이언트의 요청에 따라 작업을 수행하거나 정보를 검색하고, 그 결과를 클라이언트에 돌려준다. 서버를 구성하는 기계에서 작업을 신속하고 정확하게 처리하기 위해서는 데이터를 안전하게 관리해야 할 뿐 아니라 신속하게 데이터에 접근하여 데이터의 읽기나 쓰기를 실행할 필요가 있다. 이런 필요를 충족시키기 위해 개발된 것이 \uline{㉠분산 파일 관리 체계}이다. 분산 파일 관리 체계는 데이터의 불법 변조를 방지하고, 기계적 결함에 대응하여 데이터의 유실로부터 안전을 도모하며, 서버의 부하를 줄여 신속한 연산을 수행하도록 기계적 성능과 저장 용량을 확보하는 데 유리한 전략이다. 단일 기계로 구성된 서버에서 데이터를 관리할 경우에는 하드웨어적으로 램이나 CPU의 성능이 우수한 기계를 확보하는 데 큰 비용이 들지만, 분산 파일 관리 체계에서는 다수의 저렴한 저성능 기…(중략)
}
\begin{kfQBlock}{10}{객관식}
\kfQStem{<보기>를 참고해 ⓐ의 운영 원리를 이해한 내용으로 적절하지 \underline{않은} 것은?}
\begin{kfEvidence}클러스터 X는 마스터-슬레이브 복제를 사용한다. 읽기 요청은 마스터와 슬레이브 모두에서 처리하고, 쓰기 요청은 마스터에서만 처리한 뒤 슬레이브로 반영한다. \\\relax
클러스터 Y는 피어 투 피어 복제를 사용한다. 읽기와 쓰기 요청을 바쁘지 않은 피어에서 처리하고, 수정된 데이터는 다른 피어들로 복제한다.\end{kfEvidence}
\begin{kfChoices}
\kfChoice Y에서는 데이터 일관성을 위해 쓰기 요청을 반드시 하나의 중앙 피어에서만 처리해야 한다.
\kfChoice Y에서는 데이터 일관성을 위해 쓰기 요청을 반드시 하나의 중앙 피어에서만 처리해야 한다.
\kfChoice X에서는 쓰기 요청이 마스터에 집중되므로 쓰기 경로 병목 가능성을 점검해야 한다.
\kfChoice X에서는 마스터 반영 후 슬레이브 전파 지연 동안 일시적 불일치가 발생할 수 있다.
\end{kfChoices}
\end{kfQBlock}

\kfPassageNote{문항 11 / 긍부정 반전}{%
[16\textasciitilde{}20] 다음 글을 읽고 물음에 답하시오.

(가)

헌 누더기 입은 무리가 남자인지 여자인지    어린 자식 등에 업고 자란 자식 손에 끌고    울면서 눈물 씻고 엎어지며 오는 모양    차마 보지 못할네라 나직이 묻는 말씀    어디로서 쫓아오며 어디로 가려는고    주려들 가는 사람인가 가게 되면 얻어 먹나    아무 데도 한가지라 날 따라 도로 가면    자네 원님 가서 보고 안접(安接)하게 하여줌세    겨우겨우 대답하되 우리 곳은 당진(唐津)이라    여러 해 흉년 들어 살길이 없는 중에    도망한 자 신구환(新舊還)을 있는 자에 물리니    제 것도 못 바치며 남의 곡식 어찌할꼬    못 바치면 매 맞으니 매 맞고 더욱 살까    정처 없이 가게 되면 죽을 줄 알건마는    아니 가고 어찌하리 굶고 맞고 죽을 지경    차라리 구렁*에나 염려 없이 뭇치이면   도리어 편할지라 이런 고로 가노메라    급히 급히 넘어가자 이 백성들 살려보세    둘째 령(嶺)을 …(중략)
}
\begin{kfQBlock}{11}{객관식}
\kfQStem{(나)에 대한 설명으로 적절하지 \underline{않은} 것은?}
\begin{kfChoices}
\kfChoice <2수>: 강산을 즐기느라 임금에게 가지 못하는 상황에 대한 미안함을 드러내고 있다.
\kfChoice <2수>: 강산을 즐기느라 임금에게 가지 못하는 상황에 대한 미안함을 드러내고 있다.
\kfChoice <1수>: 돌아온 고향에서 변해 버린 인사(人事)에 대한 슬픔을 나타내고 있다.
\kfChoice <4수>: 세속의 어지러운 소식을 모른 체하며 살고 싶은 심정을 표현하고 있다.
\end{kfChoices}
\end{kfQBlock}

\kfPassageNote{문항 12 / 비교 대상}{%
[26\textasciitilde{}31] 다음 글을 읽고 물음에 답하시오.

(가)   십 년 종사 후에 고향으로 도라오니   산천 의구하되 인사(人事)는 달라졌구나   아마도 세간의 존멸을 못내 슬허 하노라 <제1수>

산화(山花)는 믈의 픠고 물새는 산의 운다   일신이 한가하야 산수간의 누어시니   세상의 어즈러은 긔별을 나는 몰라 하노라 <제4수>

거믄고 빗기 들고 산수를 희롱하니   청풍은 건듯 블고 명월도 도라온다   하물며 유신(有信)한 갈매기는 오명 가명 하나니 <제5수>

거믄고 흥진(興盡)커던 조대(釣臺)로 내려가니   도화 뜬 말근 믈 뛰노나니 고기로다   아이야 밋기 다지 마라 취적(取適)*이나 하오리라 <제7수>

* 신교, ｢귀산음(歸山吟)｣ -   * 취적 : 낚시질의 참뜻이 세상 생각을 잊고자 하는 데 있음.

(나)   백수(白首)에 산수 구경 늦은 줄 알지마는   평생 품은 뜻을 이루고야 말리라 여겨   병자년 봄에 봄옷을 새로 입고   죽장망혜(竹杖芒鞋)로 노계 깊은 골에 …(중략)
}
\begin{kfQBlock}{12}{객관식}
\kfQStem{(가)\textasciitilde{}(다)의 공통점으로 가장 적절한 것은?}
\begin{kfChoices}
\kfChoice 구체적인 경험을 바탕으로 지향하는 삶의 모습을 드러내고 있다.
\kfChoice 구체적인 경험을 바탕으로 지향하는 삶의 모습을 드러내고 있다.
\kfChoice 과거의 삶을 후회하며 이상적 세계에 대한 동경을 드러내고 있다.
\kfChoice 역사적 사실을 언급하며 상황에 대한 비판적 시각을 드러내고 있다.
\end{kfChoices}
\end{kfQBlock}

\kfPassageNote{문항 13 / 주체·객체}{%
\#\#\# [18～21] 다음 글을 읽고 물음에 답하시오.

제1회 봄놀이   오작교에선 선랑(仙郞)이 봄바람에 취하고   버드나무 언덕에선 가인(佳人)이 그네를 뛰네

‘광한루기’는 작품 전체의 제목이다. 광한루가 없었더라면 [A] 이도린이 놀러 가지 않았을 것이요, 이도린이 놀러 가지 않았더라면 춘향이 이도린을 만날 수 없었을 것이요, 춘향이 이도린을 만나지 못했더라면 8회로 구성된 한 편의 작품이 무엇을 바탕으로 탄생할 수 있었겠는가. 광한루 하나가 공중에 솟구쳐 있었기에 이도린이 놀러 갈 수밖에 없었고, 춘향이 이도린을 만날 수밖에 없었으며, 8회로 구성된 한 편의 작품이 만들어질 수밖에 없었다.   (중략)   그네 뛰는 모습을 이도린이 보고 자기도 모르게 눈앞이 어질어질하여 김한에게 말했다.   “너는 저런 것을 본 적이 있느냐? 저것이 금이냐, 옥이냐? 아니면 귀신이냐? 그것도 아니면 선녀냐? 너는 저것을 아느냐?”   김한이 대답했다.   “금도 아니고 옥도 아닙니다. 낙…(중략)
}
\begin{kfQBlock}{13}{객관식}
\kfQStem{윗글에서 ‘김한’의 역할을 이해한 것으로 가장 적절한 것은?}
\begin{kfChoices}
\kfChoice 춘향에게 이도린과의 만남은 거듭된 우연으로 이루어진 인연임을 알려 주어, 두 사람을 만나게 하는 매개자 역할을 한다.
\kfChoice 춘향에게 이도린과의 만남은 거듭된 우연으로 이루어진 인연임을 알려 주어, 두 사람을 만나게 하는 매개자 역할을 한다.
\kfChoice 이도린에게 눈앞에 보이는 것이 금과 옥이 아니라고 알려 주어, 이도린의 무지를 일깨우는 비판자 역할을 한다.
\kfChoice 이도린에게 춘향이 선녀 같은 아가씨라고 말하여, 이도린이 춘향의 고귀한 신분을 알게 하는 조력자 역할을 한다.
\end{kfChoices}
\end{kfQBlock}

\kfPassageNote{문항 14 / 내용·방법}{%
[18 \textasciitilde{} 23] 다음 글을 읽고 물음에 답하시오.

(가)   삼 년을 임을 떠나 해도(海島)에 유배되니   ㉠ 내 언제 무심하여 임에게 득죄했나   임이 언제 박정(薄情)하여 날 대접 소홀히 했나   내 얼굴 고왔던지 질투하는 건 뭇 여자로다   유한한* 이내 몸을 음란하다 이르로세   (중략)   긴 소매 들고 앉아 옛 잘못을 헤아리니   우직하기 본성이오 망령됨도 내 죄로되   근본을 생각하면 임 위한 정성일세   일월 같은 우리 임이 거의 아니 굽어볼까   날 살리신 이 은혜를 결초(結草)하기 생각하나   광주리의 가을 부채 어느 날 다시 날꼬   황금을 못 얻으니 장문부*를 어이 사리*   \textit{마름과 연(蓮)으로 옷을 짓고 부용(芙蓉)으로 치마 지어}   \textit{상자 안에 두어신들 눌 위하여 단장할꼬}   \textit{고향에 돌아갈 꿈 벽해(碧海)를 밟아 건너}   \textit{옥루(玉樓) 높은 곳에 밤마다 임을 모셔}   \textit{일당우불에 수답이 여향하니}   가까이 다가앉아 귀신을 묻던 가태부 이러한가…(중략)
}
\begin{kfQBlock}{14}{객관식}
\kfQStem{㉠ \textasciitilde{} ㉤의 표현상의 특징으로 적절하지 \underline{않은} 것은?}
\begin{kfChoices}
\kfChoice ㉢ : 과장법을 사용하여 임을 향한 사랑을 포기해야 하는 것에 대한 화자의 절망감을 강조하고 있다.
\kfChoice ㉢ : 과장법을 사용하여 임을 향한 사랑을 포기해야 하는 것에 대한 화자의 절망감을 강조하고 있다.
\kfChoice ㉠ : 대구적 표현을 사용하여 운율감을 조성하고 있다.
\kfChoice ㉡ : 감각적 심상을 활용하여 화자의 그리움을 부각하고 있다.
\end{kfChoices}
\end{kfQBlock}

\kfPassageNote{문항 15 / 내용 왜곡}{%
[26\textasciitilde{}31] 다음 글을 읽고 물음에 답하시오.

(가)   십 년 종사 후에 고향으로 도라오니   산천 의구하되 인사(人事)는 달라졌구나   아마도 세간의 존멸을 못내 슬허 하노라 <제1수>

산화(山花)는 믈의 픠고 물새는 산의 운다   일신이 한가하야 산수간의 누어시니   세상의 어즈러은 긔별을 나는 몰라 하노라 <제4수>

거믄고 빗기 들고 산수를 희롱하니   청풍은 건듯 블고 명월도 도라온다   하물며 유신(有信)한 갈매기는 오명 가명 하나니 <제5수>

거믄고 흥진(興盡)커던 조대(釣臺)로 내려가니   도화 뜬 말근 믈 뛰노나니 고기로다   아이야 밋기 다지 마라 취적(取適)*이나 하오리라 <제7수>

* 신교, ｢귀산음(歸山吟)｣ -   * 취적 : 낚시질의 참뜻이 세상 생각을 잊고자 하는 데 있음.

(나)   백수(白首)에 산수 구경 늦은 줄 알지마는   평생 품은 뜻을 이루고야 말리라 여겨   병자년 봄에 봄옷을 새로 입고   죽장망혜(竹杖芒鞋)로 노계 깊은 골에 …(중략)
}
\begin{kfQBlock}{15}{객관식}
\kfQStem{ⓐ와 ⓑ에 대한 이해로 가장 적절한 것은?}
\begin{kfChoices}
\kfChoice ⓐ는 하늘의 모습을 물에서 보게 된 것에 대한, ⓑ는 산의 모습이 평소와 달리 보이는 것에 대한 반응이다.
\kfChoice ⓐ는 하늘의 모습을 물에서 보게 된 것에 대한, ⓑ는 산의 모습이 평소와 달리 보이는 것에 대한 반응이다.
\kfChoice ⓐ는 하늘과 물의 변함없는 모습을 본 것에 대한, ⓑ는 선명하게 드러난 산의 모습을 본 것에 대한 반응이다.
\kfChoice ⓐ는 하늘이 물의 모습을 닮아 변해 가는 것에 대한, ⓑ는 산이 주변의 모습을 닮아 변해 가는 것에 대한 반응이다.
\end{kfChoices}
\end{kfQBlock}

\kfPassageNote{문항 16 / 과잉 해석}{%
[38\textasciitilde{}42] 다음 글을 읽고 물음에 답하시오.

(가)   장마 가뭄에 피해 입은 백성이 관찰사 가을 순행 기다림은   가을걷이 부족함을 채워줄까 해서인데 지나는 곳마다 죄를 묻는 폐단 있네   무논 재해도 감췄는데 목화밭이야 거론할까   백 묘(畝)나 되는 벌건 땅에 백지징세 하는구나   인자한 우리 임금 곡식 한 묶음도 모래 덮일까 염려하는데   불쌍한 백성 논밭에다 좁은 길 넓히란다   각읍 관리 독촉하니 채찍 몽둥이 낭자하다   허다한 관인들이 대호(大戶) 소호(小戶)에 분담시켜   사방(四方) 부근 십 리 안에 닭과 개가 멸종하네   부자는 괜찮지만 가련한 이 가난한 자로다   해는 기울고 이정*은 저녁밥 재촉할 때   텅 빈 부엌에서 우는 아낙 발 구르며 하는 말이   방아품에 얻은 양식 한두 되 있건마는   채소도 있건마는 그릇은 누구에게 빌릴꼬   앞뒷집 돌아보니 섣달그믐에 시루 빌리는 격이로다   한 마을 닭과 개 다 먹어 치우고 집집마다 또 거둔단 말인가   대호…(중략)
}
\begin{kfQBlock}{16}{객관식}
\kfQStem{<보기>를 참고하여 (가)를 감상한 내용으로 적절하지 \underline{않은} 것은? [3점]}
\begin{kfChoices}
\kfChoice ‘유생’들이 ‘과거장’에서 ‘재주’를 ‘겨루는’ 것은 의로운 선비가 되기 위해 과거에 통과하기를 바라는 유생들의 기대를 드러낸다고 볼 수 있겠군.
\kfChoice ‘유생’들이 ‘과거장’에서 ‘재주’를 ‘겨루는’ 것은 의로운 선비가 되기 위해 과거에 통과하기를 바라는 유생들의 기대를 드러낸다고 볼 수 있겠군.
\kfChoice ‘이 놀이’를 ‘다시’ 하게 되면 백성들이 ‘못 살겠’다고 한 것은 지배 계층의 유흥을 위해 수탈을 당하는 백성들의 현실을 드러낸다고 볼 수 있겠군.
\kfChoice 백성들이 ‘집과 논밭’을 ‘다 팔고서’ 떠나는 것은 가렴주구로 인해 유랑의 길을 떠나야 하는 백성들의 고통스러운 현실을 드러낸다고 볼 수 있겠군.
\end{kfChoices}
\end{kfQBlock}

\kfPassageNote{문항 17 / 범위 오류}{%
\#\#\# [42 \textasciitilde{} 45] 다음 글을 읽고 물음에 답하시오.

[앞부분 줄거리] 정 소저는 계모 박 씨의 모함을 의심 없이 받아들인 아버지 정공 때문에 위기에 처하고, 집에서 나와 숨어 다니던 중 도적을 만나 강물에 몸을 던진다. 이때, 정혼자 조무(용홍)와 동생 조성이 정 소저를 우연히 발견하여 구출한다.

소저가 매우 놀라며 말하였다.   “내가 외가로 가지 않고 구차하게 길가에서 분주하게 다닌 것은 조숙모에게 부끄럽고, 아버지의 허물을 드러내고 싶지 않아서였다. 뜻밖에 저 공자들을 만나니 내가 차마 사실을 말하여 부끄러움을 더하겠는가? 은인의 덕이 산과 바다 같으나 차마 근본을 아뢰게 되어 저 집에서 우리 집의 허물을 알게 되면 매우 부끄럽게 될 것이다. 모름지기 너는 다만 대답하기를 내가 타향에서 떠돌아다니다가 서울의 친척을 찾으러 왔다가 도적을 만나 물에 빠져 죽을 뻔했다고 말하여라. 조 공자가 이미 우리가 여자인 줄을 알았으니 남녀는 구별이 있는 것이다. 생명을 구해준 은혜…(중략)
}
\begin{kfQBlock}{17}{객관식}
\kfQStem{[A]와 [B]에 대한 이해로 가장 적절하지 \underline{않은} 것은?}
\begin{kfChoices}
\kfChoice [A]와 [B]는 모두 과거에 일어난 일을 상대에게 언급하고 있다.
\kfChoice [A]와 [B]는 모두 과거에 일어난 일을 상대에게 언급하고 있다.
\kfChoice [A]는 [B]와 달리 상대의 행동에 변화를 촉구하고 있다.
\kfChoice [B]는 [A]와 달리 상대에게 다른 인물의 말을 전하고 있다.
\end{kfChoices}
\end{kfQBlock}

\kfPassageNote{문항 18 / 시점·화자}{%
\#\#\# [18 \textasciitilde{} 23] 다음 글을 읽고 물음에 답하시오.

(가)   녹양방초 언덕에 소 먹이는 저 아이야   인간 영욕을 아는가 모르는가   인생 백년이 풀 끝에 이슬이라   삼만 육천 일이 그 아니 보잘것없는가   하물며 장수 단명이 운명이어니 사생(死生)을 정할쏘냐   여관 같은 세상에 하루살이같이 나왔다가   ㉠ 공명도 못 이루고 초목같이 썩어지면   ⓐ 공산 백골이 그 아니 느꺼우냐   하늘의 뜻을 이어 법칙을 세움은 옛 성인의 사업이요   아름다운 이름을 후세에 전함은 대장부의 할 일이라   생애는 유한하고 사일(死日)은 무궁하니   유한한 생애로 썩지 않을 이름을   영구히 전하여 ⓑ 천지와 함께 무궁하려고   (중략)   어와 그 뉘신고 어떠한 사람인고   형용이 초췌하니 초나라 대부 굴원이신가   잔혼이 영락하니 학사 유자후이신가*   눈썹을 찡그리시니 근심이 많으신가   발끝으로 서시니 어디를 바라는고   아름다운 기약을 바라는가 이별의 슬픔이 중하신가   ⓒ 해…(중략)
}
\begin{kfQBlock}{18}{객관식}
\kfQStem{(다)에 대한 이해로 적절하지 \underline{않은} 것은?}
\begin{kfChoices}
\kfChoice 사람들이 ‘팽조나 노담의 장수를 부러워하’고 ‘요절한 이를 슬퍼하’는 것은 ‘몸이 장애물이 되’어 ‘삶은 낮, 죽음은 밤과 같다’고 여겨 삶과 죽음을 다른 것으로 인식하기 때문이다.
\kfChoice 사람들이 ‘팽조나 노담의 장수를 부러워하’고 ‘요절한 이를 슬퍼하’는 것은 ‘몸이 장애물이 되’어 ‘삶은 낮, 죽음은 밤과 같다’고 여겨 삶과 죽음을 다른 것으로 인식하기 때문이다.
\kfChoice 사람들이 ‘보석으로 치장하’는 것과 ‘길거리에서 구걸을 하’는 것을 ‘영예나 치욕’으로 여기는 것은 정신이 몸에 갇혀 ‘바깥의 상황에’ 영향을 받기 때문이다.
\kfChoice ‘나와 남이 서로 통하지 않게 되어’ ‘싸움이 번지고 혼란이 야기’된 것은 사람의 정신이 그 몸에 얽매여 ‘만물은 본래 하나’라는 사실을 잊게 되었기 때문이다.
\end{kfChoices}
\end{kfQBlock}



\end{document}
